% !TeX root = ../Book4.tex
% 纲要标题：# 第二编　复形、消解与导出函子
% 纲要位置：计划纲要.md，第 2561 行。

\BookPart{复形、消解与导出函子}{b4:part:02}
\BookPartGuide
短正合列经过 Hom 或张量积以后，可能不再正合。第一编已经说明这种失败发生在哪一端；本编进一步把失去的正合性表示为可以计算的对象。为此，先用投射模或内射模组成的正合列代替原来的模，再对整列施加函子。所得复形的同调记录了新的关系，由此得到 Ext、Tor 和一般导出函子。

\ProseParagraphBreak
读者需要掌握第一编的加性函子、核与余核、交换范畴和蛇引理，以及第二卷第三、六、七章的模正合列、张量积、投射模、内射模和平坦模。一般伴随的若干应用会在需要时回引。第三卷的局部环理论不参与本编基础构造；它在 Koszul 消解\footnote{科斯居尔（Jean-Louis Koszul，1921-2018），法国数学家，研究 Lie 代数上同调与微分几何，并提出以其命名的复形。}和同调维数的讨论中提供应用。

\ProseParagraphBreak
\chapref{b4:ch:05}从复形开始，构造上同调长正合列，并用 Euler 特征\footnote{欧拉（Leonhard Euler，1707-1783），瑞士数学家，在分析、数论和图论等领域留下奠基性成果。}说明为什么同调保留了交错和。\chapref{b4:ch:06}解释什么样的替换不会改变同调，同时区分同伦等价与拟同构。映射锥把一个映射是否为拟同构转化为一个复形是否正合；双复形则使两种逐步计算可以放在同一张分次排列中。这里统一使用升次的上链记号，模的投射消解另用递降下标，并明确给出二者的转换。

\ProseParagraphBreak
\chapref{b4:ch:07}建立消解的存在、比较和相容构造。比较定理控制更换消解产生的差别，马蹄引理使短正合列可以同时消解。\chapref{b4:ch:08}据此定义导出函子，证明其长正合列和泛性质，再研究 Ext、Tor 的两个变量。计算 \(\operatorname{Ext}_R^n(M,N)\) 时，可以对第一变量 \(M\) 作投射消解，或对第二变量 \(N\) 作内射消解；Tor 则可以对任一变量作投射消解。平衡性保证这些相应的计算得到同一结果，证明使用有限对角上的双复形。

\ProseParagraphBreak
\chapref{b4:ch:05}至\chapref{b4:ch:08}是后两编反复调用的复形、消解与导出计算基础。若主要目的是学习第三卷的局部环同调方法，可先读至\secref{b4:sec:0805}。该节已经证明正次数的维数平移，并说明合冲模的选择与导出不变量的自然性有何区别。\chapref{b4:ch:09}继续把 Ext 解释为扩张，讨论乘积，并建立主理想整环上的万有系数定理及 Künneth 公式。\chapref{b4:ch:10}从 Ext、Tor 的消失研究消解长度，再说明何时可以用更方便的零调对象代替内射对象。其基本论证接续\chapref{b4:ch:08}。复合函子谱序列的适用条件在本编末列出，证明见\chapref{b4:ch:15}。

\ProseParagraphBreak
本编反复使用 \(\mathbb Z/n\mathbb Z\) 的两项自由消解。最初它给出非零同调与连接同态的直接计算，随后用来求 Ext 和 Tor，再解释短正合扩张及系数变换。每次使用都应留意当前所求的是复形的同调、函子的导出对象，还是扩张的等价类。它们可以联系起来，却不是同一个定义。

\ProseParagraphBreak
本编对应 Weibel\cite[Chapters 1--4]{Weibel1994HomologicalAlgebra}：复形与同伦、消解与导出函子、Ext 和 Tor 及万有系数、同调维数分别见该书四章。各章另给节或定理定位；该书 Chapter 3 的逆极限导出函子在本卷\chapref{b4:ch:13}讨论。

% !TeX root = ../../Book4.tex
% 纲要标题：## 第6章　复形、同调与 Euler–Poincaré 不变量
% 纲要标签：b4:ch:05
% 纲要位置：计划纲要.md，第 3216 行。

\chapter{复形、同调与 Euler–Poincaré 不变量}
\label{b4:ch:05}
\BookChapterNumber{b4:ch:05}

\begin{chapterabstract}
本章建立复形、Hom 复形、张量复形与同调，说明长正合列和连接同态的构造，并以 Grothendieck 群表述有界复形的 Euler 特征。
\end{chapterabstract}

\begin{learninggoals}
\begin{itemize}
\item 从给定微分计算圈、边界及同调，并正确转换链指标、上链指标和移位。
\item 区分复形映射与任意分次映射，核对 Hom 与张量复形的微分平方。
\item 构造并计算短正合列的连接同态，解释长正合列的自然性。
\item 用短正合关系构造 \(K_0\)，说明 Euler–Poincaré 等式中的有限性条件。
\end{itemize}
\end{learninggoals}

考虑整数上的两步计算：先乘 \(2\)，再取模 \(2\) 的余类。第二步会把第一步的结果全部送到零，但它还会把哪些元素送到零？答案正好是偶数，故这条序列在中间位置没有留下缺口。若只保留第一步而把末项改成零，中间位置的核变成整个 \(\mathbb Z\)，原来的像仍只有 \(2\mathbb Z\)，缺口便是 \(\mathbb Z/2\mathbb Z\)。复形把连续运算的复合为零作为起点；同调则逐处测量「成为零的元素」与「由前一步产生的元素」之间的差别。这个简单计算也提示，比较两个复形不能只比较它们的各项。

\ProseParagraphBreak
本章需要\chapref{b4:ch:04}的交换范畴、核与余核及蛇引理。模的算例另用第二卷第6.1节的张量积泛性质；有限长度版本用第二卷第12.2节已经证明的长度可加性，有限维向量空间则已足够呈现基本计算。复形的基本运算与长正合列可参阅 Weibel\cite[Sections 1.1--1.3, Theorem 1.3.1, Addendum 1.3.3, Proposition 1.3.4]{Weibel1994HomologicalAlgebra}；\(K_0\) 的关系与 Euler 特征见\cite[Chapter II, Definition 6.1.1, Proposition 6.6]{Weibel2013KBook}。移位和映射锥在文献中有不同符号，本书始终采用 \(C[1]^n=C^{n+1}\)、\(\mathrm{d}_{C[1]}=-\mathrm{d}_C\)。

% !TeX root = ../../Book4.tex
% 纲要标题：### 6.1 复形
% 纲要标签：b4:sec:0501
% 纲要位置：计划纲要.md，第 3218 行。

\section{复形}
\label{b4:sec:0501}

相邻映射的复合为零，只保证像包含在核中。同调把两者的差额组织成逐次数的不变量，链与上链的编号则须先固定转换规则。

% 纲要内容（仅作写作计划，尚非教材正文）：
%
% * 链复形与上链复形
% * 圈（cycle）、边界（boundary）及同调
% * 指标转换与次数平移约定
%

\subsection{链复形与上链复形}
\label{b4:subsec:0501:01}

一个序列若处处正合，核与像就完全一致；但许多计算只能保证相邻两步的复合为零。先保留这一较弱的条件，才有可能记录正合性在哪个次数失败，以及核中还剩下什么。

\begin{definition}\label{b4:def:cochain-complex}
设 \(\mathcal A\) 是交换范畴。给定对象 \(C^n\) 及态射 \(\mathrm{d}_C^n:C^n\to C^{n+1}\)（\(n\in\mathbb Z\)），若每个 \(n\) 都满足 \(\mathrm{d}_C^{n+1}\mathrm{d}_C^n=0\)，则称 \(C=(C^n,\mathrm{d}_C^n)\) 为 \(\mathcal A\) 中的\term[shanglian-fuxing]{上链复形}[cochain complex]，\(\mathrm{d}_C^n\) 称为微分。类似地，对象 \(C_n\) 及态射 \(\partial_n:C_n\to C_{n-1}\) 满足 \(\partial_{n-1}\partial_n=0\) 时，组成\term[lian-fuxing]{链复形}[chain complex]。
\end{definition}
\symidx{\(\mathrm{d}_C^n,\ \partial_n\)：上链复形与链复形的微分}

复形条件在每个次数只保证
\[
\mathrm{d}_C^n\mathrm{d}_C^{n-1}=0
\quad\Longleftrightarrow\quad
\operatorname{im}\mathrm{d}_C^{n-1}\subseteq\ker\mathrm{d}_C^n;
\]
正合还要求这里的像与核是同一个子对象。同调将记录这个包含还差多少。链与上链的差别在指标方向：链微分降一次，上链微分升一次。它们的核、像和正合性问题相同。为统一后面的 Hom 复形、移位和映射锥，本编主要用上链指标；讨论消解或 Tor 时也会使用链指标，并明确作转换。

\begin{example}\label{b4:exm:two-term-complex}
在左 \(R\)-模范畴中，任一同态 \(u:M\to N\) 都给出仅在次数 \(0,1\) 非零的复形 \(M\xrightarrow{u}N\)。此时 \(\mathrm{d}^0=u\)，下一微分是 \(\mathrm{d}^1=0_{N,0}:N\to0\)，所以所需条件为
\[
\mathrm{d}^1\mathrm{d}^0=0_{N,0}\circ u=0.
\]
其余相邻复合也均为零。不要求 \(u^2=0\)，因为 \(u\) 的目标是 \(N\)，而源是 \(M\)，一般根本不能把 \(u\) 与自身复合。相反，若每个次数均放同一模 \(M\)，并在每一步都用同一自同态 \(a\)，便确实需要 \(a^2=0\)。例如 \(R=k[\varepsilon]/(\varepsilon^2)\) 上反复乘 \(\varepsilon\) 给出复形。
\end{example}

\begin{definition}\label{b4:def:bounded-complex}
设 \(C\) 是交换范畴中的上链复形。若存在整数 \(a\)，使 \(n<a\) 时 \(C^n=0\)，则称 \(C\) 下有界；若存在 \(b\)，使 \(n>b\) 时 \(C^n=0\)，则称其上有界；同时满足二者时称其\term[youjie-fuxing]{有界复形}[bounded complex]。对象 \(M\) 放在次数零、其余项为零所得复形记为 \(M[0]\)。
\end{definition}
\symidx{\(M[0]\)：集中在零次的复形}

这里的上下指整数次数的大小，与箭头在版面上朝哪个方向无关。

\subsection{圈、边界与同调}
\label{b4:subsec:0501:02}

\begin{definition}\label{b4:def:cohomology}
设 \(C\) 是交换范畴 \(\mathcal A\) 中的上链复形。记
\[
Z^n(C)\coloneqq\ker \mathrm{d}_C^n,\qquad
B^n(C)\coloneqq\operatorname{im}\mathrm{d}_C^{n-1}.
\]
记像分解为 \(\mathrm{d}_C^{n-1}=i^n\pi^{n-1}\)，其中 \(\pi^{n-1}:C^{n-1}\to B^n(C)\) 满，\(i^n:B^n(C)\to C^n\) 单。由 \(\mathrm{d}_C^ni^n\pi^{n-1}=0\) 及 \(\pi^{n-1}\) 满，得到 \(\mathrm{d}_C^ni^n=0\)。因此，核嵌入 \(z^n:Z^n(C)\to C^n\) 的泛性质唯一给出 \(j^n:B^n(C)\to Z^n(C)\)，使
\[
z^nj^n=i^n.
\]
因为 \(i^n\) 单，\(j^n\) 也单。这就是余边对象到余圈对象的典范嵌入，于是定义
\[
H^n(C)\coloneqq\operatorname{coker}\bigl(j^n:B^n(C)\longrightarrow Z^n(C)\bigr).
\]
它称为次数 \(n\in\mathbb Z\) 的\term[shangtongdiao]{上同调对象}[cohomology object]。当 \(R\) 为含幺环、\(C\) 为幺性左 \(R\) 模复形时，\(Z^n(C)\) 的元素称\term[yuquan]{余圈}[cocycle]，\(B^n(C)\) 的元素称\term[yubian]{余边}[coboundary]，且 \(H^n(C)=Z^n(C)/B^n(C)\)。对 \(\mathcal A\) 中的链复形 \((C_n,\partial_n)\)，相应定义 \(Z_n(C)=\ker\partial_n\)、\(B_n(C)=\operatorname{im}\partial_{n+1}\)，余核 \(H_n(C)\coloneqq\operatorname{coker}\bigl(B_n(C)\longrightarrow Z_n(C)\bigr)\) 称为次数 \(n\) 的\term[tongdiao]{同调对象}[homology object]；对模链复形，\(Z_n(C)\) 与 \(B_n(C)\) 的元素分别称\term[quan]{圈}[cycle]与\term[bianjie]{边界}[boundary]。
\end{definition}
\symidx{\(Z_n(C),B_n(C),H_n(C)\)：链复形的圈、边界与同调对象}
\symidx{\(Z^n(C),B^n(C),H^n(C)\)：余圈、余边与上同调对象}

在模范畴中，两个余圈代表同一上同调类，恰好表示它们之差是余边：能够由前一步微分产生的改变被忽略了。在一般交换范畴中，相同的思想由上述余核表达，不能预先假定对象带有元素。

\begin{definition}\label{b4:def:acyclic-complex}
设 \(C\) 是交换范畴中的复形。若所有次数的同调对象均为零，则称 \(C\) 为\term[lingtiao-fuxing]{零调复形}[acyclic complex]，也称正合复形，即在每一个次数处都正合。对上链复形，这等价于 \(\operatorname{im}\mathrm{d}_C^{n-1}=\ker\mathrm{d}_C^n\) 对每个 \(n\in\mathbb Z\) 成立，其中等号指同一个子对象；链复形作相应的指标转换。
\end{definition}

\begin{example}\label{b4:exm:integer-three-complex}
取交换群复形
\[
C^0=\mathbb Z\xrightarrow{\ \mathrm{d}^0\ }C^1=\mathbb Z^2
\xrightarrow{\ \mathrm{d}^1\ }C^2=\mathbb Z,
\qquad \mathrm{d}^0(a)=(2a,0),\quad \mathrm{d}^1(x,y)=3y.
\]
复合为零，且
\[
\begin{aligned}
H^0(C)&=0,\\
H^1(C)&=\mathbb Z(1,0)/2\mathbb Z(1,0)\cong\mathbb Z/2\mathbb Z,\\
H^2(C)&=\mathbb Z/3\mathbb Z.
\end{aligned}
\]
中间的核由 \((1,0)\) 生成，前一步只得到其偶数倍；末项没有下一步限制，却须除去 \(3\mathbb Z\)。这两个商分别记录不同次数的非正合性。
\end{example}

\subsection{指标转换与次数平移约定}
\label{b4:subsec:0501:03}

把链复形改记为上链复形时，本书取
\[
C^n=C_{-n},\qquad \mathrm{d}_C^n=\partial_{-n},\qquad H^n(C)=H_{-n}(C).
\]
因为 \(C^{n+1}=C_{-(n+1)}=C_{-n-1}\)，原来的 \(\partial_{-n}:C_{-n}\to C_{-n-1}\) 正好成为 \(\mathrm{d}_C^n:C^n\to C^{n+1}\)。这只是重编号，没有额外的负号；负号出现在下面的移位中。

\begin{definition}\label{b4:def:complex-shift}
设 \(C\) 是交换范畴中的上链复形，\(r\in\mathbb Z\)。规定
\[
C[r]^n=C^{n+r},\qquad \mathrm{d}_{C[r]}^n=(-1)^r\cdot\mathrm{d}_C^{n+r},
\]
所得复形称为 \(C\) 的 \(r\) 次\term[yiwei]{移位}[shift]。对复形之间的零次态射，移位只改指标，不增加符号。
\end{definition}
\symidx{\(C[r]\)：复形的移位，\(C[r]^n=C^{n+r}\)}

无论是否乘上这个整体符号，相邻微分的复合都仍为零；采用 \((-1)^r\) 并不是复形条件本身所强制的。核与像不因微分乘 \(-1\) 改变，所以
\[
H^n(C[r])=H^{n+r}(C),\qquad (C[r])[s]=C[r+s].
\]
原来次数为 \(m\) 的对象，在 \(C[r]\) 中出现于次数 \(m-r\)，因为 \(C[r]^{m-r}=C^m\)。特别地，\(C[1]^{-1}=C^0\)、\(C[1]^0=C^1\)，所以 \([1]\) 把原来的各项向较小次数移动一格。例如 \(M[0][1]\) 的唯一非零项在次数 \(-1\)，不是次数 \(1\)。若沿用链记号，同一约定写成 \((C[r])_n=C_{n-r}\)、\(\partial_{C[r]}=(-1)^r\partial_C\)。阅读采用相反移位记号的文献时，须同时转换次数和微分。

\ProseParagraphBreak
移位的负号使复形能够相容地组成映射锥，并使同伦在移位后仍满足同一形式的公式。此处先固定约定，\secref{b4:sec:0602}将通过微分平方和连接同态的计算说明它的作用。

% !TeX root = ../../Book4.tex
% 纲要标题：### 6.2 复形的运算
% 纲要标签：b4:sec:0502
% 纲要位置：计划纲要.md，第 3224 行。

\section{复形的运算}
\label{b4:sec:0502}

复形之间的态射必须与微分相容，Hom 和张量积的微分还必须消去两个方向的混合项。逐次的核与余核可以沿用交换范畴的构造，分次符号却需要另作核验。

% 纲要内容（仅作写作计划，尚非教材正文）：
%
% * 复形映射、核、余核与短正合列
% * Hom 复形及张量复形的初步定义
% * 分次符号规则
%

\subsection{复形映射、核、余核与短正合列}
\label{b4:subsec:0502:01}

\begin{definition}\label{b4:def:complex-map}
设 \(C,D\) 是交换范畴 \(\mathcal A\) 中的上链复形。态射族 \(f^n:C^n\to D^n\) 若满足
\[
\mathrm{d}_D^nf^n=f^{n+1}\mathrm{d}_C^n\qquad(n\in\mathbb Z),
\]
则称为\term[fuxing-yingshe]{复形映射}[map of complexes]，记为 \(f:C\to D\)。以逐次复合作为复合，得到复形范畴 \(\operatorname{Ch}(\mathcal A)\)。
\end{definition}
\symidx{\(\operatorname{Ch}(\mathcal A)\)：上链复形及复形映射组成的范畴}

记余圈嵌入为 \(z_X^n:Z^n(X)\to X^n\)（\(X=C,D\)）。交换等式给出 \(\mathrm{d}_D^nf^nz_C^n=f^{n+1}\mathrm{d}_C^nz_C^n=0\)，所以核的泛性质唯一产生
\[
Z^n(f):Z^n(C)\longrightarrow Z^n(D),
\qquad z_D^nZ^n(f)=f^nz_C^n.
\]
另一方面，\(f^n\mathrm{d}_C^{n-1}=\mathrm{d}_D^{n-1}f^{n-1}\) 保证 \(f^n\) 将前一微分的像送入目标的余边对象，因而 \(Z^n(f)\) 将 \(B^n(C)\) 的嵌入送入 \(B^n(D)\)。记 \(q_X^n:Z^n(X)\to H^n(X)\) 为余核态射，则 \(q_D^nZ^n(f)\) 杀死 \(B^n(C)\) 的嵌入，沿 \(q_C^n\) 唯一下降为
\[
H^n(f):H^n(C)\longrightarrow H^n(D),
\qquad H^n(f)q_C^n=q_D^nZ^n(f).
\]
在模范畴中，这两次因子化正是先将余圈 \(z\) 送到 \(f^n(z)\)，再取类 \([z]\mapsto[f^n(z)]\)。恒等映射和复合的保持也由这两种唯一性得到，因此 \(H^n\) 是函子。

\begin{proposition}\label{b4:prop:complex-category-abelian}
若 \(\mathcal A\) 是交换范畴，则 \(\operatorname{Ch}(\mathcal A)\) 也是交换范畴。核、余核、像与有限双积均逐次数计算；特别地，复形序列 \(0\to A\to B\to C\to0\) 短正合，当且仅当每个次数的对象序列短正合。
\end{proposition}
\begin{proof}
给定复形 \(C,D\)，有
\[
\operatorname{Hom}_{\operatorname{Ch}(\mathcal A)}(C,D)
\subseteq\prod_{n\in\mathbb Z}\operatorname{Hom}_{\mathcal A}(C^n,D^n).
\]
由于 \(\mathcal A\) 局部小，右端是集合指标的集合直积，故左端也是集合。这证明了 \(\operatorname{Ch}(\mathcal A)\) 的局部小性。态射逐次相加，零复形和逐次双积给出加性结构。

对 \(f:C\to D\)，令 \(K^n=\ker f^n\)，核嵌入为 \(k^n:K^n\to C^n\)。微分 \(\mathrm{d}_C^n\) 接在这个嵌入之后，得到从 \(K^n\) 到 \(C^{n+1}\) 的态射，且
\[
f^{n+1}\mathrm{d}_C^nk^n=\mathrm{d}_D^nf^nk^n=0.
\]
因此这条态射唯一地通过 \(f^{n+1}\) 的核分解，产生 \(\mathrm{d}_K^n:K^n\to K^{n+1}\)，满足 \(k^{n+1}\mathrm{d}_K^n=\mathrm{d}_C^nk^n\)。再计算
\[
k^{n+2}\mathrm{d}_K^{n+1}\mathrm{d}_K^n
=\mathrm{d}_C^{n+1}\mathrm{d}_C^nk^n=0,
\]
由 \(k^{n+2}\) 单，得到 \(\mathrm{d}_K^{n+1}\mathrm{d}_K^n=0\)。若复形映射 \(h:T\to C\) 满足 \(fh=0\)，其逐次核分解在复合 \(k^{n+1}\) 后由 \(h\) 的微分相容性互相匹配；消去这个单态射，便知分解本身是复形映射。故所得 \(K\to C\) 满足复形范畴中的核泛性质。

对余核 \(q^n:D^n\to Q^n\)，有 \(q^{n+1}\mathrm{d}_D^nf^n=q^{n+1}f^{n+1}\mathrm{d}_C^n=0\)。所以 \(q^{n+1}\mathrm{d}_D^n:D^n\to Q^{n+1}\) 杀死 \(f^n\)，沿其余核 \(q^n\) 唯一下降，产生 \(\mathrm{d}_Q^n:Q^n\to Q^{n+1}\)，满足
\[
\mathrm{d}_Q^nq^n=q^{n+1}\mathrm{d}_D^n.
\]
于是 \(\mathrm{d}_Q^{n+1}\mathrm{d}_Q^nq^n=q^{n+2}\mathrm{d}_D^{n+1}\mathrm{d}_D^n=0\)，消去满态射 \(q^n\) 得到复形条件。任何杀死 \(f\) 的复形映射都逐次通过 \(q^n\) 唯一分解；所得分解的微分相容性也在复合 \(q^n\) 后成立，再由其满性得到。因此 \(D\to Q\) 是复形范畴中的余核。

像是余核的核，余像是核的余核，所以二者也逐次计算。记 \(P=\operatorname{coim}f\)、\(J=\operatorname{im}f\)，结构复形映射为 \(p:C\to P\)、\(i:J\to D\)。每个次数的典范比较 \(a^n:P^n\to J^n\) 满足 \(i^na^np^n=f^n\)，且因 \(\mathcal A\) 交换而可逆。由 \(p,i,f\) 的微分相容性，得到
\[
\begin{aligned}
i^{n+1}\mathrm{d}_J^na^np^n
&=\mathrm{d}_D^nf^n=f^{n+1}\mathrm{d}_C^n\\
&=i^{n+1}a^{n+1}\mathrm{d}_P^np^n.
\end{aligned}
\]
消去单态射 \(i^{n+1}\) 与满态射 \(p^n\)，可知 \(\mathrm{d}_J^na^n=a^{n+1}\mathrm{d}_P^n\)，故这些比较组成复形映射。等式两侧再复合相应的逆，得到
\[
\mathrm{d}_P^n(a^n)^{-1}=(a^{n+1})^{-1}\mathrm{d}_J^n.
\]
所以逐次的逆也组成复形映射，比较确为复形同构。最后的正合性判据来自逐次核和像的识别。
\end{proof}

后面讨论长正合列与导出函子，还需要同调保持 Hom 群的加法。对复形映射 \(f,g:C\to D\)，记它们在余圈对象上诱导的态射为 \(Z^n(f),Z^n(g)\)。核的唯一分解给出 \(Z^n(f+g)=Z^n(f)+Z^n(g)\)。记余核态射为 \(q_X^n:Z^n(X)\to H^n(X)\)（\(X=C,D\)），则
\[
\begin{aligned}
H^n(f+g)q_C^n
&=q_D^n\bigl(Z^n(f)+Z^n(g)\bigr)\\
&=\bigl(H^n(f)+H^n(g)\bigr)q_C^n.
\end{aligned}
\]
由于 \(q_C^n\) 满，得到 \(H^n(f+g)=H^n(f)+H^n(g)\)。所以每个 \(H^n:\operatorname{Ch}(\mathcal A)\to\mathcal A\) 都是加性函子。

\ProseParagraphBreak
逐次分裂的短正合列不必有复形截面。微分可能把各次数所选补对象送出补对象；\secref{b4:sec:0503}中的非零连接同态能够检测这类不分裂现象。不过，连接同态全部为零一般也不足以推出存在复形截面。

\subsection{Hom 复形与张量复形}
\label{b4:subsec:0502:02}

给两个复形逐项取 Hom 或张量积，并不能只把所有微分原样相加：两个方向的混合项必须抵消。以下分别固定产生这种抵消的符号。

\ProseParagraphBreak
次数增加 \(m\) 的分次同态需要同时给出每个次数上的分量 \(f^p:C^p\to D^{p+m}\)，不要求只有有限多个分量非零，因此这些同态组成直积。例如取每个 \(p\in\mathbb Z\) 处都有 \(C^p=\mathbb Z\)，并令微分为零，则恒等映射 \(1_C=(1_{C^p})_{p\in\mathbb Z}\) 有无穷多个非零分量。它属于 \(\prod_p\operatorname{Hom}_{\mathbb Z}(C^p,C^p)\)，却不属于相应直和。

\ProseParagraphBreak
Hom 微分应把次数 \(m\) 的分次同态送到次数 \(m+1\)，而在次数零，微分为零应当恰好表示 \(\mathrm{d}_Df=f\mathrm{d}_C\)。因此取目标微分后的复合与源微分前的复合之差；为了让两次微分中的混合项抵消，第二项的系数随次数增加而变号，得到下面的带符号公式。

\begin{definition}\label{b4:def:hom-complex}
设 \(R\) 是含幺环，\(C,D\) 是左 \(R\)-模的上链复形。定义交换群复形
\[
\operatorname{Hom}_R^m(C,D)\coloneqq
\prod_{p\in\mathbb Z}\operatorname{Hom}_R(C^p,D^{p+m}),
\qquad
(\mathrm{d}f)^p=\mathrm{d}_D^{p+m}f^p-(-1)^m f^{p+1}\mathrm{d}_C^p.
\]
它称为\term[Hom-fuxing]{Hom 复形}[Hom complex]。这里 \(m\in\mathbb Z\)，\(f=(f^p)_{p\in\mathbb Z}\in\operatorname{Hom}_R^m(C,D)\) 是一整族使次数增加 \(m\) 的同态，不要求其与微分交换。
\end{definition}
\symidx{\(\operatorname{Hom}_R^\bullet(C,D),\ \operatorname{Hom}_R^m(C,D)\)：Hom 复形及其第 \(m\) 次项}

\(\mathrm{d}f\) 的次数已是 \(m+1\)，所以再次施行微分时使用的符号是 \((-1)^{m+1}\)。暂略各分量的指标，展开为
\[
\begin{aligned}
\mathrm{d}(\mathrm{d}f)
={}&\mathrm{d}_D^2f-(-1)^m\cdot\mathrm{d}_Df\mathrm{d}_C\\
&-(-1)^{m+1}\cdot\mathrm{d}_Df\mathrm{d}_C
+(-1)^{2m+1}f\mathrm{d}_C^2=0.
\end{aligned}
\]
次数零的余圈正是复形映射。若 \(\mathcal A\) 为一般局部小交换范畴，也可把上述 Hom 换为其态射交换群，得到同样的交换群复形。

\begin{definition}\label{b4:def:tensor-complex}
设 \(R\) 是含幺环，\(C\) 是右 \(R\)-模的上链复形，\(D\) 是左 \(R\)-模的上链复形。定义交换群复形
\[
\begin{gathered}
(C\otimes_R D)^n\coloneqq\bigoplus_{p+q=n}C^p\otimes_R D^q,\\
\mathrm{d}(x\otimes y)=\mathrm{d}_Cx\otimes y+(-1)^p x\otimes \mathrm{d}_Dy,\\
x\in C^p,\quad y\in D^q,\quad p,q\in\mathbb Z.
\end{gathered}
\]
它称为\term[zhangliang-fuxing]{张量复形}[tensor product of complexes]。若 \(R\) 交换且两者均为 \(R\)-模复形，结果也带自然的 \(R\)-模结构。
\end{definition}
\symidx{\(C\otimes_RD\)：复形的张量积}

每个张量项由整数对 \((p,q)\) 标记，取 \(p+q=n\) 的全部项便是总次数 \(n\)。\(\mathrm{d}_C\) 将 \(p\) 增一，\(\mathrm{d}_D\) 将 \(q\) 增一，所以二者都将总次数增一。虽然同一个总次数中可能有无限多个非零项，直和的定义只允许每个元素使用其中有限多个；微分后仍只有有限多个分量。公式与张量积的平衡关系相容，因为微分均为 \(R\)-线性。

\ProseParagraphBreak
混合项有两种产生次序：先作用 \(\mathrm{d}_C\)，则 \(x\) 的次数已经变为 \(p+1\)，接着作用 \(\mathrm{d}_D\) 时带系数 \((-1)^{p+1}\)；反过来先作用 \(\mathrm{d}_D\)，带系数 \((-1)^p\)，再作用 \(\mathrm{d}_C\)。因此它们相消，完整计算为
\[
\mathrm{d}^2(x\otimes y)=\mathrm{d}_C^2x\otimes y
+\bigl((-1)^{p+1}+(-1)^p\bigr)\cdot\mathrm{d}_Cx\otimes \mathrm{d}_Dy
+x\otimes \mathrm{d}_D^2y=0.
\]
这证明了复形条件。张量与同调能否交换是另一问题：后面的\cref{b4:thm:kunneth-pid}将表明，在主理想整环上即使复形各项自由，同调的张量积仍可能漏掉 Tor 项。

\subsection{分次符号规则}
\label{b4:subsec:0502:03}

在本书所用的分次张量结构中，一个次数为 \(r\) 的齐次因子跨过次数为 \(s\) 的齐次因子时，带系数 \((-1)^{rs}\)；相应齐次映射的张量运算也遵守这一约定。这称为\term[fenci-fuhao-guize]{分次符号规则}[graded sign rule]。它适用于具体的交换或张量操作，不能只因公式中出现两个次数就任意加号；是否与微分相容，仍须从所定义的映射核验。

\begin{proposition}\label{b4:prop:graded-leibniz}
设 \(R\) 为含幺环，\(C,D,E\) 是幺性左 \(R\) 模的上链复形，\(r,s\in\mathbb Z\)，\(f\in\operatorname{Hom}_R^r(C,D)\)、\(g\in\operatorname{Hom}_R^s(D,E)\)。以 \((gf)^p=g^{p+r}f^p\)（\(p\in\mathbb Z\)）定义复合，则
\[
\mathrm{d}(gf)=(\mathrm{d}g)f+(-1)^s g(\mathrm{d}f).
\]
这里 \(gf=g\circ f\)，即先作用 \(f\)，再作用 \(g\)；因子在式中仍依次写作 \(g,f\)。在通常的分次乘法公式 \(\mathrm{d}(ab)=(\mathrm{d}a)b+(-1)^{|a|}a\cdot\mathrm{d}b\) 中取 \(a=g\)、\(b=f\)，左因子的次数便是 \(s\)，所以这里出现 \((-1)^s\)。
若 \(R\) 交换且 \(C,D\) 为 \(R\)-模复形，则
\[
\tau:C\otimes_R D\longrightarrow D\otimes_R C,
\qquad \tau(x\otimes y)=(-1)^{pq}y\otimes x
\quad(x\in C^p,\ y\in D^q)
\]
是复形同构，且 \(\tau_{D,C}\tau_{C,D}=1\)。
\end{proposition}
\begin{proof}
第一式右端展开为
\[
\begin{aligned}
(\mathrm{d}g)f+(-1)^s g(\mathrm{d}f)
={}&(\mathrm{d}_Eg-(-1)^sg\mathrm{d}_D)f\\
&+(-1)^sg(\mathrm{d}_Df-(-1)^rf\mathrm{d}_C)\\
={}&\mathrm{d}_Egf-(-1)^{r+s}gf\mathrm{d}_C.
\end{aligned}
\]
这正是左端。第二式的两种复合在纯张量上分别为
\[
\begin{aligned}
\tau\mathrm{d}(x\otimes y)
={}&(-1)^{(p+1)q}y\otimes\mathrm{d}_Cx\\
&+(-1)^{p+p(q+1)}\cdot\mathrm{d}_Dy\otimes x,\\
\mathrm{d}\tau(x\otimes y)
={}&(-1)^{pq}\cdot\mathrm{d}_Dy\otimes x\\
&+(-1)^{pq+q}y\otimes\mathrm{d}_Cx.
\end{aligned}
\]
\(y\otimes\mathrm{d}_Cx\) 的两份系数相等，因为 \((p+1)q=pq+q\)；\(\mathrm{d}_Dy\otimes x\) 的两份系数相等，因为 \(p+p(q+1)=pq+2p\)。因此 \(\tau\) 与微分相容。再交换一次乘上 \((-1)^{2pq}=1\)，所以映射可逆。
\end{proof}

例如两个只在次数一非零的 \(k\)-向量空间复形交换时，交换映射带负号。若删掉该号，在存在非零微分的复形上就可能不再与微分交换。特征二时正负号数值相同，但定义仍沿用同一公式，以便不同特征的论证可以比较。

% !TeX root = ../../Book4.tex
% 纲要标题：### 6.3 上同调长正合列
% 纲要标签：b4:sec:0503
% 纲要位置：计划纲要.md，第 3230 行。

\section{上同调长正合列}
\label{b4:sec:0503}

在模复形的短正合列中，右端余圈可以提升到中项，但其原像未必仍是余圈。对这个提升作微分，便得到连接同态及其带来的上同调长正合列。

% 纲要内容（仅作写作计划，尚非教材正文）：
%
% * 连接同态的构造
% * 长正合列及其自然性
% * 显式计算而非仅引用存在性
%

\subsection{连接同态的构造}
\label{b4:subsec:0503:01}

先在模范畴中考察复形的短正合列
\[
0\longrightarrow A\xrightarrow{i}B\xrightarrow{p}C\longrightarrow0.
\]
若 \(c\in C^n\) 是余圈，将它提升到 \(b\in B^n\) 后，\(b\) 未必仍是余圈。然而 \(p(\mathrm{d}b)=\mathrm{d}c=0\)，故 \(\mathrm{d}b\) 来自 \(A^{n+1}\)。连接同态记录的正是这个提升失败的程度。
\[
\begin{tikzcd}[column sep=large,row sep=large]
A^n\arrow[r,hook,"i^n"]\arrow[d,"\mathrm{d}_A^n"']&B^n\arrow[r,two heads,"p^n"]\arrow[d,"\mathrm{d}_B^n"]&C^n\arrow[d,"\mathrm{d}_C^n"]\\
A^{n+1}\arrow[r,hook,"i^{n+1}"']&B^{n+1}\arrow[r,two heads,"p^{n+1}"']&C^{n+1}
\end{tikzcd}
\]

\begin{proposition}\label{b4:prop:connecting-map-construction}
设 \(\mathcal A\) 是交换范畴，\(0\to A\xrightarrow{i}B\xrightarrow{p}C\to0\) 是上链复形的短正合列。每个整数 \(n\) 都有典范态射
\[
\delta^n:H^n(C)\longrightarrow H^{n+1}(A).
\]
在模范畴中，它由下列规则给出：取 \(c\in Z^n(C)\)，选 \(b\in B^n\) 使 \(p(b)=c\)，唯一写成 \(\mathrm{d}_Bb=i(a)\)，则 \(a\in Z^{n+1}(A)\)，并规定 \(\delta^n[c]=[a]\)。
选定原像 \(b\) 后，元素的去向为：
\[
\begin{tikzcd}[column sep=large,row sep=large]
&b\arrow[r,mapsto,"p^n"]\arrow[d,mapsto,"\mathrm{d}_B^n"']&c\arrow[d,mapsto,"\mathrm{d}_C^n"]\\
a\arrow[r,mapsto,"i^{n+1}"']&\mathrm{d}_Bb\arrow[r,mapsto,"p^{n+1}"']&0
\end{tikzcd}
\]
\end{proposition}
\begin{proof}
固定整数 \(n\)。在此证明中将中间复形暂记为 \(E\)，把短正合列写成 \(0\to A\xrightarrow{i}E\xrightarrow{p}C\to0\)。要运用蛇引理构造连接态射，需要一张竖直箭头的核为 \(H^n\)、余核为 \(H^{n+1}\) 的图。为此，对每个上链复形 \(X\) 定义
\[
Q^n(X)\coloneqq\operatorname{coker}\bigl(B^n(X)\hookrightarrow X^n\bigr),
\qquad q_X:X^n\twoheadrightarrow Q^n(X).
\]
在模范畴中，\(Q^n(X)\) 就是 \(X^n/B^n(X)\)。记余圈嵌入为 \(z_X^r:Z^r(X)\hookrightarrow X^r\)。微分的像包含在 \(Z^{n+1}(X)\) 中，并在 \(B^n(X)\) 上为零，因而唯一诱导
\[
\bar{\mathrm{d}}_X:Q^n(X)\longrightarrow Z^{n+1}(X),
\qquad
\mathrm{d}_X^n=z_X^{n+1}\bar{\mathrm{d}}_Xq_X.
\]

先识别这条态射的核。令 \(k_X:K_X\hookrightarrow Q^n(X)\) 为 \(\bar{\mathrm{d}}_X\) 的核。因为 \(\bar{\mathrm{d}}_Xq_Xz_X^n=0\)，存在唯一 \(u_X:Z^n(X)\to K_X\) 使 \(k_Xu_X=q_Xz_X^n\)。由 \(z_X^{n+1}\) 单，对任意 \(x:T\to X^n\)，条件 \(\bar{\mathrm{d}}_Xq_Xx=0\) 等价于 \(\mathrm{d}_X^nx=0\)，也就是 \(x\) 唯一通过 \(z_X^n\) 分解。若还给定 \(y:T\to K_X\) 满足 \(q_Xx=k_Xy\)，则该分解通过 \(u_X\) 所得的态射与 \(y\) 相等，因为 \(k_X\) 单。这说明 \(Z^n(X)\) 是 \(q_X\) 沿 \(k_X\) 的拉回，故由\cref{b4:lem:epi-pullback}，\(u_X\) 满。又因 \(k_X\) 单，\(u_Xv=0\) 等价于 \(q_Xz_X^nv=0\)，后者恰表示 \(z_X^nv\) 通过 \(B^n(X)\hookrightarrow X^n\) 分解。消去 \(z_X^n\) 可知，\(u_X\) 的核就是 \(B^n(X)\hookrightarrow Z^n(X)\)。由\cref{b4:prop:abelian-normal-factorization}，满态射 \(u_X\) 是其核的余核，故由上同调的定义得到典范识别
\[
\ker\bar{\mathrm{d}}_X\cong H^n(X).
\]
另一方面，\(q_X\) 满，故 \(\bar{\mathrm{d}}_X\) 的像与微分在 \(Z^{n+1}(X)\) 中的像相同，即 \(B^{n+1}(X)\)。因此
\[
\operatorname{coker}\bar{\mathrm{d}}_X
\cong\operatorname{coker}\bigl(B^{n+1}(X)\hookrightarrow Z^{n+1}(X)\bigr)
=H^{n+1}(X).
\]

短正合列中的复形映射保持余圈与余边，因而诱导交换图
\[
\begin{tikzcd}[column sep=large]
Q^n(A)\arrow[r,"Q^n(i)"]\arrow[d,"\bar{\mathrm{d}}_A"']&Q^n(E)\arrow[r,"Q^n(p)"]\arrow[d,"\bar{\mathrm{d}}_E"]&Q^n(C)\arrow[r]\arrow[d,"\bar{\mathrm{d}}_C"]&0\\
Z^{n+1}(A)\arrow[r,"Z^{n+1}(i)"']&Z^{n+1}(E)\arrow[r,"Z^{n+1}(p)"']&Z^{n+1}(C)&
\end{tikzcd}
\]
下排左端单，因为 \(i^{n+1}\) 与余圈嵌入均单。为验证中项正合，设 \(e:T\to Z^{n+1}(E)\) 满足 \(Z^{n+1}(p)e=0\)。由 \(i^{n+1}=\ker p^{n+1}\)，存在唯一 \(a:T\to A^{n+1}\) 使 \(i^{n+1}a=z_E^{n+1}e\)。复形映射的相容性给出
\[
i^{n+2}\mathrm{d}_A^{n+1}a
=\mathrm{d}_E^{n+1}z_E^{n+1}e=0.
\]
由 \(i^{n+2}\) 单，\(a\) 是余圈，故它唯一通过 \(Z^{n+1}(A)\) 分解，得到 \(e\) 在下排左端的原像。反向由 \(pi=0\) 给出。

上排右端满：复合 \(q_Cp^n:E^n\to Q^n(C)\) 满，且它在 \(B^n(E)\) 上为零，故满足
\[
q_Cp^n=Q^n(p)q_E.
\]
原复合满，推出 \(Q^n(p)\) 满。

上排中项的正合性用\cref{b4:lem:generalized-element-calculus}检验。以下每次提升都继续预合成所得满态射，并保留原符号；同一等式中的广义元素具有共同定义域。设 \(\xi:T\to Q^n(E)\) 满足 \(Q^n(p)\xi=0\)。由 \(q_E\) 满，在满态射后取 \(e:T'\to E^n\)，使 \(q_Ee=\xi\)。此时 \(q_Cp^ne=0\)，而 \(\ker q_C=B^n(C)=\operatorname{im}\mathrm{d}_C^{n-1}\)，故经进一步的满态射预合成后可取 \(c':T''\to C^{n-1}\)，使 \(p^ne=\mathrm{d}_C^{n-1}c'\)。再由 \(p^{n-1}\) 满，继续预合成并取 \(e'\) 使 \(p^{n-1}e'=c'\)。于是
\[
p^n\bigl(e-\mathrm{d}_E^{n-1}e'\bigr)=0.
\]
由于 \(i^n=\ker p^n\)，存在唯一 \(a\) 使 \(e-\mathrm{d}_E^{n-1}e'=i^na\)。\(q_E\) 消去余边，所以
\[
\xi=q_Ee=q_Ei^na=Q^n(i)q_Aa.
\]
这给出了 \(\xi\) 在共同满态射后的原像；反向仍由 \(pi=0\) 给出，故上排在中项正合。两排满足蛇引理的假设，已识别的竖直核与余核因此给出所需 \(\delta^n\)。

在模中，蛇引理的提升规则正是陈述中的公式。由 \(i(\mathrm{d}a)=\mathrm{d}_E^2e=0\) 和 \(i\) 单得到 \(\mathrm{d}a=0\)。改变提升 \(e\) 为 \(e+i(a_0)\)，\(a\) 改变为 \(a+\mathrm{d}a_0\)，故其上同调类不变。改变余圈代表 \(c\) 为 \(c+\mathrm{d}c_0\) 时，提升 \(c_0\) 为 \(e_0\)，改取 \(e+\mathrm{d}e_0\) 即可，二次微分为零使 \(a\) 不变。加法也与各步相容，得到所声称的同态公式。
\end{proof}

在一般交换范畴中，最后一段仍可作为广义元素的计算方法，只须在需要提升时先预合成满态射。蛇引理已经保证连接态射存在且与选择无关，这种计算说明了它如何作用于代表元。

\subsection{长正合列及其自然性}
\label{b4:subsec:0503:02}

\begin{theorem}[上同调长正合列]\label{b4:thm:cohomology-long-exact}
设 \(\mathcal A\) 是交换范畴，\(0\to A\xrightarrow{i}B\xrightarrow{p}C\to0\) 是 \(\mathcal A\) 中上链复形的短正合列。以 \(\delta^n:H^n(C)\to H^{n+1}(A)\)（\(n\in\mathbb Z\)）表示其典范连接态射。它们组成\term[shangtongdiao-chang-zhenghelie]{上同调长正合列}[long exact cohomology sequence]
\[
\cdots\longrightarrow H^n(A)\xrightarrow{H^n(i)}H^n(B)
\xrightarrow{H^n(p)}H^n(C)\xrightarrow{\delta^n}H^{n+1}(A)
\longrightarrow\cdots.
\]
给定上述复形短正合列到另一条 \(\mathcal A\) 中的复形短正合列 \(0\to A'\to B'\to C'\to0\) 的态射 \((a,b,c):(A,B,C)\to(A',B',C')\)，记后者的连接态射为 \(\delta'^n\)。诱导的上同调态射与长正合列的所有箭头交换，尤其
\[
H^{n+1}(a)\,\delta^n=\delta'^n H^n(c).
\]
把相邻的四项一起排列，得到自然性的交换图：
\[
\begin{tikzcd}[column sep=large,row sep=large]
H^n(A)\arrow[r,"H^n(i)"]\arrow[d,"H^n(a)"']&H^n(B)\arrow[r,"H^n(p)"]\arrow[d,"H^n(b)"]&H^n(C)\arrow[r,"\delta^n"]\arrow[d,"H^n(c)"]&H^{n+1}(A)\arrow[d,"H^{n+1}(a)"]\\
H^n(A')\arrow[r]&H^n(B')\arrow[r]&H^n(C')\arrow[r,"\delta'^n"']&H^{n+1}(A')
\end{tikzcd}
\]
\end{theorem}
\begin{proof}
对\cref{b4:prop:connecting-map-construction}中的图应用蛇引理，得到
\[
H^n(A)\to H^n(B)\to H^n(C)\to
H^{n+1}(A)\to H^{n+1}(B)\to H^{n+1}(C)
\]
在四个内项处正合。这六项中的箭头与原复形所诱导的箭头相同，因为核、像和余核上的态射都是由原来的 \(i,p\) 唯一诱导的。让 \(n\) 遍历整数，相邻六项列的重叠部分相同，因而拼成所述长列，每项都落在某个六项列的内项位置。

短正合列的态射诱导 \(Q^n\) 与 \(Z^{n+1}\) 两排之间的交换图。\cref{b4:prop:snake-naturality}立即给出连接箭头的交换；其余箭头交换由上同调的函子性得到。这也说明连接同态与次数的排列没有额外选择。
\end{proof}

链指标下，同一定理得到\term[tongdiao-chang-zhenghelie]{同调长正合列}[long exact homology sequence]，连接同态为 \(H_n(C)\to H_{n-1}(A)\)。升次与降次的区别只来自 \(H^n=H_{-n}\)，不表示有两种不同的构造。

\begin{corollary}\label{b4:cor:acyclic-two-of-three}
在交换范畴中，复形短正合列 \(0\to A\to B\to C\to0\) 的三个复形若有两个零调，则第三个也零调。若 \(A\) 零调，则 \(H^n(B)\to H^n(C)\) 对每个 \(n\) 均为同构；若 \(C\) 零调，则 \(H^n(A)\to H^n(B)\) 均为同构。
\end{corollary}
\begin{proof}
把相应零项代入长正合列即可。例如 \(H^n(A)=H^{n+1}(A)=0\) 时，相关三项成为 \(0\to H^n(B)\to H^n(C)\to0\)，所以中间箭头既单又满；其他情形使用同一列在相应次数的正合性。
\end{proof}

\subsection{连接同态的显式计算}
\label{b4:subsec:0503:03}

\begin{example}\label{b4:exm:connecting-nonzero}
设 \(k\) 是域，\(B\) 是次数 \(0,1\) 的复形 \(k\xrightarrow{1}k\)，令 \(A\) 只在次数 \(1\) 等于 \(k\)，以恒等映射嵌入 \(B^1\)，再令 \(C=B/A\)。于是 \(C\) 只在次数零等于 \(k\)，每个次数的短正合列都分裂。可是
\[
H^0(C)=k\xrightarrow{\delta^0}H^1(A)=k
\]
是恒等映射：把 \(c\) 提升为 \(B^0\) 中同一个数，其微分仍为 \(c\)。如果复形短正合列有复形截面，便能把余圈提升为余圈，从而 \(\delta^0=0\)，矛盾。逐次分裂因此弱于复形分裂。
\end{example}

\begin{example}\label{b4:exm:connecting-divisibility}
固定整数 \(m\ge2\)，令 \(D\) 为交换群复形 \(\mathbb Z\xrightarrow{\times m}\mathbb Z\)，非零次数为 \(0,1\)。逐次乘 \(m\) 给出短正合列
\[
0\longrightarrow D\xrightarrow{\times m}D\longrightarrow D/mD\longrightarrow0.
\]
模 \(m\) 后微分为零，故 \(H^0(D/mD)=H^1(D/mD)=\mathbb Z/m\)；原复形则有 \(H^0(D)=0\)、\(H^1(D)=\mathbb Z/m\)。将 \(\bar a\in H^0(D/mD)\) 提升为整数 \(a\)，其微分为 \(ma\)，通过左端嵌入乘 \(m\) 取原像得到 \(a\)。因此 \(\delta^0(\bar a)=\bar a\)，长列的非零部分为
\[
0\longrightarrow\mathbb Z/m\xrightarrow{1}\mathbb Z/m
\xrightarrow{0}\mathbb Z/m\xrightarrow{1}\mathbb Z/m\longrightarrow0.
\]
中间的零映射来自乘 \(m\)，并不破坏正合性，因为它的前一个箭头满、后一个箭头单。
\end{example}

更一般地，只要逐次乘 \(m\) 在模复形 \(D\) 上单射，同一短正合列就成立。连接同态要求先取 \(\mathrm{d}b\)，再用已知的整除关系 \(\mathrm{d}b=ma\) 求 \(a\)，最后取其上同调类；直接把 \(\mathrm{d}b\) 模 \(m\) 约化只会得到零，丢掉了需要记录的信息。

% !TeX root = ../../Book4.tex
% 纲要标题：### 6.4 Euler–Poincaré 等式与 Grothendieck 群
% 纲要标签：b4:sec:0504
% 纲要位置：计划纲要.md，第 3236 行。

\section{Euler–Poincaré 等式与 Grothendieck 群}
\label{b4:sec:0504}

边界对象在相邻次数中以相反符号出现，使复形各项的交错和等于同调的交错和。Grothendieck 群把这一计算推广到没有数值维数的交换范畴。

% 纲要内容（仅作写作计划，尚非教材正文）：
%
% * 有限长度或有限维情形的交错和
% * Grothendieck 群 K_0 中的可加不变量
% * 有界复形与其同调的 Euler 特征相等
%
% ---
%

\subsection{有限长度与有限维情形的交错和}
\label{b4:subsec:0504:01}

求出所有同调群有时很费力，但一个交错和却能在求核取商之前计算。原因是每个边界对象在两个相邻次数中各出现一次，且符号相反。

\begin{definition}\label{b4:def:euler-characteristic}
设 \(k\) 是域，\(C\) 是由有限维 \(k\)-向量空间组成的有界复形。规定
\[
\chi(C)\coloneqq\sum_n(-1)^n\dim_k C^n,
\]
称其为\term[Euler-tezheng]{Euler 特征}[Euler characteristic]。设 \(R\) 为含幺环；若 \(C\) 是由有限长度幺性左 \(R\) 模组成的有界上链复形，以 \(\ell_R(C^n)\) 表示第 \(n\) 项的合成长度，则把维数替换为此长度所得的有限交错和称为长度的 Euler 特征。
\end{definition}
\symidx{\(\chi(C)\)：有界有限维复形的 Euler 特征}
\symidx{\(\ell_R(M)\)：模的合成长度}

有界性保证只有有限多个非零项；每项有限维或有限长度，则保证每个被相加的数值都是有限整数。这是两个不同的要求。这里有限长度指具有有限合成列，其长度是合成因子的个数；长度的良定义与短正合列可加性已在第二卷命题12.2.5“长度的可加性”证明。有限维情形只需秩—零化度公式。两种情形都有
\[
0\to U\to V\to W\to0
\quad\Longrightarrow\quad
\lambda(V)=\lambda(U)+\lambda(W),
\]
其中 \(\lambda\) 分别表示维数或长度。因而对有界复形的短正合列，逐次相加即可得到 \(\chi(B)=\chi(A)+\chi(C)\)。

\begin{example}\label{b4:exm:graph-euler}
取一个有限图，顶点数为 \(v\)，边数为 \(e\)，连通分支数为 \(c\)。给每条边定向，在域 \(k\) 上定义链复形 \(C_1=k^e\xrightarrow{\partial}C_0=k^v\)，其余项为零。若用上链指标 \(C^{-n}=C_n\)，则 \((-1)^{-n}=(-1)^n\)，所以其 Euler 特征仍为 \(v-e\)。以 \(\varepsilon_w\) 表示顶点 \(w\) 的基向量，把一条边送到其终点与起点的基向量之差。每条边界在所属连通分支中的坐标和为零，故像包含在各分支的零和子空间中。反过来，在一个分支中选定基准顶点 \(w_0\)，从 \(w_0\) 到任意顶点 \(w\) 的路径按行进方向给各边加上正负号，其边界就是 \(\varepsilon_w-\varepsilon_{w_0}\)。这些差生成该分支的全部零和向量：若 \(\sum_w a_w=0\)，则
\[
\sum_w a_w\varepsilon_w
=\sum_{w\ne w_0}a_w(\varepsilon_w-\varepsilon_{w_0}).
\]
因此两个子空间相等，每个分支贡献的秩是其顶点数减一，总计 \(\operatorname{rank}\partial=v-c\)，从而
\[
\dim_kH_0=c,\qquad \dim_kH_1=e-v+c,
\qquad v-e=\dim_kH_0-\dim_kH_1.
\]
例如三角形的 \(H_1\) 为一维，而加上一条没有构成新回路的悬挂边不改变这个维数。
\end{example}

\subsection{\texorpdfstring{Grothendieck 群 \(K_0\) 中的可加不变量}{Grothendieck 群 K0 中的可加不变量}}
\label{b4:subsec:0504:02}

维数把所有一维向量空间看成同一种贡献；一般交换范畴中可能有几种不能互相替代的简单对象。为了保留这种差别，可以只规定短正合列应满足的加法关系。Euler 交错和还需要取对象类的差，因此不能只在交换幺半群中相加；下面从自由交换群出发，允许有限整数线性组合，包括负系数。

\begin{definition}\label{b4:def:grothendieck-group}
设 \(\mathcal A\) 是本质小的交换范畴，即其对象同构类组成一个集合。取以各同构类 \([M]\) 为基的自由交换群，再除去全部关系
\[
[B]-[A]-[C]
\qquad\bigl(0\to A\to B\to C\to0\;\text{短正合}\bigr)
\]
所生成的子群，所得群记为 \(K_0(\mathcal A)\)，称为 \(\mathcal A\) 的\term[Grothendieck-qun]{Grothendieck 群}[Grothendieck group]。
\end{definition}
\symidx{\(K_0(\mathcal A)\)：交换范畴的 Grothendieck 群}

本质小的条件保证上述自由群以集合为基。对于全部模组成的大范畴，需要先限制大小或选取所关心的本质小子范畴；不能不加说明地对真类取自由群。本节的交换范畴版本对全部短正合列施加关系。环的 \(K_0(R)\) 则以有限生成投射模的同构类为生成元；它们之间态射在模范畴中的核、余核未必仍为有限生成投射模，因而此范畴通常不是交换范畴。三项都是有限生成投射模的短正合列必因商模投射而分裂，所以这里使用的关系是 \([A\oplus C]=[A]+[C]\)。对象范围和所用的正合列均须区分，不能仅凭同一个符号将两种群等同。\Cref{b4:def:B-grothendieck-group}会以指定正合结构统一这两种定义；\cref{b4:thm:B-euler-isomorphism}进一步证明投射模与完美复形的 Euler 类给出相同的 Grothendieck 群。

\begin{proposition}\label{b4:prop:grothendieck-universal}
设 \(\mathcal A\) 是本质小交换范畴，\(G\) 是交换群。对 \(\mathcal A\) 的对象同构类赋值 \(\lambda(M)\in G\)，若每条短正合列 \(0\to A\to B\to C\to0\) 均满足 \(\lambda(B)=\lambda(A)+\lambda(C)\)，则存在唯一群同态 \(\bar\lambda:K_0(\mathcal A)\to G\)，使 \(\bar\lambda([M])=\lambda(M)\)。若 \(\mathcal B\) 也是本质小交换范畴，则每个正合函子 \(F:\mathcal A\to\mathcal B\) 诱导群同态 \(K_0(F):K_0(\mathcal A)\to K_0(\mathcal B)\)，\([M]\mapsto[F(M)]\)。
\end{proposition}
\begin{proof}
由自由交换群的泛性质，生成元上的赋值唯一延拓为群同态。可加性恰好说定义 \(K_0\) 的每个关系都映到零，故唯一地通过商群分解。第二项中正合函子把定义关系送到定义关系，因此同样可以下降。关系 \([0]=[0]+[0]\) 还说明 \([0]=0\)，不必另外添加零对象关系。
\end{proof}

\begin{example}\label{b4:exm:kzero-vectors}
设 \(k\) 为域。有限维 \(k\)-向量空间范畴满足 \([V]=(\dim_kV)[k]\)，因为基给出 \(V\cong k^{\dim V}\)，且直和对应分裂短正合列。维数诱导 \(K_0\to\mathbb Z\)，而 \(1\mapsto[k]\) 是其逆，故 \(K_0\cong\mathbb Z\)。若改为有限维 \(k\times k\)-模，令 \(e_1=(1,0)\)、\(e_2=(0,1)\)，则 \(M=e_1M\oplus e_2M\)。记两种简单模为 \(S_1=k\times0\)、\(S_2=0\times k\)，并令 \(a=\dim_k e_1M\)、\(b=\dim_k e_2M\)。在各分量中选基给出
\[
e_1M\cong S_1^{\oplus a},\qquad
e_2M\cong S_2^{\oplus b},\qquad
[M]=a[S_1]+b[S_2].
\]
模同态保持这两个分量，短正合列逐分量正合，所以两坐标维数诱导 \(K_0\to\mathbb Z^2\)。它与群同态 \((a,b)\mapsto a[S_1]+b[S_2]\) 互为逆：两种简单模分别映到 \((1,0)\)、\((0,1)\)，而上式把任意 \([M]\) 恢复出来。因此 \(K_0\cong\mathbb Z^2\)。总维数只是这两个坐标之和，不能区别两种简单模。
\end{example}

\subsection{有界复形与同调的 Euler 特征}
\label{b4:subsec:0504:03}

\begin{theorem}[Euler–Poincaré 等式]\label{b4:thm:euler-poincare}
设 \(\mathcal A\) 是本质小交换范畴，\(C\) 是其中的有界复形。则在 \(K_0(\mathcal A)\) 中有
\[
\chi_{K_0}(C)\coloneqq\sum_n(-1)^n[C^n]
=\sum_n(-1)^n[H^n(C)].
\]
这称为\term[Euler-Poincare-dengshi]{Euler–Poincaré 等式}[Euler–Poincaré formula]。因此任何取值于交换群的短正合列可加不变量，都使复形与其同调的交错和相等。
\end{theorem}
\symidx{\(\chi_{K_0}(C)\)：取值于 Grothendieck 群的 Euler 类}
\begin{proof}
将微分分解为到像的满态射再接像的包含态射：
\[
\mathrm{d}^n:C^n\xrightarrow{e^n}B^{n+1}(C)
\xrightarrow{i^{n+1}}C^{n+1},
\qquad \mathrm{d}^n=i^{n+1}e^n.
\]
因为 \(i^{n+1}\) 为单态射，\(\ker e^n=\ker\mathrm{d}^n=Z^n(C)\)。因此对每个 \(n\)，有两条短正合列
\[
0\to Z^n(C)\to C^n\xrightarrow{e^n}B^{n+1}(C)\to0,
\qquad
0\to B^n(C)\to Z^n(C)\to H^n(C)\to0.
\]
第二条来自同调的定义。由这两条列得
\[
[C^n]=[B^n(C)]+[H^n(C)]+[B^{n+1}(C)].
\]
取整数 \(a\le b\)，使 \(C^n=0\) 对 \(n<a\) 或 \(n>b\) 成立。于是 \(B^a(C)=\operatorname{im}\mathrm{d}^{a-1}=0\)，且 \(B^{b+1}(C)=\operatorname{im}\mathrm{d}^b=0\)。将上式乘 \((-1)^n\) 并从 \(a\) 到 \(b\) 求和，其中一组边界项重新编号为
\[
\begin{aligned}
\sum_{n=a}^{b}(-1)^n[B^{n+1}(C)]
&=-\sum_{j=a+1}^{b+1}(-1)^j[B^j(C)]\\
&=-\sum_{j=a}^{b}(-1)^j[B^j(C)].
\end{aligned}
\]
最后一个等号用了上述两个端点为零。这组有限和恰与另一组边界项相消，剩下的正是所需公式。最后对该等式应用\cref{b4:prop:grothendieck-universal}的群同态即可。
\end{proof}

\begin{corollary}\label{b4:cor:euler-quasi-invariance}
设 \(k\) 为域、\(R\) 为含幺环。对有限维 \(k\) 向量空间的有界上链复形，或有限长度幺性左 \(R\) 模的有界上链复形，复形各项的维数或长度交错和等于其上同调的相应交错和。若同一类中的复形映射 \(f:C\to D\) 在每个次数诱导上同调同构，则 \(C,D\) 的相应 Euler 特征相等；此类有界零调复形的相应 Euler 特征为零。
\end{corollary}
\begin{proof}
边界、余圈及上同调仍为有限维向量空间或有限长度模，所以刚才的两条短正合列中，各项的维数或长度均有定义。以 \(\lambda\) 表示这一数值，由可加性得
\[
\lambda(C^n)=\lambda(B^n(C))+\lambda(H^n(C))+\lambda(B^{n+1}(C)).
\]
按前一证明的有限重新编号消去边界项，便得到所需等式。上同调同构使等式右端逐项相同；零调时每项为零。
\end{proof}

有限性不能只当作便于书写的条件。例如令 \(C^n=k\)（\(n\ge0\)）、\(C^n=0\)（\(n<0\)），并取 \(\mathrm{d}^{2r}=1_k\)、\(\mathrm{d}^{2r+1}=0\)（\(r\ge0\)），即
\[
0\to k\xrightarrow{1}k\xrightarrow{0}k\xrightarrow{1}k\xrightarrow{0}\cdots,
\]
其中第一个 \(k\) 位于零次。这是零调复形，但 \(1-1+1-1+\cdots\) 不是这里定义的有限交错和。\(K_0(\mathcal A)\) 的群运算只给出有限求和，并未定义无穷级数运算，所以不能将相邻项成对抵消后宣称这个复形的 Euler 类为零。若要给无界复形赋予更广的 Euler 特征，必须另给允许无限求和的结构和条件，或使所求交错和仍只有有限个非零项；本定理并不提供这样的延拓。


\begin{chaptersummary}
同调对象是边界对象嵌入圈对象的余核；对上链复形，所用的是余边对象与余圈对象。在模范畴中，这就是相邻微分的“核除以像”。复形映射在这些对象之间诱导态射，从而使同调具有函子性。复形范畴的正合性逐次判断，但同调函子一般不保持短正合列；缺少的连接箭头由提升余圈、施加微分和取商构造。Hom 复形使用直积，张量复形使用直和；微分中的符号使连续作用产生的混合项相消，分次复合满足带符号的 Leibniz 律，交换底环上的带符号张量交换与微分相容。

\ProseParagraphBreak
有界复形的边界对象在交错和中成对抵消，因此 Euler 特征可以直接由复形各项计算，也可以由同调计算。Grothendieck 群把这种计算表达为普遍的短正合关系，保留了数值维数可能忘掉的信息。下一章将比较复形之间的同伦与拟同构，说明哪些改变保持同调、哪些改变实际上更强。
\end{chaptersummary}

\BookExercises

\begin{exercise}\label{b4:ex:05-periodic}
在每个整数次数放 \(\mathbb Z/12\mathbb Z\)，每个微分均为乘 \(6\)。验证得到复形并计算各阶上同调。把 \(6\) 改为 \(4\) 后还能得到复形吗？
\end{exercise}
\begin{solution}
\(6^2\equiv0\pmod {12}\)，故是复形。乘 \(6\) 的核为 \(\{\bar0,\bar2,\bar4,\bar6,\bar8,\overline{10}\}\)，由 \(\bar2\) 生成，同构于 \(\mathbb Z/6\mathbb Z\)。像为 \(\{\bar0,\bar6\}\)，在此同构下对应 \(\mathbb Z/6\mathbb Z\) 中由 \(\bar3\) 生成的二阶子群。因此每阶上同调为
\[
\ker(\times6)/\operatorname{im}(\times6)
\cong(\mathbb Z/6\mathbb Z)/\langle\bar3\rangle
\cong\mathbb Z/3\mathbb Z.
\]
而 \(4^2\equiv4\not\equiv0\pmod {12}\)，连续两次微分不为零，不能谈其同调。
\end{solution}

\begin{exercise}\label{b4:ex:05-integer-matrices}
计算次数 \(0,1,2\) 的交换群复形
\(\mathbb Z\xrightarrow{a\mapsto(2a,4a)}\mathbb Z^2
\xrightarrow{(x,y)\mapsto6x-3y}\mathbb Z\)
的上同调。指出中间同调的一个生成元。
\end{exercise}
\begin{solution}
复合为 \(12a-12a=0\)。第一个映射单，故 \(H^0=0\)。第二个映射的核是 \(\mathbb Z(1,2)\)，而前一步的像是 \(2\mathbb Z(1,2)\)，故 \(H^1\cong\mathbb Z/2\mathbb Z\)，生成元由 \((1,2)\) 表示。末项上同调是 \(\mathbb Z/3\mathbb Z\)。
\end{solution}

\begin{exercise}\label{b4:ex:05-shift-support}
设 \(M\) 是模，\(M[0]\) 集中在次数零。分别写出 \(M[0][2]\)、\(M[0][-3]\) 的非零次数。若 \(C\) 的非零项仅在 \([-2,4]\)，\(C[3]\) 的非零项可能位于哪里？
\end{exercise}
\begin{solution}
\(M[0][2]\) 只在 \(-2\) 次非零，\(M[0][-3]\) 只在 \(3\) 次非零。\(C[3]^n=C^{n+3}\)，所以可能非零的范围为 \(-5\le n\le1\)。一般移位 \([r]\) 使支集向左移动 \(r\)，并使微分乘 \((-1)^r\)。
\end{solution}

\begin{exercise}\label{b4:ex:05-hom-support}
设 \(M,N\) 是左 \(R\)-模，\(a,b\in\mathbb Z\)。确定 \(\operatorname{Hom}_R^\bullet(M[0][a],N[0][b])\) 的全部项和微分。
\end{exercise}
\begin{solution}
源仅在 \(p=-a\) 非零，目标仅在 \(-b\) 非零，故 \(p+m=-b\) 等价于 \(m=a-b\)。该次数为 \(\operatorname{Hom}_R(M,N)\)，其余项为零，所有微分为零。因而它是 \(\operatorname{Hom}_R(M,N)[0][b-a]\)：移位 \([b-a]\) 把原来的零次项移到 \(-(b-a)=a-b\) 次，与上述计算一致。
\end{solution}

\begin{exercise}\label{b4:ex:05-tensor-integers}
设 \(m,n\) 是非零整数。令 \(P=(\mathbb Z\xrightarrow{m}\mathbb Z)\)、\(Q=(\mathbb Z\xrightarrow{n}\mathbb Z)\) 均放在次数 \(-1,0\)。写出 \(P\otimes Q\) 的微分并计算上同调。
\end{exercise}
\begin{solution}
在次数 \(-1\) 按 \(P^0\otimes Q^{-1}\)、\(P^{-1}\otimes Q^0\) 排列两项。于是复形为
\[
\mathbb Z\xrightarrow{t\mapsto(mt,-nt)}\mathbb Z^2
\xrightarrow{(u,v)\mapsto nu+mv}\mathbb Z.
\]
第一微分的第二分量来自 \(1\otimes\mathrm{d}_Q\)，其符号为 \((-1)^{-1}=-1\)，因为输入的第一因子在 \(P^{-1}\) 中；第二微分的两个输入分别有第一因子次数 \(0\) 和 \(-1\)，但后一输入只有 \(\mathrm{d}_P\otimes1\) 非零，故得到 \(nu+mv\)。令 \(g=\gcd(m,n)>0\)。第一箭头单，第二箭头的像为 \(g\mathbb Z\)，核由 \((m/g,-n/g)\) 生成，前一像为该生成元的 \(g\) 倍。因此 \(H^{-2}=0\)，\(H^{-1}\cong H^0\cong\mathbb Z/g\mathbb Z\)。
\end{solution}

\begin{exercise}\label{b4:ex:05-tensor-wrong-sign}
设 \(k\) 是特征不为二的域，\(C=D=(k\xrightarrow{1}k)\) 均在次数 \(0,1\)。若把张量微分误写成 \(\mathrm{d}(x\otimes y)=\mathrm{d}x\otimes y+x\otimes \mathrm{d}y\)，证明其平方不为零。
\end{exercise}
\begin{solution}
取次数零的基向量 \(e\in C^0\)、\(f\in D^0\)。其二次微分为 \(2\cdot\mathrm{d}_Ce\otimes \mathrm{d}_Df\)，这是次数二的一维空间中的非零向量。正确公式使两条混合路径的系数分别为 \(-1\) 和 \(1\)，从而相消。
\end{solution}

\begin{exercise}\label{b4:ex:05-homology-not-mono}
找出一个逐次单射的交换群复形映射，使其诱导的某个上同调映射不是单射。再找出逐次满射却不在某个上同调次数满射的例子。
\end{exercise}
\begin{solution}
令 \(D=(\mathbb Z\xrightarrow{2}\mathbb Z)\) 在次数 \(0,1\)。逐次乘 \(2\) 是单射，而 \(H^1(D)=\mathbb Z/2\mathbb Z\) 上诱导零映射，不单。逐次商映射 \(D\to D/2D\) 满，但 \(H^0(D)=0\)、\(H^0(D/2D)=\mathbb Z/2\mathbb Z\)，所以零次上同调映射不满。这不违背函子性：函子保持恒等、复合及同构，一般不保持单态射或满态射；\(H^n\) 的加性也不自动保证这种保持性。
\end{solution}

\begin{exercise}\label{b4:ex:05-prescribed-connecting}
给定左 \(R\)-模同态 \(u:M\to N\)，构造逐次分裂的复形短正合列，其连接同态 \(H^0(C)\to H^1(A)\) 恰为 \(u\)。
\end{exercise}
\begin{solution}
令 \(B=(M\xrightarrow{u}N)\) 在次数 \(0,1\)，令 \(A\) 仅在次数一等于 \(N\)，并嵌入 \(B\)，则 \(C=B/A\) 仅在次数零等于 \(M\)。两次数的短列都显然分裂。\(m\in C^0\) 的提升为同一 \(m\in B^0\)，其微分为 \(u(m)\in N\)，所以 \(\delta^0=u\)。因此连接同态从代数可能性上没有额外限制：任意给定模同态都可由这样的复形短正合列实现。
\end{solution}

\begin{exercise}\label{b4:ex:05-modfour-connecting}
取 \(D=(\mathbb Z\xrightarrow{2}\mathbb Z)\) 在次数 \(0,1\)。计算 \(0\to D\xrightarrow{4}D\to D/4D\to0\) 的连接同态，并写出非零部分的长正合列。
\end{exercise}
\begin{solution}
\(H^0(D)=0\)、\(H^1(D)=\mathbb Z/2\mathbb Z\)。在商复形中，\(H^0(D/4D)=\{\bar0,\bar2\}\cong\mathbb Z/2\mathbb Z\)，\(H^1(D/4D)=(\mathbb Z/4\mathbb Z)/2(\mathbb Z/4\mathbb Z)\cong\mathbb Z/2\mathbb Z\)。提升 \(\bar2\) 为整数二，微分为四，除以左端嵌入的四得到一，所以连接同态是上述生成元下的同构。长列为
\[
0\to\mathbb Z/2\mathbb Z\xrightarrow{1}\mathbb Z/2\mathbb Z
\xrightarrow{0}\mathbb Z/2\mathbb Z\xrightarrow{1}\mathbb Z/2\mathbb Z\to0.
\]
\end{solution}

\begin{exercise}\label{b4:ex:05-split-complex-les}
设交换范畴中的复形短正合列 \(0\to A\to B\to C\to0\) 有复形截面 \(s:C\to B\)。证明所有连接同态为零，并说明 \(H^n(B)\) 的分解依赖于什么选择。
\end{exercise}
\begin{solution}
记短正合列的箭头为 \(i:A\to B\)、\(p:B\to C\)。由 \(ps=1_C\) 得 \(H^n(p)H^n(s)=1\)，所以 \(H^n(p)\) 满；长正合列中的 \(\delta^nH^n(p)=0\) 因而给出 \(\delta^n=0\)。截面确定复形同构 \(F_s=(i,s):A\oplus C\to B\)，取上同调便得到 \(H^n(B)\cong H^n(A)\oplus H^n(C)\)。

若截面改为 \(s'=s+it\)，其中 \(t:C\to A\) 是复形映射，则同一对象的两套直和坐标由
\[
F_{s'}^{-1}F_s=
\begin{pmatrix}1_A&-t\\0&1_C\end{pmatrix}
\]
联系，因为 \(F_{s'}\begin{psmallmatrix}1&-t\\0&1\end{psmallmatrix}=(i,s'-it)=F_s\)。在模范畴中，这就是说 \(b=i(a)+s(c)=i(a-t(c))+s'(c)\)。由 \(H^n\) 的加性，上同调中的坐标变换为 \(\begin{psmallmatrix}1&-H^n(t)\\0&1\end{psmallmatrix}\)；模情形即 \((a,c)\mapsto(a-H^n(t)c,c)\)。所以此直和同构依赖所选的截面。
\end{solution}

\begin{exercise}\label{b4:ex:05-euler-ranks}
设域 \(k\) 上的复形仅在次数 \(0,1,2\) 非零，维数依次为 \(3,5,4\)，且两个微分的秩依次为 \(2,2\)。计算各阶同调维数和 Euler 特征。说明给定秩为何与复形条件相容。
\end{exercise}
\begin{solution}
维数分别为 \(3-2=1\)、\((5-2)-2=1\)、\(4-2=2\)。Euler 特征为 \(3-5+4=2=1-1+2\)。第二微分的核维数为 \(5-2=3\)，第一微分要求的像维数为 \(2\)，故可以把这二维像选在三维核中。再选第一微分为到该像的满射，就得到符合给定秩且复合为零的两张线性映射。
\end{solution}

\begin{exercise}\label{b4:ex:05-finite-groups-kzero}
证明有限交换群范畴的 Grothendieck 群是以素数为指标的自由交换群，其中 \([\mathbb Z/p\mathbb Z]\) 对应第 \(p\) 个基向量。
\end{exercise}
\begin{solution}
有限交换群有合成列：对非零群取非零的极小子群 \(H\)。任取 \(0\ne x\in H\)，极小性给出 \(H=\langle x\rangle\)；若 \(|x|=ab\) 且 \(a,b>1\)，则 \(\langle ax\rangle\) 是阶为 \(b\) 的真非零子群，矛盾。因此 \(H\) 的阶为素数，再对商群按群阶归纳。短正合关系因此把任意 \([M]\) 写成 \([\mathbb Z/p\mathbb Z]\) 的有限整数和。另一方面，令 \(v_p(|M|)\) 为群阶中素数 \(p\) 的指数；短正合列满足 \(|B|=|A|\cdot|C|\)，故这些赋值诱导 \(K_0\to\bigoplus_p\mathbb Z\)。它把 \([\mathbb Z/p\mathbb Z]\) 送到对应基向量，与前述生成映射互为逆，因而没有额外关系。
\end{solution}

\begin{exercise}\label{b4:ex:05-countable-swindle}
令 \(\mathcal A\) 为域 \(k\) 上至多可数维向量空间范畴。证明它本质小且为交换范畴，并计算 \(K_0(\mathcal A)\)。为何结果与有限维范畴不同？
\end{exercise}
\begin{solution}
同构类由 \(0,1,2,\ldots,\aleph_0\) 表示，故本质小；核、余核及有限双积仍至多可数维，故是交换范畴。令 \(V\) 有可数无限基，则对任意 \(W\in\mathcal A\)，都有 \(V\oplus W\cong V\)。分裂正合关系给出 \([V]+[W]=[V]\)，从而 \([W]=0\)，故 \(K_0=0\)。这里使用的是一个已经属于该范畴的无限维对象 \(V\) 及一个分裂短正合关系，并未在 \(K_0\) 中作无限求和。有限维对象不能满足 \(V\oplus k\cong V\)，而可数无限维的 \(V\) 能吸收任意至多可数维的 \(W\)，这使后者的类全部为零。整数维数不能延拓为这一范畴上的可加整数值函数，否则 \([k]=0\) 将与维数一矛盾。
\end{solution}

\begin{exercise}\label{b4:ex:05-kzero-more-than-dimension}
设 \(k\) 为域，\(R=k\times k\)，\(S_1=k\times0\)、\(S_2=0\times k\)。令复形 \(C\) 在次数零为 \(S_1\)、次数一为 \(S_2\)，微分为零。比较它的 \(k\)-维数 Euler 特征与 \(K_0\) 中的 Euler 特征。
\end{exercise}
\begin{solution}
数值为 \(1-1=0\)。在\cref{b4:exm:kzero-vectors}的同构 \(K_0\cong\mathbb Z^2\) 下，\(\chi_{K_0}(C)=(1,-1)\ne0\)。总维数对应群同态
\[
\dim_k:\mathbb Z^2\longrightarrow\mathbb Z,
\qquad(a,b)\longmapsto a+b,
\]
而 \((1,-1)\) 是其核中的非零元素。因此数值 Euler 特征把 \(K_0\) 的信息进一步压缩，总维数为零并不表示各种简单对象的贡献分别抵消。
\end{solution}

\begin{exercise}\label{b4:ex:05-euler-shift-tensor}
对有限维 \(k\)-向量空间的有界复形 \(C,D\)，证明 \(\chi(C[r])=(-1)^r\chi(C)\) 及 \(\chi(C\otimes_kD)=\chi(C)\chi(D)\)。这里是否需要先知道张量复形的同调公式？
\end{exercise}
\begin{solution}
第一式令 \(j=n+r\) 重新编号即可。第二式由有限求和得到
\[
\sum_n(-1)^n\sum_{p+q=n}\dim C^p\dim D^q
=\left(\sum_p(-1)^p\dim C^p\right)
 \left(\sum_q(-1)^q\dim D^q\right).
\]
这两式都不需要张量复形的同调公式：计算只用了各项的直和分解、维数对直和的可加性，以及 \(\dim_k(C^p\otimes_kD^q)=\dim_k C^p\,\dim_k D^q\)。Künneth 公式还进一步确定张量复形的同调，是更强的结论。
\end{solution}

\begin{exercise}\label{b4:ex:05-euler-zero-insufficient}
给出有界有限维复形，使 Euler 特征为零但同调不全为零；再说明仅给定 Euler 特征为什么不能恢复各阶同调。
\end{exercise}
\begin{solution}
取次数 \(0,1\) 的 \(k\xrightarrow0k\)。它有 \(H^0=H^1=k\)，而 Euler 特征为零。次数 \(0,1\) 的 \(k\xrightarrow1k\) 则同调全零且有相同 Euler 特征。交错和会把相邻次数的维数抵消，所以不能保留全部逐次信息。
\end{solution}

% !TeX root = ../../Book4.tex
% 纲要标题：## 第7章　链同伦、映射锥与拟同构
% 纲要标签：b4:ch:06
% 纲要位置：计划纲要.md，第 3342 行。

\chapter{链同伦、映射锥与拟同构}
\label{b4:ch:06}
\BookChapterNumber{b4:ch:06}

\begin{chapterabstract}
链同伦保证两个复形映射在同调上诱导相同映射；同伦等价则比拟同构要求更多。本章通过映射锥和映射柱比较这些概念，再用有限消元证明双复形总化的比较结果，并说明无界总化时直和与直积的差别。
\end{chapterabstract}

\begin{learninggoals}
\begin{itemize}
\item 构造链同伦，判断正合复形何时可缩，并区别同调相同与映射同伦。
\item 写出映射锥、映射柱及其收缩同伦，检验移位和连接箭头的符号。
\item 用锥的正合性识别拟同构，说明拟同构何时能加强为同伦等价。
\item 正确总化双复形，使用有限消元证明正合性，并识别无界总化的不同结果。
\end{itemize}
\end{learninggoals}

把一条短正合列视为只在三个连续次数非零的复形，正合性说明它的同调处处为零。然而，要把恒等链映射写成「微分与一个降次映射的复合之和」，需要逐项找到能退回前一项的映射。例如，对交换群短正合列 \(0\rightarrow\mathbb Z\xrightarrow{2}\mathbb Z\rightarrow\mathbb Z/2\mathbb Z\rightarrow0\)，这样的降次映射会迫使商映射有截面，而这个截面并不存在。可缩正是复形与零复形同伦等价，因此一定零调；这个例子表明，零调并不保证可缩。同调计算无法区分该复形与零复形，二者却不同伦等价。对一般复形映射，也可以追问它在同调上成为同构以后，能否在复形层面找到同伦逆。映射锥把前一种可逆性转化为一个新复形的正合性，使两种要求可以直接比较。

\ProseParagraphBreak
本章使用\chapref{b4:ch:05}的复形运算及上同调长正合列。同伦与映射锥可参阅 Weibel\cite[Section 1.4, Lemma 1.4.5; Section 1.5, Lemma 1.5.3, Corollary 1.5.4]{Weibel1994HomologicalAlgebra}；双复形、总化与符号见\cite[Example 1.2.4, Sign Trick 1.2.5, Total Complexes 1.2.6]{Weibel1994HomologicalAlgebra}。本章使用的锥矩阵和移位按本书约定书写，与上述文献比较时需留意分量次序与符号的转换。

% !TeX root = ../../Book4.tex
% 纲要标题：### 7.1 链同伦
% 纲要标签：b4:sec:0601
% 纲要位置：计划纲要.md，第 3344 行。

\section{链同伦}
\label{b4:sec:0601}

同调看不见的复形映射未必都能用同伦消去。先研究可显式写成 \(\mathrm{d}h+h\mathrm{d}\) 的差别，便能区分零伦、可缩与仅仅正合。

% 纲要内容（仅作写作计划，尚非教材正文）：
%
% * 同伦定义与同调不变性
% * 零伦映射及可缩复形
% * 正合复形未必可缩
%

\subsection{同伦与同调不变性}
\label{b4:subsec:0601:01}

两个复形映射可能逐项不同，却在同调上完全相同。最容易控制的一类差别是能够写成“先作微分再作某个降次映射，加上先降次再作微分”。这类差别在余圈上自动变成余边。

\ProseParagraphBreak
\cref{b4:def:cochain-complex,b4:def:complex-map}在交换范畴中定义了复形与复形映射；在加性范畴中，同样可以取对象族 \(C^n\) 和满足 \(\mathrm{d}_C^{n+1}\mathrm{d}_C^n=0\) 的态射 \(\mathrm{d}_C^n:C^n\to C^{n+1}\)，并要求复形映射与微分交换。这些定义只用到零态射和复合。同伦还用到态射的加减法，因而也可在加性范畴中讨论；核、像及同调对象则仍在交换范畴中使用。

\begin{definition}\label{b4:def:chain-homotopy}
设 \(\mathcal A\) 是加性范畴，\(C,D\) 为 \(\mathcal A\) 中的上链复形，\(f,g:C\to D\) 是复形映射。若有态射族 \(h^n:C^n\to D^{n-1}\)（\(n\in\mathbb Z\)），使
\[
f^n-g^n=\mathrm{d}_D^{n-1}h^n+h^{n+1}\mathrm{d}_C^n
\quad(n\in\mathbb Z),
\]
则称 \(f\) 与 \(g\) \term[lian-tonglun]{链同伦}[chain homotopic]，记为 \(f\simeq g\)，并称 \(h\) 为一个链同伦。链指标中同一个公式写为 \(f_n-g_n=\partial_{D,n+1}h_n+h_{n-1}\partial_{C,n}\)，其中 \(h_n:C_n\to D_{n+1}\)。
\end{definition}
\symidx{\(f\simeq g\)：复形映射之间的链同伦关系}

上链指标下同伦降一次，链指标下升一次，两者都与微分方向相反。在次数 \(n\)，两个复合
\[
\mathrm{d}_D^{n-1}h^n:C^n\longrightarrow D^n,
\qquad h^{n+1}\mathrm{d}_C^n:C^n\longrightarrow D^n
\]
分别先降次再升次、先升次再降次，因而都能与 \(f^n-g^n\) 相加比较。这里保留“链同伦”这一通用名称，不因为指标升次而另换概念。

\ProseParagraphBreak
对模复形，\cref{b4:def:hom-complex}中的 Hom 复形恰好解释了这个公式的来源。态射族 \(h\) 是次数 \(-1\) 的元素，其微分为
\[
\mathrm{d}_{\operatorname{Hom}}h
=\mathrm{d}_Dh-(-1)^{-1}h\mathrm{d}_C
=\mathrm{d}_Dh+h\mathrm{d}_C.
\]
所以 \(f\simeq g\) 就是说零次余圈 \(f-g\) 是 Hom 复形中的余边。

\begin{proposition}\label{b4:prop:homotopy-cohomology}
设 \(\mathcal A\) 是加性范畴。在任意两个固定的 \(\mathcal A\) 中的上链复形之间，同伦是复形映射的等价关系，且与加法、左右复合相容。若 \(\mathcal A\) 还是交换范畴，\(f,g:C\to D\) 为同伦的复形映射，则对每个 \(n\in\mathbb Z\)，有 \(H^n(f)=H^n(g)\)。
\end{proposition}
\begin{proof}
零同伦、\(-h\) 及同伦之和分别给出自反、对称和传递性。若 \(u:C'\to C\)、\(v:D\to D'\) 为复形映射，则
\[
v(f-g)u=\mathrm{d}_{D'}(vhu)+(vhu)\mathrm{d}_{C'};
\]
这证明复合相容性。两个差的同伦相加给出加法相容性。

再设 \(\mathcal A\) 是交换范畴，记 \(z_C^n,z_D^n\) 为余圈的核嵌入，\(j_D^n:B^n(D)\to Z^n(D)\) 为余边到余圈的嵌入，并写
\[
\mathrm{d}_D^{n-1}=z_D^nj_D^n\pi_D^{n-1},
\qquad \pi_D^{n-1}:D^{n-1}\twoheadrightarrow B^n(D).
\]
设 \(f_Z^n,g_Z^n:Z^n(C)\to Z^n(D)\) 是 \(f,g\) 在余圈上诱导的态射。把同伦公式右复合 \(z_C^n\)，含 \(h^{n+1}\mathrm{d}_C^n\) 的项消失，再消去单态射 \(z_D^n\)，得到
\[
f_Z^n-g_Z^n=j_D^n\pi_D^{n-1}h^nz_C^n.
\]
同调商映射 \(q_D^n:Z^n(D)\to H^n(D)\) 满足 \(q_D^nj_D^n=0\)，故 \(q_D^n(f_Z^n-g_Z^n)=0\)。由 \(q_C^n:Z^n(C)\to H^n(C)\) 为满态射，得到 \(H^n(f)-H^n(g)=0\)。在模中，这就是余圈 \(z\) 的两个像之差 \(f(z)-g(z)=\mathrm{d}_Dh(z)\) 是余边。
\end{proof}

因此对任意加性范畴 \(\mathcal A\)，都可以把同伦的映射视为同一个态射，得到\term[tonglun-fanchou]{同伦范畴}[homotopy category] \(K(\mathcal A)\)。对象仍是复形，态射则是复形映射的同伦类；上述相容性保证态射的复合与加法不依赖代表。\symidx{\(K(\mathcal A)\)：复形的同伦范畴}例如投射模构成的加性范畴也有同伦范畴，不必要求它本身是交换范畴。

\ProseParagraphBreak
若 \(C,D\) 是左 \(R\)-模复形，则
\[
\begin{aligned}
Z^0\operatorname{Hom}_R^\bullet(C,D)
&=\{C\to D\text{ 的复形映射}\},\\
B^0\operatorname{Hom}_R^\bullet(C,D)
&=\{\mathrm{d}_Dh+h\mathrm{d}_C\mid h\text{ 的次数为 }-1\}.
\end{aligned}
\]
右边第二个集合恰是零伦映射所成的子群。将零次余圈模掉这个子群，便是将两个同伦的复形映射视为相同，故
\[
\operatorname{Hom}_{K(R\text{-}\mathbf{Mod})}(C,D)
=H^0\operatorname{Hom}_R^\bullet(C,D),
\]
这把同伦类的定义与 Hom 复形的零次同调联系起来。

\subsection{零伦映射及可缩复形}
\label{b4:subsec:0601:02}

\begin{definition}\label{b4:def:contractible-complex}
设 \(\mathcal A\) 为加性范畴，\(C,D\) 为其中的上链复形，\(f:C\to D\) 为复形映射。若 \(f\) 与零映射同伦，则称 \(f\) 为\term[linglun-yingshe]{零伦映射}[null-homotopic map]。若 \(1_C\simeq0\)，则称 \(C\) 为\term[kesuo-fuxing]{可缩复形}[contractible complex]；满足 \(1_{C^n}=\mathrm{d}_C^{n-1}h^n+h^{n+1}\mathrm{d}_C^n\) 的态射族 \(h^n:C^n\to C^{n-1}\)（\(n\in\mathbb Z\)）称为收缩同伦。若存在复形映射 \(g:D\to C\)，使 \(gf\simeq1_C\)、\(fg\simeq1_D\)，则称 \(f\) 为\term[tonglun-dengjia]{同伦等价}[homotopy equivalence]。
\end{definition}

同伦等价所需的 \(g\) 是同伦逆，两个复合只要求与恒等映射同伦；复形同构的严格逆当然满足这一要求，反向却不必成立。可缩复形在交换范畴中必零调，因为恒等同态在同调上等于零。反向所缺的条件是各次数的短正合列能否分裂：分裂把 \(C^n\) 分成当前的余圈与下一次余边的选定提升，微分便只负责把后一个部分送到下一项。

\begin{proposition}\label{b4:prop:contractible-split-exact}
设 \(\mathcal A\) 是交换范畴，\(C\) 为其中的上链复形。它可缩，当且仅当它零调，且对每个 \(n\in\mathbb Z\)，短正合列
\[
0\longrightarrow Z^n(C)\xrightarrow{z_C^n}C^n
\xrightarrow{\widetilde{\mathrm{d}}_C^n} Z^{n+1}(C)\longrightarrow0
\]
均分裂；其中 \(z_C^n:Z^n(C)\to C^n\) 为核态射，\(\widetilde{\mathrm{d}}_C^n\) 由 \(\mathrm{d}_C^n=z_C^{n+1}\widetilde{\mathrm{d}}_C^n\) 定义，并在零调假设下为满态射。
\end{proposition}
\begin{proof}
若 \(1_C=\mathrm{d}h+h\mathrm{d}\)，则已知 \(C\) 零调。\(\widetilde{\mathrm{d}}_C^n\) 是微分的像分解中到 \(Z^{n+1}(C)=B^{n+1}(C)\) 的满态射；这里满的是这个分量，不是到整个 \(C^{n+1}\) 的微分。收缩同伦 \(h^{n+1}\) 将下一次余圈降回 \(C^n\)，因而截面的候选是
\[
s^{n+1}\coloneqq h^{n+1}z_C^{n+1}:Z^{n+1}(C)\longrightarrow C^n.
\]
把次数 \(n+1\) 的同伦公式右复合 \(z_C^{n+1}\)，由 \(\mathrm{d}_C^{n+1}z_C^{n+1}=0\) 得
\[
z_C^{n+1}=\mathrm{d}_C^nh^{n+1}z_C^{n+1}
=z_C^{n+1}\widetilde{\mathrm{d}}_C^nh^{n+1}z_C^{n+1}.
\]
消去单态射 \(z_C^{n+1}\)，得到
\[
\widetilde{\mathrm{d}}_C^n\bigl(h^{n+1}z_C^{n+1}\bigr)
=1_{Z^{n+1}(C)}.
\]
所以 \(s^{n+1}\) 确是 \(\widetilde{\mathrm{d}}_C^n\) 的截面。

反过来，选各短列的截面 \(s^{n+1}:Z^{n+1}(C)\to C^n\)。分裂短正合列给出的同构具体为
\[
\alpha^n=(z_C^n,s^{n+1}):Z^n(C)\oplus Z^{n+1}(C)\xrightarrow{\sim}C^n.
\]
第一分量由微分消去，第二分量经微分成为下一项的余圈嵌入，故在这些双积坐标中，微分为 \(\mathrm{d}(z,w)=(w,0)\)。定义 \(h^n(z,w)=(0,z)\in Z^{n-1}(C)\oplus Z^n(C)\)，则 \(\mathrm{d}h(z,w)=(z,0)\)、\(h\mathrm{d}(z,w)=(0,w)\)，相加等于恒等。这里的坐标公式表示双积的分量态射，因而计算在一般交换范畴中也成立。
\end{proof}

例如在次数 \(n,n+1\) 放同一对象 \(M\)，微分为 \(1_M\)，其余各项为零。只取 \(h^{n+1}=1_M\)，其余同伦分量为零，就得到收缩同伦。它的两个非零项可以很大，同调却为零；可缩所说的“消去”是通过显式同伦实现的，并不是逐项零对象。

\subsection{正合复形的非可缩例子}
\label{b4:subsec:0601:03}

\begin{example}\label{b4:exm:acyclic-not-contractible}
把交换群短正合列 \(0\to\mathbb Z\xrightarrow{2}\mathbb Z\to\mathbb Z/2\mathbb Z\to0\) 放在次数 \(0,1,2\)，便得到零调复形。若它可缩，记次数一的商映射为 \(\mathrm{d}^1\)，由于 \(\mathrm{d}^2=0\)，次数二的同伦公式就是 \(1_{\mathbb Z/2\mathbb Z}=\mathrm{d}^1h^2\)。这要求 \(h^2:\mathbb Z/2\mathbb Z\to\mathbb Z\) 为商映射的截面。但任何这样的群同态都为零，因此不存在截面，复形不可缩。
\end{example}

\begin{example}\label{b4:exm:acyclic-free-periodic}
令 \(R=\mathbb Z/4\mathbb Z\)，在每个整数次数取自由模 \(R\)，微分均为乘 \(2\)。因为核与像同为 \(2R\)，复形正合。每个 \(R\)-线性 \(h^n:R\to R\) 都是乘某个 \(a_n\in R\)。若有收缩同伦，便须
\[
1=2a_n+2a_{n+1}\quad\text{在}\;R\;\text{中成立},
\]
右端是偶数剩余类，矛盾。这说明即使各项自由，无界正合复形也未必可缩。

\ProseParagraphBreak
此时余圈 \(2R\cong R/(2)\) 不是投射模。否则满射 \(R\xrightarrow{2}2R\) 会有截面 \(s:2R\to R\)；由 \(2s(2)=s(2\cdot2)=0\) 可知 \(s(2)\in2R\)，所以乘 \(2\) 又将它送到零，无法得到所要求的 \(2\)。因此各项自由并没有保证上述余圈短正合列分裂。
\end{example}

作为对照，对任意含幺环 \(R\)，各项均为投射幺性左 \(R\)-模的上有界零调复形必可缩。取 \(b\) 使 \(n>b\) 时 \(C^n=0\)，则 \(Z^b(C)=C^b\) 投射；于是 \(0\to Z^{b-1}(C)\to C^{b-1}\to Z^b(C)\to0\) 分裂，且 \(Z^{b-1}(C)\) 作为投射模的直和因子仍投射。如此向较低次数归纳，所有余圈短正合列都分裂，再用上述命题得到可缩性。周期反例没有可供这一步归纳开始的最高次数。后面比较投射消解时，也将从一个有限端开始逐次提升同伦，不能在没有起点的无界复形上直接照搬。

% !TeX root = ../../Book4.tex
% 纲要标题：### 7.2 映射锥
% 纲要标签：b4:sec:0602
% 纲要位置：计划纲要.md，第 3350 行。

\section{映射锥}
\label{b4:sec:0602}

一个复形映射也可以成为新复形的微分组成部分。映射锥将它在上同调上诱导的映射的核与余核信息集中到长正合列中，柱构造则把相关比较写成具体映射。

% 纲要内容（仅作写作计划，尚非教材正文）：
%
% * 锥与柱的显式构造
% * 锥的上同调长正合列
% * 移位和锥的符号检验
%

\subsection{锥与柱的显式构造}
\label{b4:subsec:0602:01}

上同调长正合列从短正合列出发，而许多问题最初只有一个映射 \(f:C\to D\)。要把它转成上同调之间的比较，可以构造一个逐次分裂的复形短列，左端为 \(D\)，右端为 \(C[1]\)。于是中项的次数 \(n\) 取 \(D^n\oplus C[1]^n=D^n\oplus C^{n+1}\)，两个对角微分分别取 \(\mathrm{d}_D^n\)、\(\mathrm{d}_{C[1]}^n=-\mathrm{d}_C^{n+1}\)。再把给定的 \(f^{n+1}:C^{n+1}\to D^{n+1}\) 放在非对角位置，便得到下面的映射锥。移位使这个非对角项的源、目标次数能够相接，所得长正合列将同时比较两端的上同调与 \(H(f)\)。

\begin{definition}\label{b4:def:mapping-cone}
设 \(f:C\to D\) 是加性范畴中的上链复形映射。规定
\[
\operatorname{Cone}(f)^n=D^n\oplus C^{n+1},\qquad
 \mathrm{d}_{\operatorname{Cone}(f)}^n=
 \begin{pmatrix}\mathrm{d}_D^n&f^{n+1}\\0&-\mathrm{d}_C^{n+1}\end{pmatrix}.
\]
所得复形称为 \(f\) 的\term[yingshe-zhui]{映射锥}[mapping cone]。在模中写作 \(\mathrm{d}(y,x)=(\mathrm{d}_Dy+f(x),-\mathrm{d}_Cx)\)。
\end{definition}
\symidx{\(\operatorname{Cone}(f)\)：复形映射的映射锥}

右下角的负号正是\cref{b4:def:complex-shift}中 \(C[1]\) 的微分符号。矩阵平方的非对角项为 \(\mathrm{d}_Df-f\mathrm{d}_C=0\)，对角项为两个微分的平方，因此确为复形。若把右下角的负号删去，非对角项通常会变成 \(\mathrm{d}_Df+f\mathrm{d}_C\)，不能保证为零。

映射柱要把 \(f\) 分解成 \(C\xrightarrow{j}\operatorname{Cyl}(f)\xrightarrow{r}D\)，其中 \(j\) 逐次分裂单，\(r\) 为同伦等价，并保持 \(rj=f\)。为给 \(C\) 留出嵌入，同时保留 \(D\) 的同伦信息，可以在 \(D\) 上添入由 \(C^n\oplus C^{n+1}\) 组成的可缩辅助复形。下面的三个分量正是 \(D\) 与这份辅助数据；定义后将用坐标变换识别它们。

\begin{definition}\label{b4:def:mapping-cylinder}
设 \(f:C\to D\) 是加性范畴中的上链复形映射。规定
\[
\operatorname{Cyl}(f)^n=D^n\oplus C^n\oplus C^{n+1},
\qquad
 \mathrm{d}(y,c,a)=(\mathrm{d}_Dy+f(a),\ \mathrm{d}_Cc-a,\ -\mathrm{d}_Ca).
\]
所得复形称为 \(f\) 的\term[yingshe-zhu]{映射柱}[mapping cylinder]。元组公式在一般加性范畴中表示相应双积矩阵。
\end{definition}
\symidx{\(\operatorname{Cyl}(f)\)：复形映射的映射柱}

这三项可以从一个坐标变换理解。令 \(z=y+f(c)\)，即把 \((y,c,a)\) 换成 \((z,c,a)\)，逆变换为 \((z,c,a)\mapsto(z-f(c),c,a)\)。因为 \(f\) 与微分交换，新坐标中的微分为
\[
\begin{gathered}
\mathrm{d}_Dy+f(a)+f(\mathrm{d}_Cc-a)
=\mathrm{d}_D(y+f(c))=\mathrm{d}_Dz,\\
\mathrm{d}(z,c,a)=(\mathrm{d}_Dz,\,\mathrm{d}_Cc-a,\,-\mathrm{d}_Ca).
\end{gathered}
\]
第一分量因而与后两分量分开：映射柱同构于 \(D\oplus K\)，其中 \(K^n=C^n\oplus C^{n+1}\)，微分为 \((c,a)\mapsto(\mathrm{d}_Cc-a,-\mathrm{d}_Ca)\)。辅助复形 \(K\) 可缩，收缩为 \((c,a)\mapsto(0,-c)\)。映射柱是在 \(D\) 上添入这一份同伦上为零的复形；同时 \(c\mapsto(0,c,0)\) 为 \(C\) 留出逐次分裂的嵌入。在新坐标中，该嵌入是 \(c\mapsto(f(c),c,0)\)，而到 \(D\) 的投影正是原坐标中的 \(y+f(c)\)。

\begin{proposition}\label{b4:prop:cylinder-factorization}
设 \(\mathcal A\) 为加性范畴，\(C,D\) 为其中的上链复形，\(f:C\to D\) 为复形映射。令 \(j:C\to\operatorname{Cyl}(f)\)、\(r:\operatorname{Cyl}(f)\to D\)、\(s:D\to\operatorname{Cyl}(f)\) 的逐次公式为
\[
j(c)=(0,c,0),\qquad r(y,c,a)=y+f(c),\qquad s(y)=(y,0,0).
\]
则 \(rj=f\)、\(rs=1_D\)，\(r,s\) 互为同伦逆。\(j:C\to\operatorname{Cyl}(f)\) 逐次为分裂单态射，且 \(q^n(y,c,a)=(y,a)\) 给出 \(j\) 的余核复形 \(\operatorname{Cone}(f)\)。在交换范畴中因而有逐次分裂短正合列
\[
0\longrightarrow C\xrightarrow{j}\operatorname{Cyl}(f)
\xrightarrow{q}\operatorname{Cone}(f)\longrightarrow0.
\]
\end{proposition}
\begin{proof}
将微分公式再用一次，三个分量分别为 \(\mathrm{d}_D^2y+\mathrm{d}_Df(a)-f\mathrm{d}_C(a)\)、\(\mathrm{d}_C^2c-\mathrm{d}_Ca+\mathrm{d}_Ca\)、\(\mathrm{d}_C^2a\)，均为零。由 \(f\) 与微分交换，\(j,r,s\) 都是复形映射；两个复合等式直接成立。

还需证明 \(1-sr\) 零伦。按同伦约定，比较恒等态射与 \(sr\)，要找到 \(h\) 使 \(1-sr=\mathrm{d}h+h\mathrm{d}\)。定义降次态射 \(h(y,c,a)=(0,0,-c)\)，则
\[
\mathrm{d}h(y,c,a)=(-f(c),c,\mathrm{d}_Cc),\qquad
h\mathrm{d}(y,c,a)=(0,0,-\mathrm{d}_Cc+a),
\]
相加得 \((-f(c),c,a)=(1-sr)(y,c,a)\)。因此 \(sr\simeq1\)。

\(j^n\) 的逐次左逆是 \(\pi_C^n(y,c,a)=c\)，满足 \(\pi_C^nj^n=1_{C^n}\)。逐次投影到第一、第三个双积分量得到 \(q\)，其逐次截面为 \(\sigma^n(y,a)=(y,0,a)\)，满足 \(q^n\sigma^n=1\)。其目标上的微分为 \((\mathrm{d}_Dy+f(a),-\mathrm{d}_Ca)\)，正是映射锥的微分，且 \(qj=0\)。若复形映射 \(u:\operatorname{Cyl}(f)\to E\) 满足 \(uj=0\)，则双积泛性质给出唯一的态射族 \(v^n:\operatorname{Cone}(f)^n\to E^n\)，使 \(u^n=v^nq^n\)。由于每个 \(q^n\) 都是分裂满态射，等式
\[
\begin{aligned}
\bigl(\mathrm{d}_E^nv^n-v^{n+1}\mathrm{d}_{\operatorname{Cone}(f)}^n\bigr)q^n
&=\mathrm{d}_E^nu^n-u^{n+1}\mathrm{d}_{\operatorname{Cyl}(f)}^n\\
&=0
\end{aligned}
\]
可以在右边消去满态射 \(q^n\)，得到 \(\mathrm{d}_E^nv^n=v^{n+1}\mathrm{d}_{\operatorname{Cone}(f)}^n\)。所以逐次分解组成复形映射 \(v\)，\(q\) 满足 \(j\) 在复形范畴中的余核泛性质。
\end{proof}

这里逐次分裂是指每个次数上的短正合列分裂；复形分裂还要求分裂态射与微分交换。上面的 \(\pi_C\) 与 \(\sigma\) 通常不满足这个要求：\(\pi_C\mathrm{d}(y,c,a)=\mathrm{d}_Cc-a\)，而 \(\mathrm{d}_C\pi_C(y,c,a)=\mathrm{d}_Cc\)；\(\mathrm{d}\sigma(y,a)\) 的中间分量为 \(-a\)，而 \(\sigma\mathrm{d}(y,a)\) 的中间分量为零。可缩则是单个复形的恒等映射零伦的条件；这里可缩的是辅助复形 \(K\)，映射柱本身仍保留 \(D\) 的同伦信息。

\subsection{锥的上同调长正合列}
\label{b4:subsec:0602:02}

\begin{proposition}\label{b4:prop:cone-long-exact}
设 \(\mathcal A\) 是交换范畴，\(C,D\) 为其中的上链复形，\(f:C\to D\) 是复形映射。嵌入 \(y\mapsto(y,0)\) 与投影 \((y,x)\mapsto x\) 给出逐次分裂短正合列
\[
0\longrightarrow D\longrightarrow\operatorname{Cone}(f)
\longrightarrow C[1]\longrightarrow0.
\]
其连接态射在 \(H^n(C[1])=H^{n+1}(C)\) 的识别下就是 \(H^{n+1}(f)\)。因此有自然上同调长正合列
\[
\cdots\to H^n(C)\xrightarrow{H^n(f)}H^n(D)
\to H^n\operatorname{Cone}(f)\to H^{n+1}(C)
\xrightarrow{H^{n+1}(f)}H^{n+1}(D)\to\cdots.
\]
\end{proposition}
\begin{proof}
微分矩阵表明嵌入与投影均为复形映射，前者为后者的核，后者为前者的余核；这些泛性质由逐次双积分解得到。先应用\cref{b4:thm:cohomology-long-exact}，所得原始片段为
\[
H^n(D)\longrightarrow H^n\operatorname{Cone}(f)
\longrightarrow H^n(C[1])\xrightarrow{\delta^n}H^{n+1}(D).
\]
再用 \(H^n(C[1])=H^{n+1}(C)\) 识别移位项，并计算 \(\delta^n\)。投影的逐次截面 \(x\mapsto(0,x)\) 不必为复形映射，因为其微分还产生 \(f(x)\) 分量。对余圈子对象 \(Z^{n+1}(C)\) 使用这一截面，再作锥微分，得到进入 \(D^{n+1}\) 的态射 \(f^{n+1}|_{Z^{n+1}(C)}\)。由\cref{b4:prop:connecting-map-construction}，它在上同调上诱导的正是连接态射 \(H^{n+1}(f)\)。在模中，取余圈 \(x\in Z^{n+1}(C)\)，它代表 \(H^n(C[1])\) 中的同一个类。计算为
\[
\begin{aligned}
x&\longmapsto(0,x),\\
\mathrm{d}_{\operatorname{Cone}(f)}^n(0,x)
&=(f^{n+1}(x),-\mathrm{d}_C^{n+1}x)\\
&=(f^{n+1}(x),0),
\end{aligned}
\]
故 \(\delta^n[x]=[f^{n+1}(x)]\)，没有额外负号。最终长列中的前一个箭头 \(H^n(C)\to H^n(D)\) 来自上一次数的连接 \(\delta^{n-1}:H^{n-1}(C[1])\to H^n(D)\)，在移位识别下正是 \(H^n(f)\)。
\end{proof}

\begin{example}\label{b4:exm:cone-single-map}
把左 \(R\)-模 \(M,N\) 放在次数零，同态 \(u:M\to N\) 的锥只有次数 \(-1,0\) 非零，恰为 \(M\xrightarrow{u}N\)。所以
\[
H^{-1}\operatorname{Cone}(u)=\ker u,\qquad
H^0\operatorname{Cone}(u)=\operatorname{coker}u.
\]
例如 \(\mathbb Z\xrightarrow{m}\mathbb Z\) 的锥在 \(m\ne0\) 时只剩零次上同调 \(\mathbb Z/m\mathbb Z\)。因此 \(\operatorname{Cone}(u)\) 零调，当且仅当 \(\ker u=0\)、\(\operatorname{coker}u=0\)，也就是 \(u\) 为模同构。下一节的\cref{b4:prop:cone-quasi-isomorphism}将把这一原型推广到任意复形映射：锥零调恰好检测它是否在所有次数诱导上同调同构。
\end{example}

\subsection{移位和锥的符号检验}
\label{b4:subsec:0602:03}

锥的不同符号约定可以彼此同构，但必须明确写出比较映射。例如有些文献把 \(C^{n+1}\) 放在前面，并在 \(f\) 前加负号；这时交换两个分量并改变其中一个分量的符号，才能与本书定义比较。

\begin{proposition}\label{b4:prop:cone-sign-compatibility}
设 \(\mathcal A\) 为加性范畴，\(C,D\) 为其中的上链复形，\(f,g:C\to D\) 为复形映射。以下元组公式表示相应双积上的态射矩阵。
\begin{enumerate}
\item 若态射族 \(h^n:C^n\to D^{n-1}\) 满足 \(f-g=\mathrm{d}_Dh+h\mathrm{d}_C\)，则次数 \(n\) 上的公式 \((y,x)\mapsto(y+h^{n+1}(x),x)\) 给出 \(\operatorname{Cone}(f)\cong\operatorname{Cone}(g)\)。这个同构由所给的 \(h\) 确定，不同同伦一般会给出不同同构。
\item 对整数 \(r\)，映射 \((y,x)\mapsto(y,(-1)^r x)\) 给出 \(\operatorname{Cone}(f[r])\cong\operatorname{Cone}(f)[r]\)。
\item \(\operatorname{Cone}(1_C)\) 可缩，收缩同伦为 \((y,x)\mapsto(0,y)\)。
\end{enumerate}
\end{proposition}
\begin{proof}
第一项的映射与微分交换等价于
\[
\mathrm{d}_Dy+\mathrm{d}_Dh(x)+g(x)=\mathrm{d}_Dy+f(x)-h(\mathrm{d}_Cx),
\]
这恰为同伦公式；逆映射把 \(h\) 换成 \(-h\)。第二项中 \(f[r]^n=f^{n+r}\)，不另加符号。源 \(\operatorname{Cone}(f[r])^n\) 与目标 \(\operatorname{Cone}(f)[r]^n\) 的底对象都是 \(D^{n+r}\oplus C^{n+r+1}\)，差别只在微分。暂略已经移到 \(n+r\) 的指标，源的微分矩阵为
\[
\begin{pmatrix}(-1)^r\cdot\mathrm{d}_D&f\\0&-(-1)^r\cdot\mathrm{d}_C\end{pmatrix},
\]
目标的微分矩阵为 \((-1)^r\begin{pmatrix}\mathrm{d}_D&f\\0&-\mathrm{d}_C\end{pmatrix}\)。令 \(\epsilon=(-1)^r\)，比较态射 \(t^n=\operatorname{diag}(1,\epsilon)\) 在第二分量补上符号。用 \(\epsilon^2=1\) 计算，两种复合的矩阵均为
\[
t^{n+1}\mathrm{d}_{\operatorname{Cone}(f[r])}^n
=\mathrm{d}_{\operatorname{Cone}(f)[r]}^nt^n
=\begin{pmatrix}\epsilon\cdot\mathrm{d}_D&f\\0&-\mathrm{d}_C\end{pmatrix}.
\]
而 \((t^n)^2=1\)，故给出复形同构。第三项令 \(h(y,x)=(0,y)\)，则 \(\mathrm{d}h(y,x)=(y,-\mathrm{d}y)\)、\(h\mathrm{d}(y,x)=(0,\mathrm{d}y+x)\)，其和为 \((y,x)\)。
\end{proof}

移位后的同伦也须改符号：若 \(h\) 是 \(f\simeq g\) 的同伦，则 \(h_r^n=(-1)^r h^{n+r}\) 是 \(f[r]\simeq g[r]\) 的同伦。这从两边微分均乘 \((-1)^r\) 直接得出。当 \(r\) 为奇数时，若只移动指标而不乘这个符号，同伦公式右边会多出整体负号；\(r\) 为偶数时则无此差异。

% !TeX root = ../../Book4.tex
% 纲要标题：### 7.3 拟同构
% 纲要标签：b4:sec:0603
% 纲要位置：计划纲要.md，第 3258 行。

\section{拟同构}
\label{b4:sec:0603}

同调上的可逆性弱于在复形层面构造同伦逆。映射锥提供拟同构的判据，而非分裂正合列说明这两种可逆性为何会分离。

% 纲要内容（仅作写作计划，尚非教材正文）：
%
% * 拟同构的定义
% * 与同伦等价的区别
% * 非分裂正合列产生的基本例子
%

\subsection{拟同构的定义}
\label{b4:subsec:0603:01}

\begin{definition}\label{b4:def:quasi-isomorphism}
设 \(\mathcal A\) 是交换范畴，\(f:C\to D\) 是复形映射。若每个整数 \(n\) 的 \(H^n(f):H^n(C)\to H^n(D)\) 都是同构，则称 \(f\) 为\term[ni-tonggou]{拟同构}[quasi-isomorphism]。
\end{definition}

拟同构必须是已经给定的复形映射。“两个复形的同调对象逐次同构”只给出对象之间的比较，未必指定了能够实现这些同构的复形映射。更不能由此断言原来的某个映射就是拟同构。

\begin{proposition}\label{b4:prop:cone-quasi-isomorphism}
设 \(\mathcal A\) 是交换范畴，\(f:C\to D\) 是 \(\mathcal A\) 中上链复形之间的映射。则 \(f\) 是拟同构，当且仅当 \(\operatorname{Cone}(f)\) 零调。对 \(\mathcal A\) 中任意可复合的复形映射 \(f:C\to D\)、\(g:D\to E\)，若 \(f,g,gf\) 的任意两个为拟同构，则第三个也是拟同构。
\end{proposition}
\begin{proof}
锥的长正合列中有
\[
H^n(C)\xrightarrow{H^n(f)}H^n(D)\to
H^n\operatorname{Cone}(f)\to H^{n+1}(C)
\xrightarrow{H^{n+1}(f)}H^{n+1}(D).
\]
若左右两个 \(H(f)\) 都是同构，中间对象到右侧的箭头既单又为零，故中间对象为零。若所有锥同调均为零，则相关正合段给出 \(H^n(f)\) 既单又满，故为同构。最后，对 \(H^n(f),H^n(g)\) 使用同构的复合与消去性质，再让 \(n\) 遍历全部整数即可。
\end{proof}

\begin{proposition}\label{b4:prop:short-exact-quasi-comparison}
设交换范畴中两条复形短正合列之间给定态射。如果三个竖直复形映射中有两个为拟同构，则第三个也为拟同构。
\end{proposition}
\begin{proof}
由\cref{b4:thm:cohomology-long-exact}得到两条长正合列及其间交换图。对待证映射 \(H^n\) 所在的项，取其前后各两项组成五项正合图。其余四个竖直态射均是另两个复形在某个次数的同调态射，所以都是同构。\cref{b4:thm:five-lemma}证明中间态射也是同构。
\end{proof}

\subsection{与同伦等价的区别}
\label{b4:subsec:0603:02}

同伦等价必为拟同构，因为同伦逆在同调上成为真正的逆；这是\cref{b4:prop:homotopy-cohomology}的直接应用。拟同构只要求同调上可逆，同伦等价还要求构造反向的复形映射及两个同伦，故要求更强。

\begin{proposition}\label{b4:prop:field-complex-splitting}
设 \(k\) 是域，\(C\) 是任意 \(k\)-向量空间复形。把各个 \(H^n(C)\) 组成微分全零的复形 \(H(C)\)，则 \(C\) 与 \(H(C)\) 同伦等价。因此 \(k\)-向量空间复形之间的每个拟同构都是同伦等价。所用同伦等价一般依赖补空间的选择。
\end{proposition}
\begin{proof}
在每个次数选择补空间
\[
Z^n(C)=B^n(C)\oplus V^n,\qquad C^n=Z^n(C)\oplus W^n.
\]
微分的限制
\[
\mathrm{d}_C^n|_{W^n}:W^n\xrightarrow{\sim}B^{n+1}(C)
\]
是同构，因为其核与 \(W^n\) 的交为零且其像就是全部边界。令 \(q^n:Z^n(C)\to H^n(C)\) 为商映射，则其限制 \(\rho^n=q^n|_{V^n}\) 也是同构。\(H^n(C)\) 本身是商空间；通过 \((\rho^n)^{-1}\)，才能将它识别为选定的代表子空间 \(V^n\subseteq C^n\)。

利用分解 \(C^n=B^n(C)\oplus V^n\oplus W^n\)，定义
\[
i^n:H^n(C)\longrightarrow C^n,\qquad
i^n(\eta)=(\rho^n)^{-1}(\eta),
\]
以及
\[
p^n:C^n\longrightarrow H^n(C),\qquad
p^n(b+v+w)=\rho^n(v).
\]
\(i^n\) 的像在余圈中，而 \(p^{n+1}\) 在微分的像 \(B^{n+1}(C)\) 上为零，故它们组成复形映射 \(i:H(C)\to C\)、\(p:C\to H(C)\)，且 \(pi=1_{H(C)}\)。再定义
\[
h^n(b+v+w)
=\bigl(\mathrm{d}_C^{n-1}|_{W^{n-1}}\bigr)^{-1}(b).
\]
逐个检查这三个分量，得到
\[
\mathrm{d}_C^{n-1}h^n+h^{n+1}\mathrm{d}_C^n
=1_{C^n}-i^np^n.
\]
于是 \(ip\simeq1_C\)，这给出所需同伦等价。

对拟同构 \(f:C\to D\)，给两者分别作上述选择。\(i_C\) 选取同调类的余圈代表，而 \(p_D\) 在余圈上的限制就是取同调类，因此 \(p_Dfi_C=H(f)\) 是微分为零复形之间的同构。取 \(g=i_C H(f)^{-1}p_D\)。由同伦对复合的相容性，
\[
gf\simeq gf i_Cp_C=i_Cp_C\simeq1_C,
\qquad
fg\simeq i_Dp_Dfg=i_Dp_D\simeq1_D.
\]
所以 \(g\) 是 \(f\) 的同伦逆。选择补空间改变了 \(i,p,h\)，所以这不是无须选择的复形分解。
\end{proof}

上述证明使用的是向量空间每个短正合列都分裂。对一般模，核和像可能没有补模；不能只把 \(k\) 改成任意环便保留结论。

\subsection{非分裂正合列产生的基本例子}
\label{b4:subsec:0603:03}

\begin{example}\label{b4:exm:resolution-quasi-not-homotopy}
令 \(P\) 为次数 \(-1,0\) 的复形 \(\mathbb Z\xrightarrow{2}\mathbb Z\)，令 \(M=(\mathbb Z/2)[0]\)。商映射给出增广 \(\varepsilon:P\to M\)，且 \(H^{-1}(P)=0\)、\(H^0(P)=\mathbb Z/2\)，所以它是拟同构。

\ProseParagraphBreak
但任意复形映射 \(g:M\to P\) 的零次分量 \(\mathbb Z/2\to\mathbb Z\) 都为零，故 \(g=0\)。在集中于零次的复形 \(M\) 上，所有可能的降次同伦也为零，因此 \(\varepsilon g=0\) 不可能同伦于 \(1_M\)。所以 \(\varepsilon\) 不是同伦等价。
\end{example}

\begin{example}\label{b4:exm:additive-not-quasi-preserving}
对上一例逐次应用加性函子 \(\operatorname{Hom}_{\mathbb Z}(\mathbb Z/2,-)\)。由于 \(\operatorname{Hom}_{\mathbb Z}(\mathbb Z/2,\mathbb Z)=0\)，所得 \(F(P)\) 是零复形；而 \(F(M)\) 集中于零次且等于 \(\mathbb Z/2\)。因此 \(F(\varepsilon)\) 不是拟同构。
\end{example}

加性函子会保持已经给定的链同伦，因为它保持同伦公式中的加法与复合，却未必保持拟同构。后面先取合适消解再施加函子，正是为了处理这种差别。对非分裂短正合列形成的零调复形，零复形到它的映射也是拟同构；若它不是可缩复形，就同样不会成为同伦等价。

% !TeX root = ../../Book4.tex
% 纲要标题：### 7.4 双复形的预备
% 纲要标签：b4:sec:0604
% 纲要位置：计划纲要.md，第 3264 行。

\section{双复形的预备}
\label{b4:sec:0604}

映射锥把拟同构的判定化为一个复形的正合性。若分别为张量积的两个变量选取消解，两个次数会同时出现；比较这些消解时，知道每列的上同调相同，还须说明合成后的复形是否也有相同的上同调。下面沿每条对角线合并两个次数，并用映射锥把逐列比较归结为有限消元。对角线上有无限多项时，消元可能需要无限步，直和与直积也可能给出不同的上同调。

% 纲要内容（仅作写作计划，尚非教材正文）：
%
% * 横、纵微分的反交换约定
% * 直和总复形与直积总复形
% * 无界情形的差异暂不隐去
% * 在后续谱序列所用反交换约定下取总微分 \(\mathrm{d}=\mathrm{d}_h+\mathrm{d}_v\)；若从交换方格起步，明确在哪个方向加入 \((-1)^p\) 符号后再总化
%
% ---
%

\subsection{横纵微分的反交换约定}
\label{b4:subsec:0604:01}

给两个复形作张量积时，横向应用第一个复形的微分，纵向应用第二个复形的微分。若直接把两个微分相加，平方里还会出现交叉项；反交换条件恰好使这两项抵消。

\begin{definition}\label{b4:def:double-complex}
设 \(\mathcal A\) 是加性范畴。给定 \(\mathcal A\) 中的对象 \(E^{p,q}\)（\(p,q\in\mathbb Z\)）及态射
\[
\mathrm{d}_h:E^{p,q}\to E^{p+1,q},\qquad
 \mathrm{d}_v:E^{p,q}\to E^{p,q+1},
\]
若每个位置都满足 \(\mathrm{d}_h^2=\mathrm{d}_v^2=0\) 且 \(\mathrm{d}_h\mathrm{d}_v+\mathrm{d}_v\mathrm{d}_h=0\)，则称 \(E\) 为\term[shuang-fuxing]{双复形}[double complex]。这里两个方向的方格反交换。对 \(\mathcal A\) 中的双复形 \(E,F\)，若态射族 \(u^{p,q}:E^{p,q}\to F^{p,q}\) 分别与两个方向的微分交换，则称其为双复形映射。
\end{definition}

例如右 \(R\)-模复形 \(C\) 与左 \(R\)-模复形 \(D\) 给出
\[
E^{p,q}=C^p\otimes_R D^q,\qquad
 \mathrm{d}_h=\mathrm{d}_C\otimes1,\qquad \mathrm{d}_v=(-1)^p(1\otimes \mathrm{d}_D).
\]
先横后纵时，纵向符号变为 \((-1)^{p+1}\)；先纵后横时仍为 \((-1)^p\)，所以两次复合之和为零。这个例子会在 Tor 的两个变量比较中使用。

\ProseParagraphBreak
若最初给出的横纵微分 \(\partial_h,\partial_v\) 满足交换关系 \(\partial_h\partial_v=\partial_v\partial_h\)，则取 \(\mathrm{d}_h=\partial_h\)、\(\mathrm{d}_v|_{E^{p,q}}=(-1)^p\partial_v\) 才得到本书意义下的双复形。后文总化时不再另添一次相同的符号。

\subsection{直和总复形与直积总复形}
\label{b4:subsec:0604:02}

\begin{definition}\label{b4:def:total-complex}
设 \(E\) 是交换范畴 \(\mathcal A\) 中的双复形。若相应的可数余积或积存在，对每个 \(n\in\mathbb Z\) 分别规定
\[
\operatorname{Tot}_{\oplus}(E)^n=\bigoplus_{p+q=n}E^{p,q},\qquad
\operatorname{Tot}_{\Pi}(E)^n=\prod_{p+q=n}E^{p,q},
\qquad \mathrm{d}=\mathrm{d}_h+\mathrm{d}_v.
\]
所得复形分别称为 \(E\) 的\term[zhihe-zongfuxing]{直和总复形}[direct-sum total complex]与\term[zhiji-zongfuxing]{直积总复形}[product total complex]。若每条对角线 \(p+q=n\) 仅有有限个非零对象，则二者都是有限双积，典范相同，简记为 \(\operatorname{Tot}(E)\)。
\end{definition}
\symidx{\(\operatorname{Tot}_{\oplus},\operatorname{Tot}_{\Pi}\)：直和与直积总复形}

直和中每个输入分量至多产生两个输出分量；直积中每个输出分量也只依赖两个输入分量，因此公式都定义了态射。反交换关系给出 \(\mathrm{d}^2=\mathrm{d}_h^2+\mathrm{d}_h\mathrm{d}_v+\mathrm{d}_v\mathrm{d}_h+\mathrm{d}_v^2=0\)。一般交换范畴未必有无限直和或直积；当每条对角线有限时，有限双积已经足够。例如支持在 \(p\ge a,q\ge b\) 或支持在 \(p\le a,q\le b\) 的双复形均满足此条件。

\begin{lemma}[有限对角双复形的消元]\label{b4:lem:bounded-double-complex}
设 \(\mathcal A\) 是交换范畴，\(E\) 是其中的双复形，并且对每个 \(n\in\mathbb Z\)，对角线 \(p+q=n\) 仅有有限个非零项。若每个纵列 \((E^{p,\bullet},\mathrm{d}_v)\) 正合，则 \(\operatorname{Tot}(E)\) 正合。若改为每个横行 \((E^{\bullet,q},\mathrm{d}_h)\) 正合，同一结论成立。
\end{lemma}
\begin{proof}
先用模的元组记号写出有限消元过程。固定总次数 \(n\)，设 \(z=(z_p)\) 为余圈，其中 \(z_p\in E^{p,n-p}\)。从最小的可能非零横指标 \(p\) 开始。该位置的总余圈等式为 \(\mathrm{d}_vz_p+\mathrm{d}_hz_{p-1}=0\)，而较小分量已为零，所以 \(\mathrm{d}_vz_p=0\)。纵列正合，故可取 \(y_p\in E^{p,n-p-1}\) 使 \(\mathrm{d}_vy_p=z_p\)。从 \(z\) 减去 \(\mathrm{d}y_p\)，第 \(p\) 个分量消失，至多改变第 \(p+1\) 个分量。

继续沿横指标增加的方向操作。只有总次数 \(n\) 和 \(n-1\) 两条对角线中的位置会被使用，它们的非零位置并集有限，所以过程在有限步后停止。若某一步所需的 \(E^{p,n-p-1}\) 为零，纵列正合已迫使当前余圈分量为零，无须选择。最终 \(z=\mathrm{d}(\sum_p y_p)\)，所以每个余圈是余边。

在一般交换范畴中，从任意满足 \(\mathrm dz=0\) 的态射 \(z:T\to\operatorname{Tot}(E)^n\) 开始。纵列正合，使纵微分诱导的态射 \(E^{p,n-p-1}\to Z^{n-p}(E^{p,\bullet})\) 为满态射，因而可用\cref{b4:lem:generalized-element-calculus}在满态射后提升 \(z_p\) 到 \(y_p\)。每进行一步提升，都把 \(z\) 与先前选出的全部 \(y_p\) 一同预合成到新的测试对象上，并仍用原符号表示。步骤有限，这些满态射的复合仍是一个满态射 \(e:T'\to T\)；最终所得 \(y:T'\to\operatorname{Tot}(E)^{n-1}\) 满足
\[
\mathrm dy=ze.
\]
该引理的正合性判据于是证明总复形正合。横纵互换仍满足反交换关系，对角有限性不变，故得到横行版本。
\end{proof}

\begin{corollary}\label{b4:cor:double-complex-comparison}
设 \(\mathcal A\) 是交换范畴，\(u:E\to F\) 是其中的双复形映射，且 \(E,F\) 的每条对角线 \(p+q=n\) 均只有有限个非零项。若 \(u\) 在每个纵列上是拟同构，则 \(\operatorname{Tot}(u)\) 是拟同构。逐横行拟同构也有同样结论。
\end{corollary}
\begin{proof}
对每个纵列取映射锥，组成双复形
\[
G^{p,q}=F^{p,q}\oplus E^{p,q+1},\quad
 \mathrm{d}_v^G=\begin{pmatrix}\mathrm{d}_v^F&u\\0&-\mathrm{d}_v^E\end{pmatrix},\quad
 \mathrm{d}_h^G=\begin{pmatrix}\mathrm{d}_h^F&0\\0&-\mathrm{d}_h^E\end{pmatrix}.
\]
两个微分的反交换来自原反交换关系及 \(\mathrm{d}_h^Fu=u\mathrm{d}_h^E\)；新双复形仍逐对角有限。每个纵列是拟同构的锥，由\cref{b4:prop:cone-quasi-isomorphism}正合。上一引理使 \(\operatorname{Tot}(G)\) 正合，而其微分矩阵正是 \(\operatorname{Cone}(\operatorname{Tot}(u))\) 的矩阵。再次用锥判据即可。横行版本交换两个指标进行同一构造。
\end{proof}

\subsection{无界总复形的差异}
\label{b4:subsec:0604:03}

每条对角线有限是上一个证明真正使用的条件；总次数固定，并不自动使一次消元只走有限步。下面的例子同时说明两种无限总化为什么不同。

\begin{example}\label{b4:exm:sum-product-total-difference}
设 \(k\) 是域。对每个整数 \(p\ge0\)，在 \((p,-p)\) 放一维空间 \(k e_p\)，在 \((p,1-p)\) 放 \(k f_p\)，其中 \(e_p,f_p\) 为所选基向量，其余位置为零。定义
\[
\mathrm{d}_v(e_p)=f_p,\qquad \mathrm{d}_h(e_p)=-f_{p+1},\qquad
 \mathrm{d}_v(f_p)=\mathrm{d}_h(f_p)=0.
\]
每个非零纵列都是同构 \(k e_p\to k f_p\)，因而正合；两个微分的所有二次复合为零，故满足反交换条件。总复形只有次数零和一。用各一维分量的基向量 \(e_p,f_p\) 表示坐标，两种总微分均写成
\[
(\alpha_0,\alpha_1,\alpha_2,\ldots)\longmapsto
(\alpha_0,\alpha_1-\alpha_0,\alpha_2-\alpha_1,\ldots).
\]
直积情形中，对任意目标坐标序列 \((\beta_p)\)，唯一的原像由 \(\alpha_p=\sum_{i=0}^p\beta_i\) 给出，所以微分是同构，总复形正合。

\ProseParagraphBreak
直和情形中，每个序列须有有限支集。微分仍单，但像恰好是坐标总和为零的有限支集序列：像的坐标和通过相邻相消为零；反过来，若 \(\sum_i\beta_i=0\)，上述部分和 \(\alpha_p\) 最终为零，确实属于直和。坐标求和映射到 \(k\) 是满射，其核正是微分的像，因此
\[
H^0\operatorname{Tot}_{\oplus}(E)=0,\qquad
H^1\operatorname{Tot}_{\oplus}(E)\cong k,
\]
而直积总复形的上同调全零。一个纵列处处正合的双复形，在直和总化后仍可能出现非零上同调。
\end{example}

解差分方程得到的部分和未必有有限支集。例如次数一的基向量 \(f_0\) 在直积中有原像 \((1,1,1,\ldots)\)，在直和中却没有原像。逐列消元所需的无限多个分量因此未必能组成直和总复形中的一个元素。

\subsection{总微分与总化的符号}
\label{b4:subsec:0604:04}

为便于后续计算，把两种起点区别如下。若横纵微分已经反交换，总微分就是 \(\mathrm{d}_h+\mathrm{d}_v\)；若原方格交换，则应先把纵微分改为 \((-1)^p\partial_v\)，总微分才是 \(\partial_h+(-1)^p\partial_v\)。这两句话描述同一个构造，不能把后一个符号重复施加到前一种约定上。

\begin{example}\label{b4:exm:double-square-sign}
设 \(R\) 交换，\(a,b\in R\)。在四个位置 \((0,0),(1,0),(0,1),(1,1)\) 都放 \(R\)，原横箭头均为乘 \(a\)，原纵箭头均为乘 \(b\)。它们交换。把第 \(p\) 列的纵箭头乘 \((-1)^p\)，并把总次数一按 \((1,0),(0,1)\) 排列，总复形成为
\[
R\xrightarrow{t\mapsto(at,bt)}R^2
\xrightarrow{(u,v)\mapsto-bu+av}R.
\]
复合为 \(-bat+abt=0\)。若漏掉负号，复合一般为 \(2abt\)，所以仅知道原方格交换不足以直接相加。该复形是两个两项复形的张量总化，也是后面两元 Koszul 计算中会反复遇到的形式。
\end{example}

当两套消解形成的双复形逐对角有限时，逐列拟同构的映射锥也逐对角有限，有限消元便把每列的正合性传到总复形，从而得到总复形的拟同构。无限对角线上的比较则还取决于总化方式，以及所需的原像能否组成该总复形中的态射。


\begin{chaptersummary}
同伦把两个复形映射的差写成 \(\mathrm{d}h+h\mathrm{d}\)，因此它们在同调上诱导相同映射。可缩意味着恒等映射零伦，比同调消失更强；对零调复形而言，所缺的条件是每个次数的余圈短正合列都分裂。映射锥把核、余核与连接态射的比较合成一条上同调长正合列，拟同构正好对应锥零调。映射柱则把任意映射分成逐次分裂单态射和同伦等价。张量与 Hom 复形在两个变量都保持链同伦，因而保持同伦等价；保持拟同构则不能只靠这些同伦等式。

\ProseParagraphBreak
双复形的总微分依赖横纵反交换的约定。当每条对角线有限时，逐列或逐行正合可通过有限消元传到总复形；同样的论证给出逐列拟同构的总化比较。无界时，提升得到的序列可能没有有限支集，因而直和与直积总化会产生不同同调。这些区别将在消解、导出函子以及后面的谱序列中继续起作用。
\end{chaptersummary}

\BookExercises

\begin{exercise}\label{b4:ex:06-scalar-homotopy}
固定非零整数 \(m\)，令 \(P=(\mathbb Z\xrightarrow{m}\mathbb Z)\) 位于次数 \(-1,0\)。对整数 \(a,b\)，逐次乘 \(a,b\) 得到复形自映射。证明它们同伦当且仅当 \(m\mid(a-b)\)。
\end{exercise}
\begin{solution}
唯一可能非零的同伦分量是 \(h^0:P^0\to P^{-1}\)，设其为乘 \(t\)。两个次数的同伦公式分别给出 \(a-b=tm\) 与 \(a-b=mt\)，是同一条件。故存在同伦恰好等价于整除条件。
\end{solution}

\begin{exercise}\label{b4:ex:06-homology-zero-not-null}
令 \(P=(\mathbb Z\xrightarrow{2}\mathbb Z)\) 位于次数 \(-1,0\)，\(D=\mathbb Z[0][1]\)。定义 \(f:P\to D\) 的 \(-1\) 次分量为恒等，零次分量为零。证明所有 \(H^n(f)\) 为零，但 \(f\) 不零伦。
\end{exercise}
\begin{solution}
按移位约定，\(D=\mathbb Z[0][1]\) 只有次数 \(-1\) 的项 \(\mathbb Z\)，其微分为零。于是复形映射等式成立。源只有 \(H^0=\mathbb Z/2\mathbb Z\)，目标只有 \(H^{-1}=\mathbb Z\)，所以每阶诱导映射都为零。若 \(f=\mathrm{d}h+h\mathrm{d}\)，唯一可能的分量 \(h^0:\mathbb Z\to\mathbb Z\) 为乘 \(t\)，在次数 \(-1\) 的等式要求 \(1=2t\)，不可能。因此同调相同的两个映射未必同伦。
\end{solution}

\begin{exercise}\label{b4:ex:06-tensor-homotopy}
设 \(R\) 为含幺环，\(C,C'\) 是幺性右 \(R\)-模的上链复形，\(D,D'\) 是幺性左 \(R\)-模的上链复形。设 \(f,g:C\to C'\) 由 \(h\) 给出同伦，\(u,v:D\to D'\) 由 \(k\) 给出同伦。在直和总化的张量复形上，写出 \(f\otimes u\) 与 \(g\otimes v\) 之间的同伦。再取 \(R=\mathbb Z\)、\(C=C'=(\mathbb Z\xrightarrow1\mathbb Z)\)、\(D=D'=(\mathbb Z\xrightarrow1\mathbb Z)\)，两个复形均位于次数 \(0,1\)，取 \(f=g=1_C,u=1_D,v=0\)。计算所得收缩，并检验省去第二因子同伦的次数符号会发生什么。
\end{exercise}
\begin{solution}
把差分成
\[
f\otimes u-g\otimes v=(f-g)\otimes u+g\otimes(u-v).
\]
由\cref{b4:prop:tensor-hom-preserve-homotopy}的两个同伦前后复合相应复形映射，得到总次数为 \(-1\) 的同伦
\[
H(x\otimes y)=h(x)\otimes u(y)+(-1)^p g(x)\otimes k(y)
\quad(x\in C^p).
\]
两个同伦等式相加即得 \(\mathrm{d}H+H\mathrm{d}=f\otimes u-g\otimes v\)。

在所给例子中取 \(h=0\)，取 \(k:D^1\to D^0\) 为恒等，其余分量为零。设 \(e_0,e_1\) 与 \(v_0,v_1\) 分别是两复形的标准基。则非零的收缩值为
\[
H(e_0\otimes v_1)=e_0\otimes v_0,\qquad
H(e_1\otimes v_1)=-e_1\otimes v_0.
\]
它满足 \(\mathrm{d}H+H\mathrm{d}=1\)，所以张量复形可缩。若省去 \((-1)^p\)，得到 \(H_0(x\otimes y)=x\otimes k(y)\)，则
\[
(\mathrm{d}H_0+H_0\mathrm{d})(e_0\otimes v_1)
=e_0\otimes v_1+2e_1\otimes v_0,
\]
不等于 \(e_0\otimes v_1\)。第二因子的降次映射越过第一因子时，次数符号保证这些混合项抵消。
\end{solution}

\begin{exercise}\label{b4:ex:06-additive-contractibility}
证明加性函子逐次作用于复形时保持可缩性。对交换群上的函子
\[
F=\operatorname{Hom}_{\mathbb Z}(\mathbb Z/2\mathbb Z,-),
\]
举例说明它未必保持正合复形。
\end{exercise}
\begin{solution}
由 \(1_C=\mathrm{d}h+h\mathrm{d}\)，加性和复合保持性给出 \(1_{F(C)}=F(\mathrm{d})F(h)+F(h)F(\mathrm{d})\)，所以 \(F(h)\) 是收缩同伦。对正合复形 \(\mathbb Z\xrightarrow2\mathbb Z\to\mathbb Z/2\mathbb Z\)（次数 \(0,1,2\)）应用 \(F\)，得到次数 \(0,1,2\) 的复形
\[
0\longrightarrow0\longrightarrow\mathbb Z/2\mathbb Z,
\]
所以 \(H^2(F(C))=\mathbb Z/2\mathbb Z\ne0\)。
\end{solution}

\begin{exercise}\label{b4:ex:06-no-natural-contraction}
设 \(k\) 是域。把 \(0\to k\xrightarrow{x\mapsto(x,0)}k^2\xrightarrow{(x,y)\mapsto y}k\to0\) 放在次数 \(0,1,2\)。证明此复形可缩，但不存在与其所有复形自同构都交换的收缩同伦。
\end{exercise}
\begin{solution}
短列分裂，故可缩。若 \(h\) 是收缩同伦，次数二的等式要求 \(h^2(t)=(at,t)\)，其中 \(a\in k\)。对任意 \(b\in k\)，在中项用剪切 \(u_b(x,y)=(x+by,y)\)，两端用恒等，得到复形自同构。若 \(h\) 与它交换，就需 \(u_bh^2(t)=h^2(t)\)，即 \((a+b)t=at\)。取 \(b=1,t=1\) 矛盾。因此，每个选定的分裂都能给出收缩，却没有一个收缩对所有这些自同构保持不变；收缩存在不意味着可以自然地选取。
\end{solution}

\begin{exercise}\label{b4:ex:06-cone-euler}
设 \(C,D\) 是有限维 \(k\)-向量空间的有界复形，\(f:C\to D\)。证明 \(\chi(\operatorname{Cone}(f))=\chi(D)-\chi(C)\)。若右端为零，能否断言 \(f\) 是拟同构？
\end{exercise}
\begin{solution}
锥的次数 \(n\) 为 \(D^n\oplus C^{n+1}\)，交错求和即得公式。不能反推拟同构：取 \(C=D=k[0]\)、\(f=0\)，两者 Euler 特征相等，但 \(H^0(f)=0:k\to k\) 不可逆，锥在 \(-1,0\) 两次各有一个 \(k\) 同调。
\end{solution}

\begin{exercise}\label{b4:ex:06-cone-residue}
令 \(R=\mathbb Z/12\mathbb Z\)，把乘 \(4\) 看成集中于零次的 \(R\)-模复形之间的映射。计算其锥的同调，并判断该映射是否为拟同构。
\end{exercise}
\begin{solution}
锥为次数 \(-1,0\) 的 \(R\xrightarrow4R\)。核为 \(\{0,3,6,9\}\cong\mathbb Z/4\mathbb Z\)，余核为 \(R/4R\cong\mathbb Z/4\mathbb Z\)，分别是 \(H^{-1}\)、\(H^0\)。锥不正合，故乘四不是拟同构。
\end{solution}

\begin{exercise}\label{b4:ex:06-cone-comparison-integers}
令 \(P=(\mathbb Z\xrightarrow2\mathbb Z)\)、\(Q=(\mathbb Z\xrightarrow6\mathbb Z)\) 均在次数 \(-1,0\)。取 \(f^{-1}=1\)、\(f^0=3\)。写出映射锥并直接计算同调，再与锥长正合列核对。
\end{exercise}
\begin{solution}
\(6\cdot1=3\cdot2\)，故 \(f\) 是复形映射。锥在次数 \(-2,-1,0\) 为
\[
\mathbb Z\xrightarrow{t\mapsto(t,-2t)}\mathbb Z^2
\xrightarrow{(u,v)\mapsto6u+3v}\mathbb Z.
\]
第一映射单，第二映射的核恰为第一映射的像，余核为 \(\mathbb Z/3\mathbb Z\)。另一方面，\(H^0(f):\mathbb Z/2\mathbb Z\to\mathbb Z/6\mathbb Z\) 将一送到三，是单射且余核 \(\mathbb Z/3\mathbb Z\)；锥长正合列给出相同结果。
\end{solution}

\begin{exercise}\label{b4:ex:06-cylinder-homotopies}
设 \(f,g:C\to D\) 为模复形映射。证明 \(f\simeq g\) 当且仅当存在复形映射 \(F:\operatorname{Cyl}(1_C)\to D\)，使 \(F(0,c,0)=f(c)\)、\(F(y,0,0)=g(y)\)。写出 \(F\) 与同伦的关系。
\end{exercise}
\begin{solution}
在 \(\operatorname{Cyl}(1_C)^n=C^n\oplus C^n\oplus C^{n+1}\) 中，第一份 \(C^n\) 是映射 \(1_C\) 的目标，第二份是它的源，第三份属于移位复形。任何这样的分次映射都形如 \(F(y,c,a)=g(y)+f(c)+t(a)\)，其中 \(t\) 降一次。将映射柱的微分代入，与微分交换的额外条件为 \(\mathrm{d}_Dt+t\mathrm{d}_C=g-f\)。因此若 \(h\) 满足 \(f-g=\mathrm{d}_Dh+h\mathrm{d}_C\)，取 \(t=-h\)；反之令 \(h=-t\) 即得同伦。这里第三分量把两个端部映射之间的同伦直接保存下来。
\end{solution}

\begin{exercise}\label{b4:ex:06-tensor-breaks-quasi}
对\cref{b4:exm:resolution-quasi-not-homotopy}中的拟同构 \(P\to(\mathbb Z/2\mathbb Z)[0]\) 张量 \((\mathbb Z/2\mathbb Z)[0]\)。证明所得映射不是拟同构，并解释这与张量保持同伦等价并不冲突。
\end{exercise}
\begin{solution}
\(P\otimes\mathbb Z/2\mathbb Z\) 是次数 \(-1,0\) 的 \(\mathbb Z/2\mathbb Z\xrightarrow0\mathbb Z/2\mathbb Z\)，所以负一次同调为 \(\mathbb Z/2\mathbb Z\)。目标张量积仍集中于零次，负一次同调为零，故不是拟同构。这里 \(\mathbb Z/2\mathbb Z\) 不是平坦的 \(\mathbb Z\)-模：对单射 \(\mathbb Z\xrightarrow2\mathbb Z\) 张量后得到零映射 \(\mathbb Z/2\mathbb Z\to\mathbb Z/2\mathbb Z\)，单射性丢失，产生了额外的负一次同调。原映射没有同伦逆；张量保持给定的同伦等式，并不保证保持任意拟同构。
\end{solution}

\begin{exercise}\label{b4:ex:06-homotopy-not-isomorphism}
给出域 \(k\) 上同伦等价但不逐次同构的两个复形。它们可以有不同的总维数吗？
\end{exercise}
\begin{solution}
取 \(C=(k\xrightarrow1k)\) 和零复形。\(C\) 可缩，所以两者之间的零映射互为同伦逆，但 \(C\) 有两个非零项，无法逐次同构于零复形。总维数分别为二和零。在同伦范畴中，\(C\) 与零对象同构；这种同构只要求复合等于恒等的同伦类，不要求每个次数的对象同构。同伦等价保持同调；对各项均有限维的有界 \(k\)-向量空间复形，它也保持 Euler 特征，但不保持各项维数之和。
\end{solution}

\begin{exercise}\label{b4:ex:06-coprime-square-contraction}
将\cref{b4:exm:double-square-sign}取 \(R=\mathbb Z,a=2,b=3\)。证明总复形可缩，并写出一个收缩同伦。
\end{exercise}
\begin{solution}
复形为 \(\mathbb Z\xrightarrow{(2,3)}\mathbb Z^2\xrightarrow{(-3,2)}\mathbb Z\)，次数 \(0,1,2\)。取 \(h^1(u,v)=-u+v\)、\(h^2(t)=(-t,-t)\)。则 \(h^1\mathrm{d}^0=1\)、\(\mathrm{d}^1h^2=1\)，而
\[
\mathrm{d}^0h^1(u,v)+h^2\mathrm{d}^1(u,v)
=(2(-u+v),3(-u+v))+(3u-2v,3u-2v)=(u,v).
\]
所以 \(\mathrm{d}h+h\mathrm{d}=1\)。互素性允许这里出现整数线性组合等于一。
\end{solution}

\begin{exercise}\label{b4:ex:06-first-quadrant-edge}
设交换范畴中的双复形 \(E\) 支持在 \(p,q\ge0\)，且各纵列只可能在 \(q=0\) 有同调。令 \(Z^{p,0}=\ker(\mathrm{d}_v:E^{p,0}\to E^{p,1})\)。证明这些对象随横微分组成复形，并证明自然嵌入 \(Z^{\bullet,0}\to\operatorname{Tot}(E)\) 是拟同构。
\end{exercise}
\begin{solution}
由反交换，\(\mathrm{d}_v\mathrm{d}_hz=-\mathrm{d}_h\mathrm{d}_vz=0\)，故横微分保持 \(Z^{p,0}\)。把它作为仅有 \(q=0\) 一行的双复形。嵌入在每个纵列上诱导同调同构：零次由 \(E^{p,-1}=0\) 得 \(H^0_v=Z^{p,0}\)，其他次数由假设为零。当总次数 \(n\ge0\) 时，\(p+q=n\) 且 \(p,q\ge0\) 迫使 \(0\le p\le n\)；当 \(n<0\) 时没有非零项。因此第一象限保证逐对角有限，应用\cref{b4:cor:double-complex-comparison}即可。
\end{solution}

\begin{exercise}\label{b4:ex:06-third-quadrant-edge}
设交换范畴中的双复形 \(E\) 支持在 \(p,q\le0\)，且各横行只在 \(p=0\) 可能有同调。证明横同调 \(H_h^0(E^{\bullet,q})\) 随纵微分组成复形，并且从总复形到这一复形的自然投影是拟同构。
\end{exercise}
\begin{solution}
因为 \(E^{1,q}=0\)，零次横同调为 \(E^{0,q}/\operatorname{im}\mathrm{d}_h\)。反交换保证 \(\mathrm{d}_v\) 将横边界送到横边界，所以诱导商上的微分。把这些商置于 \(p=0\) 一列，其他列为零；从 \(E\) 到它的商映射在每个横行上是拟同构。当总次数 \(n\le0\) 时，\(p+q=n\) 且 \(p,q\le0\) 迫使 \(n\le p\le0\)；当 \(n>0\) 时没有非零项。因此第三象限同样逐对角有限，由横行版本的总化比较推论得到结论。
\end{solution}

\begin{exercise}\label{b4:ex:06-infinite-diagonal-size}
设 \(k\) 是域，双复形在每个 \((p,-p)\)、\(p\in\mathbb Z\) 放 \(k\)，其余位置为零，所有微分为零。比较两种总复形及自然映射 \(\operatorname{Tot}_{\oplus}\to\operatorname{Tot}_{\Pi}\)，并指出其为何不是拟同构。
\end{exercise}
\begin{solution}
两者都集中于总次数零，分别为 \(\bigoplus_{p\in\mathbb Z}k\) 与 \(\prod_{p\in\mathbb Z}k\)。自然映射是有限支集序列的包含，不满，例如全一序列不在像中。因此零次同调映射不满，不能成为拟同构。即使没有微分，总化所选的积或余积也会改变对象。
\end{solution}

\begin{exercise}\label{b4:ex:06-hom-homotopies}
固定非零整数 \(m\)，令交换群复形 \(P=(\mathbb Z\xrightarrow{m}\mathbb Z)\) 位于次数 \(-1,0\)，令 \(D=\mathbb Z[0]\)。写出 \(\operatorname{Hom}_{\mathbb Z}^\bullet(P,D)\) 的非零项及微分，计算每个整数 \(r\) 的 \(H^r\operatorname{Hom}_{\mathbb Z}^\bullet(P,D)\)。再用\cref{b4:prop:tensor-hom-preserve-homotopy}识别 \(\operatorname{Hom}_{K(\mathbf{Ab})}(P,D[r])\)：在 \(r=1\) 时用具体复形映射代表同伦类，并说明它何时零伦，核对 Hom 余边与移位同伦之间的符号。
\end{exercise}
\begin{solution}
Hom 复形仅有次数零和一非零，分别来自 \(P^0\to D^0\) 与 \(P^{-1}\to D^0\)。两项均为 \(\mathbb Z\)，而次数零的微分为 \(\mathrm{d}_D\beta-\beta\mathrm{d}_P\)，所以复形为
\[
\mathbb Z\xrightarrow{-m}\mathbb Z
\quad\text{，位于次数 }0,1.
\]
因此零次同调为零，一次同调为 \(\mathbb Z/m\mathbb Z\)，其余次数的同调全零。命题给出的同构说明这些分别也是 \(P\to D[r]\) 的同伦类群。

当 \(r=1\) 时，目标 \(D[1]\) 只有次数 \(-1\) 非零。对整数 \(a\)，令 \(\alpha_a^{-1}:P^{-1}\to D[1]^{-1}\) 为乘 \(a\)，其他分量为零；这给出全部复形映射。两个映射 \(\alpha_a,\alpha_b\) 同伦，当且仅当存在乘 \(t\) 的同伦分量 \(h^0:P^0\to D[1]^{-1}\)，使 \(a-b=tm\)。所以类由 \(a\bmod m\) 决定，且 \(\alpha_a\) 零伦当且仅当 \(m\mid a\)。若 \(a=mt\)，对应的次数零 Hom 元素 \(\beta\) 为乘 \(-t\)，于是 \(\mathrm{d}_{\mathrm{Hom}}\beta=a\)，且 \(h=(-1)^1\beta\) 为乘 \(t\)，与上述零伦同伦一致。
\end{solution}

% !TeX root = ../../Book4.tex
% 纲要标题：## 第8章　投射消解与内射消解
% 纲要标签：b4:ch:07
% 纲要位置：计划纲要.md，第 3372 行。

\chapter{投射消解与内射消解}
\label{b4:ch:07}
\BookChapterNumber{b4:ch:07}

\begin{chapterabstract}
本章构造投射与内射消解，证明比较定理和马蹄引理，说明消解的提升与唯一性应在同伦意义下理解，并给出短消解、Koszul 消解和 bar 消解的模型。
\end{chapterabstract}

\begin{learninggoals}
\begin{itemize}
\item 区分增广复形与去掉增广项的消解复形，说明投射和内射的次数方向。
\item 从核或余核逐次构造消解，并指出一般交换范畴所需的存在性假设。
\item 构造比较映射及两个提升之间的同伦，解释消解独立性。
\item 用马蹄引理构造相容消解，辨别逐次分裂与复形分裂。
\item 识别短消解的适用条件，并核验 Koszul 与 bar 微分的正合性来源。
\end{itemize}
\end{learninggoals}

设整数 \(n\ge2\)。映射 \(\mathbb Z\xrightarrow{n}\mathbb Z\twoheadrightarrow\mathbb Z/n\mathbb Z\) 把商模的生成元和关系都留在一条正合列里。给定从中间的自由模 \(\mathbb Z\) 到另一个模 \(M\) 的同态，它何时能经过商模 \(\mathbb Z/n\mathbb Z\)？答案取决于前一项的像是否被送到零。从自由模出发可以先选择生成元的像，再由前一项检查关系。对一般模，关系组成的核还可能有自己的关系，因此须继续取自由满射，逐层记录。自由模之所以适合这种构造，不仅因为可以指定基元的像，还因为从它出发的态射能沿满射提升。后面的比较定理反复使用的是这一提升性质，不需要依赖某组基；因此将自由模推广为具备同样性质的投射对象，就得到投射消解。接下来还须证明：换一种选择得到的消解，仍会给出相同的同调计算。

\ProseParagraphBreak
本章使用\chapref{b4:ch:04}的交换范畴与蛇引理、\chapref{b4:ch:05}的同调长正合列，以及\chapref{b4:ch:06}的链同伦。模论中的投射、内射及其存在性可复习第二卷第7.1—7.2节；其中定理7.2.8已完整证明任意环的模范畴有足够多内射对象。抽象交换范畴的结论总要另外声明足够多投射或内射对象，不能只由交换范畴公理推得。

\ProseParagraphBreak
Weibel 的《An Introduction to Homological Algebra》\cite[Sections 2.2--2.3, Comparison Theorem 2.2.6, Horseshoe Lemma 2.2.8, Comparison Theorem 2.3.7]{Weibel1994HomologicalAlgebra}给出本章前三节的基本构造。该书使用链指标；本章为展示投射消解保留非负链指标 \(P_n\)，同时以 \(P^{-n}=P_n\) 明确它与本书上链约定的转换。

% !TeX root = ../../Book4.tex
% 纲要标题：### 8.1 消解的构造
% 纲要标签：b4:sec:0701
% 纲要位置：计划纲要.md，第 3374 行。

\section{消解的构造}
\label{b4:sec:0701}

生成元给出自由满射，关系组成其核；继续为核取生成元，便得到消解。自由模可以沿满射提升映射，后面的比较定理只需要这一性质，因而可以把自由模放宽为投射对象。相反方向的构造使用内射对象的延拓性质；这两个方向都依赖范畴中相应对象是否足够多。


% 纲要内容（仅作写作计划，尚非教材正文）：
%
% * 模的自由及投射消解
% * 内射消解
% * 足够多投射或内射对象的假设
%

\subsection{模的自由及投射消解}
\label{b4:subsec:0701:01}

一个模的生成元给出自由满射，生成元之间的关系则成为这个满射的核。若核还不是容易处理的模，就继续为核选择生成元。消解把这项重复操作写成一个正合复形；它所保留的不是某一组最优生成元，而是所有关系都被逐层记录这一事实。

\begin{definition}\label{b4:def:projective-resolution}
设 \(\mathcal A\) 是交换范畴，\(P\) 为其对象。若对每个满态射 \(e:X\to Y\) 及态射 \(f:P\to Y\)，都存在态射 \(\widetilde f:P\to X\)，使 \(e\widetilde f=f\)，则称 \(P\) 为\term[tousheduixiang]{投射对象}[projective object]。给定对象 \(M\in\mathcal A\) 及增广链复形
\[
\cdots\longrightarrow P_2\xrightarrow{\mathrm{d}_2}P_1
\xrightarrow{\mathrm{d}_1}P_0\xrightarrow{\varepsilon}M\longrightarrow0
\]
若此序列正合，且每个 \(P_n\)（\(n\ge0\)）都投射，则称它为 \(M\) 的\term[toushexiaojie]{投射消解}[projective resolution]。若 \(R\) 为含幺环、\(\mathcal A=R\text{-}\mathbf{Mod}\) 为幺性左 \(R\) 模范畴，且各 \(P_n\) 自由，则称为\term[ziyouxiaojie]{自由消解}[free resolution]。
\end{definition}

正合性说明这条序列确实记录了 \(M\) 及各层关系；各项的投射性则保证从它出发的映射可以沿满态射逐层提升。后一条件将在\secref{b4:sec:0702}的比较定理中使用，不是正合性的另一个说法。

\ProseParagraphBreak
增广项 \(M\) 不属于非负次数的复形 \(P_\bullet\)：它的零次项是 \(P_0\)，从 \(P_0\) 向更低次数的微分为零。由增广诱导的商同构是
\[
H_0(P_\bullet)=P_0/\operatorname{im}\mathrm{d}_1
=P_0/\ker\varepsilon\xrightarrow{\;\overline\varepsilon\;}M.
\]
它是同构，因为 \(\varepsilon\) 满，且核正好是 \(\operatorname{im}\mathrm{d}_1\)；正次数同调则由正合性全部为零。写成上链复形时约定 \(P^{-n}=P_n\)，并在正次数补零，得到
\[
P^\bullet:\quad\cdots\longrightarrow P^{-2}\longrightarrow P^{-1}
\longrightarrow P^0\longrightarrow0.
\]
与它比较的是 \(M[0]\)，其中 \(M[0]^0=M\)，其余次数为零。增广给出复形映射 \(\varepsilon:P^\bullet\to M[0]\)：零次分量为原来的 \(\varepsilon\)，其余分量为零；\(\varepsilon\mathrm{d}^{-1}=0\) 保证复形映射条件。由刚才的同调计算，它是拟同构。

\begin{proposition}\label{b4:prop:free-resolution-exists}
设 \(R\) 是含幺环。每个幺性左 \(R\) 模都有自由消解，因而也有投射消解。每个幺性右 \(R\) 模也有相应的右模消解。
\end{proposition}
\begin{proof}
设 \(M\) 是幺性左 \(R\) 模，令 \(K_0=M\)。以 \(K_0\) 的底集合为基取自由模 \(P_0\)，把基向量 \([m]\) 送到 \(m\)，得到满射 \(p_0:P_0\to K_0\)。归纳地令 \(K_{n+1}=\ker p_n\)，取以 \(K_{n+1}\) 的底集合为基的自由模 \(P_{n+1}\) 及满射 \(p_{n+1}:P_{n+1}\to K_{n+1}\)。这里 \(K_1\) 是原生成元之间的关系模，\(K_2\) 是为 \(K_1\) 选择的生成元之间的关系模，后续各 \(K_n\) 以同样方式记录上一层的关系。

取增广 \(\varepsilon=p_0\)，对 \(n\ge1\) 定义
\[
\mathrm{d}_n:P_n\xrightarrow{p_n}K_n\hookrightarrow P_{n-1}.
\]
后一个箭头是单射，所以 \(\ker\mathrm{d}_n=\ker p_n=K_{n+1}\)；而 \(p_{n+1}\) 满射，所以 \(\operatorname{im}\mathrm{d}_{n+1}=K_{n+1}\)。因此每个正次数上像等于核，零次上也有 \(\operatorname{im}\mathrm{d}_1=K_1=\ker\varepsilon\)，增广满射，所构造的序列正合。

各项自由还保证它们投射。具体地，给定满射 \(e:X\to Y\) 及线性映射 \(f:R^{(S)}\to Y\)，按本卷的集合指标选择约定，为每个标准基向量 \(b_s\) 选择 \(x_s\in X\)，使 \(e(x_s)=f(b_s)\)。自由模的泛性质把 \(b_s\mapsto x_s\) 唯一延拓为线性映射 \(\widetilde f:R^{(S)}\to X\)；\(e\widetilde f\) 与 \(f\) 在基上相同，故 \(e\widetilde f=f\)。选择用在这一族原像上。对反环 \(R^{\mathrm{op}}\) 使用同一构造便得到右模版本。
\end{proof}

这个构造通常很大。对 \(\mathbb Z/m\mathbb Z\)，其中 \(m>1\)，较合适的消解是 \(0\to\mathbb Z\xrightarrow{\times m}\mathbb Z\to\mathbb Z/m\mathbb Z\to0\)。设 \(k\) 是域。对 \(R=k[\epsilon]/(\epsilon^2)\) 上的 \(k=R/(\epsilon)\)，则有无限消解
\[
\cdots\xrightarrow{\epsilon}R\xrightarrow{\epsilon}R
\xrightarrow{\epsilon}R\longrightarrow k\longrightarrow0.
\]
每个元素唯一写成 \(a+b\epsilon\)，其中 \(a,b\in k\)，并且 \(\epsilon(a+b\epsilon)=a\epsilon\)。因此乘 \(\epsilon\) 的核是所有 \(b\epsilon\)，像也是所有 \(a\epsilon\)，二者均为 \((\epsilon)\)；增广的核同样为 \((\epsilon)\)，所以该序列正合。这个有限生成模的消解仍有无限多个非零项，有限生成本身不使逐层取关系的过程停下来。

\subsection{内射消解}
\label{b4:subsec:0701:02}

投射性约束映射的出发对象：从它出发的映射能够沿满态射提升。内射性则约束映射的到达对象：映入它的映射能够沿单态射延拓。因而内射消解从嵌入原对象开始，逐层处理余核。

\begin{definition}\label{b4:def:injective-resolution}
设 \(\mathcal A\) 是交换范畴，\(I\) 为其对象。若对每个单态射 \(e:X\to Y\) 及态射 \(f:X\to I\)，都存在态射 \(\widetilde f:Y\to I\)，使 \(\widetilde f e=f\)，则称 \(I\) 为\term[neisheduixiang]{内射对象}[injective object]。给定对象 \(N\in\mathcal A\) 及增广上链复形
\[
0\longrightarrow N\xrightarrow{\eta}I^0\xrightarrow{\mathrm{d}^0}I^1
\xrightarrow{\mathrm{d}^1}I^2\longrightarrow\cdots
\]
若此序列正合且各 \(I^n\)（\(n\ge0\)）内射，则称它为 \(N\) 的\term[neishexiaojie]{内射消解}[injective resolution]。
\end{definition}

这里的增广项 \(N\) 不属于 \(I^\bullet\)；后者位于非负上链次数，增广 \(\eta\) 诱导 \(N\cong\ker\mathrm{d}^0=H^0(I^\bullet)\)，正次数上同调为零。构造时，先选到内射对象的单态射 \(j^0:N\hookrightarrow I^0\)，令 \(C^0=\operatorname{coker}j^0\)，再选到内射对象的单态射 \(j^1:C^0\hookrightarrow I^1\)。其余核记为 \(C^1\)，如此继续，并令
\[
\mathrm{d}^n:I^n\xrightarrow{q^n}C^n\xrightarrow{j^{n+1}}I^{n+1},
\qquad C^n=\operatorname{coker}j^n,
\]
其中 \(n\ge0\)，\(q^n\) 是余核映射，各 \(I^n\) 都内射。因为 \(j^{n+1}\) 单，\(\ker\mathrm{d}^n=\ker q^n=\operatorname{im}j^n\)；对 \(n\ge1\)，因为前一步的 \(q^{n-1}\) 满，\(\operatorname{im}\mathrm{d}^{n-1}=\operatorname{im}j^n\)。零次的核是 \(\operatorname{im}j^0\)，所以这条增广序列正合。

\ProseParagraphBreak
% 跨卷引用目标：b2:thm:enough-injectives（7.2.8）。
对任意含幺环 \(R\)，第二卷定理7.2.8已经证明每个幺性左 \(R\) 模都能嵌入内射模。这个构造先用内射交换群 \(\mathbb Q/\mathbb Z\) 分离非零元素：对 \(0\ne n\in N\)，总有加法群同态 \(\chi:N\to\mathbb Q/\mathbb Z\) 使 \(\chi(n)\ne0\)。再通过余诱导把交换群中的延拓变为 \(R\) 模中的延拓，得到内射左模
\[
J=\operatorname{Hom}_{\mathbb Z}(R,\mathbb Q/\mathbb Z),
\qquad (rf)(s)=f(sr).
\]
每个 \(\chi\) 对应左 \(R\) 线性映射 \(n\mapsto(s\mapsto\chi(sn))\)。单个映射未必分离所有元素，把全部这样的映射组成积映射，便得到嵌入
\[
N\longrightarrow\prod_{\chi\in\operatorname{Hom}_{\mathbb Z}(N,\mathbb Q/\mathbb Z)}J,
\qquad n\longmapsto\bigl(s\longmapsto\chi(sn)\bigr)_\chi.
\]
若 \(n\ne0\)，取分离它的 \(\chi\)，在 \(s=1\) 处这一分量非零，所以积映射单射。各 \(J\) 的内射性和内射模积的内射性都已在第二卷证明。本章把这个结果反复用于余核模，就得到内射消解；幺性右模的情形对反环应用同一结果。

\begin{example}\label{b4:exm:integer-injective-resolution}
作为 \(\mathbb Z\) 模，\(\mathbb Z\) 有内射消解
\[
0\longrightarrow\mathbb Z\longrightarrow\mathbb Q
\longrightarrow\mathbb Q/\mathbb Z\longrightarrow0.
\]
% 跨卷引用目标：b2:thm:divisible-injective（7.2.5）。
正合性由商群定义给出。\(\mathbb Q\) 中每个元素都能除以任意非零整数；在 \(\mathbb Q/\mathbb Z\) 中，对 \(q+\mathbb Z\) 和整数 \(m\ne0\)，有 \(m(q/m+\mathbb Z)=q+\mathbb Z\)。所以两项都可除，由第二卷定理7.2.5的可除群判据，它们都是内射模。

它与 \(\mathbb Z\) 自身的长度零投射消解适用于 Hom 的不同变量：对投射模 \(P\)，\(\operatorname{Hom}_{\mathbb Z}(P,-)\) 正合；对内射模 \(I\)，\(\operatorname{Hom}_{\mathbb Z}(-,I)\) 正合。对一般含幺环 \(R\) 及幺性左模 \(M,N\)，\secref{b4:sec:0802}将用第一变量的投射消解或第二变量的内射消解计算 \(\operatorname{Ext}_R^n(M,N)\)，并证明两种计算相同。
\end{example}

\subsection{足够多投射或内射对象的假设}
\label{b4:subsec:0701:03}

\begin{definition}\label{b4:def:enough-projectives-injectives}
设 \(\mathcal A\) 是交换范畴。若对每个对象 \(M\) 都存在投射对象 \(P\) 及满态射 \(P\twoheadrightarrow M\)，就说 \(\mathcal A\) 有\term[zugouduo-tousheduixiang]{足够多投射对象}[enough projectives]；若对每个对象 \(M\) 都存在内射对象 \(I\) 及单态射 \(M\hookrightarrow I\)，就说它有\term[zugouduo-neisheduixiang]{足够多内射对象}[enough injectives]。
\end{definition}

在前一条件下，对核反复选择投射满态射，就得到投射消解；在后一条件下，对余核反复选择内射单态射，就得到内射消解。这是\cref{b4:prop:free-resolution-exists}的范畴形式，正合性仍由每一步的核或余核保证。为每个对象选定一条消解，并不会自动得到函子：对象间的态射还需要提升到所选消解之间，提升也未必唯一。下一节的比较定理将证明提升存在且在链同伦意义下唯一；施加加性函子后再取同调，链同伦的提升给出相同的态射。这正是后续同调不变量能够成为函子的原因。

\ProseParagraphBreak
% 跨卷引用目标：b1:thm:fga-structure（9.3.1）。
交换范畴本身的公理没有包括这两项条件。例如有限交换群组成交换范畴，却没有非零投射对象。事实上，按第一卷定理9.3.1的有限生成交换群结构定理，非零有限交换群可分解为
\[
P\cong\bigoplus_i\mathbb Z/p_i^{a_i}\mathbb Z,
\qquad a_i\ge1,
\]
其中 \(p_i\) 为素数。假如 \(P\) 投射，任取一个非零直和因子 \(T=\mathbb Z/p^a\mathbb Z\)。对从 \(T\) 出发的映射，先与 \(P\to T\) 的投影复合，用 \(P\) 的投射性提升，再限制到 \(T\)，故 \(T\) 也投射。于是自然满射
\[
\mathbb Z/p^{a+1}\mathbb Z\longrightarrow\mathbb Z/p^a\mathbb Z
\]
应有截面 \(s\)。但 \(s(\bar1)\) 必须是 \(\bar1\) 的原像，可写成 \(1+p^at\pmod{p^{a+1}}\)；它不被 \(p\) 整除，所以阶为 \(p^{a+1}\)。另一方面，同态性要求
\[
p^a s(\bar1)=s(p^a\bar1)=0,
\]
所以它的阶又必须整除 \(p^a\)，矛盾。因此有限交换群范畴没有非零投射对象，不能把模范畴中的存在性结论无条件搬到它的交换子范畴。

% !TeX root = ../../Book4.tex
% 纲要标题：### 8.2 Ext
% 纲要位置：计划纲要.md，第 2656 行。

\section{Ext}
\label{b4:sec:0802}

Hom 的两个变量具有相反的箭头方向，分别用投射或内射消解得到的计算必须比较。平衡性使两种模型计算同一组 Ext，并保留两变量的连接态射。


% 纲要内容（仅作写作计划，尚非教材正文）：
%
% * Hom 的右导出函子
% * 反变与协变变量
% * Ext 长正合列及具体计算
%

\subsection{Hom 的右导出函子}
\label{b4:subsec:0802:01}

\begin{definition}\label{b4:def:ext}
设 \(R\) 是含幺环，\(M,N\) 是左 \(R\) 模。对 \(N\) 取内射消解 \(N\to I^\bullet\)，定义交换群
\[
\operatorname{Ext}_R^n(M,N)\coloneqq
H^n\bigl(\operatorname{Hom}_R(M,I^\bullet)\bigr),\qquad n\ge0.
\]
这就是左正合函子 \(\operatorname{Hom}_R(M,-)\) 的右导出函子。
\end{definition}
\symidx{\(\operatorname{Ext}_R^n(M,N)\)：Ext 群}

零次是 \(\operatorname{Hom}_R(M,N)\)。在一般非交换环上，Ext 首先是交换群，不自动是左 \(R\) 模；若 \(R\) 交换，标量乘法与所有微分相容，便得到自然的 \(R\) 模结构。下标的环也不能略去，例如 \(\operatorname{Ext}_{\mathbb Z}^1(\mathbb Z/2,\mathbb Z/2)\) 与 \(\operatorname{Ext}_{\mathbb F_2}^1(\mathbb F_2,\mathbb F_2)\) 不同。

\ProseParagraphBreak
固定 \(N\) 时，\(\operatorname{Hom}_R(-,N)\) 是从左模反范畴到交换群的左正合函子。左模的投射对象正是反范畴的内射对象，所以投射消解也能给出一套右导出函子。两套计算是否一致，需要证明，不能由名字相同得出。

\subsection{反变与协变变量}
\label{b4:subsec:0802:02}

\begin{theorem}[Ext 的平衡性]\label{b4:thm:ext-balanced}
设 \(R\) 是含幺环，\(M,N\) 是幺性左 \(R\) 模，\(P_\bullet\to M\) 为投射消解，\(N\to I^\bullet\) 为内射消解。对整数 \(n\ge0\) 取
\[
\operatorname{Hom}_R(P_\bullet,N)^n=\operatorname{Hom}_R(P_n,N),
\qquad \delta^n(f)=f\,\mathrm{d}_{P,n+1}.
\]
则有自然同构
\[
H^n\bigl(\operatorname{Hom}_R(P_\bullet,N)\bigr)
\cong\operatorname{Ext}_R^n(M,N).
\]
它对 \(M\) 反变，对 \(N\) 协变，并与两变量的连接态射相容。
\end{theorem}
\begin{proof}
令 \(D^{p,q}=\operatorname{Hom}_R(P_p,I^q)\)，其中 \(p,q\ge0\)，定义
\[
\mathrm{d}_h(f)=f \mathrm{d}_P,\qquad \mathrm{d}_v(f)=(-1)^p \mathrm{d}_I f.
\]
两次水平或竖直微分为零；一水平一竖直的两种复合相差一个负号，故 \(\mathrm{d}_h \mathrm{d}_v+\mathrm{d}_v\mathrm{d}_h=0\)。总次数为 \(p+q\)，每条对角线有限，因而总复形只用有限直和。

固定 \(p\)，\(P_p\) 的投射性使 \(\operatorname{Hom}_R(P_p,-)\) 正合，所以竖列只在 \(q=0\) 有上同调 \(\operatorname{Hom}_R(P_p,N)\)。固定 \(q\)，\(I^q\) 的内射性使 \(\operatorname{Hom}_R(-,I^q)\) 正合，所以横行只在 \(p=0\) 有上同调 \(\operatorname{Hom}_R(M,I^q)\)。增广分别给出两条复形映射
\[
\operatorname{Hom}_R(P_\bullet,N)\longrightarrow\operatorname{Tot}D,
\qquad
\operatorname{Hom}_R(M,I^\bullet)\longrightarrow\operatorname{Tot}D.
\]
在第一条的 \(q=0\) 边，竖直微分为零；在第二条的 \(p=0\) 边，水平微分为零。因此它们确实与总微分相容。\cref{b4:cor:double-complex-comparison}分别沿竖列和横行说明这两条映射都是拟同构。该结果通过逐个消去有限对角线上的分量证明。

投射、内射比较映射使上述两条映射对 \(M,N\) 自然；不同选择同伦，故同调态射确定。对短正合列选择马蹄消解后，所有总复形也组成逐次短正合列；两侧增广给出它们之间的交换图。\cref{b4:thm:cohomology-long-exact}的自然性于是保证同构与连接态射相容。
\end{proof}

为了方便表示上同调类，本定理使用无符号的 \(\delta(f)=f\mathrm{d}\)。若把 \(P_n\) 放在上链次数 \(-n\)，按一般内部 Hom 规则会得到 \(\delta(f)=(-1)^{n+1}f\mathrm{d}\)；每个次数 \(n\) 乘 \((-1)^{n(n+1)/2}\) 给出这两种复形的同构。本节以后用前一种计算约定。

\ProseParagraphBreak
给定 \(u:M'\to M\)，其投射消解提升经预合成给出 \(u^*:\operatorname{Ext}_R^n(M,N)\to\operatorname{Ext}_R^n(M',N)\)；给定 \(v:N\to N'\)，后合成给出 \(v_*\)。预合成与后合成交换，所以 Ext 是两个变量的函子，而非仅对每个固定模分别定义的一组交换群。

\subsection{Ext 长正合列及具体计算}
\label{b4:subsec:0802:03}

\begin{theorem}\label{b4:thm:ext-long-exact}
设 \(R\) 是含幺环。左 \(R\) 模短正合列 \(0\to A\to B\to C\to0\) 对任意左模 \(N\) 给出自然长正合列
\[
\begin{aligned}
0&\longrightarrow\operatorname{Hom}_R(C,N)\longrightarrow\operatorname{Hom}_R(B,N)
\longrightarrow\operatorname{Hom}_R(A,N)\\
&\longrightarrow\operatorname{Ext}_R^1(C,N)\longrightarrow\operatorname{Ext}_R^1(B,N)
\longrightarrow\operatorname{Ext}_R^1(A,N)\longrightarrow\cdots.
\end{aligned}
\]
对任意左模 \(M\)，左模短正合列 \(0\to N'\to N\to N''\to0\) 给出自然长正合列
\[
\begin{aligned}
0&\longrightarrow\operatorname{Hom}_R(M,N')\longrightarrow\operatorname{Hom}_R(M,N)
\longrightarrow\operatorname{Hom}_R(M,N'')\\
&\longrightarrow\operatorname{Ext}_R^1(M,N')\longrightarrow\operatorname{Ext}_R^1(M,N)
\longrightarrow\operatorname{Ext}_R^1(M,N'')\longrightarrow\cdots.
\end{aligned}
\]
\end{theorem}
\begin{proof}
第一列取 \(N\) 的内射消解，逐次施加反变正合函子 \(\operatorname{Hom}_R(-,I^q)\)，得到 \(0\to\operatorname{Hom}(C,I)\to\operatorname{Hom}(B,I)\to\operatorname{Hom}(A,I)\to0\)，再取同调长正合列。第二列取 \(M\) 的投射消解，对给定序列逐次施加正合函子 \(\operatorname{Hom}_R(P_p,-)\)，取长正合列，再用\cref{b4:thm:ext-balanced}识别各项。增广在零次给出通常 Hom 映射，连接态射与选择独立且自然。
\end{proof}

\begin{example}\label{b4:exm:ext-cyclic}
设 \(R\) 是主理想整环，\(0\ne a\in R\)，\(N\) 是 \(R\) 模。把 \(R/(a)\) 的两项消解送入 Hom，得到 \(N\xrightarrow{a}N\)，故
\[
\operatorname{Ext}_R^0(R/(a),N)=\{\,n\in N\mid an=0\,\},
\qquad\operatorname{Ext}_R^1(R/(a),N)=N/aN,
\]
二次及以上为零。特别地，\(\operatorname{Ext}_{\mathbb Z}^1(\mathbb Z/m,\mathbb Z/n)\cong\mathbb Z/\gcd(m,n)\mathbb Z\)，其中 \(m,n>0\)。若改取 \(N=\mathbb Q\)，乘 \(m\) 为同构，零次及正次均为零；这与 \(\mathbb Q\) 内射一致。
\end{example}

在 \(R=k[\epsilon]/(\epsilon^2)\) 上，对\secref{b4:sec:0701}给出的 \(k\) 的周期消解施加 \(\operatorname{Hom}_R(-,k)\)，所有微分为零。因此 \(\operatorname{Ext}_R^n(k,k)\cong k\) 对每个 \(n\ge0\) 成立。相同的底层向量空间，在底环为 \(k\) 时正次 Ext 全为零；记录底环是计算的一部分。

% !TeX root = ../../Book4.tex
% 纲要标题：### 8.3 马蹄引理
% 纲要标签：b4:sec:0703
% 纲要位置：计划纲要.md，第 3287 行。

\section{马蹄引理}
\label{b4:sec:0703}

短正合列的两端已有消解时，中项的消解仍须保存原扩张的不分裂信息。马蹄构造把这部分信息放进微分的非对角项，使三份消解逐次相容。


% 纲要内容（仅作写作计划，尚非教材正文）：
%
% * 短正合列与相容消解
% * 投射及内射版本
% * 消解上的长正合列
%

\subsection{短正合列与相容消解}
\label{b4:subsec:0703:01}

假设已有短正合列 \(0\to A\xrightarrow{i}B\xrightarrow{q}C\to0\)，又已分别选好 \(A,C\) 的投射消解。希望构造 \(B\) 的消解时，最自然的逐次对象是 \(P_n^A\oplus P_n^C\)。困难在于微分通常不能取两个微分的直和：那样得到的零次同调是 \(A\oplus C\)，未必是 \(B\)。原列的不分裂信息必须进入微分的非对角项。

\begin{lemma}\label{b4:lem:horseshoe-step}
设 \(\mathcal A\) 是交换范畴，\(0\to A\xrightarrow{i}B\xrightarrow{q}C\to0\) 正合，\(p_A:P_A\to A\)、\(p_C:P_C\to C\) 满，且 \(P_C\) 投射。选 \(t:P_C\to B\) 使 \(qt=p_C\)，令 \(p_B=(ip_A,t):P_A\oplus P_C\to B\)。则 \(p_B\) 满，且
\[
0\longrightarrow\ker p_A\longrightarrow\ker p_B
\longrightarrow\ker p_C\longrightarrow0
\]
正合。
\end{lemma}
\begin{proof}
在两行分别为 \(0\to P_A\to P_A\oplus P_C\to P_C\to0\) 和给定短正合列的图中，竖直态射取 \(p_A,p_B,p_C\)。图由 \(qt=p_C\)、\(qi=0\) 交换：
\[
\begin{tikzcd}[column sep=2.5em,row sep=large]
0\arrow[r]&P_A\arrow[r]\arrow[d,two heads,"p_A"']&P_A\oplus P_C\arrow[r]\arrow[d,"p_B"]&P_C\arrow[r]\arrow[d,two heads,"p_C"]&0\\
0\arrow[r]&A\arrow[r,"i"']&B\arrow[r,"q"']&C\arrow[r]&0
\end{tikzcd}
\]
由\cref{b4:thm:snake-lemma}得到正合列
\[
\begin{aligned}
0\longrightarrow\ker p_A\longrightarrow\ker p_B
&\longrightarrow\ker p_C\longrightarrow\operatorname{coker}p_A\\
&\longrightarrow\operatorname{coker}p_B
\longrightarrow\operatorname{coker}p_C\longrightarrow0.
\end{aligned}
\]
因 \(p_A,p_C\) 均满，两端余核为零。于是 \(\ker p_B\to\ker p_C\) 满，且 \(\operatorname{coker}p_B=0\)，这正是结论。
\end{proof}

\subsection{马蹄引理的投射与内射版本}
\label{b4:subsec:0703:02}

\begin{theorem}[马蹄引理]\label{b4:thm:horseshoe}
设 \(\mathcal A\) 是交换范畴，\(0\to A\to B\to C\to0\) 是短正合列。给定 \(A,C\) 的投射消解 \(P^A_\bullet,P^C_\bullet\)，存在 \(B\) 的投射消解 \(P^B_\bullet\)，使增广相容，且
\[
0\longrightarrow P^A_\bullet\longrightarrow P^B_\bullet
\longrightarrow P^C_\bullet\longrightarrow0
\]
逐次分裂正合，并可取 \(P_n^B=P_n^A\oplus P_n^C\)。对偶地，给定 \(A,C\) 的内射消解，可取 \(I_B^n=I_A^n\oplus I_C^n\)，得到相容且逐次分裂的短正合列 \(0\to I_A^\bullet\to I_B^\bullet\to I_C^\bullet\to0\)。此结果称为\term[matiyinli]{马蹄引理}[horseshoe lemma]。
\end{theorem}
\begin{proof}
这里的分裂逐次数成立，截面未必与微分交换。选定双积表示后，中间微分的非对角块正是需要另外构造的部分；只有存在复形截面时，才能用它选取双积表示，使这些非对角块同时为零。

先将\cref{b4:lem:horseshoe-step}用于 \(P_0^A\to A\)、\(P_0^C\to C\)，得到 \(P_0^B\to B\) 及其核列。记这三个核为 \(K_1^A,K_1^B,K_1^C\)。给定消解中的下一项分别满射到 \(K_1^A,K_1^C\)，于是再用引理得到 \(P_1^B\to K_1^B\)，其核又组成短正合列。归纳继续。把 \(P_n^B\to K_n^B\) 与 \(K_n^B\hookrightarrow P_{n-1}^B\) 复合，就是中间消解的微分；每一步定义保证像等于下一个核。

各次包含与投影通过核图构造，因而与微分交换。其对象列是双积的分裂列；\(P_n^A\oplus P_n^C\) 投射，因为两个分量的态射可以分别提升后相加。这样得到全部结论。内射版本在 \(\mathcal A^{\mathrm{op}}\) 应用已证投射版本，并反转箭头与链指标；有限双积在反范畴仍是双积，故逐次分裂性不变。
\end{proof}

选定逐次双积表示以后，相容消解可以排列如下：
\[
\begin{tikzcd}[column sep=large,row sep=large]
P_1^A\arrow[r]\arrow[d,"\mathrm{d}_1^A"']&P_1^B\arrow[r]\arrow[d,"\mathrm{d}_1^B"]&P_1^C\arrow[d,"\mathrm{d}_1^C"]\\
P_0^A\arrow[r]\arrow[d]&P_0^B\arrow[r]\arrow[d]&P_0^C\arrow[d]\\
A\arrow[r]&B\arrow[r]&C
\end{tikzcd}
\]
从 \(P_0^A,P_0^B,P_0^C\) 向下的箭头为增广；上面两排来自逐次分裂的消解短正合列。微分具体为
\[
\mathrm{d}_n^B=\begin{pmatrix}\mathrm{d}_n^A&h_n\\0&\mathrm{d}_n^C\end{pmatrix},
\qquad \mathrm{d}_{n-1}^Ah_n+h_{n-1}\mathrm{d}_n^C=0.
\]
第二个等式在 \(n\ge2\) 时成立，是 \((\mathrm{d}^B)^2=0\) 的非对角分量，显示原来的扩张如何影响消解。零次的相容关系则由增广给出。

\subsection{消解上的长正合列}
\label{b4:subsec:0703:03}

\begin{proposition}\label{b4:prop:horseshoe-functor-sequence}
设 \(\mathcal A,\mathcal B\) 为交换范畴，\(F:\mathcal A\to\mathcal B\) 为加性函子，\(0\to A\to B\to C\to0\) 为 \(\mathcal A\) 中的短正合列。给定两端 \(A,C\) 的投射消解，并按马蹄引理构造相容的消解短正合列 \(0\to P^A_\bullet\to P^B_\bullet\to P^C_\bullet\to0\)。此列经 \(F\) 逐次作用后仍是逐次分裂的复形短正合列，因此其同调对象有自然连接态射及长正合列。两端取内射消解及相应的马蹄构造时，同样成立。
\end{proposition}
\begin{proof}
分裂短正合列由态射等式 \(ri=1\)、\(ps=1\)、\(ir+sp=1\) 表示。加性函子保持这些等式，故不必假定 \(F\) 正合。再用\cref{b4:thm:cohomology-long-exact}，在链情形反转指标，即得到长正合列。

为证明自然性，给定短正合列间的交换图
\[
\begin{tikzcd}[column sep=large,row sep=large]
0\arrow[r]&A\arrow[r,"i"]\arrow[d,"a"']
&B\arrow[r,"q"]\arrow[d,"b"]
&C\arrow[r]\arrow[d,"c"]&0\\
0\arrow[r]&A'\arrow[r,"i'"']
&B'\arrow[r,"q'"']&C'\arrow[r]&0.
\end{tikzcd}
\]
源端的马蹄消解记为 \(P^A,P^B,P^C\)，目标端记为 \(Q^A,Q^B,Q^C\)。选定分别提升 \(a,c\) 的端点比较映射 \(u:P^A\to Q^A\)、\(v:P^C\to Q^C\)。用双积表示两套微分，分别记非对角项为 \(h_n\)、\(h'_n\)。设各端点增广为 \(\varepsilon_{P^A},\varepsilon_{P^C},\varepsilon_{Q^A},\varepsilon_{Q^C}\)，中间增广写成
\[
\begin{aligned}
\varepsilon_P&=(i\varepsilon_{P^A},t_P):
P_0^A\oplus P_0^C\longrightarrow B,\\
\varepsilon_Q&=(i'\varepsilon_{Q^A},t_Q):
Q_0^A\oplus Q_0^C\longrightarrow B',
\end{aligned}
\]
其中 \(qt_P=\varepsilon_{P^C}\)、\(q't_Q=\varepsilon_{Q^C}\)。我们构造中间比较映射
\[
w_n=\begin{pmatrix}u_n&y_n\\0&v_n\end{pmatrix}.
\]
零次需要满足 \(\varepsilon_Qw_0=b\varepsilon_P\)。第一列已由 \(u_0\) 的增广条件满足；第二列要求
\[
i'\varepsilon_{Q^A}y_0=bt_P-t_Qv_0:
P_0^C\longrightarrow B'.
\]
右端与 \(q'\) 的复合为
\[
q'(bt_P-t_Qv_0)
=c\varepsilon_{P^C}-\varepsilon_{Q^C}v_0=0.
\]
因 \(i'\) 是 \(q'\) 的核，右端唯一分解为 \(i'z_0\)，其中 \(z_0:P_0^C\to A'\)。由 \(P_0^C\) 投射，沿满态射 \(\varepsilon_{Q^A}:Q_0^A\to A'\) 提升 \(z_0\)，得到 \(y_0:P_0^C\to Q_0^A\)。这就满足了中间的增广条件。

已构造到次数 \(n-1\) 后，次数 \(n\) 的链映射条件只剩非对角块
\[
\mathrm{d}_n^{Q^A}y_n=r_n,
\qquad r_n=u_{n-1}h_n+y_{n-1}\mathrm{d}_n^{P^C}-h'_nv_n.
\]
这里 \(r_n:P_n^C\to Q_{n-1}^A\)，其余三块由 \(u,v\) 是链映射及零块自动满足。当 \(n=1\) 时，源、目标增广与微分的相容关系分别给出
\[
i\varepsilon_{P^A}h_1+t_P\mathrm{d}_1^{P^C}=0,
\qquad
i'\varepsilon_{Q^A}h'_1+t_Q\mathrm{d}_1^{Q^C}=0.
\]
结合已构造的 \(y_0\)，可得
\[
\begin{aligned}
i'\varepsilon_{Q^A}r_1
&=-bt_P\mathrm{d}_1^{P^C}
 +(bt_P-t_Qv_0)\mathrm{d}_1^{P^C}
 +t_Q\mathrm{d}_1^{Q^C}v_1\\
&=0,
\end{aligned}
\]
最后一步使用 \(v_0\mathrm{d}_1^{P^C}=\mathrm{d}_1^{Q^C}v_1\)。因 \(i'\) 单，\(\varepsilon_{Q^A}r_1=0\)。当 \(n\ge2\) 时，利用两套微分的非对角关系和上一次数的链映射等式，得到
\[
\begin{aligned}
\mathrm{d}_{n-1}^{Q^A}r_n
={}&-u_{n-2}h_{n-1}\mathrm{d}_n^{P^C}\\
&+\bigl(u_{n-2}h_{n-1}
 +y_{n-2}\mathrm{d}_{n-1}^{P^C}
 -h'_{n-1}v_{n-1}\bigr)\mathrm{d}_n^{P^C}\\
&+h'_{n-1}\mathrm{d}_n^{Q^C}v_n
=0.
\end{aligned}
\]
其中含 \(y_{n-2}\) 的项由 \((\mathrm{d}^{P^C})^2=0\) 消失，其余项由 \(v\) 的链映射等式抵消。因此 \(r_n\) 分解到目标核；目标消解正合，使 \(Q_n^A\) 满射到该核，再由 \(P_n^C\) 投射取提升 \(y_n\)。归纳给出消解短正合列之间的严格交换图。

对图施加 \(F\)，再用复形长正合列的自然性，即有
\[
H_{n-1}(F(u))\,\delta_P
=\delta_Q\,H_n(F(v)).
\]
若固定这两套消解而改变比较映射或中间提升，则所得三个比较链映射仍分别提升 \(a,b,c\)。由\cref{b4:prop:comparison-homotopy-unique}，它们分别链同伦；加性函子保持这些同伦，故诱导的同调态射不变。

若改变端点消解或马蹄消解，则把上述构造用于原短正合列的恒等图。三个比较映射都提升恒等，由\cref{b4:cor:resolutions-homotopy-equivalent}在施加 \(F\) 后诱导同调同构。比较映射的同伦唯一性保证这些同构不依赖提升，并满足传递律；刚才的自然性等式保证它们也与连接态射交换。因此，不同消解所得的长正合列在这些典范比较同构下相容，而不是在字面上具有相同的同调对象。内射情形在反范畴作同一构造，转回后得到上同调长正合列及其相容性。
\end{proof}

连接态射有可操作的含义：把商复形中的圈提升到中间复形，对提升取微分，再把所得像读成左端复形中的圈。马蹄引理保证即使函子不正合，仍有逐次正合的复形列来执行这一过程；下一章的 Ext 与 Tor 正是由此获得正合列。

% !TeX root = ../../Book4.tex
% 纲要标题：### 8.4 特殊消解
% 纲要标签：b4:sec:0704
% 纲要位置：计划纲要.md，第 3392 行。

\section{特殊消解}
\label{b4:sec:0704}

消解的长度与形式会随底环改变：主理想整环给出短自由消解，正则列给出 Koszul 消解，群或代数乘法给出 bar 消解。每种公式都需要独立检查其正合性。


% 纲要内容（仅作写作计划，尚非教材正文）：
%
% * 主理想整环模的短消解
% * Koszul 消解简介
% * 群代数和代数的 bar 消解预告
%
% ---
%

\subsection{主理想整环模的短消解}
\label{b4:subsec:0704:01}

% 跨卷引用目标：b2:thm:pid-free-submodule（4.3.1）。
第二卷定理4.3.1证明了主理想整环上有限秩自由模的子模自由。有限秩证明逐个加入基向量，每一步用一个主理想描述新增加的部分。无限基未必能按自然数排成一列，也没有供有限归纳结束的最后一项。这里使用选择公理的良序形式，把任意基按序数排列：有直接前驱的阶段只增加一个基向量，没有直接前驱的极限阶段则取此前各阶段的并。自由模中的向量只有有限个非零坐标，因而取并仍只使用有限线性组合。

\begin{theorem}\label{b4:thm:pid-free-submodule}
设 \(R\) 是主理想整环，\(F\) 是任意秩自由 \(R\) 模。每个子模 \(N\subseteq F\) 都自由。因此每个 \(R\) 模 \(M\) 都有长度至多为一的自由消解 \(0\to F_1\to F_0\to M\to0\)。
\end{theorem}
\begin{proof}
用选择公理良序排列 \(F\) 的一组基，写成 \((e_\alpha)_{\alpha<\kappa}\)。对 \(\beta\le\kappa\)，令 \(F_\beta\) 由 \(\alpha<\beta\) 的基向量生成，\(N_\beta=N\cap F_\beta\)。在增加 \(e_\alpha\) 的阶段，取其坐标映射
\[
\pi_\alpha:N_{\alpha+1}\longrightarrow R.
\]
其核恰为 \(N_\alpha\)，像 \(I_\alpha\) 是 \(R\) 的一个理想，故模的第一同构定理给出
\[
N_{\alpha+1}/N_\alpha\cong I_\alpha.
\]
若 \(I_\alpha\ne0\)，取非零生成元 \(a_\alpha\)，并在 \(N_{\alpha+1}\) 选一个 \(e_\alpha\) 坐标为 \(a_\alpha\) 的元素 \(x_\alpha\)。任意 \(y\in N_{\alpha+1}\) 的该坐标是 \(ra_\alpha\)，所以 \(y-rx_\alpha\in N_\alpha\)，这证明 \(N_{\alpha+1}=N_\alpha+Rx_\alpha\)。若 \(rx_\alpha\in N_\alpha\)，看 \(e_\alpha\) 坐标得 \(ra_\alpha=0\)。因为 \(R\) 是整环且 \(a_\alpha\ne0\)，必有 \(r=0\)，所以两部分交为零，即
\[
N_{\alpha+1}=N_\alpha\oplus Rx_\alpha.
\]
若 \(I_\alpha=0\)，则 \(N_{\alpha+1}=N_\alpha\)。

在非零极限序数 \(\lambda\) 处，给定非零 \(y\in N_\lambda\)，它的非零坐标指标构成有限集，因而有最大元 \(\gamma<\lambda\)。极限序数没有直接前驱，所以 \(\gamma+1<\lambda\)，且 \(y\in N_{\gamma+1}\)。零向量也属于此前各阶段，故
\[
N_\lambda=\bigcup_{\beta<\lambda}N_\beta.
\]
从 \(N_0=0\) 开始，上述后继阶段的分解与极限阶段的取并表明，每个 \(N_\beta\) 都由已选出的 \(x_\alpha\)（\(\alpha<\beta\)）生成，特别是它们生成 \(N=N_\kappa\)。若有有限关系
\[
r_1x_{\alpha_1}+\cdots+r_mx_{\alpha_m}=0,
\qquad \alpha_1<\cdots<\alpha_m,
\]
看 \(e_{\alpha_m}\) 坐标：前面各项在该坐标为零，故 \(r_ma_{\alpha_m}=0\)，从而 \(r_m=0\)。再对余下的有限项重复，得到全部系数为零。因此这些 \(x_\alpha\) 构成 \(N\) 的基。

最后取自由满射 \(F_0\twoheadrightarrow M\)，其核 \(F_1\subseteq F_0\) 由刚证明的结论自由，于是
\[
0\longrightarrow F_1\longrightarrow F_0\longrightarrow M\longrightarrow0
\]
是所需的长度至多为一的自由消解。
\end{proof}

有限生成 \(M\) 时可先取有限秩 \(F_0\)，并由有限秩版本得到有限秩 \(F_1\)。若 \(M=R/(a)\)、\(a\ne0\)，其短消解就是乘 \(a\) 的映射；整环假设在这里保证该映射单射。

\subsection{Koszul 消解简介}
\label{b4:subsec:0704:02}

设 \(R\) 为交换含幺环，\(x\in R\)。两项链复形 \(K(x)=(0\to R\xrightarrow{x}R\to0)\) 把两个 \(R\) 分别置于次数一、零，微分 \(\mathrm{d}_1\) 是乘 \(x\)。其零次同调为余核 \(R/(x)\)，一次同调为核 \(\{\,r\in R\mid xr=0\,\}\)。因此它成为 \(R/(x)\) 的消解，准确地要求乘 \(x\) 单射。

\begin{definition}\label{b4:def:koszul-preview}
设 \(R\) 是交换含幺环，\(r\ge0\) 为整数，\(x_1,\ldots,x_r\in R\)。取自由 \(R\) 模 \(R^r\) 的基 \(e_1,\ldots,e_r\)，令 \(K_n=\bigwedge^nR^r\)（\(0\le n\le r\)），其余次数取零，并用 \(R\) 线性延拓定义微分
\[
\mathrm{d}(e_{i_1}\wedge\cdots\wedge e_{i_n})
=\sum_{j=1}^n(-1)^{j-1}x_{i_j}
 e_{i_1}\wedge\cdots\widehat{e_{i_j}}\cdots\wedge e_{i_n}.
\]
所得链复形记为 \(K_\bullet(x_1,\ldots,x_r;R)\)，称为\term[koszul-fuxing]{Koszul 复形}[Koszul complex]。
\end{definition}
\symidx{\(K_\bullet(x_1,\ldots,x_r;R)\)：Koszul 复形}

例如，当 \(i<j\) 时，
\[
\mathrm{d}(e_i\wedge e_j)=x_ie_j-x_je_i,
\qquad
\mathrm{d}^2(e_i\wedge e_j)=x_ix_j-x_jx_i=0.
\]
% 跨卷引用目标：b3:thm:koszul-regular-resolution（15.3.2）。
一般情形也如此：删去两个基向量的两种次序符号相反，而 \(R\) 的交换性使两种次序中的系数相同，所以 \(\mathrm{d}^2=0\)。外幂及楔积的定义与基已在第二卷第九章建立。第三卷定理15.3.2证明了正则序列给出 Koszul 消解；下面利用同调长正合列给出第二证明，考察每增加一个元素时同调如何变化。

\begin{proposition}\label{b4:prop:koszul-resolution}
% 跨卷引用目标：b3:def:regular-sequence（15.1.1）。
设 \(R\) 是交换含幺环，\(r\ge0\) 为整数，\(x_1,\ldots,x_r\in R\)。假设 \((x_1,\ldots,x_r)\ne R\)，且对每个 \(1\le i\le r\)，乘 \(x_i\) 在 \(R/(x_1,\ldots,x_{i-1})\) 上单射，即按第三卷定义15.1.1，\(x_1,\ldots,x_r\) 是一个 \(R\) 正则序列。则 \(R\) 上的 Koszul 复形 \(K_\bullet(x_1,\ldots,x_r;R)\) 连同零次增广 \(R\to R/(x_1,\ldots,x_r)\) 是 \(R/(x_1,\ldots,x_r)\) 的有限自由消解。
\end{proposition}
\begin{proof}
\(r=0\) 时取集中于零次的 \(R\)，结论成立。对正整数 \(r\) 归纳。\(r=1\) 就是前面的两项计算。设 \(K'=K_\bullet(x_1,\ldots,x_{r-1};R)\)。楔积的每个基元要么不含 \(e_r\)，要么唯一写成前 \(r-1\) 个基向量的楔积再楔上 \(e_r\)。因此
\[
K_n=K'_n\oplus(K'_{n-1}\wedge e_r).
\]
若 \(a\in K'_n\)、\(b\in K'_{n-1}\)，则在 \(b\) 的每个楔积基元之后添入的 \(e_r\) 位于第 \(n\) 个位置，删去它的符号为 \((-1)^{n-1}\)，所以
\[
\mathrm{d}(a+b\wedge e_r)=\mathrm{d}'a+(-1)^{n-1}x_rb+(\mathrm{d}'b)\wedge e_r.
\]
令 \(Q_n=K'_{n-1}\)、\(\mathrm{d}_{Q,n}=\mathrm{d}'_{n-1}\)。这里重新编号次数并保留原微分，与前章微分带负号的移位约定不同。以
\[
\iota_n(a)=a,\qquad q_n(a+b\wedge e_r)=b
\]
定义包含与投影。上述微分公式给出
\(\mathrm{d}\iota=\iota\mathrm{d}'\) 及
\(q_{n-1}\mathrm{d}=\mathrm{d}_{Q,n}q_n\)，因此它们组成短正合复形列
\[
0\longrightarrow K'_\bullet\xrightarrow{\iota}K_\bullet
\xrightarrow{q}Q_\bullet\longrightarrow0.
\]
每个次数的投影都以 \(b\mapsto b\wedge e_r\) 为截面，但这个截面与微分的差为左端中的 \((-1)^{n-1}x_rb\)，一般并不组成复形截面。这与马蹄微分的非对角项一样，把逐次分裂与复形分裂区分开来。

对 \(n\ge1\)，\(Q_n\) 中的圈 \(b\in K'_{n-1}\) 满足 \(\mathrm{d}'b=0\)。按照同调长正合列中连接同态的构造，先在 \(K_n\) 中以 \(b\wedge e_r\) 提升它，再取微分，得到
\[
\mathrm{d}(b\wedge e_r)=(-1)^{n-1}x_rb.
\]
所得元素完全落在子复形 \(K'\) 中，故在识别 \(H_n(Q)=H_{n-1}(K')\) 后，连接同态为
\[
\delta_n:H_{n-1}(K')\longrightarrow H_{n-1}(K'),
\qquad [b]\longmapsto(-1)^{n-1}[x_rb].
\]
归纳假设给出 \(H_j(K')=0\)（\(j>0\)）及 \(H_0(K')=R/(x_1,\ldots,x_{r-1})\)。当 \(n\ge2\) 时，长正合列中的片段为
\[
H_n(K')\longrightarrow H_n(K)\longrightarrow H_n(Q)=H_{n-1}(K'),
\]
两端为零，所以 \(H_n(K)=0\)。只有 \(n=1\) 时右端仍为 \(H_0(K')\)，对应的片段为
\[
0\longrightarrow H_1(K)\longrightarrow R/(x_1,\ldots,x_{r-1})
\xrightarrow{x_r}R/(x_1,\ldots,x_{r-1})
\longrightarrow H_0(K)\longrightarrow0.
\]
正则性在这一归纳步骤中保证中间乘 \(x_r\) 的映射单射，故一次同调也消失，而零次同调为所需商环。Koszul 复形的定义、上述短正合列及连接同态公式本身都不要求正则性；这一条件控制的是正次数同调的消失。
\end{proof}

例如，设 \(k\) 为域，令 \(R=k[x,y]\)，则 \(k=R/(x,y)\) 有消解
\[
0\longrightarrow R\xrightarrow{c\mapsto(-yc,xc)}R^2
\xrightarrow{(a,b)\mapsto xa+yb}R\longrightarrow k\longrightarrow0.
\]
把此消解与 \(k\) 在 \(R\) 上张量后，\(x,y\) 在 \(k\) 上的作用均为零，因此张量后的微分全为零。下一章将由此计算出
\[
\operatorname{Tor}_1^R(k,k)\cong k^2,
\qquad \operatorname{Tor}_2^R(k,k)\cong k.
\]
这些非零同调来自张量后的复形；原来的 Koszul 消解在正次数仍然正合。这也说明 \(-\otimes_Rk\) 没有保持该消解的正合性，与\cref{b4:ex:06-tensor-breaks-quasi}中用 \(\mathbb Z/2\mathbb Z\) 张量破坏拟同构的计算一样。Tor 函子正是用这些新出现的同调记录张量函子未保持正合性的信息。

\subsection{群代数与代数的 bar 消解}
\label{b4:subsec:0704:03}

群环和一般代数的自然模往往没有短消解，却有一种由“插入一个符号”控制正合性的自由消解。这里介绍其形式；后面群上同调与 Hochschild 上同调的章节将说明它为何适合相应问题。

\begin{example}\label{b4:exm:homogeneous-bar}
设 \(G\) 是群，\(k\) 是交换含幺环。对 \(n\ge0\)，令 \(B_n\) 为以 \((g_0,\ldots,g_n)\in G^{n+1}\) 为基的自由 \(k\) 模，并让 \(G\) 同时左乘所有坐标。任意元组唯一写成
\[
(g_0,\ldots,g_n)
=g_0\cdot(1,g_0^{-1}g_1,\ldots,g_0^{-1}g_n).
\]
因此首坐标为 \(1\) 的元组构成各轨道的唯一代表，每个轨道上的自由 \(k\) 模同构于左正则模 \(kG\)，从而
\[
B_n\cong\bigoplus_{(h_1,\ldots,h_n)\in G^n}kG
\]
是自由幺性左 \(kG\) 模，其中 \(n=0\) 时右边只有一个直和项。对 \(n\ge1\) 令
\[
\mathrm{d}_n(g_0,\ldots,g_n)
=\sum_{i=0}^n(-1)^i(g_0,\ldots,\widehat{g_i},\ldots,g_n),
\qquad\varepsilon(g_0)=1.
\]
若 \(i<j\)，先删去第 \(j\) 项再删第 \(i\) 项的系数为 \((-1)^{i+j}\)，先删第 \(i\) 项再删原来的第 \(j\) 项的系数为 \((-1)^{i+j-1}\)：第二次删除时该项的编号已变为 \(j-1\)。两项相消，故 \(\mathrm{d}_{n-1}\mathrm{d}_n=0\)（\(n\ge2\)）；次数一还满足 \(\varepsilon\mathrm{d}_1(g_0,g_1)=1-1=0\)。

在底层 \(k\) 模上，用 \(k\) 线性延拓定义
\[
s_n(g_0,\ldots,g_n)=(1,g_0,\ldots,g_n)\quad(n\ge0),
\qquad s_{-1}(1)=(1).
\]
删除新插入的首项给出原元组，其余删除项与 \(s\mathrm{d}\) 抵消，因此
\[
\begin{aligned}
\mathrm{d}_{n+1}s_n+s_{n-1}\mathrm{d}_n&=1_{B_n}\quad(n\ge1),\\
\mathrm{d}_1s_0+s_{-1}\varepsilon&=1_{B_0},
\qquad \varepsilon s_{-1}=1_k.
\end{aligned}
\]
底层增广 \(k\) 模复形于是可缩，特别是正合。所有微分与增广都是 \(kG\) 线性的，核、像在遗忘 \(G\) 作用后仍是同一底层子模，故底层等式 \(\ker\mathrm{d}_n=\operatorname{im}\mathrm{d}_{n+1}\) 也就是 \(kG\) 模中的正合性；增广处同样成立。这便是平凡幺性左 \(kG\) 模 \(k\) 的自由消解。

不过，对 \(g\in G\) 一般有
\[
\begin{aligned}
s_n\bigl(g\cdot(g_0,\ldots,g_n)\bigr)&=(1,gg_0,\ldots,gg_n),\\
g\cdot s_n(g_0,\ldots,g_n)&=(g,gg_0,\ldots,gg_n).
\end{aligned}
\]
首坐标的差别说明上述 \(k\) 线性收缩通常不是 \(kG\) 线性的。正合性可由底层核、像检测，可缩性却要求收缩属于正在讨论的模范畴，因而不能由这个收缩断言增广复形在 \(kG\) 模范畴中可缩。
\end{example}

代数的 bar 构造把元组换成张量串。设 \(k\) 为域，\(A\) 为含幺 \(k\) 代数，令 \(A^e=A\otimes_k A^{\mathrm{op}}\)。对 \(n\ge0\)，取 \(B_n=A^{\otimes_k(n+2)}\)，其幺性左 \(A^e\) 模作用为
\[
\begin{aligned}
&(a\otimes b^{\mathrm{op}})\cdot(a_0\otimes\cdots\otimes a_{n+1})\\
&\hspace{2em}=aa_0\otimes a_1\otimes\cdots\otimes a_n\otimes a_{n+1}b.
\end{aligned}
\]
也就是说，首尾因子分别承担左、右 \(A\) 作用。对 \(n\ge1\)，微分交错地乘合相邻因子：
\[
\mathrm{d}_n(a_0\otimes\cdots\otimes a_{n+1})
=\sum_{i=0}^n(-1)^i
 a_0\otimes\cdots\otimes a_ia_{i+1}\otimes\cdots\otimes a_{n+1},
\]
增广 \(\varepsilon:B_0=A\otimes_k A\to A\) 为乘法。结合律使两次乘合的项成对抵消，故得到复形。在左端插入 \(1\)，即取
\[
\begin{aligned}
s_n(a_0\otimes\cdots\otimes a_{n+1})
&=1\otimes a_0\otimes\cdots\otimes a_{n+1},\\
s_{-1}(a)&=1\otimes a.
\end{aligned}
\]
便给出底层向量空间的收缩：乘合新插入的首因子产生原张量，其余项与 \(s\mathrm{d}\) 抵消。收缩只需 \(k\) 线性，而微分与增广都为 \(A^e\) 线性，故与群的情形一样，可以由底层正合性得到 \(A^e\) 模的正合性。

\ProseParagraphBreak
自由性则由首尾因子与中间因子的分离得到。取 \(A^{\otimes_k n}\) 的一组 \(k\) 基 \(\mathcal B_n\)，其中 \(A^{\otimes_k0}=k\)。映射
\[
A^e\otimes_k A^{\otimes_k n}\longrightarrow B_n,
\qquad
(a\otimes b^{\mathrm{op}})\otimes v\longmapsto a\otimes v\otimes b
\]
是左 \(A^e\) 模同构，因此
\[
B_n\cong A^e\otimes_k A^{\otimes_k n}
\cong\bigoplus_{v\in\mathcal B_n}A^e.
\]
由此得到 \(A\) 作为幺性左 \(A^e\) 模的自由消解。域的假设在此保证中间张量因子有向量空间基；在一般交换底环上不能直接由同一公式断言各项自由。双模约定、归一化及由此得到的上同调计算见\chapref{b4:ch:12}。


\begin{chaptersummary}
投射消解反复把关系写成投射对象的商，内射消解反复处理嵌入后的商对象。模范畴中的存在性使每个模都能进入这两种构造；一般交换范畴必须先检查相应假设。比较定理给出不同消解之间的同伦等价，加性函子仍保持这些等价；当目标为交换范畴时，取同调便得到典范比较同构。

\ProseParagraphBreak
马蹄消解的每次对象是双积，微分却可能含非对角项。逐次分裂性保证加性函子作用后的列仍逐次分裂；若目标为交换范畴，再取同调就得到长正合列。计算时则宜尽量选适合对象的消解：主理想整环给出两项自由消解，正则列给出有限 Koszul 消解，而 bar 构造用统一的较大模型换取态射上的便利。
\end{chaptersummary}

\BookExercises

\begin{exercise}\label{b4:ex:07-noetherian-finite-resolution}
设 \(R\) 是左 Noether 的含幺环\footnote{埃米·诺特（Amalie Emmy Noether，1882-1935），德国数学家，建立现代抽象代数中理想的链条件理论，并提出连接对称性与守恒律的定理。}，\(M\) 是有限生成幺性左 \(R\)-模。证明可以给 \(M\) 选择每次秩都有限的自由消解。该证明是否保证只有有限个非零项？
\end{exercise}
\begin{solution}
有限生成性给出 \(R^{r_0}\to M\) 的满射。因 \(R\) 左 Noether，有限秩自由模的子模有限生成，故核有有限自由满射。重复此论证即可。这里控制的是每个 \(P_n\) 的秩，并没有保证某个关系核 \(K_n\) 变为零，所以每项有限秩并不等于消解只有有限个非零项。\(k[\epsilon]/(\epsilon^2)\) 上的 \(k\) 有逐项秩一的无限消解，下一章还会用它的正次导出函子排除有限长度消解。
\end{solution}

\begin{exercise}\label{b4:ex:07-nonnoetherian-relations}
设 \(k\) 是域，\(R=k[x_1,x_2,\ldots]\)，\(\mathfrak m=(x_1,x_2,\ldots)\)。证明循环模 \(R/\mathfrak m\) 的自然自由满射 \(R\to R/\mathfrak m\) 的核不是有限生成的。
\end{exercise}
\begin{solution}
若 \(\mathfrak m=(f_1,\ldots,f_s)\)，每个 \(f_i\) 只有有限个单项式，每个单项式又只含有限个变量，故有限多个 \(f_i\) 所涉及的变量总并仍有限。取其中未出现的 \(x_j\)，把除 \(x_j\) 外的所有变量置零，得到 \(R\to k[x_j]\) 的同态。各 \(f_i\) 的常数项为零，故其像均为零，而 \(x_j\) 的像非零，与它属于所生成理想矛盾。有限生成对象的关系未必有限生成。
\end{solution}

\begin{exercise}\label{b4:ex:07-cyclic-injective-resolution}
对整数 \(m>1\)，用 \(\mathbb Q/\mathbb Z\) 写出 \(\mathbb Z/m\mathbb Z\) 的两项内射消解，并写清嵌入。
\end{exercise}
\begin{solution}
所需序列是
\[
0\longrightarrow\mathbb Z/m\mathbb Z
\xrightarrow{\bar a\mapsto a/m+\mathbb Z}\mathbb Q/\mathbb Z
\xrightarrow{\times m}\mathbb Q/\mathbb Z\longrightarrow0.
\]
若 \(a'-a=mt\)，则 \(a'/m-a/m=t\in\mathbb Z\)，所以嵌入不依赖剩余类代表。又 \(a/m\in\mathbb Z\) 当且仅当 \(m\mid a\)，故它单射。乘 \(m\) 的核恰为这些分数类；乘 \(m\) 满，因为有理数代表元可以除以 \(m\)。两项可除，因而内射。
\end{solution}

\begin{exercise}\label{b4:ex:07-cyclic-comparison}
设 \(m,n>0\)，\(f:\mathbb Z/m\mathbb Z\to\mathbb Z/n\mathbb Z\) 把 \(\bar1\) 送到 \(\bar c\)。写出它在标准两项自由消解间的提升，并解释 \(c\) 换成 \(c+nt\) 时两个提升的同伦。
\end{exercise}
\begin{solution}
良定义要求 \(n\mid mc\)。零次映射取乘 \(c\)，一次取乘 \(mc/n\)，则 \(n(mc/n)=cm\)，方块交换。记使用 \(c+nt\) 的新提升为 \(u\)，使用 \(c\) 的原提升为 \(v\)。则 \(u-v\) 在零次为乘 \(nt\)，一次为乘 \(mt\)。唯一可能非零的链同伦分量是 \(s_0:P_0\to Q_1\)。取 \(s_0:\mathbb Z\to\mathbb Z\) 为乘 \(t\)，则两个次数的同伦公式分别为
\[
(u-v)_0=\mathrm{d}_1^Qs_0=\times nt,
\qquad
(u-v)_1=s_0\mathrm{d}_1^P=\times mt.
\]
因此它给出所需同伦。
\end{solution}

\begin{exercise}\label{b4:ex:07-projective-source-necessary}
若把比较定理中的“源为投射消解”改为“源为任意正合增广复形”，结论是否仍成立？用 \(\mathbb Z/2\mathbb Z\) 构造反例。
\end{exercise}
\begin{solution}
令源只有零次 \(P_0=\mathbb Z/2\mathbb Z\)，增广为恒等；目标取标准自由消解 \(\mathbb Z\xrightarrow2\mathbb Z\to\mathbb Z/2\mathbb Z\)。若恒等有提升，零次便是 \(\mathbb Z/2\mathbb Z\to\mathbb Z\) 的同态且复合到 \(\mathbb Z/2\mathbb Z\) 为恒等。但任何这样的同态均为零，因为整数无二挠。这排除了提升。失败恰好发生在比较定理的第一步：源的零次项不是投射对象，恒等态射无法沿目标的自由满射提升；即使目标是一份投射消解，也不能弥补源缺少的投射性。
\end{solution}

\begin{exercise}\label{b4:ex:07-add-contractible}
设 \(P_\bullet\to M\) 是投射消解，\(Q\) 是投射对象。在相邻非负次数加入复形 \(Q\xrightarrow1Q\)，其余次数为零。证明所得直和仍是 \(M\) 的投射消解，并写出它到原消解的同伦等价。
\end{exercise}
\begin{solution}
设新增复形 \(T\) 的非零项在次数 \(r+1,r\)，其中 \(r\ge0\)，微分 \(T_{r+1}\to T_r\) 为 \(1_Q\)。在新增零次项上把增广取为零。\(T\) 的同调为零，每次对象仍为投射双积，故增广同调不变。包含 \(i\) 和投影 \(p\) 满足 \(pi=1\)；\(1-ip\) 仅在新增复形上为恒等。取 \(s_r:T_r\to T_{r+1}\) 为 \(1_Q\)，其余同伦分量取零，就有 \(\mathrm{d}s+s\mathrm{d}=1-ip\)。因此 \(i,p\) 是同伦逆。
\end{solution}

\begin{exercise}\label{b4:ex:07-projective-augmented-split}
设 \(M\) 本身投射。证明任何投射消解 \(P_\bullet\to M\) 的增广复形都可缩，并解释这里的收缩为什么一般依赖选择。
\end{exercise}
\begin{solution}
\(\varepsilon:P_0\to M\) 有截面 \(s_{-1}:M\to P_0\)，满足 \(\varepsilon s_{-1}=1_M\)，所以短正合列
\[
0\longrightarrow K_1\longrightarrow P_0\xrightarrow{\varepsilon}M\longrightarrow0
\]
分裂，\(K_1\) 是 \(P_0\) 的投射直和因子。继而
\[
0\longrightarrow K_2\longrightarrow P_1\longrightarrow K_1\longrightarrow0
\]
也分裂，\(K_2\) 因而投射。如此递归，依次得到 \(P_n\cong K_{n+1}\oplus K_n\)，其中 \(K_0=M\)。微分把第二因子恒等送到前一项的第一因子，在另一因子上为零。将各第一因子恒等提升到后一项的第二因子，连同增广复形的次数 \(-1\) 上的 \(s_{-1}:M\to P_0\)，即得 \(\mathrm{d}s+s\mathrm{d}=1\)；在次数 \(-1\) 上，该式就是 \(\varepsilon s_{-1}=1_M\)。截面没有被投射性唯一指定，所以收缩一般依赖选择。
\end{solution}

\begin{exercise}\label{b4:ex:07-horseshoe-p-square}
令 \(p\) 为素数，对 \(0\to\mathbb Z/p\mathbb Z\xrightarrow{\bar a\mapsto p a}\mathbb Z/p^2\mathbb Z\to\mathbb Z/p\mathbb Z\to0\) 构造马蹄消解。验证中间微分可取 \(\begin{pmatrix}p&-1\\0&p\end{pmatrix}\)，增广为 \((a,b)\mapsto pa+b\pmod{p^2}\)。
\end{exercise}
\begin{solution}
增广满射，与微分复合为 \(p^2a\)，故为零；矩阵行列式为 \(p^2\ne0\)，所以它在 \(\mathbb Z^2\) 上单射。若 \(pa+b=p^2t\)，则 \(b=p(pt-a)\)，而
\[
\begin{pmatrix}p&-1\\0&p\end{pmatrix}
\binom{t}{pt-a}=\binom{a}{b}.
\]
因此其像恰为增广核。左端消解以第一坐标包含，右端以第二坐标投影，两个对角块均为乘 \(p\)。非对角项 \(-1\) 使增广复合中的两项 \(-pb\) 与 \(pb\) 抵消，因而把两端消解相容地拼成给定的非分裂扩张，而不是它们的直和。若去掉非对角项 \(-1\)，未增广复形就成为两端消解的直和，零次同调为 \((\mathbb Z/p\mathbb Z)^2\)。而原增广与新微分的复合为 \(p^2a+pb\pmod{p^2}\)，不恒为零，所以它已不再是给定扩张 \(\mathbb Z/p^2\mathbb Z\) 的增广消解。
\end{solution}

\begin{exercise}\label{b4:ex:07-koszul-repeat}
设 \(k\) 是域。在 \(R=k[x]\) 上计算 \(K_\bullet(x,x;R)\) 的全部同调，说明第二个元素的正则性为何失败。
\end{exercise}
\begin{solution}
微分为 \(R\xrightarrow{(-x,x)}R^2\xrightarrow{(x,x)}R\)。左端单射，故 \(H_2=0\)；中间核为 \(R(1,-1)\)，而
\[
\operatorname{im}\mathrm{d}_2=R(-x,x)=Rx(1,-1),
\]
因为乘以单位 \(-1\) 不改变生成的子模。因此
\[
H_1=R(1,-1)/Rx(1,-1)\cong R/(x).
\]\(H_0=R/(x)\)。第二个 \(x\) 在 \(R/(x)\) 上为零，乘法不是单射，正是一次同调不消失的原因。
\end{solution}

\begin{exercise}\label{b4:ex:07-koszul-nonreduced}
设 \(k\) 是域。在 \(R=k[x,y]\) 上写出 \(R/(x,y^2)\) 的两步自由消解。末端商环含有非零幂零元，是否妨碍该消解正合？
\end{exercise}
\begin{solution}
\(x\) 在 \(R\) 上非零因子，\(y^2\) 在 \(R/(x)\cong k[y]\) 上也非零因子。因此消解为
\[
0\to R\xrightarrow{(-y^2,x)}R^2\xrightarrow{(x,y^2)}R\to R/(x,y^2)\to0.
\]
末端 \(\bar y\ne0\) 而 \(\bar y^2=0\)，但正则性是在逐步取商以前检验乘法的单射性，最终商环仍可能非约化。
\end{solution}

\begin{exercise}\label{b4:ex:07-bar-contraction-equivariance}
在\cref{b4:exm:homogeneous-bar}中，假设 \(k\ne0\)，并取非平凡元素 \(g\in G\)。比较 \(s_0(g)\) 与 \(g\cdot s_0(1)\)，并在次数零验证收缩等式。
\end{exercise}
\begin{solution}
作用的两种复合为
\[
s_0\bigl(g\cdot(1)\bigr)=s_0(g)=(1,g),
\qquad
 g\cdot s_0(1)=g\cdot(1,1)=(g,g).
\]
两者是不同基元，故 \(s_0\) 不是 \(G\) 等变态射。又 \(\mathrm{d}_1s_0(g)=(g)-(1)\)，\(s_{-1}\varepsilon(g)=(1)\)，两者之和为 \((g)\)。底层 \(k\) 模上的收缩足以证明正合性，却不能推出平凡 \(kG\) 模投射。
\end{solution}

\begin{exercise}\label{b4:ex:07-finitely-generated-no-injectives}
证明有限生成交换群范畴没有非零内射对象，因而没有足够多内射对象。
\end{exercise}
\begin{solution}
若 \(I\) 在此范畴内射，给定整数 \(n\ge1\) 和 \(a\in I\)，先定义 \(n\mathbb Z\to I\)，使 \(n\mapsto a\)。由单态射 \(n\mathbb Z\hookrightarrow\mathbb Z\) 和内射性，将其延拓为 \(\varphi:\mathbb Z\to I\)。令 \(b=\varphi(1)\)，便有 \(a=\varphi(n)=nb\)，故 \(I\) 可除。有限生成交换群的结构定理给出 \(I\cong\mathbb Z^r\oplus T\)，\(T\) 有限。乘二满射排除 \(r>0\)，乘 \(|T|\) 满射又迫使 \(T=0\)。所以只有零对象内射，而 \(\mathbb Z\) 不能嵌入零对象。
\end{solution}

\begin{exercise}\label{b4:ex:07-cyclic-injective-comparison}
设 \(m,n>1\) 为整数，\(f:\mathbb Z/m\mathbb Z\to\mathbb Z/n\mathbb Z\) 为交换群同态，并设 \(f(\bar1)=\bar c\)，其中 \(c\in\mathbb Z\)。分别用
\[
0\longrightarrow\mathbb Z/r\mathbb Z
\xrightarrow{\bar a\mapsto a/r+\mathbb Z}\mathbb Q/\mathbb Z
\xrightarrow{\times r}\mathbb Q/\mathbb Z\longrightarrow0,
\qquad r=m,n,
\]
作为两者的内射消解。写出 \(f\) 在消解之间的延拓，使两个分量均为乘某个整数的映射。若把 \(c\) 换为 \(c+nt\)（\(t\in\mathbb Z\)），写出两个延拓之间的链同伦。
\end{exercise}
\begin{solution}
同态的良定义要求 \(n\mid mc\)，故 \(\ell=mc/n\) 是整数。记两份内射消解为 \(I_m^\bullet,I_n^\bullet\)。嵌入把 \(\bar1\) 分别送到 \(1/m+\mathbb Z\) 和 \(1/n+\mathbb Z\)。要延拓 \(f\)，零次乘数应使 \(\ell/m+\mathbb Z=c/n+\mathbb Z\)，取 \(\ell=mc/n\) 正好满足这个要求。再由微分相容等式 \(n\ell=mc\) 确定一次分量取乘 \(c\)，所以定义
\[
F^0=\times\ell,\qquad F^1=\times c.
\]
对 \(\bar a\in\mathbb Z/m\mathbb Z\)，有 \(\ell a/m+\mathbb Z=ca/n+\mathbb Z\)，所以 \(F^0\) 延拓 \(f\)。又 \(n\ell=mc\)，故 \((\times n)F^0=F^1(\times m)\)，两个分量组成上链复形映射。

更换代表元后，\(\ell\) 变为 \(\ell+mt\)，两个延拓 \(F',F\) 的差在零次为乘 \(mt\)，在一次为乘 \(nt\)。取唯一可能非零的同伦分量
\[
h^1:I_m^1\longrightarrow I_n^0,\qquad h^1=\times t,
\]
则 \(h^1\mathrm d_{I_m}^0=\times mt\)、\(\mathrm d_{I_n}^0h^1=\times nt\)，正好给出 \(F'-F=\mathrm dh+h\mathrm d\)。
\end{solution}

% !TeX root = ../../Book4.tex
% 纲要标题：## 第9章　导出函子、Ext 与 Tor
% 纲要标签：b4:ch:08
% 纲要位置：计划纲要.md，第 3400 行。

\chapter{导出函子、Ext 与 Tor}
\label{b4:ch:08}
\BookChapterNumber{b4:ch:08}

\begin{chapterabstract}
本章将 Hom 和张量积不保持正合的缺口组织为 Ext 与 Tor，建立导出函子的消解计算、长正合列、消去性与维数平移，并给出典型模上的计算。
\end{chapterabstract}

\begin{learninggoals}
\begin{itemize}
\item 写出左、右导出函子的次数约定，并说明零次为何需要相应单侧正合性。
\item 通过消去性构造泛 \(\delta\)-函子的延拓，区分上同调与同调泛性的方向。
\item 在两个变量中计算 Ext 和 Tor，证明平衡性并正确使用左右模约定。
\item 从长正合列解释 Hom、张量积的正合性失败，使用一次消失判据。
\item 定义相对于消解的合冲模，推导正次数平移及零次的核、余核公式。
\end{itemize}
\end{learninggoals}

取整数 \(n\ge2\)，对短正合列 \(0\rightarrow\mathbb Z\xrightarrow{n}\mathbb Z\rightarrow\mathbb Z/n\mathbb Z\rightarrow0\) 施加 \(\operatorname{Hom}_{\mathbb Z}(-,\mathbb Z)\)。由于从有限群到 \(\mathbb Z\) 的同态只能为零，所得序列在自然识别下为
\[
0\longrightarrow0\longrightarrow\mathbb Z
\xrightarrow{n}\mathbb Z.
\]
Hom 的左正合性保证前三个位置正合，却不能使最后的乘 \(n\) 映射满射；其余核为 \(\mathbb Z/n\mathbb Z\)。投射消解把这一缺口识别为 \(\operatorname{Ext}_{\mathbb Z}^1(\mathbb Z/n\mathbb Z,\mathbb Z)\)，同时要求我们证明它不随消解的选择而变。张量积丢失的正合性也可用消解测量，这便得到 Tor。

\ProseParagraphBreak
本章需要\chapref{b4:ch:07}的比较定理和马蹄引理：比较映射的存在及同伦唯一性使消解定义的群与态射不依赖选择，马蹄引理的逐次分裂短正合列则在施加加性函子后仍正合，从而产生长正合列。\secref{b4:sec:0604}的有限对角双复形比较把按两个变量作消解的结果联系起来，用于证明 Ext 和 Tor 的平衡性。第二卷第6.1—6.3节给出张量积与 Hom 的基本性质，第7.3节证明投射模平坦。以下讨论 Ext 时两个变量都是同一含幺环的左模；讨论 Tor 时第一变量为右模，第二变量为左模。除明确写出的交换情形外，不增加底环交换性的假设。

\ProseParagraphBreak
导出函子也适用于模范畴以外的交换范畴。例如，对拓扑空间 \(X\) 上的交换群层取全局截面是左正合函子，其右导出给出层上同调。层上同调的具体计算见\secref{b4:sec:C03}。

\ProseParagraphBreak
本章可与 Weibel 的《An Introduction to Homological Algebra》\cite[Section 2.1, Definition 2.1.4; Sections 2.4--2.7]{Weibel1994HomologicalAlgebra}对照。左、右导出函子的泛性分别见\cite[Theorem 2.4.7; Section 2.5.1]{Weibel1994HomologicalAlgebra}，消去性的一般判据见\cite[Exercise 2.4.5]{Weibel1994HomologicalAlgebra}；Tor 与 Ext 的平衡性分别见\cite[Theorem 2.7.2, Theorem 2.7.6]{Weibel1994HomologicalAlgebra}。关于整数模的 Tor 与平坦性，另见\cite[Sections 3.1--3.2]{Weibel1994HomologicalAlgebra}。本章的双复形证明在每个总次数只需有限次消元。

% !TeX root = ../../Book4.tex
% 纲要标题：### 9.1 导出函子
% 纲要标签：b4:sec:0801
% 纲要位置：计划纲要.md，第 2650 行。

\section{导出函子}
\label{b4:sec:0801}

左正合或右正合函子丢失的信息，可以在施加函子后的消解同调中读出。比较定理保证所得对象不依赖消解的选择，并给出它们在态射上的作用；马蹄引理给出连接态射和长正合列；消去性进一步说明这套导出函子具有怎样的泛性质。


% 纲要内容（仅作写作计划，尚非教材正文）：
%
% * 左、右导出及次数约定
% * δ-函子与泛性质
% * 消去性与泛 δ-函子
%

\subsection{左、右导出及次数约定}
\label{b4:subsec:0801:01}

一个左正合函子已经正确保留了核，却可能把满态射变成非满态射。内射消解把这种失败放入一个复形，正次数上同调便记录原函子没有保留的信息。投射消解对于右正合函子起相反方向的作用。

\ProseParagraphBreak
例如，固定整数 \(n\ge2\)，对交换群取左正合函子 \(F=\operatorname{Hom}_{\mathbb Z}(\mathbb Z/n,-)\)。\cref{b4:exm:integer-injective-resolution}已经给出内射消解
\[
 0\longrightarrow\mathbb Z\longrightarrow\mathbb Q
 \longrightarrow\mathbb Q/\mathbb Z\longrightarrow0.
\]
施加 \(F\) 后，\(F(\mathbb Q)=0\)，而 \(F(\mathbb Q/\mathbb Z)\) 是 \(\mathbb Q/\mathbb Z\) 中被 \(n\) 零化的子群，即 \((\tfrac1n\mathbb Z)/\mathbb Z\cong\mathbb Z/n\)。所得复形为 \(0\to\mathbb Z/n\)，位于零次和一次。因此原函子在 \(\mathbb Z\) 上取值为零，消解却在一次留下 \(\mathbb Z/n\)。这正是 \(\mathbb Q\to\mathbb Q/\mathbb Z\) 经过 \(F\) 后失去满射性的部分；下面将这种计算定义为导出函子。

\begin{definition}\label{b4:def:derived-functors}
设 \(\mathcal A,\mathcal B\) 是交换范畴，\(F:\mathcal A\to\mathcal B\) 是加性函子。若 \(\mathcal A\) 有足够多内射对象，对每个对象 \(A\in\mathcal A\) 选内射消解 \(A\to I_A^\bullet\)，定义
\[
R^nF(A)\coloneqq H^n(F(I_A^\bullet)),\qquad n\ge0,
\]
称为 \(F\) 的第 \(n\) 个\term[youdaochuhanzi]{右导出函子}[right derived functor]。若 \(\mathcal A\) 有足够多投射对象，对每个对象 \(A\in\mathcal A\) 选投射消解 \(P^A_\bullet\to A\)，定义
\[
L_nF(A)\coloneqq H_n(F(P^A_\bullet)),\qquad n\ge0,
\]
称为第 \(n\) 个\term[zuodaochuhanzi]{左导出函子}[left derived functor]。对态射 \(f:A\to A'\)，选择延拓或提升 \(f\) 的消解比较映射，逐次应用 \(F\) 后取相应次数的同调，所得态射分别定义 \(R^nF(f)\) 或 \(L_nF(f)\)。
\end{definition}

由\cref{b4:prop:comparison-homotopy-unique,b4:cor:resolutions-homotopy-equivalent}，提升同一个态射的比较映射同伦，因而诱导相同的同调态射；恒等和复合也由比较映射的恒等和复合诱导。于是上述定义给出函子，更换消解只引起满足传递律的自然同构。它们还是加性函子：若 \(f,g:A\to A'\) 分别由比较映射 \(\widetilde f,\widetilde g\) 提升，则 \(\widetilde f+\widetilde g\) 提升 \(f+g\)。由于 \(F\) 加性，取同调后有
\[
R^nF(f+g)=R^nF(f)+R^nF(g),\qquad
L_nF(f+g)=L_nF(f)+L_nF(g).
\]
若把投射消解写成 \(P^{-n}=P_n\)，则 \(L_nF(A)=H^{-n}(F(P^\bullet))\)。下标 \(n\) 是非负同调次数，不是把右导出函子也改放到负次数。

\begin{proposition}\label{b4:prop:derived-degree-zero}
设 \(\mathcal A,\mathcal B\) 是交换范畴。若 \(\mathcal A\) 有足够多内射对象且 \(F:\mathcal A\to\mathcal B\) 加性、左正合，则有自然同构 \(R^0F\cong F\)，且内射对象 \(I\) 满足 \(R^nF(I)=0\) 对 \(n>0\)。若有足够多投射对象且 \(F\) 加性、右正合，则 \(L_0F\cong F\)，且投射对象的正次数左导出函子消失。
\end{proposition}
\begin{proof}
由 \(0\to A\to I^0\to I^1\) 正合及左正合性，\(F(A)\) 就是 \(F(I^0)\to F(I^1)\) 的核，即零次上同调；构造与态射相容。内射对象可取自身集中于零次的内射消解，故正次消失。投射情形用 \(P_1\to P_0\to A\to0\) 的余核及零次链消解。
\end{proof}

\cref{b4:tab:derived-conventions}集中列出两种导出的消解与次数约定，其中左、右正合性保证零次导出与原函子相同。表中的连接态射对应短正合列 \(0\to A\to B\to C\to0\)：右导出的连接升一次，适用于 \(n\ge0\)；左导出的连接降一次，适用于 \(n\ge1\)。模范畴中的 Ext 属于右导出，Tor 属于左导出；Ext 的第一变量反变，还需在相应长正合列中反转对象的顺序。

\begin{table}[htbp]
\centering
\caption{左、右导出函子的约定}
\label{b4:tab:derived-conventions}
\begin{tabular}{@{}lllll@{}}
\toprule
类型 & \(F\) 的条件 & 所选消解 & 记号 & 连接态射 \\
\midrule
右导出 & 加性、左正合 & 内射上链消解 & \(R^nF\)
& \(R^nF(C)\to R^{n+1}F(A)\) \\
左导出 & 加性、右正合 & 投射链消解 & \(L_nF\)
& \(L_nF(C)\to L_{n-1}F(A)\) \\
\bottomrule
\end{tabular}
\end{table}


\cref{b4:prop:horseshoe-functor-sequence}给出短正合列的导出长正合列。下一小节把这项性质与消解的具体选取分开表述。

\subsection{\texorpdfstring{\(\delta\)-函子与泛性质}{\AlgebraPdfDelta-函子与泛性质}}
\label{b4:subsec:0801:02}

\begin{definition}\label{b4:def:delta-functor}
设 \(\mathcal A,\mathcal B\) 是交换范畴。一族加性函子 \(T^n:\mathcal A\to\mathcal B\)，\(n\ge0\)，若对每个短正合列 \(0\to A\to B\to C\to0\) 配有态射 \(\delta^n:T^n(C)\to T^{n+1}(A)\)，使
\[
0\longrightarrow T^0(A)\longrightarrow T^0(B)\longrightarrow T^0(C)
\xrightarrow{\delta^0}T^1(A)\longrightarrow\cdots
\]
正合，且连接态射对短正合列之间的态射自然，就称 \(T^\bullet\) 连同这些连接态射为\term[delta-hanzi]{上同调 \(\delta\)-函子}[cohomological delta functor]。若 \(T^\bullet,U^\bullet:\mathcal A\to\mathcal B\) 均为此类函子，且逐次自然变换 \(a^n:T^n\to U^n\) 与每个连接态射交换，则称 \(a^\bullet\) 为它们之间的态射。
\end{definition}

对短正合列 \(0\to A\to B\to C\to0\)，\(\delta\)-函子态射 \(a^\bullet:T^\bullet\to U^\bullet\) 所要求的相容性就是下列方块交换：
\[
\begin{tikzcd}[column sep=large,row sep=large]
T^n(C)\arrow[r,"\delta_T^n"]\arrow[d,"a_C^n"']&T^{n+1}(A)\arrow[d,"a_A^{n+1}"]\\
U^n(C)\arrow[r,"\delta_U^n"']&U^{n+1}(A)
\end{tikzcd}
\]
同调版本由 \(T_n\) 和连接态射 \(T_n(C)\to T_{n-1}(A)\) 组成，长正合列在右端以 \(T_0(C)\to0\) 结束。右导出函子组成上同调 \(\delta\)-函子，左导出函子组成同调 \(\delta\)-函子。

\begin{definition}\label{b4:def:universal-delta-functor}
设 \(\mathcal A,\mathcal B\) 为交换范畴，\(T^\bullet:\mathcal A\to\mathcal B\) 为上同调 \(\delta\)-函子。若对任意上同调 \(\delta\)-函子 \(U^\bullet:\mathcal A\to\mathcal B\)，每个自然变换 \(T^0\to U^0\) 都唯一延拓为 \(\delta\)-函子态射 \(T^\bullet\to U^\bullet\)，则称 \(T^\bullet\) 为\term[fan-delta-hanzi]{泛上同调 \(\delta\)-函子}[universal cohomological delta functor]。对同调 \(\delta\)-函子 \(T_\bullet:\mathcal A\to\mathcal B\)，若对任意同调 \(\delta\)-函子 \(U_\bullet:\mathcal A\to\mathcal B\)，每个 \(U_0\to T_0\) 都唯一延拓为 \(\delta\)-函子态射 \(U_\bullet\to T_\bullet\)，则称 \(T_\bullet\) 具有同调泛性。
\end{definition}

这里方向不能省略。右导出函子从自身的零次项出发控制到其他理论的映射；左导出函子接收其他理论到自身零次项的映射。它们分别由余核式和核式的归纳构造产生。相同零次函子的两个泛 \(\delta\)-函子必自然同构：双向延拓零次恒等，再用唯一性证明两个复合都是恒等。

\subsection{\texorpdfstring{消去性与泛 \(\delta\)-函子}{消去性与泛 \AlgebraPdfDelta-函子}}
\label{b4:subsec:0801:03}

\begin{definition}\label{b4:def:effaceable}
设 \(T:\mathcal A\to\mathcal B\) 是交换范畴间的加性函子。若每个对象 \(A\) 都有单态射 \(j:A\to E\) 使 \(T(j)=0\)，则称 \(T\) \term[xiaoqvxing]{可消去}[effaceable]。若每个对象 \(A\) 都有满态射 \(P\to A\) 使其在 \(T\) 下为零，则称 \(T\) 可对偶消去。
\end{definition}

消去的是一个诱导态射，不一定要求 \(T(E)=0\)。导出函子的消去性更强：嵌入内射对象便使其正次取值本身为零。

\begin{theorem}[消去性与泛性]\label{b4:thm:derived-universal}
设 \(\mathcal A,\mathcal B\) 为交换范畴，\(T^\bullet:\mathcal A\to\mathcal B\) 为上同调 \(\delta\)-函子。若每个 \(T^n\)、\(n>0\)，均可消去，则它是泛上同调 \(\delta\)-函子。对偶地，正次均可对偶消去的同调 \(\delta\)-函子具有同调泛性。特别地，对加性函子 \(F:\mathcal A\to\mathcal B\)，若 \(\mathcal A\) 有足够多内射对象且 \(F\) 左正合，则 \(R^\bullet F\) 具有上同调泛性；若 \(\mathcal A\) 有足够多投射对象且 \(F\) 右正合，则 \(L_\bullet F\) 具有同调泛性。
\end{theorem}
\begin{proof}
给定 \(a^0:T^0\to U^0\)。归纳假定次数小于 \(n\) 的自然变换已定义，并与已有连接态射相容。对 \(A\) 取单态射 \(j:A\to E\)，使 \(T^n(j)=0\)，记 \(Q=E/A\)。正合性给出满态射
\[
\delta_T:T^{n-1}(Q)\longrightarrow T^n(A),
\qquad \ker\delta_T=\operatorname{im}\bigl(T^{n-1}(E)\to T^{n-1}(Q)\bigr).
\]
复合 \(\delta_Ua_Q^{n-1}\) 在这个核上为零：由低次自然性，其在 \(T^{n-1}(E)\) 上的值等于先用 \(a_E^{n-1}\)，再沿 \(U\) 的相邻两个态射复合，而后者为零。因此它唯一下降为
\[
a_A^n:T^n(A)\longrightarrow U^n(A),
\qquad a_A^n\delta_T=\delta_Ua_Q^{n-1}.
\]
这里 \(\delta_T\) 是满态射，所以等式唯一决定 \(a_A^n\)；任何与连接态射相容的延拓都必须取此值。

检验选择独立性。若另选 \(j':A\to E'\)，则 \(j''=(j,j'):A\to E\oplus E'\) 仍单，且被 \(T^n\) 零化。记 \(Q''=\operatorname{coker}j''\)。向 \(E\) 的投影 \(\pi_E\) 诱导下图，其中两行短正合：
\[
\begin{tikzcd}[column sep=small,row sep=large]
0\arrow[r]&A\arrow[r,"j''"]\arrow[d,"1_A"']
&E\oplus E'\arrow[r]\arrow[d,"\pi_E"]
&Q''\arrow[r]\arrow[d,"\bar\pi_E"]&0\\
0\arrow[r]&A\arrow[r,"j"']&E\arrow[r]&Q\arrow[r]&0.
\end{tikzcd}
\]
向 \(E'\) 的投影给出同样的比较图。用下标 \(j,j''\) 区分两种选择所得的候选态射。低次自然性及连接态射的自然性给出
\[
a_{A,j}^n\delta_T''=\delta_U''a_{Q''}^{n-1}
=a_{A,j''}^n\delta_T''.
\]
上行的 \(\delta_T''\) 为满态射，故 \(a_{A,j}^n=a_{A,j''}^n\)。对另一投影重复此论证，得到 \(a_{A,j'}^n=a_{A,j''}^n\)，从而构造不依赖选择。

检验对任意 \(f:A\to B\) 的自然性。选分别消去 \(T^n(A),T^n(B)\) 的 \(j_A:A\to E_A\)、\(j_B:B\to E_B\)。用 \(j_f=(j_A,j_B\circ f):A\to E_A\oplus E_B\) 作为 \(A\) 的选择，记 \(Q_f=\operatorname{coker}j_f\)、\(Q_B=\operatorname{coker}j_B\)。向 \(E_B\) 的投影诱导
\[
\begin{tikzcd}[column sep=small,row sep=large]
0\arrow[r]&A\arrow[r,"j_f"]\arrow[d,"f"']
&E_A\oplus E_B\arrow[r]\arrow[d,"\pi_B"]
&Q_f\arrow[r]\arrow[d,"\bar f"]&0\\
0\arrow[r]&B\arrow[r,"j_B"']&E_B\arrow[r]&Q_B\arrow[r]&0.
\end{tikzcd}
\]
已证的选择独立性允许用上行定义 \(a_A^n\)。将低次自然性与连接态射的自然性用于此图，两种复合满足
\[
\bigl(U^n(f)a_A^n\bigr)\delta_{T,f}
=\delta_{U,B}a_{Q_B}^{n-1}T^{n-1}(\bar f)
=\bigl(a_B^nT^n(f)\bigr)\delta_{T,f}.
\]
消去满态射 \(\delta_{T,f}\)，即得 \(U^n(f)a_A^n=a_B^nT^n(f)\)。

最后检验与任意 \(0\to A\xrightarrow{i}B\to C\to0\) 的连接态射相容。选择 \(j:B\to E\) 消去 \(T^n(B)\)，则 \(ji:A\to E\) 消去 \(T^n(A)\)。原短正合列映入 \(0\to A\to E\to E/A\to0\)，其中左端为恒等。定义式已保证后一列的相容；把它与 \(C\to E/A\) 的低次自然性复合，就得到原列的相容。归纳完成存在、自然性及唯一性。

把论证同时用于 \(\mathcal A^{\mathrm{op}}\) 与 \(\mathcal B^{\mathrm{op}}\)，单态射变满态射，余核下降变核分解，自然变换的方向也反转，遂得同调版本。最后，\cref{b4:prop:derived-degree-zero}使内射嵌入、投射满射消去所有正次导出函子，故可应用已证判据。
\end{proof}

这一泛性质确定的不仅是各个同调对象，也包括所有态射和连接态射。以后从另一种消解或另一种构造得到同一理论时，可以用零次的自然同构来识别整套长正合列。

% !TeX root = ../../Book4.tex
% 纲要标题：### 9.2 Ext
% 纲要标签：b4:sec:0802
% 纲要位置：计划纲要.md，第 3408 行。

\section{Ext}
\label{b4:sec:0802}

Hom 的两个变量具有相反的箭头方向，分别用投射或内射消解得到的计算必须比较。平衡性使两种模型计算同一组 Ext，并保留两变量的连接态射。


% 纲要内容（仅作写作计划，尚非教材正文）：
%
% * Hom 的右导出函子
% * 反变与协变变量
% * Ext 长正合列及具体计算
%

\subsection{Hom 的右导出函子}
\label{b4:subsec:0802:01}

\begin{definition}\label{b4:def:ext}
设 \(R\) 是含幺环，\(M,N\) 是幺性左 \(R\) 模。对 \(N\) 取内射消解 \(N\to I^\bullet\)，定义交换群
\[
\operatorname{Ext}_R^n(M,N)\coloneqq
H^n\bigl(\operatorname{Hom}_R(M,I^\bullet)\bigr),\qquad n\ge0.
\]
这就是左正合函子 \(\operatorname{Hom}_R(M,-)\) 的右导出函子。
\end{definition}
\symidx{\(\operatorname{Ext}_R^n(M,N)\)：Hom 函子的右导出函子值}

内射消解的开头是 \(0\to N\xrightarrow{\eta}I^0\xrightarrow{\mathrm{d}_I^0}I^1\)。因此
\[
\operatorname{Ext}_R^0(M,N)
=\ker\bigl(\operatorname{Hom}_R(M,I^0)
\xrightarrow{\mathrm{d}_I^0\circ-}\operatorname{Hom}_R(M,I^1)\bigr)
\cong\operatorname{Hom}_R(M,N).
\]
这里同构由 \(f\mapsto\eta f\) 给出：映入 \(I^0\) 后被微分零化的态射，恰好唯一经过核 \(\eta:N\hookrightarrow I^0\)。

\ProseParagraphBreak
在一般非交换环上，Ext 首先是交换群，不自动是左 \(R\) 模。即使在零次 Hom 中，试图定义 \((rf)(m)=r f(m)\) 也可能破坏 \(R\) 线性，因为
\[
(rf)(sm)=rs f(m),\qquad s(rf)(m)=sr f(m),
\]
两者未必相等。若 \(R\) 交换，这种标量乘法保持线性，并与 Hom 复形的所有微分相容，便得到自然的 \(R\) 模结构。下标的环也不能略去，例如 \(\operatorname{Ext}_{\mathbb Z}^1(\mathbb Z/2\mathbb Z,\mathbb Z/2\mathbb Z)\) 与 \(\operatorname{Ext}_{\mathbb F_2}^1(\mathbb F_2,\mathbb F_2)\) 不同。

\ProseParagraphBreak
固定 \(N\) 时，\(\operatorname{Hom}_R(-,N)\) 是从左模反范畴到交换群的左正合函子。左模的投射对象正是反范畴的内射对象；原来的投射消解
\[
\cdots\longrightarrow P_1\longrightarrow P_0\longrightarrow M\longrightarrow0
\]
在反范畴中成为 \(0\to M\to P_0\to P_1\to\cdots\)，即其中的一条内射消解。施加这个反变 Hom 后，\(\operatorname{Hom}_R(P_n,N)\) 的次数从 \(n\) 向 \(n+1\) 增加，所以它给出另一套右导出函子。两套计算是否一致，需要证明，不能由名字相同得出。

\subsection{反变与协变变量}
\label{b4:subsec:0802:02}

为了方便表示上同调类，以下把投射消解 \(P_\bullet\to M\) 送入 Hom 时，采用无符号的微分 \(\delta^n(f)=f\mathrm{d}_{P,n+1}\)。这与\cref{b4:def:hom-complex}的约定可以相互转换：若把 \(P_n\) 放在上链次数 \(-n\)，把 \(N\) 放在零次，则内部 Hom 的微分是 \(\delta_{\mathrm{Hom}}^n(f)=(-1)^{n+1}\cdot f\mathrm{d}_{P,n+1}\)。令 \(c_n=(-1)^{n(n+1)/2}\)，在次数 \(n\) 作重标度 \(\rho_n(f)=c_nf\)，则
\[
c_{n+1}=(-1)^{n+1}c_n,
\qquad
\rho_{n+1}\delta_{\mathrm{Hom}}^n(f)
=c_{n+1}(-1)^{n+1}\cdot f\mathrm{d}_{P,n+1}
=\delta^n\rho_n(f).
\]
因此 \(\rho\) 是从内部 Hom 复形到无符号复形的同构。例如 \(c_0,c_1,c_2,c_3\) 依次为 \(1,-1,-1,1\)，恰好吸收微分在相邻次数的符号。本节以后均采用无符号的计算约定。

\begin{theorem}[Ext 的平衡性]\label{b4:thm:ext-balanced}
设 \(R\) 是含幺环，\(M,N\) 是幺性左 \(R\) 模，\(P_\bullet\to M\) 为投射消解，\(N\to I^\bullet\) 为内射消解。对整数 \(n\ge0\) 取
\[
\operatorname{Hom}_R(P_\bullet,N)^n=\operatorname{Hom}_R(P_n,N),
\qquad \delta^n(f)=f\,\mathrm{d}_{P,n+1}.
\]
则有自然同构
\[
H^n\bigl(\operatorname{Hom}_R(P_\bullet,N)\bigr)
\xrightarrow[\sim]{\beta_n}\operatorname{Ext}_R^n(M,N).
\]
它对 \(M\) 反变，对 \(N\) 协变，并与两变量的连接态射相容。
\end{theorem}
\begin{proof}
两种计算将通过同一个双复形比较。令 \(D^{p,q}=\operatorname{Hom}_R(P_p,I^q)\)，其中 \(p,q\ge0\)：固定 \(p\) 沿内射消解计算，固定 \(q\) 则沿投射消解计算。反变 Hom 把 \(P_{p+1}\to P_p\) 变为 \(D^{p,q}\to D^{p+1,q}\)，所以这里的两个指标都按上链次数增加。定义
\[
\mathrm{d}_h(f)=f\mathrm{d}_{P,p+1},\qquad
\mathrm{d}_v(f)=(-1)^p\cdot\mathrm{d}_I^qf.
\]
两次水平或竖直微分为零；一水平一竖直的两种复合中，符号分别为 \((-1)^p\) 和 \((-1)^{p+1}\)，故 \(\mathrm{d}_h\mathrm{d}_v+\mathrm{d}_v\mathrm{d}_h=0\)。总次数为 \(p+q\)。对 \(n\ge0\)，\(p+q=n\) 只有 \(n+1\) 对非负整数解；对 \(n<0\) 则没有解，故总复形逐对角有限。

固定 \(p\)，\(P_p\) 的投射性使 \(\operatorname{Hom}_R(P_p,-)\) 正合。把内射消解的开头送入它，得到
\[
0\longrightarrow\operatorname{Hom}_R(P_p,N)
\longrightarrow D^{p,0}\longrightarrow D^{p,1},
\]
故纵列零次上同调为 \(\ker(D^{p,0}\to D^{p,1})\cong\operatorname{Hom}_R(P_p,N)\)，正次数上同调为零。固定 \(q\)，\(I^q\) 的内射性同样给出横行零次上同调 \(\ker(D^{0,q}\to D^{1,q})\cong\operatorname{Hom}_R(M,I^q)\)，正次数上同调为零。因此需要比较的两个 Hom 复形分别由下边缘的纵核和左边缘的横核组成，而不是双复形中整条底行或整条左列。

记 \(\eta:N\to I^0\)、\(\varepsilon:P_0\to M\) 为内射、投射消解的增广。它们分别给出复形映射
\[
\begin{aligned}
\operatorname{Hom}_R(P_\bullet,N)&\longrightarrow\operatorname{Tot}D,
&f&\longmapsto\eta f\in D^{p,0},\\
\operatorname{Hom}_R(M,I^\bullet)&\longrightarrow\operatorname{Tot}D,
&g&\longmapsto g\varepsilon\in D^{0,q}.
\end{aligned}
\]
第一条由 \(N\to I^0\) 后合成而来，\(\mathrm{d}_I^0\eta=0\) 使其像的竖直微分为零；第二条由 \(P_0\to M\) 预合成而来，\(\varepsilon\mathrm{d}_{P,1}=0\) 使其像的水平微分为零。其余微分正是边缘复形自身的微分，所以两条映射与总微分相容。将边缘复形分别放在 \(q=0\) 和 \(p=0\)，刚才的列、行上同调计算说明这两个双复形映射分别逐列、逐行拟同构。\cref{b4:cor:double-complex-comparison}的有限对角条件已经满足，故两条总化映射都是拟同构。

将 \(P_p\) 放在上链次数 \(-p\)，总复形到内部 Hom 复形的同构在分量 \(D^{p,q}\) 上为
\[
 S^{p,q}(f)=c_p(-1)^{pq}f.
\]
沿水平方向，其系数之比为 \((-1)^{p+q+1}\)；沿竖直方向，其系数之比为 \((-1)^p\)，恰好把上面的总微分变成内部 Hom 微分。左边缘 \(p=0\) 的系数为 \(1\)，下边缘 \((p,q)=(n,0)\) 的系数为 \(c_n\)。因此 \(\beta_n\) 将无符号余圈 \(b:P_n\to N\) 对应于内部 Hom 余圈 \(c_nb\)，这一系数来自两条边缘映射的比较。

投射、内射比较映射使上述两条映射对 \(M,N\) 自然；不同选择同伦，故同调态射确定。还需比较连接态射：两种模型不仅应得到同构的交换群，也应得到同一组长正合列，否则不能无歧义地使用 Ext 的连接态射。分别在两个变量中核对这一点。给定第一变量的短正合列 \(0\to A\to B\to C\to0\)，取投射马蹄消解
\[
0\longrightarrow P_{A,\bullet}\longrightarrow P_{B,\bullet}
\longrightarrow P_{C,\bullet}\longrightarrow0,
\]
它在每个次数分裂。固定 \(N\) 的内射消解，对 \(X=A,B,C\) 令 \(T_X=\operatorname{Tot}\operatorname{Hom}_R(P_{X,\bullet},I^\bullet)\)。Hom 的第一变量反变，把 \(A,B,C\) 的次序反转为 \(C,B,A\)，因此得到的短正合列方向是
\[
0\longrightarrow T_C\longrightarrow T_B\longrightarrow T_A\longrightarrow0.
\]
两侧边缘复形也分别组成 \(C,B,A\) 顺序的短正合列：投射计算一侧用马蹄消解的逐次分裂性，内射计算一侧用各 \(I^q\) 的内射性。边缘到总复形的增广映射给出短正合列之间的交换图。由\cref{b4:thm:cohomology-long-exact}的自然性，两种计算中的连接态射
\[
\operatorname{Ext}_R^n(A,N)\longrightarrow\operatorname{Ext}_R^{n+1}(C,N)
\]
在平衡同构下相互对应。

给定第二变量的短正合列 \(0\to N'\to N\to N''\to0\)，取逐次分裂的内射马蹄消解 \(0\to I_{N'}^\bullet\to I_N^\bullet\to I_{N''}^\bullet\to0\)。固定 \(M\) 的投射消解，令 \(T_V=\operatorname{Tot}\operatorname{Hom}_R(P_\bullet,I_V^\bullet)\)，其中 \(V=N',N,N''\)。这次 Hom 保持箭头方向，得到
\[
0\longrightarrow T_{N'}\longrightarrow T_N\longrightarrow T_{N''}\longrightarrow0.
\]
投射计算一侧的边缘短正合列使用各 \(P_p\) 的投射性，内射计算一侧则使用马蹄消解的逐次分裂性。增广仍给出短正合列之间的交换图，故连接态射
\[
\operatorname{Ext}_R^n(M,N'')\longrightarrow\operatorname{Ext}_R^{n+1}(M,N')
\]
也与平衡同构相容。
\end{proof}

给定 \(u:M'\to M\)，其投射消解提升经预合成给出 \(u^*:\operatorname{Ext}_R^n(M,N)\to\operatorname{Ext}_R^n(M',N)\)；给定 \(v:N\to N'\)，后合成给出 \(v_*\)。预合成与后合成交换，所以得到双函子
\[
\operatorname{Ext}_R^n:
(R\text{-}\mathbf{Mod})^{\mathrm{op}}\times R\text{-}\mathbf{Mod}
\longrightarrow\mathbf{Ab}.
\]
它同时记录第一变量的反变作用与第二变量的协变作用。

\subsection{Ext 长正合列及具体计算}
\label{b4:subsec:0802:03}

\begin{theorem}\label{b4:thm:ext-long-exact}
设 \(R\) 是含幺环。幺性左 \(R\) 模短正合列 \(0\to A\to B\to C\to0\) 对任意幺性左模 \(N\) 给出自然长正合列
\[
\begin{aligned}
0&\longrightarrow\operatorname{Hom}_R(C,N)\longrightarrow\operatorname{Hom}_R(B,N)
\longrightarrow\operatorname{Hom}_R(A,N)\\
&\longrightarrow\operatorname{Ext}_R^1(C,N)\longrightarrow\operatorname{Ext}_R^1(B,N)
\longrightarrow\operatorname{Ext}_R^1(A,N)\longrightarrow\cdots.
\end{aligned}
\]
其一般中间段为
\[
\begin{aligned}
\cdots\longrightarrow\operatorname{Ext}_R^n(C,N)
&\longrightarrow\operatorname{Ext}_R^n(B,N)
\longrightarrow\operatorname{Ext}_R^n(A,N)\\
&\longrightarrow\operatorname{Ext}_R^{n+1}(C,N)\longrightarrow\cdots
\qquad(n\ge1).
\end{aligned}
\]
对任意幺性左模 \(M\)，幺性左模短正合列 \(0\to N'\to N\to N''\to0\) 给出自然长正合列
\[
\begin{aligned}
0&\longrightarrow\operatorname{Hom}_R(M,N')\longrightarrow\operatorname{Hom}_R(M,N)
\longrightarrow\operatorname{Hom}_R(M,N'')\\
&\longrightarrow\operatorname{Ext}_R^1(M,N')\longrightarrow\operatorname{Ext}_R^1(M,N)
\longrightarrow\operatorname{Ext}_R^1(M,N'')\longrightarrow\cdots.
\end{aligned}
\]
其一般中间段为
\[
\begin{aligned}
\cdots\longrightarrow\operatorname{Ext}_R^n(M,N')
&\longrightarrow\operatorname{Ext}_R^n(M,N)
\longrightarrow\operatorname{Ext}_R^n(M,N'')\\
&\longrightarrow\operatorname{Ext}_R^{n+1}(M,N')\longrightarrow\cdots
\qquad(n\ge1).
\end{aligned}
\]
\end{theorem}
\begin{proof}
第一列取 \(N\) 的内射消解，逐次施加反变正合函子 \(\operatorname{Hom}_R(-,I^q)\)，得到 \(0\to\operatorname{Hom}(C,I)\to\operatorname{Hom}(B,I)\to\operatorname{Hom}(A,I)\to0\)，再取同调长正合列。第二列取 \(M\) 的投射消解，对给定序列逐次施加正合函子 \(\operatorname{Hom}_R(P_p,-)\)，取长正合列，再用\cref{b4:thm:ext-balanced}识别各项。增广在零次给出通常 Hom 映射，连接态射与选择独立且自然。
\end{proof}

\begin{example}\label{b4:exm:ext-cyclic}
设 \(R\) 是主理想整环，\(0\ne a\in R\)，\(N\) 是幺性 \(R\) 模。\(R/(a)\) 的两项自由消解为 \(0\to R\xrightarrow{\times a}R\to R/(a)\to0\)。送入 Hom 后得到
\[
0\longrightarrow N\xrightarrow{\times a}N\longrightarrow0,
\]
两个 \(N\) 分别位于上链次数零、一；零次上同调是乘 \(a\) 的核，一次上同调是其余核。因此
\[
\operatorname{Ext}_R^0(R/(a),N)=\{\,n\in N\mid an=0\,\},
\qquad\operatorname{Ext}_R^1(R/(a),N)=N/aN,
\]
二次及以上为零。特别地，\(\operatorname{Ext}_{\mathbb Z}^1(\mathbb Z/m\mathbb Z,\mathbb Z/n\mathbb Z)\cong\mathbb Z/\gcd(m,n)\mathbb Z\)，其中 \(m,n>0\)。若改取 \(N=\mathbb Q\)，从这个具体复形可见乘 \(m\) 为同构，所以零次及正次均为零。另一方面，\(\mathbb Q\) 的内射性也保证正次数 Ext 消失，与计算相符。
\end{example}

设 \(k\) 为域，在 \(R=k[\epsilon]/(\epsilon^2)\) 上，对\secref{b4:sec:0701}给出的 \(k=R/(\epsilon)\) 的周期消解施加 \(\operatorname{Hom}_R(-,k)\)，所有微分为零。因此 \(\operatorname{Ext}_R^n(k,k)\cong k\) 对每个 \(n\ge0\) 成立。相同的底层向量空间，在底环为 \(k\) 时正次 Ext 全为零；记录底环是计算的一部分。

% !TeX root = ../../Book4.tex
% 纲要标题：### 9.3 Tor
% 纲要标签：b4:sec:0803
% 纲要位置：计划纲要.md，第 3414 行。

\section{Tor}
\label{b4:sec:0803}

张量积的两个模必须取相反的侧别，任一变量的投射消解都应得到同一组 Tor。总复形的比较同时解释平衡性和短正合列中的连接映射。


% 纲要内容（仅作写作计划，尚非教材正文）：
%
% * 张量积的左导出函子
% * 平衡性
% * Tor 长正合列及具体计算
%

\subsection{张量积的左导出函子}
\label{b4:subsec:0803:01}

\begin{definition}\label{b4:def:tor}
设 \(R\) 是含幺环，\(M\) 是幺性右 \(R\) 模，\(N\) 是幺性左 \(R\) 模。对 \(M\) 取右模投射消解 \(P_\bullet\to M\)，定义交换群
\[
\operatorname{Tor}_n^R(M,N)\coloneqq H_n(P_\bullet\otimes_R N),\qquad n\ge0.
\]
这就是右正合函子 \(-\otimes_R N\) 的左导出函子。
\end{definition}
\symidx{\(\operatorname{Tor}_n^R(M,N)\)：张量积函子的左导出函子值}

投射消解在零次给出 \(P_1\to P_0\to M\to0\)。张量的右正合性使
\[
P_1\otimes_RN\longrightarrow P_0\otimes_RN
\longrightarrow M\otimes_RN\longrightarrow0
\]
正合，因此
\[
\operatorname{Tor}_0^R(M,N)
=\operatorname{coker}(P_1\otimes_RN\to P_0\otimes_RN)
\cong M\otimes_RN.
\]
这一零次识别只用右正合性，不要求 \(N\) 平坦。

\ProseParagraphBreak
左右侧别不可交换省略。若 \(P_\bullet\to M\) 和 \(Q_\bullet\to N\) 分别为右模、左模投射消解，下一小节的平衡性给出
\[
H_n(P_\bullet\otimes_RN)\cong H_n(M\otimes_RQ_\bullet).
\]
它没有交换 \(M,N\)，只改变在哪个变量取消解。只有在交换环上，左右模有自然识别，才能进一步比较 \(\operatorname{Tor}_n^R(M,N)\) 与 \(\operatorname{Tor}_n^R(N,M)\)，得到对称性。

\subsection{平衡性}
\label{b4:subsec:0803:02}

\begin{theorem}[Tor 的平衡性]\label{b4:thm:tor-balanced}
设 \(R\) 是含幺环，\(M\) 是幺性右 \(R\) 模，\(N\) 是幺性左 \(R\) 模。若 \(Q_\bullet\to N\) 是左 \(R\) 模投射消解，则对每个整数 \(n\ge0\)，有对两个变量自然的同构
\[
\operatorname{Tor}_n^R(M,N)\cong H_n(M\otimes_RQ_\bullet).
\]
因而可以在任意一侧取投射消解计算 Tor。
\end{theorem}
\begin{proof}
与 Ext 的平衡性一样，取两份消解组成双复形，再比较其两条边缘。这里使行、列正合的性质是投射模的平坦性。取右模投射消解 \(P_\bullet\to M\)，令 \(D_{p,q}=P_p\otimes_RQ_q\)，\(p,q\ge0\)，以
\[
\mathrm{d}_h=\mathrm{d}_P\otimes1,\qquad \mathrm{d}_v=(-1)^p1\otimes \mathrm{d}_Q
\]
为微分。它们反交换，所以总链微分平方为零。第二卷定理7.3.4证明了投射左模平坦。投射右 \(R\) 模正是投射左 \(R^{\mathrm{op}}\) 模，将该定理用于 \(R^{\mathrm{op}}\)，并经 \(X\otimes_{R^{\mathrm{op}}}P\cong P\otimes_RX\) 识别相应的张量函子，便知 \(P\otimes_R-\) 正合，即投射右模平坦。固定 \(p\) 时，\(P_p\) 是投射右 \(R\) 模，因而 \(P_p\otimes_R-\) 在左模范畴上正合；竖列只在零次有同调 \(P_p\otimes_RN\)。固定 \(q\) 时，\(Q_q\) 是投射左 \(R\) 模，因而 \(-\otimes_RQ_q\) 在右模范畴上正合；横行只在零次有同调 \(M\otimes_RQ_q\)。

记两份消解的增广为 \(\varepsilon_P:P_0\to M\)、\(\varepsilon_Q:Q_0\to N\)。第一条总化映射在 \(q=0\) 上取 \(1\otimes\varepsilon_Q\)，在 \(q>0\) 上取零；第二条在 \(p=0\) 上取 \(\varepsilon_P\otimes1\)，在 \(p>0\) 上取零。由增广零化对应的一次微分，它们与总链微分交换，给出
\[
\operatorname{Tot}D\longrightarrow P_\bullet\otimes_RN,
\qquad
\operatorname{Tot}D\longrightarrow M\otimes_RQ_\bullet.
\]
把两条边缘分别放在 \(q=0\) 和 \(p=0\)，刚才的纵列、横行同调计算说明这两条映射分别逐列、逐行拟同构。对于 \(n\ge0\)，总次数对角线 \(p+q=n\) 只有 \((0,n),(1,n-1),\ldots,(n,0)\) 共 \(n+1\) 个位置；\(n<0\) 时没有位置。因此\cref{b4:cor:double-complex-comparison}的链版本适用，两条总化映射均为拟同构。取同调并组合两同构即得结论。比较映射与增广交换，比较选择同伦，故所得同构对两变量自然。

还需要使两种计算给出相同的连接态射。对于第一变量的短正合列，取逐次分裂的右模马蹄消解，并固定第二变量的左模投射消解。三个总复形及边缘 \(P_\bullet\otimes_RN\) 的短正合列使用马蹄的逐次分裂性；另一边缘 \(M\otimes_RQ_\bullet\) 的短正合列则使用各 \(Q_q\) 平坦。对于第二变量的短正合列，改用左模马蹄消解并固定 \(P_\bullet\)，两条边缘的正合性分别来自各 \(P_p\) 平坦和马蹄的逐次分裂性。两种情况下，增广映射都给出总复形短正合列到边缘短正合列的交换图。由同调长正合列的自然性，平衡同构便与两变量的连接态射相容。
\end{proof}

若 \(R\) 交换，对左模定义右作用 \(n\cdot r\coloneqq rn\)，对右模定义左作用 \(r\cdot m\coloneqq mr\)，便可自然识别左右模；同样识别消解中的各项。这时交换两因子的映射才有如下张量积类型，并且满足平衡关系：
\[
\tau_{p,q}:P_p\otimes_RQ_q\longrightarrow Q_q\otimes_RP_p,
\qquad x\otimes y\longmapsto(-1)^{pq}y\otimes x.
\]
这些 \(\tau_{p,q}\) 合为保持总次数的映射 \(\tau\)。在交换后以 \(Q_q\) 为第一因子的总复形中，微分为 \(\mathrm{d}'=\mathrm{d}_Q\otimes1+(-1)^q1\otimes\mathrm{d}_P\)。对 \(x\in P_p\)、\(y\in Q_q\)，有
\[
\begin{aligned}
\mathrm{d}'\tau(x\otimes y)
&=(-1)^{pq}\cdot\mathrm{d}_Qy\otimes x
 +(-1)^{pq+q}y\otimes\mathrm{d}_Px,\\
\tau\mathrm{d}(x\otimes y)
&=(-1)^{(p-1)q}y\otimes\mathrm{d}_Px
 +(-1)^{p+p(q-1)}\cdot\mathrm{d}_Qy\otimes x.
\end{aligned}
\]
其中 \(p+p(q-1)=pq\)，而 \((p-1)q\) 与 \(pq+q\) 相差 \(2q\)，故两个分量的系数分别相同，得到 \(\mathrm{d}'\tau=\tau\mathrm{d}\)。因此 \(\tau\) 给出总复形同构，取同调并用平衡性，得到对两变量自然的同构 \(\operatorname{Tor}_n^R(M,N)\cong\operatorname{Tor}_n^R(N,M)\)。一般非交换环上，仅凭给定的左右模结构，不能作上述自然识别或定义这样的交换映射；平衡性仍只表示可在任一变量取消解。

\subsection{Tor 长正合列及具体计算}
\label{b4:subsec:0803:03}

\begin{theorem}\label{b4:thm:tor-long-exact}
设 \(R\) 是含幺环，\(M\) 是幺性右 \(R\) 模，\(N\) 是幺性左 \(R\) 模。幺性右模短正合列 \(0\to M'\to M\to M''\to0\) 给出
\[
\begin{aligned}
\cdots&\longrightarrow\operatorname{Tor}_1^R(M,N)
\longrightarrow\operatorname{Tor}_1^R(M'',N)
\longrightarrow M'\otimes_RN\\
&\longrightarrow M\otimes_RN\longrightarrow M''\otimes_RN\longrightarrow0,
\end{aligned}
\]
其一般中间段为
\[
\begin{aligned}
\cdots\longrightarrow\operatorname{Tor}_n^R(M',N)
&\longrightarrow\operatorname{Tor}_n^R(M,N)
\longrightarrow\operatorname{Tor}_n^R(M'',N)\\
&\xrightarrow{\delta}\operatorname{Tor}_{n-1}^R(M',N)
\longrightarrow\cdots\qquad(n\ge1).
\end{aligned}
\]
幺性左 \(R\) 模短正合列 \(0\to N'\to N\to N''\to0\) 同样给出自然长正合列，其一般中间段为
\[
\begin{aligned}
\cdots\longrightarrow\operatorname{Tor}_n^R(M,N')
&\longrightarrow\operatorname{Tor}_n^R(M,N)
\longrightarrow\operatorname{Tor}_n^R(M,N'')\\
&\xrightarrow{\delta}\operatorname{Tor}_{n-1}^R(M,N')
\longrightarrow\cdots\qquad(n\ge1).
\end{aligned}
\]
\end{theorem}
\begin{proof}
在第一变量取马蹄消解，得到逐次分裂的短正合列 \(0\to P_n'\to P_n\to P_n''\to0\)。加性函子 \(-\otimes_RN\) 保持分裂短正合列，故张量后仍为复形短正合列，可以取同调长正合列。这里保证左端单射的是分裂性；张量仅有右正合性并不足以得到它。第二变量可用平衡性改在该侧取马蹄消解，或固定右模投射消解并用各 \(P_n\) 平坦。由\cref{b4:thm:tor-balanced}证明中的连接态射相容性，两种构造对应同一长正合列；零次识别给出所列张量尾端。
\end{proof}

\begin{example}\label{b4:exm:tor-cyclic}
设 \(R\) 是主理想整环，\(0\ne a\in R\)，\(N\) 是幺性 \(R\) 模。将 \(R/(a)\) 的两项消解与 \(N\) 张量，得到链次数一、零上的 \(N\xrightarrow{\times a}N\)，故零次同调是乘 \(a\) 的余核，一次同调是其核：
\[
\operatorname{Tor}_0^R(R/(a),N)=N/aN,
\qquad \operatorname{Tor}_1^R(R/(a),N)=\{\,n\in N\mid an=0\,\},
\]
且高次为零。对 \(R=\mathbb Z\)、\(N=\mathbb Z/n\mathbb Z\)、\(a=m\)，其中 \(m,n>0\)，一次 Tor 同构于 \(\mathbb Z/\gcd(m,n)\mathbb Z\)。与\cref{b4:exm:ext-cyclic}比较，对同一个乘法映射 \(a:N\to N\)，\(\operatorname{Ext}_R^1(R/(a),N)\) 取其余核 \(N/aN\)，\(\operatorname{Tor}_1^R(R/(a),N)\) 则取其核。有限循环群算例给出同构的交换群，并不使两者成为同一构造。
\end{example}

设 \(k\) 为域，令 \(R=k[x,y]\)，并把 \(k\) 看作商模 \(R/(x,y)\)。\cref{b4:prop:koszul-resolution}给出的两变量消解中，\(P_0=R\) 记录一个生成元，\(P_1=R^2\) 的微分为 \((a,b)\mapsto xa+yb\)，记录两个关系 \(x,y\)，而 \(P_2=R\) 的微分为 \(c\mapsto(-yc,xc)\)，记录这两个关系之间的关系 \(x(-y)+yx=0\)。与 \(k\) 张量后，\(x,y\) 都作用为零，所以微分全为零，各项 \(k,k^2,k\) 直接成为同调：
\[
\operatorname{Tor}_0^R(k,k)\cong k,\qquad
\operatorname{Tor}_1^R(k,k)\cong k^2,\qquad
\operatorname{Tor}_2^R(k,k)\cong k,
\]
更高次为零。这里一次和二次 Tor 的维数分别记录了这个消解中的两个关系以及它们之间的一条关系。

% !TeX root = ../../Book4.tex
% 纲要标题：### 9.4 一次 Ext、Tor 与正合性判据
% 纲要标签：b4:sec:0804
% 纲要位置：计划纲要.md，第 3420 行。

\section{一次 Ext、Tor 与正合性判据}
\label{b4:sec:0804}

一次 Ext 或 Tor 的消失能检测投射性、内射性或平坦性，但测试对象必须遍历相应的全部模。整数模与多项式理想的例子可以把正合性的失败具体算出来。


% 纲要内容（仅作写作计划，尚非教材正文）：
%
% * Ext 检验投射性和内射性
% * Tor 检验平坦性
% * 回到第二卷反例解释失败机制
%

\subsection{Ext 与投射性、内射性}
\label{b4:subsec:0804:01}

\begin{theorem}\label{b4:thm:ext-projective-injective}
设 \(R\) 是含幺环，\(P,I\) 是幺性左 \(R\) 模。下列判据成立：
\begin{enumerate}
\item \(P\) 投射，当且仅当对所有幺性左 \(R\) 模 \(N\)，\(\operatorname{Ext}_R^1(P,N)=0\)；这也等价于对所有这样的 \(N\) 及所有 \(n>0\)，\(\operatorname{Ext}_R^n(P,N)=0\)。
\item \(I\) 内射，当且仅当对所有幺性左 \(R\) 模 \(M\)，\(\operatorname{Ext}_R^1(M,I)=0\)；这也等价于对所有这样的 \(M\) 及所有 \(n>0\)，\(\operatorname{Ext}_R^n(M,I)=0\)。
\end{enumerate}
\end{theorem}
\begin{proof}
投射或内射时用集中于零次的相应消解，正次 Ext 消失。反过来，设对所有左模 \(N\) 都有 \(\operatorname{Ext}_R^1(P,N)=0\)。任取满射 \(q:B\to C\)，令 \(A=\ker q\)，第二变量长正合列给出
\[
\operatorname{Hom}_R(P,B)\xrightarrow{q\circ-}\operatorname{Hom}_R(P,C)
\longrightarrow\operatorname{Ext}_R^1(P,A)=0.
\]
所以每个 \(P\to C\) 都是某个 \(P\to B\) 与 \(q\) 的复合，即 \(P\) 有投射提升性质。

同样，设对所有左模 \(M\) 都有 \(\operatorname{Ext}_R^1(M,I)=0\)。任取单射 \(j:A\to B\)，令 \(C=\operatorname{coker}j\)，第一变量长正合列给出
\[
\operatorname{Hom}_R(B,I)\xrightarrow{-\circ j}\operatorname{Hom}_R(A,I)
\longrightarrow\operatorname{Ext}_R^1(C,I)=0.
\]
所以每个 \(A\to I\) 都能延拓为 \(B\to I\)，即 \(I\) 内射。
\end{proof}

“对所有模”的量词有实质作用。例如 \(\operatorname{Ext}_{\mathbb Z}^1(\mathbb Z/2\mathbb Z,\mathbb Q)=0\)，并不能推出 \(\mathbb Z/2\mathbb Z\) 投射：这次消失来自第二变量 \(\mathbb Q\) 的内射性，而换成第二变量 \(\mathbb Z\) 时，Ext 就是 \(\mathbb Z/2\mathbb Z\)。又如 \(\operatorname{Ext}_{\mathbb Z}^1(\mathbb Z,\mathbb Z)=0\)，不能推出第二个 \(\mathbb Z\) 内射：这次消失来自第一变量 \(\mathbb Z\) 的投射性，换成第一变量 \(\mathbb Z/2\mathbb Z\) 就得到同一个非零 Ext。检验投射性时固定第一变量、遍历第二变量；检验内射性时则相反。

\subsection{Tor 与平坦性}
\label{b4:subsec:0804:02}

\begin{theorem}\label{b4:thm:tor-flat-criterion}
设 \(R\) 是含幺环，\(N\) 是幺性左 \(R\) 模。下列条件等价：\(N\) 平坦；对每个幺性右 \(R\) 模 \(M\)，\(\operatorname{Tor}_1^R(M,N)=0\)；对每个这样的 \(M\) 及每个 \(n>0\)，\(\operatorname{Tor}_n^R(M,N)=0\)。右模的判据以相反侧的幺性左模作测试。
\end{theorem}
\begin{proof}
若 \(N\) 平坦，\(-\otimes_RN\) 正合，故右模投射消解在张量后仍在正次数正合；于是所有正次 Tor 为零。全消失当然蕴含一次消失。若一次均消失，对任意右模单射 \(A\to B\)，令 \(C=B/A\)。长正合列的片段
\[
\operatorname{Tor}_1^R(C,N)\longrightarrow A\otimes_RN
\longrightarrow B\otimes_RN
\]
表明后一个箭头单射。对任意右模短正合列 \(0\to A\to B\to C\to0\)，张量积的右正合性已经保证 \(A\otimes_RN\to B\otimes_RN\to C\otimes_RN\to0\) 正合；缺少的只是最左边的单射性。现在这一步也成立，故 \(-\otimes_RN\) 正合，即 \(N\) 平坦。反环给出右模版本。
\end{proof}

不必为了判定平坦性计算所有高次 Tor；一次已经足够。相反，只对某一个模计算一次 Tor，通常只能发现反例，不能得出整体平坦性。

\subsection{非正合性的模论反例}
\label{b4:subsec:0804:03}

设 \(m>1\) 为整数。把 \(0\to\mathbb Z\xrightarrow{\times m}\mathbb Z\to\mathbb Z/m\mathbb Z\to0\) 与 \(\mathbb Z/m\mathbb Z\) 作张量，乘 \(m\) 的映射变为零。由于 \(\mathbb Z\) 投射，相应长正合列的片段是
\[
0\longrightarrow\operatorname{Tor}_1^{\mathbb Z}(\mathbb Z/m\mathbb Z,\mathbb Z/m\mathbb Z)
\longrightarrow\mathbb Z/m\mathbb Z\xrightarrow{0}\mathbb Z/m\mathbb Z.
\]
零映射的核是整个 \(\mathbb Z/m\mathbb Z\)，故 \(\operatorname{Tor}_1^{\mathbb Z}(\mathbb Z/m\mathbb Z,\mathbb Z/m\mathbb Z)\cong\mathbb Z/m\mathbb Z\)，它记录了张量后丢失的单射性。对同一列施加 \(\operatorname{Hom}_{\mathbb Z}(-,\mathbb Z)\)，得到
\[
\mathbb Z\xrightarrow{\times m}\mathbb Z
\longrightarrow\operatorname{Ext}_{\mathbb Z}^1(\mathbb Z/m\mathbb Z,\mathbb Z)
\longrightarrow0.
\]
乘 \(m\) 的映射不是满射，其余核就是 \(\operatorname{Ext}_{\mathbb Z}^1(\mathbb Z/m\mathbb Z,\mathbb Z)\cong\mathbb Z/m\mathbb Z\)。对同一个整数短正合列，张量积丢失的单射性由 Tor 的核项记录，反变 Hom 丢失的满射性由 Ext 的余核项记录。

\ProseParagraphBreak
再取域 \(k\) 上的多项式环 \(R=k[x,y]\) 及其理想 \(\mathfrak m=(x,y)\)，将 \(k\) 识别为 \(R/\mathfrak m\)。因为 \(R\) 交换，以下将这些模同时视为右 \(R\) 模。短正合列 \(0\to\mathfrak m\to R\to k\to0\) 在第一变量的 Tor 长正合列中给出
\[
\underbrace{\operatorname{Tor}_2^R(R,k)}_{=0}\longrightarrow
\operatorname{Tor}_2^R(k,k)\longrightarrow\operatorname{Tor}_1^R(\mathfrak m,k)
\longrightarrow\underbrace{\operatorname{Tor}_1^R(R,k)}_{=0}.
\]
两端因 \(R\) 投射而消失，再代入上一节算出的二次 Tor，得到
\[
\operatorname{Tor}_1^R(\mathfrak m,k)\cong
\operatorname{Tor}_2^R(k,k)\cong k.
\]
故 \(\mathfrak m\) 不平坦，尽管它是整环中的无挠子模。这具体说明主理想整环上的无挠模平坦这一结论不能直接推广到任意整环。上面的连接态射已经将二次 Tor 与核模的一次 Tor 对应；下一节说明何时能这样降低次数。

% !TeX root = ../../Book4.tex
% 纲要标题：### 9.5 维数平移与合冲模初步
% 纲要标签：b4:sec:0805
% 纲要位置：计划纲要.md，第 3426 行。

\section{维数平移与合冲模初步}
\label{b4:sec:0805}

短正合列中的投射或内射项消去高次导出群，使相邻次数之间出现自然同构。逐次取合冲模或余合冲模，便把高次计算移到更低次数。


% 纲要内容（仅作写作计划，尚非教材正文）：
%
% * 对短正合列 \(0\to\Omega M\to P\to M\to0\) 及投射模 \(P\)，由长正合列推出正次数 Ext 和 Tor 的维数平移，并注明变量侧别
% * 对内射嵌入 \(0\to N\to I\to C\to0\) 给出对偶的 Ext 平移公式；区分合冲模选择与导出不变量的自然同构
% * 迭代平移与截短消解，说明在有限步后出现投射合冲模的含义；不在此预设一般同调维数理论
% * 本节是第三卷第15–21章所用同调工具包的一部分；第11.1–11.2节在其基础上统一定义和研究各类同调维数
%
% ---
%

\subsection{投射消解中的 Ext 与 Tor 维数平移}
\label{b4:subsec:0805:01}

\begin{definition}\label{b4:def:syzygy-basic}
设 \(R\) 是含幺环，\(M\) 是幺性左 \(R\) 模。给定投射满射 \(P\to M\)，其核记为 \(\Omega M\)，称为与这一选择对应的第一\term[hechongmo]{合冲模}[syzygy module]。在固定投射消解 \(P_\bullet\xrightarrow{\varepsilon}M\) 中，令 \(\Omega^0M=M\)、\(\pi_0=\varepsilon:P_0\to M\)，并递归定义
\[
\Omega^{r+1}M=\ker\pi_r,\qquad \pi_r:P_r\twoheadrightarrow\Omega^rM.
\]
这里 \(r\ge1\) 时，消解的正合性给出 \(\operatorname{im}\mathrm{d}_r=\Omega^rM\subseteq P_{r-1}\)，所以微分 \(\mathrm{d}_r:P_r\to P_{r-1}\) 唯一因子化为满射 \(\pi_r\) 与包含映射的复合。这就确定了递归中使用的映射。右模用右投射消解作同一定义。
\end{definition}
\symidx{\(\Omega^rM\)：相对于选定投射消解的第 \(r\) 次合冲模}

\begin{theorem}[基础维数平移]\label{b4:thm:dimension-shift-basic}
设 \(R\) 是含幺环，\(0\to K\to P\to M\to0\) 是幺性左 \(R\) 模短正合列，且 \(P\) 投射。则对任意幺性左 \(R\) 模 \(N\)、幺性右 \(R\) 模 \(A\) 及整数 \(i\ge1\)，连接态射给出同构
\[
\operatorname{Ext}_R^i(K,N)\cong\operatorname{Ext}_R^{i+1}(M,N),
\qquad
\operatorname{Tor}_{i+1}^R(A,M)\cong\operatorname{Tor}_i^R(A,K).
\]
Ext 同构对 \(N\) 及所给短正合列之间的态射自然；Tor 同构对 \(A\) 及所给短正合列之间的态射自然。低次则为正合列
\[
\begin{aligned}
0&\longrightarrow\operatorname{Hom}_R(M,N)\longrightarrow\operatorname{Hom}_R(P,N)
\longrightarrow\operatorname{Hom}_R(K,N)\\
&\longrightarrow\operatorname{Ext}_R^1(M,N)\longrightarrow0,\\
0&\longrightarrow\operatorname{Tor}_1^R(A,M)\longrightarrow A\otimes_RK
\longrightarrow A\otimes_RP\\
&\longrightarrow A\otimes_RM\longrightarrow0.
\end{aligned}
\]
若原列为幺性右 \(R\) 模短正合列，则 Tor 结论改为 \(\operatorname{Tor}_{i+1}^R(M,N)\cong\operatorname{Tor}_i^R(K,N)\)，其中 \(N\) 为幺性左 \(R\) 模；该同构对 \(N\) 及所给右模短正合列之间的态射自然。
\end{theorem}
\begin{proof}
第一变量 Ext 长正合列在相关位置为
\[
\operatorname{Ext}_R^i(P,N)\longrightarrow\operatorname{Ext}_R^i(K,N)
\longrightarrow\operatorname{Ext}_R^{i+1}(M,N)
\longrightarrow\operatorname{Ext}_R^{i+1}(P,N).
\]
\(i\ge1\) 时两端由投射性消失。Tor 的相应片段是
\[
\operatorname{Tor}_{i+1}^R(A,P)\longrightarrow\operatorname{Tor}_{i+1}^R(A,M)
\longrightarrow\operatorname{Tor}_i^R(A,K)\longrightarrow\operatorname{Tor}_i^R(A,P),
\]
两端由投射蕴含平坦而消失。当 \(i=0\) 时，Ext 片段左端是 \(\operatorname{Hom}_R(P,N)\)，Tor 片段右端是 \(A\otimes_RP\)，它们一般不为零，因此连接态射未必是同构；保留这些零次项才得到所列低次正合列。自然性来自原长正合列对短正合列态射的自然性；右模情形把同一证明用于另一变量。
\end{proof}

这一操作称为\term[weishu-pingyi]{维数平移}[dimension shifting]。一条短正合列只使次数降低一步；逐次取合冲模并重复这一操作，才能把高次问题降到一次。零次仍须使用包含 Hom 或张量项的正合列。

\subsection{内射消解中的 Ext 平移与自然性}
\label{b4:subsec:0805:02}

设 \(R\) 是含幺环，\(N\) 是幺性左 \(R\) 模。给定到内射模的单射 \(j:N\to I\)，其余核 \(C=I/j(N)\) 称为与这一选择对应的第一\term[yuhechongmo]{余合冲模}[cosyzygy module]。在内射消解 \(0\to N\to I^0\xrightarrow{\mathrm{d}^0}I^1\xrightarrow{\mathrm{d}^1}\cdots\) 中，第一次取商给出
\[
C^1=I^0/N\cong\operatorname{im}\mathrm{d}^0=\ker \mathrm{d}^1\subseteq I^1.
\]
因此 \(\mathrm{d}^0\) 在这个商上诱导单射 \(C^1\to I^1\)，再取 \(C^2=I^1/C^1\)，就有短正合列 \(0\to C^1\to I^1\to C^2\to0\)。这里将 \(N\) 与其在 \(I^0\) 中的像识别，并将 \(C^1\) 与上述单射的像识别。不断取这样的余核，与投射消解中不断取核相对偶。

\begin{proposition}\label{b4:prop:injective-dimension-shift}
设 \(R\) 是含幺环，\(0\to N\to I\to C\to0\) 为幺性左 \(R\) 模短正合列，\(I\) 内射。对每个幺性左 \(R\) 模 \(M\) 和整数 \(i\ge1\)，有同构
\[
\operatorname{Ext}_R^i(M,C)\cong\operatorname{Ext}_R^{i+1}(M,N).
\]
该同构对 \(M\) 及所给短正合列之间的态射自然。零次给出
\[
\operatorname{Ext}_R^1(M,N)\cong
\operatorname{coker}\bigl(\operatorname{Hom}_R(M,I)\to\operatorname{Hom}_R(M,C)\bigr).
\]
\end{proposition}
\begin{proof}
第二变量长正合列在正次的片段为
\[
\operatorname{Ext}_R^i(M,I)\longrightarrow\operatorname{Ext}_R^i(M,C)
\longrightarrow\operatorname{Ext}_R^{i+1}(M,N)
\longrightarrow\operatorname{Ext}_R^{i+1}(M,I).
\]
\(i\ge1\) 时两端由内射性消失，故中间连接态射为同构；零次则留下
\[
\begin{aligned}
0&\longrightarrow\operatorname{Hom}_R(M,N)\longrightarrow\operatorname{Hom}_R(M,I)
\longrightarrow\operatorname{Hom}_R(M,C)\\
&\longrightarrow\operatorname{Ext}_R^1(M,N)\longrightarrow0,
\end{aligned}
\]
从而得到所述余核表达。连接态射的自然性来自长正合列的自然性。
\end{proof}

余合冲模 \(C\) 取决于所选内射嵌入，正如 \(\Omega M\) 取决于所选投射满射。这里的自然性所比较的是给定短正合列及其态射；更换消解时，比较定理提供相应的比较映射。

\begin{lemma}[Schanuel 引理]\label{b4:lem:schanuel}
设 \(R\) 是含幺环，\(0\to K\to P\to M\to0\)、\(0\to K'\to P'\to M\to0\) 为两个幺性左 \(R\) 模短正合列。若中间模 \(P,P'\) 都投射，则 \(K\oplus P'\cong K'\oplus P\)。
\end{lemma}
\begin{proof}
记两个满射为 \(\pi:P\to M\)、\(\pi':P'\to M\)，取纤维积
\[
E=P\times_M P'=\{(p,p')\in P\oplus P':\pi(p)=\pi'(p')\}.
\]
任取 \(p'\in P'\)，由 \(\pi\) 满射，可取 \(p\in P\) 使 \(\pi(p)=\pi'(p')\)，于是 \((p,p')\in E\) 映到 \(p'\)。故第二投影 \(E\to P'\) 满射，其核为 \(\{(p,0):\pi(p)=0\}\cong K\)；同理第一投影 \(E\to P\) 满射，其核为 \(\{(0,p'):\pi'(p')=0\}\cong K'\)。这样得到短正合列
\[
0\longrightarrow K\longrightarrow E\longrightarrow P'\longrightarrow0,
\qquad
0\longrightarrow K'\longrightarrow E\longrightarrow P\longrightarrow0.
\]
因右端投射，两列均分裂，故 \(E\cong K\oplus P'\cong K'\oplus P\)。
\end{proof}

这个\term[schanuel-yinli]{Schanuel 引理}[Schanuel's lemma]说明两种合冲模在加上投射直和因子后同构。\footnote{沙努埃尔（Stephen Hoel Schanuel，1933-2014），美国数学家，研究同调代数、范畴论与数论。}这不能推出 \(K\cong K'\)，两者甚至可能不同构。例如在非零环 \(R\) 上，自由模 \(M\) 可以取恒等满射，核为零；也可以取投影 \(M\oplus R\to M\)，核为非零的 \(R\)。此外，两条短正合列的分裂一般不唯一，所得直和同构也依赖于分裂的选择。

\subsection{迭代平移与截短消解}
\label{b4:subsec:0805:03}

\begin{proposition}\label{b4:prop:iterated-dimension-shift}
设 \(R\) 是含幺环，\(M,N\) 是幺性左 \(R\) 模，\(A\) 是幺性右 \(R\) 模。固定 \(M\) 的投射消解 \(P_\bullet\to M\) 及其 \(r\) 次合冲模 \(\Omega^rM\)，其中 \(r\ge0\) 为整数。则对整数 \(i\ge1\)，有
\[
\operatorname{Ext}_R^i(\Omega^rM,N)\cong\operatorname{Ext}_R^{i+r}(M,N),
\qquad
\operatorname{Tor}_i^R(A,\Omega^rM)\cong\operatorname{Tor}_{i+r}^R(A,M).
\]
若固定 \(N\) 的内射消解 \(N\to I^\bullet\)，令 \(C^0=N\)、\(C^{j+1}=\operatorname{coker}(C^j\to I^j)\)（\(j\ge0\)）。这里 \(C^0\to I^0\) 为增广单射；对 \(j\ge1\)，内射消解的正合性使 \(I^{j-1}\to I^j\) 在前一步的商 \(C^j\) 上诱导单射，其像为 \(\ker(I^j\to I^{j+1})\)。则
\[
\operatorname{Ext}_R^i(M,C^r)\cong\operatorname{Ext}_R^{i+r}(M,N).
\]
投射平移的两个同构对另一变量及所给投射消解之间的增广链映射自然；内射平移的同构对 \(M\) 及所给内射消解之间的增广上链映射自然。此外，固定消解中的 \(\Omega^rM\) 投射，当且仅当对所有幺性左 \(R\) 模 \(N\)，\(\operatorname{Ext}_R^{r+1}(M,N)=0\)。
\end{proposition}
\begin{proof}
对各短正合列 \(0\to\Omega^{j+1}M\to P_j\to\Omega^jM\to0\) 依次用\cref{b4:thm:dimension-shift-basic}。例如两步给出
\[
\operatorname{Ext}_R^{i+2}(M,N)\cong\operatorname{Ext}_R^{i+1}(\Omega M,N)
\cong\operatorname{Ext}_R^i(\Omega^2M,N).
\]
每次将次数降低一，同时将第一变量换成下一次合冲模；复合 \(r\) 次连接同构即得第一式。Tor 用相同的短正合列迭代，内射式则对 \(0\to C^j\to I^j\to C^{j+1}\to0\) 反复用\cref{b4:prop:injective-dimension-shift}。增广链映射在各核上诱导映射，增广上链映射在各余核上诱导映射；这些映射组成短正合列之间的态射，故每一步的连接同构及其复合都自然。

最后，对每个幺性左 \(R\) 模 \(N\)，第一式在 \(i=1\) 时成为
\[
\operatorname{Ext}_R^1(\Omega^rM,N)\cong\operatorname{Ext}_R^{r+1}(M,N).
\]
因此右边对所有这样的 \(N\) 消失，当且仅当左边也对所有 \(N\) 消失，而\cref{b4:thm:ext-projective-injective}说明后者当且仅当 \(\Omega^rM\) 投射。
\end{proof}

若 \(r=0\)，这一条件就是 \(M\) 本身投射。若 \(r\ge1\) 且 \(\Omega^rM\) 投射，原消解可截成
\[
0\longrightarrow\Omega^rM\longrightarrow P_{r-1}\longrightarrow\cdots
\longrightarrow P_0\longrightarrow M\longrightarrow0.
\]
这里把 \(\Omega^rM\) 作为第 \(r\) 次投射项，所谓长度不超过 \(r\)，指投射项的最高非零次数不超过 \(r\)，不把增广目标 \(M\) 计作投射项。例如长度零表示 \(M\) 投射；长度不超过一表示存在 \(0\to P_1\to P_0\to M\to0\)，其中 \(P_0,P_1\) 均投射。反过来，若有长度不超过 \(r\) 的投射消解，施加 \(\operatorname{Hom}_R(-,N)\) 后，第 \(r+1\) 次项为零，所以对每个 \(N\) 都有 \(\operatorname{Ext}_R^{r+1}(M,N)=0\)。由上述判据，任意投射消解的第 \(r\) 次合冲模都投射。高次 Ext 的全称消失与投射消解能在第 \(r\) 次截短因此等价；\chapref{b4:ch:10}再用它定义并比较各种同调维数。

\subsection{维数平移在局部代数中的应用}
\label{b4:subsec:0805:04}

第三卷的局部代数以各项有限秩的自由消解记录生成元和关系。这里已经得到其中所需的基本同调推理：把 \(0\to\Omega M\to P\to M\to0\) 送入 Ext 或 Tor 的长正合列，利用 \(P\) 的正次导出函子消失，把次数降低一步。有限性和局部性在选择极小自由消解、利用 Nakayama 引理\footnote{中山正（假名：なかやま ただし，罗马音：Tadashi Nakayama，1912-1964），日本数学家，研究环论与表示论。}时另有作用；上述维数平移本身对任意含幺环成立。

\begin{example}\label{b4:exm:local-shift-koszul}
设 \(k\) 是域，\(R=k[x,y]_{(x,y)}\)，\(\mathfrak m=(x,y)R\)，剩余域记为 \(k\)。由于 \(R\) 交换，以下各左模也通过 \(m\cdot r=r\cdot m\) 视为右 \(R\) 模。\(R\) 是整环，\(R/(x)\cong k[y]_{(y)}\) 也是整环，所以 \(x,y\) 满足\cref{b4:prop:koszul-resolution}的假设，给出自由消解
\[
0\longrightarrow R\xrightarrow{\binom{-y}{x}}R^2
\xrightarrow{(x\;y)}R\longrightarrow k\longrightarrow0.
\]
与 \(k\) 张量后，两个微分映射均为零，故 \(\operatorname{Tor}_2^R(k,k)\cong k\)。短正合列 \(0\to\mathfrak m\to R\to k\to0\) 遂给出 \(\operatorname{Tor}_1^R(\mathfrak m,k)\cong k\ne0\)。由 Tor 的平坦判据，\(\mathfrak m\) 不平坦；投射模必平坦，故它也不投射。

\ProseParagraphBreak
另一方面，在这条消解中，\(\Omega k=\operatorname{im}(R^2\to R)=\mathfrak m\)，而正合性给出
\[
\Omega^2k=\ker\bigl(R^2\xrightarrow{(x\;y)}R\bigr)
=\operatorname{im}\bigl(R\xrightarrow{\binom{-y}{x}}R^2\bigr)\cong R.
\]
最后的同构使用了最左边微分的单射性，所以 \(\Omega^2k\) 投射，而 \(\Omega k\) 不投射。这同时给出长度二的自由消解和阻止在第一步截短的一次 Tor。
\end{example}


\begin{chaptersummary}
导出函子先由消解定义，比较定理保证选择独立，马蹄引理给出长正合列，消去性进一步确定它的泛性质。Ext 对第一变量反变、第二变量协变；Tor 对两个变量协变且要求相反的模侧别。平衡性由同一个双复形的两种计算证明，因此可以选择较容易处理的变量作消解。

\ProseParagraphBreak
一次 Ext 与 Tor 足以分别检验投射、内射和平坦性，但测试变量必须遍历所有相应模。维数平移利用投射或内射项的高次消失，把高次问题逐步移到合冲模或商对象。零次仍含 Hom 或张量项，不能套用正次数同构。这些工具既用于下一章的扩张解释，也用于后续同调维数及第三卷的局部自由消解。
\end{chaptersummary}

\BookExercises

\begin{exercise}\label{b4:ex:08-derived-exactness}
设 \(\mathcal A\) 有足够多内射对象，\(F:\mathcal A\to\mathcal B\) 为交换范畴间加性左正合函子。证明 \(F\) 正合当且仅当 \(R^1F(A)=0\) 对每个 \(A\) 成立，并说明此时所有正次均消失。
\end{exercise}
\begin{solution}
若 \(F\) 正合，内射消解的增广正合性经 \(F\) 保留，故所有正次上同调为零。若一次全消失，对任意 \(0\to A\to B\to C\to0\)，长正合列使 \(F(B)\to F(C)\) 满；结合已有左正合性，\(F\) 正合。最后再用正向结论即得全消失。
\end{solution}

\begin{exercise}\label{b4:ex:08-derived-natural-transformation}
设 \(F,G:\mathcal A\to\mathcal B\) 是加性左正合函子，\(\mathcal A\) 有足够多内射对象，\(a:F\to G\) 自然。用同一内射消解构造 \(R^na:R^nF\to R^nG\)，并证明它与泛性给出的延拓相同。
\end{exercise}
\begin{solution}
对 \(A\to I_A^\bullet\)，各分量 \(a_{I_A^q}\) 因自然性与微分交换，故形成复形映射。取上同调，利用比较映射的自然性得到自然变换 \(R^na\)。对马蹄消解施加 \(a\) 给出短正合复形列之间的映射，故它与连接态射相容。零次通过核的识别正是 \(a_A\)，所以由\cref{b4:thm:derived-universal}的唯一性，等于泛延拓。
\end{solution}

\begin{exercise}\label{b4:ex:08-central-actions}
设 \(R\) 是交换环，\(M,N\) 为 \(R\) 模。证明 \(r\in R\) 通过 \(M\) 的乘法映射或 \(N\) 的乘法映射诱导的 Ext 自同态相同。对 Tor 也证明相应结论。
\end{exercise}
\begin{solution}
对 \(M\) 的投射消解，逐次乘 \(r\) 是乘法态射的提升。在 \(\operatorname{Hom}_R(P_n,N)\) 上，预合成 \(r\) 与后合成 \(r\) 相同，因为 \(f(rp)=r\cdot f(p)\)。因此在上同调上相同。Tor 的消解张量中有 \(rp\otimes n=p\otimes rn\)，故同样成立。若 \(rM=0\) 或 \(rN=0\)，对应零态射诱导零，因此该 \(r\) 零化所有相关 Ext 与 Tor。
\end{solution}

\begin{exercise}\label{b4:ex:08-ext-sum-product}
设 \(R\) 是含幺环，\((M_i)_{i\in I}\)、\(N\) 是左模。证明对每个 \(n\ge0\)，有
\[
\operatorname{Ext}_R^n(\bigoplus_iM_i,N)\cong\prod_i\operatorname{Ext}_R^n(M_i,N).
\]
指出无限指标处所用的集合论选择。
\end{exercise}
\begin{solution}
为各 \(M_i\) 选择投射消解后，逐次直和仍正合。其每项也投射：给定满射 \(X\twoheadrightarrow Y\) 及 \(\bigoplus_iP_i\to Y\)，先为每个 \(P_i\to Y\) 选择一个到 \(X\) 的提升，再由直和的泛性质合成提升。无限指标时，这里要从一族非空的提升集合中作选择。因此这些消解的直和是 \(\bigoplus_iM_i\) 的投射消解，Hom 将它变为各 Hom 复形的直积。

交换群的直积中核逐分量计算；直积也保持满射，但这是另一处选择：对一族满射 \(X_i\twoheadrightarrow Y_i\) 和给定的 \((y_i)\)，为每个 \(y_i\) 选择原像 \(x_i\)，才能组成积中的原像。因而积保持短正合列，复形的圈与边界也均可逐分量计算，得到
\[
H^n\!\left(\prod_i C_i^\bullet\right)
\cong\prod_i H^n(C_i^\bullet).
\]
将此式用于上述 Hom 复形即得结论。
\end{solution}

\begin{exercise}\label{b4:ex:08-ext-map-cyclic}
设 \(m,n\) 为正整数且 \(n\mid m\)，取商映射 \(q:\mathbb Z/m\mathbb Z\to\mathbb Z/n\mathbb Z\)。对交换群 \(N\)，在 \(\operatorname{Ext}_{\mathbb Z}^1(\mathbb Z/r\mathbb Z,N)\cong N/rN\) 的识别下，求 \(q^*\)。
\end{exercise}
\begin{solution}
分别用 \(\mathbb Z\xrightarrow{m}\mathbb Z\) 和 \(\mathbb Z\xrightarrow{n}\mathbb Z\) 作源、目标的自由消解。商映射 \(q\) 在零次提升为 \(f_0=\operatorname{id}_{\mathbb Z}\)，一次提升为 \(f_1=(m/n)\operatorname{id}_{\mathbb Z}\)；链映射条件正是
\[
n\circ f_1=f_0\circ m,\qquad n(m/n)=m.
\]
对这条链映射作 Hom，箭头反向，一次映射由预合成 \(f_1\) 给出，即乘 \(m/n\)。因此
\[
q^*:N/nN\longrightarrow N/mN,\qquad [a]\longmapsto[(m/n)a].
\]
若 \(a\) 改变 \(nb\)，像改变 \(mb\)，故良定义。这个公式表明 Ext 的反变方向不是简单地沿自然取商。
\end{solution}

\begin{exercise}\label{b4:ex:08-tor-map-cyclic}
仍设 \(m,n\) 为正整数且 \(n\mid m\)，对交换群 \(N\) 求上述 \(q\) 诱导的 \(\operatorname{Tor}_1^{\mathbb Z}(\mathbb Z/m\mathbb Z,N)\to\operatorname{Tor}_1^{\mathbb Z}(\mathbb Z/n\mathbb Z,N)\)。
\end{exercise}
\begin{solution}
沿前一题的链映射作张量，箭头保持原方向，一次分量仍为乘 \(m/n\)。在一次同调即乘法映射的核上，诱导映射为
\[
\{\,a\in N\mid ma=0\,\}\longrightarrow\{\,b\in N\mid nb=0\,\},
\qquad a\longmapsto(m/n)a.
\]
\(n(m/n)a=ma=0\)，确实进入目标。同一条消解链映射在这里给出 \(N[m]\to N[n]\)，而前一题中给出 \(N/nN\to N/mN\)：方向的区别来自张量的协变性与 Hom 第一变量的反变性。
\end{solution}

\begin{exercise}\label{b4:ex:08-integer-four-six}
计算 \(\operatorname{Hom}_{\mathbb Z}(\mathbb Z/4\mathbb Z,\mathbb Z/6\mathbb Z)\)、\(\operatorname{Ext}_{\mathbb Z}^1(\mathbb Z/4\mathbb Z,\mathbb Z/6\mathbb Z)\)、\(\operatorname{Tor}_1^{\mathbb Z}(\mathbb Z/4\mathbb Z,\mathbb Z/6\mathbb Z)\)，并分别写出具体代表。
\end{exercise}
\begin{solution}
三者的抽象群均为 \(\mathbb Z/2\mathbb Z\)。Hom 由 \(\bar1\) 的像 \(0\) 或 \(3\) 决定；Tor 是 \(\mathbb Z/6\mathbb Z\) 中乘四的核 \(\{0,3\}\)；Ext 是 \((\mathbb Z/6\mathbb Z)/4(\mathbb Z/6\mathbb Z)\)，其子群 \(4(\mathbb Z/6\mathbb Z)=\{0,2,4\}\)，可用 \(0,1\) 代表两个类。核中的元素和余核中的类承担不同角色。
\end{solution}

\begin{exercise}\label{b4:ex:08-tor-conormal}
设 \(R\) 是交换环，\(I\) 是理想。证明有典范同构 \(\operatorname{Tor}_1^R(R/I,R/I)\cong I/I^2\)，无需假定 \(I\) 平坦。
\end{exercise}
\begin{solution}
对 \(0\to I\to R\to R/I\to0\) 在另一变量取 \(R/I\)，长正合列及 \(R\) 投射给出
\[
0\longrightarrow\operatorname{Tor}_1^R(R/I,R/I)
\longrightarrow I\otimes_RR/I\longrightarrow R\otimes_RR/I.
\]
最后一个映射在纯张量上将 \(a\otimes\bar r\)（\(a\in I\)）送到 \(R\otimes_RR/I\) 中的同名张量，而
\[
a\otimes\bar r=1\otimes\overline{ar}=0,
\]
故整个映射为零。张量商同构 \(I\otimes_RR/I\to I/I^2\) 将 \(a\otimes\bar r\) 送到 \(ar+I^2\)，所以连接态射与此同构的复合给出所求典范同构。
\end{solution}

\begin{exercise}\label{b4:ex:08-truncated-cubic-ring}
设 \(k\) 是域，令 \(R=k[t]/(t^3)\)，并将 \(k\) 识别为 \(R/(t)\)。写出 \(k\) 的周期自由消解，计算 \(\operatorname{Ext}_R^n(k,k)\)、\(\operatorname{Tor}_n^R(k,k)\) 以及 \(\operatorname{Ext}_R^n(k,R)\) 的正次数。
\end{exercise}
\begin{solution}
取各 \(P_i=R\)，消解为
\[
\cdots\xrightarrow{t}R\xrightarrow{t^2}R
\xrightarrow{t}R\xrightarrow{t^2}R\xrightarrow{t}R
\longrightarrow k\longrightarrow0.
\]
从 \(P_1\to P_0\) 起，微分交替为乘 \(t\)、乘 \(t^2\)。由于
\[
\begin{aligned}
\ker(t:R\to R)&=(t^2)=\operatorname{im}(t^2:R\to R),\\
\ker(t^2:R\to R)&=(t)=\operatorname{im}(t:R\to R),
\end{aligned}
\]
每个正次数的位置均正合，零次的核则为 \((t)\)。以 \(k\) 作 Hom 目标或张量因子时，\(t\) 的作用为零，故微分全为零，得到
\[
\operatorname{Ext}_R^n(k,k)\cong k,\qquad
\operatorname{Tor}_n^R(k,k)\cong k\quad(n\ge0).
\]
以 \(R\) 作 Hom 目标时，上链微分仍交替为乘 \(t\)、乘 \(t^2\)，上述核像等式使所有正次 Ext 为零；零次是 \(\ker(t:R\to R)=(t^2)\cong k\)。
\end{solution}

\begin{exercise}\label{b4:ex:08-one-module-test-fails}
给出一个非投射模 \(M\) 与一个非零模 \(N\)，使 \(\operatorname{Ext}^n(M,N)=0\) 对全部 \(n>0\) 成立；再给出非内射 \(N\) 和非零 \(M\) 的相应例子。
\end{exercise}
\begin{solution}
在 \(\mathbb Z\) 上取 \(M=\mathbb Z/2\mathbb Z\)、\(N=\mathbb Q\)，第二变量内射使正次全消失，而 \(M\) 不是投射。再取 \(M=N=\mathbb Z\)，第一变量投射使正次全消失，\(N\) 却不内射，因为 \(2\mathbb Z\to\mathbb Z\)、\(2a\mapsto a\) 无法延拓。两个例子都说明判据中的全称测试不能缩成任意一次测试。
\end{solution}

\begin{exercise}\label{b4:ex:08-flat-kernel}
设 \(R\) 是含幺环，\(0\to A\to B\to C\to0\) 为左模短正合列，\(B,C\) 平坦。证明 \(A\) 平坦。
\end{exercise}
\begin{solution}
对任意右模 \(M\)，长正合列含有 \(\operatorname{Tor}_2^R(M,C)\to\operatorname{Tor}_1^R(M,A)\to\operatorname{Tor}_1^R(M,B)\)。两端因平坦而为零，故中间为零。再由\cref{b4:thm:tor-flat-criterion}得到 \(A\) 平坦。
\end{solution}

\begin{exercise}\label{b4:ex:08-extension-flat}
设 \(0\to A\to B\to C\to0\) 是任意含幺环上的左模短正合列，\(A,C\) 平坦。证明 \(B\) 平坦。若 \(A,B\) 平坦，结论能否改为 \(C\) 平坦？
\end{exercise}
\begin{solution}
对任意右模 \(M\)，一次 Tor 片段 \(\operatorname{Tor}_1(M,A)\to\operatorname{Tor}_1(M,B)\to\operatorname{Tor}_1(M,C)\) 两端为零，故 \(B\) 平坦。反向断言不成立：\(0\to\mathbb Z\xrightarrow2\mathbb Z\to\mathbb Z/2\mathbb Z\to0\) 前两项自由，末项与自身的一次 Tor 非零，故不平坦。
\end{solution}

\begin{exercise}\label{b4:ex:08-zero-degree-shift}
以 \(0\to2\mathbb Z\to\mathbb Z\to\mathbb Z/2\mathbb Z\to0\) 为例，说明不能把 \(\operatorname{Ext}^i(\Omega M,N)\cong\operatorname{Ext}^{i+1}(M,N)\) 延伸到 \(i=0\)。写出正确的余核公式。
\end{exercise}
\begin{solution}
取 \(N=\mathbb Z\)。分别以 \(f\mapsto f(1)\) 和 \(g\mapsto g(2)\) 识别 \(\operatorname{Hom}(\mathbb Z,\mathbb Z)\) 与 \(\operatorname{Hom}(2\mathbb Z,\mathbb Z)\)，两者均为 \(\mathbb Z\)。限制映射将 \(f(1)=a\) 的同态送到 \(g(2)=f(2)=2a\) 的同态，因此在这些坐标下为乘二。长正合列给出的正确公式是
\[
\operatorname{Ext}^1_{\mathbb Z}(\mathbb Z/2\mathbb Z,\mathbb Z)
\cong\operatorname{coker}\!\left(
\operatorname{Hom}(\mathbb Z,\mathbb Z)\longrightarrow
\operatorname{Hom}(2\mathbb Z,\mathbb Z)\right)
\cong\mathbb Z/2\mathbb Z.
\]
它不是整个 \(\operatorname{Hom}(2\mathbb Z,\mathbb Z)\)；必须先模去能够从 \(\mathbb Z\) 延拓来的同态。
\end{solution}

\begin{exercise}\label{b4:ex:08-injective-shift-integers}
设 \(m\) 为正整数。用 \(0\to\mathbb Z\to\mathbb Q\to\mathbb Q/\mathbb Z\to0\) 和内射方向的低次平移，重新求 \(\operatorname{Ext}_{\mathbb Z}^1(\mathbb Z/m\mathbb Z,\mathbb Z)\)，并给出同构的具体元素。
\end{exercise}
\begin{solution}
\(\operatorname{Hom}(\mathbb Z/m\mathbb Z,\mathbb Q)=0\)，故低次正合列给出 \(\operatorname{Ext}^1(\mathbb Z/m\mathbb Z,\mathbb Z)\cong\operatorname{Hom}(\mathbb Z/m\mathbb Z,\mathbb Q/\mathbb Z)\)。一个同态由 \(\bar1\) 的像决定，该像必须被 \(m\) 消去。若 \(z+\mathbb Z\in\mathbb Q/\mathbb Z\) 满足 \(mz\in\mathbb Z\)，写 \(a=mz\)，就有 \(z=a/m\)；反之，每个 \(a/m+\mathbb Z\) 都被 \(m\) 消去。两个整数 \(a,b\) 给出同一像，当且仅当 \(m\mid(a-b)\)。因此
\[
\mathbb Z/m\mathbb Z\longrightarrow\operatorname{Hom}(\mathbb Z/m\mathbb Z,\mathbb Q/\mathbb Z),
\qquad
\bar a\longmapsto\bigl(\bar1\mapsto a/m+\mathbb Z\bigr)
\]
是同构，再由连接态射将这些同态对应到 Ext 类。
\end{solution}

\begin{exercise}\label{b4:ex:08-syzygy-change}
设 \(R\) 为含幺环，\(P,M\) 为幺性左 \(R\)-模，\(P\) 投射。给定满射 \(p:P\twoheadrightarrow M\)，记 \(K=\ker p\)，并取任意投射左模 \(Q\)。构造另一个从投射模到 \(M\) 的满射，使其核同构于 \(K\oplus Q\)。证明对任意左 \(R\)-模 \(N\)、右 \(R\)-模 \(A\) 及 \(n\ge1\)，有
\[
\operatorname{Ext}_R^n(K\oplus Q,N)\cong\operatorname{Ext}_R^n(K,N),\qquad
\operatorname{Tor}_n^R(A,K\oplus Q)\cong\operatorname{Tor}_n^R(A,K).
\]
\end{exercise}
\begin{solution}
取 \(P\oplus Q\to M\)、\((x,q)\mapsto p(x)\)，其核为 \(K\oplus Q\)。有限双积的投射消解可逐项作双积，Hom 或张量后仍分成两部分。\(Q\) 的正次导出函子消失，故所得正次 Ext、Tor 分别等于 \(K\) 的相应值。零次却保留额外项：
\[
\begin{aligned}
\operatorname{Hom}_R(K\oplus Q,N)
&\cong\operatorname{Hom}_R(K,N)\oplus\operatorname{Hom}_R(Q,N),\\
A\otimes_R(K\oplus Q)
&\cong(A\otimes_RK)\oplus(A\otimes_RQ).
\end{aligned}
\]
这些额外项一般非零，因而正次数的同构不能直接延伸到零次。
\end{solution}

\begin{exercise}\label{b4:ex:08-truncate-independent}
设 \(R\) 为含幺环，\(M\) 为幺性左 \(R\)-模，且 \(M\) 有长度不超过 \(r\) 的投射消解，其中 \(r\ge1\)。证明在任意另一投射消解中，第 \(r\) 次合冲模投射，而第 \(r-1\) 次未必投射。
\end{exercise}
\begin{solution}
已知短消解使 \(\operatorname{Ext}_R^{r+1}(M,N)=0\) 对全部 \(N\) 成立。另一消解的迭代平移给出 \(\operatorname{Ext}_R^1(\Omega^rM,N)=0\)，故该合冲模投射。取域 \(k\)，令 \(R=k[x,y]\)、\(M=R/(x,y)\)、\(r=2\)。Koszul 自由消解的二次项为 \(R\)，而其一次合冲模为 \((x,y)\)。张量 \(k\) 后两个微分均为零，故 \(\operatorname{Tor}_2^R(k,k)\cong k\)；正次数平移给出 \(\operatorname{Tor}_1^R((x,y),k)\cong k\ne0\)。因此 \((x,y)\) 不平坦，从而不投射；二次合冲模却同构于 \(R\)，是投射模。
\end{solution}

% !TeX root = ../../Book4.tex
% 纲要标题：## 第10章　扩张、乘积与万有系数
% 纲要标签：b4:ch:09
% 纲要位置：计划纲要.md，第 3435 行。

\chapter{扩张、乘积与万有系数}
\label{b4:ch:09}
\BookChapterNumber{b4:ch:09}

\begin{chapterabstract}
本章把 Ext 与模扩张联系起来，研究扩张的乘积，并证明主理想整环上的万有系数定理与 Künneth 公式，区分自然正合列和非典范分裂。
\end{chapterabstract}

\begin{learninggoals}
\begin{itemize}
\item 从投射消解构造扩张类，并用 Baer 和与 Yoneda 乘积计算它们。
\item 说明高阶扩张等价为何允许箭头方向交替，区分扩张类与中间对象的同构类型。
\item 核对张量微分的符号，辨别平坦性及有界性在保持拟同构中的作用。
\item 应用万有系数与 Künneth 公式，保留附加的 Ext、Tor 项并辨别分裂的自然性。
\end{itemize}
\end{learninggoals}

考察以 \(\mathbb Z/2\mathbb Z\) 为两端的短正合列。若中间项也是两个二元群的直积，可以从右端选一个同态截面；若中间项是 \(\mathbb Z/4\mathbb Z\)，右端的非零元却没有阶为 \(2\) 的原像，截面便不存在。两条序列的端点相同，中间的接合方式却不同。Ext 的一阶群把这些接合方式组织起来；把两条短正合列首尾相接，又提示高阶 Ext 中的乘积。另一方面，对复形取张量后，同调可能新增 Tor 项。我们将由这些可见的差别走向扩张积、万有系数与 Künneth 公式。

\ProseParagraphBreak
扩张的拼接与复形的张量积处理的是不同数据。前者把扩张次数相加；后者先构造一个新复形，再考察其同调。\cref{b4:tab:extension-tensor-comparison}将这两种操作与它们涉及的同调比较映射并列。

\begin{table}[htbp]
\centering
\caption{扩张乘积、张量复形与同调比较}
\label{b4:tab:extension-tensor-comparison}
\begin{tabularx}{\textwidth}{@{}p{0.19\textwidth}>{\raggedright\arraybackslash}X>{\raggedright\arraybackslash}X@{}}
\toprule
构造 & 输入的数据 & 所得对象或映射 \\
\midrule
Yoneda 乘积
& \(\alpha\in\operatorname{Ext}_R^m(N,L)\)，\newline \(\beta\in\operatorname{Ext}_R^n(M,N)\)
& 拼接得到 \(\alpha\beta\in\operatorname{Ext}_R^{m+n}(M,L)\) \\
张量总复形
& 右 \(R\)-模链复形 \(C\) 与左 \(R\)-模链复形 \(D\)
& 复形 \(C\otimes_R D\)，微分包含次数符号 \((-1)^p\) \\
同调比较映射
& 两个复形的同调类的张量积，\(p+q=n\)
& \(\kappa_n\) 映到 \(H_n(C\otimes_R D)\)；Künneth 公式计算其余核 \\
\bottomrule
\end{tabularx}
\end{table}


本章使用\chapref{b4:ch:05}、\chapref{b4:ch:06}的复形和总化，以及\chapref{b4:ch:07}、\chapref{b4:ch:08}的消解、Ext、Tor 和维数平移。\secref{b4:sec:0901}以左模为对象；张量部分逐项说明模的侧别，万有系数和 Künneth 主定理则在主理想整环上讨论。为方便同调公式，后三节采用链下标 \(C_n=C^{-n}\)，不改变本卷上链微分与移位的约定。第三卷的 Koszul 复形提供了另一种计算方式，可与这里的公式对照。

\ProseParagraphBreak
扩张与乘积可参阅 Weibel 的《An Introduction to Homological Algebra》\cite[Section 3.4]{Weibel1994HomologicalAlgebra}；系数变换与复形张量积的公式见\cite[Section 3.6]{Weibel1994HomologicalAlgebra}。本章对高阶扩张给出具体的双向构造，对主理想整环上的分裂逐步证明其存在，并用链自同构说明为何不能把它称为自然分裂。

% !TeX root = ../../Book4.tex
% 纲要标题：### 10.1 Yoneda 扩张
% 纲要标签：b4:sec:0901
% 纲要位置：计划纲要.md，第 3437 行。

\section{Yoneda 扩张}
\label{b4:sec:0901}

扩张不仅包含中间模，还包含核与商的指定映射。把扩张粘接起来得到的乘积，需要与消解中提升复形映射的乘积一致。

% 纲要内容：
%
% * Ext^1 与短正合扩张
% * 高阶扩张及等价
% * Yoneda 乘积与消解乘积的一致性
%

\subsection{\texorpdfstring{\(\operatorname{Ext}^1\) 与短正合扩张}{Ext1 与短正合扩张}}
\label{b4:subsec:0901:01}

给定左 \(R\)-模 \(M,N\)，\(\operatorname{Ext}_R^1(M,N)\) 既可由消解计算，也能描述把 \(N\) 放在核位置、把 \(M\) 放在商位置的所有短正合列。这里分类的数据包括两端的指定映射；仅知道中间模的同构类型还不够。

\begin{definition}\label{b4:def:yoneda-extension}
设 \(R\) 为含幺环，\(M,N\) 为幺性左 \(R\) 模。短正合列
\[
 \xi:\quad 0\longrightarrow N\xrightarrow{i}E\xrightarrow{p}M\longrightarrow0
\]
称为 \(M\) 被 \(N\) 的一个\term[kuozhang]{扩张}[extension]。给定另一个扩张 \(\xi':0\to N\xrightarrow{i'}E'\xrightarrow{p'}M\to0\)。若存在 \(R\) 模同态 \(f:E\to E'\)，满足 \(fi=i'\)、\(p'f=p\)，则称两扩张等价。这里两端的映射固定为恒等映射；两行短正合且两端为同构，由\cref{b4:thm:five-lemma}应用于包含两端零对象的五项序列，可知 \(f\) 自动为同构。仅有 \(E\cong E'\) 还不够，所用同构必须同时保持指定的核嵌入与商映射。
\end{definition}

\begin{theorem}\label{b4:thm:ext1-extensions}
设 \(R\) 为含幺环，\(M,N\) 为幺性左 \(R\) 模。\(M\) 被 \(N\) 的短正合扩张的等价类与 \(\operatorname{Ext}_R^1(M,N)\) 自然一一对应，分裂扩张对应零元。
\end{theorem}
\begin{proof}
选定正合列 \(0\to K\xrightarrow{j}P\xrightarrow{\varepsilon}M\to0\)，其中 \(P\) 投射。\cref{b4:thm:ext-long-exact}在第一变量给出的低次片段为
\[
 \operatorname{Hom}_R(P,N)\xrightarrow{j^*}\operatorname{Hom}_R(K,N)
 \longrightarrow\operatorname{Ext}_R^1(M,N)
 \longrightarrow\operatorname{Ext}_R^1(P,N)=0.
\]
最后一项由 \(P\) 的投射性消失，且 \(j^*(a)=aj\)，所以
\[
 \operatorname{Ext}_R^1(M,N)
 \cong\operatorname{Hom}_R(K,N)\big/\{\,a j:a\in\operatorname{Hom}_R(P,N)\,\}.
\]
分母恰好由可以从 \(K\) 延拓到 \(P\) 的映射组成。对扩张 \(\xi\)，由 \(P\) 投射且 \(p:E\to M\) 满射，可将 \(\varepsilon\) 提升为 \(t:P\to E\)，使 \(pt=\varepsilon\)。随后 \(ptj=\varepsilon j=0\)，故核映射 \(i:N\to E\) 的泛性质给出唯一的 \(u:K\to N\)，使 \(iu=tj\)。提升得到 \(t\)，经过核才得到 \(u\)；后者记录了 \(P\) 中的关系在扩张里落入 \(N\) 的方式。换成另一提升 \(t'\)，有 \(p(t'-t)=0\)，因而 \(t'-t=ia\)（\(a:P\to N\)），再由 \(i\) 单射得 \(u'-u=aj\)。扩张态射与 \(i,p\) 交换，将一套提升送到另一扩张的提升，故也保持这个余类。

反过来，希望把 \(P\) 中的 \(j(k)\) 与 \(N\) 中的 \(u(k)\) 识别，即令 \([(0,j(k))]=[(u(k),0)]\)。把两边相减便得到应商掉的关系 \((u(k),-j(k))\)，所以沿 \(u:K\to N\) 作推出得到
\[
 E_u=(N\oplus P)/\{\,(u(k),-j(k)):k\in K\,\}.
\]
由 \(n\mapsto[(n,0)]\) 及 \([(n,p)]\mapsto\varepsilon(p)\) 得到扩张。第一个映射单射，因为 \(j\) 单射；第二个映射的核由 \(\varepsilon(p)=0\) 即 \(p=j(k)\) 算出，此时 \([(n,j(k))]=[(n+u(k),0)]\)，故核恰好是 \(N\) 的像。如果 \(u-v=a j\)，考虑映射 \([(n,p)]\mapsto[(n+a(p),p)]\)。它把被商掉的关系送到
\[
 (u(k),-j(k))\longmapsto(u(k)-aj(k),-j(k))=(v(k),-j(k)),
\]
因而良定义；用 \(-a\) 得到逆映射，且两端映射保持不变。对原来的 \(\xi\)，映射 \([(n,p)]\mapsto i(n)+t(p)\) 也良定义，因为 \(iu=tj\)，并给出 \(E_u\to E\) 的扩张态射，因而是同构。两个构造互逆。

若类为零，则 \(u=aj\)（\(a:P\to N\)）。令 \(t'=t-ia\)，就有 \(pt'=\varepsilon\) 及
\[
 t'j=tj-iaj=iu-iu=0.
\]
因此 \(t'\) 唯一经过余核 \(\varepsilon:P\to M\)，写为 \(s\varepsilon\)；由 \(ps\varepsilon=\varepsilon\) 且 \(\varepsilon\) 满射，得到 \(ps=1_M\)。反之，若有截面 \(s\)，可取提升 \(t=s\varepsilon\)，它在 \(K\) 上为零，故对应零类。对 \(N\to N'\) 作推出相当于把 \(u\) 后复合，对 \(M'\to M\) 作拉回相当于提升该映射到投射消解后预复合，故对应自然。
\end{proof}

扩张类的加法也有直接构造。对两个扩张先取直和，再沿对角映射 \(M\to M\oplus M\) 拉回，使两个中间元素具有同一个商像；所得核仍有两份 \(N\)。随后沿加法 \(N\oplus N\to N\) 推出，把核元素 \((n,n')\) 合并为 \(n+n'\)。所得扩张称为\term[baerhe]{Baer 和}[Baer sum]。\footnote{拜尔（Reinhold Baer，1902-1979），德裔数学家，研究群的扩张、可解群和幂零群。}加法的核由 \((n,-n)\) 组成，故中间模具体为
\[
 (E\times_M E')/\{\,(i(n),-i'(n)):n\in N\,\}.
\]

\begin{proposition}\label{b4:prop:baer-addition}
设 \(R\) 为含幺环，\(M,N\) 为幺性左 \(R\) 模。Baer 和在扩张等价类上给出交换群结构；\cref{b4:thm:ext1-extensions}的对应是群同构。
\end{proposition}
\begin{proof}
从同一个 \(P\to M\) 分别提升到 \(E,E'\)，因两套提升的商像相同，可合并为到纤维积的提升；再作推出后，在 \(K\) 上诱导的映射就是 \(u+u'\)。因而 Baer 和对应 \(\operatorname{Ext}^1\) 中的加法。前一定理给出的双射于是把商群中已成立的结合律、交换律及零元、负元传到扩张类上，也证明结果不依赖代表元。具体地，分裂扩张为零，沿 \(-1_N\) 推出得到负元。
\end{proof}

例如，对素数 \(p\)，\(\operatorname{Ext}_{\mathbb Z}^1(\mathbb Z/p\mathbb Z,\mathbb Z/p\mathbb Z)=\mathbb Z/p\mathbb Z\)。非零类的中间群都同构于 \(\mathbb Z/p^2\mathbb Z\)，但指定的核嵌入可以不同，所以这些扩张仍有 \(p-1\) 个不同的非零类。这个区别在用扩张恢复结构时不能省略。

\subsection{高阶扩张及等价}
\label{b4:subsec:0901:02}

\begin{definition}\label{b4:def:higher-yoneda-extension}
设 \(R\) 为含幺环，\(M,N\) 为幺性左 \(R\) 模，\(n\ge1\) 为整数。给定左 \(R\) 模正合列
\[
 \xi:\quad 0\longrightarrow N\longrightarrow E_{n-1}\longrightarrow\cdots
 \longrightarrow E_0\longrightarrow M\longrightarrow0
\]
称为一个 \(n\) 阶扩张。两个这样的扩张之间，两端为恒等映射的链映射称为扩张态射。由扩张态射生成的等价关系称为 \term[yonedakuozhang]{Yoneda 扩张}[Yoneda extension]的等价关系：允许有限串扩张态射，箭头方向可以交替。这些扩张的等价类族暂记为 \(\operatorname{YExt}_R^n(M,N)\)；下面的定理将证明它可由一个集合参数化。
\end{definition}
\symidx{\(\operatorname{YExt}_R^n(M,N)\)：\(n\) 阶 Yoneda 扩张的等价类}

当 \(n>1\) 时，中间各项的映射不必可逆。例如令 \(A=\mathbb Z/2\mathbb Z\)，考察两个整数模二阶扩张
\[
\begin{aligned}
0&\longrightarrow A\xrightarrow{a\mapsto(a,0)}A\oplus\mathbb Z
 \xrightarrow{(a,z)\mapsto2z}\mathbb Z
 \xrightarrow{z\mapsto\bar z}A\longrightarrow0,\\
0&\longrightarrow A\xrightarrow{1_A}A\xrightarrow{0}A
 \xrightarrow{1_A}A\longrightarrow0.
\end{aligned}
\]
第一列中间的核为 \(A\oplus0\)，像为 \(2\mathbb Z\)，所以两列均正合。中间分别取投影 \((a,z)\mapsto a\) 和取模映射 \(z\mapsto\bar z\)，两端取恒等，就得到从第一列到第二列的扩张态射：中间方块两条复合均为零，其余方块也交换。中间项却分别是无限群与有限群，不可能同构。高阶等价因而不能要求逐项同构；要由这些不必可逆的态射得到等价关系，还须允许有限箭头串及箭头反向。下面同时说明这些类确实构成集合，并给出计算它们的方法。

\begin{theorem}\label{b4:thm:yoneda-ext-comparison}
设 \(R\) 为含幺环，\(M,N\) 为幺性左 \(R\) 模，\(n\ge1\) 为整数。扩张等价类可组成集合，并有对 \(M\) 反变、对 \(N\) 协变的自然双射
\[
 \operatorname{YExt}_R^n(M,N)\simeq\operatorname{Ext}_R^n(M,N).
\]
沿 \(M\) 的对角映射拉回、沿 \(N\) 的加法映射推出所定义的高阶 Baer 和，对应右侧的加法。
\end{theorem}
\begin{proof}
固定投射消解 \(P_\bullet\xrightarrow{\varepsilon}M\)，令 \(\Omega^nM=\ker(P_{n-1}\to P_{n-2})\)，其中 \(P_{-1}=M\)，\(P_0\to P_{-1}\) 为增广。记由微分诱导的满射为 \(\pi_n:P_n\to\Omega^nM\)。\cref{b4:thm:dimension-shift-basic,b4:prop:iterated-dimension-shift}的平移公式及短正合列 \(0\to\Omega^nM\to P_{n-1}\to\Omega^{n-1}M\to0\) 给出
\[
 \operatorname{Ext}_R^n(M,N)
 \simeq\operatorname{Hom}_R(\Omega^nM,N)
 /\operatorname{res}\operatorname{Hom}_R(P_{n-1},N).
\]
也可直接从消解读出此式。例如 \(n=2\) 时，\(\Omega^2M=\ker\mathrm{d}_{P,1}=\operatorname{im}\mathrm{d}_{P,2}\)。次数二的余圈 \(b:P_2\to N\) 满足 \(b\mathrm{d}_{P,3}=0\)，而 \(\ker\pi_2=\operatorname{im}\mathrm{d}_{P,3}\)，所以唯一写成 \(b=u\pi_2\)。次数二的余边是 \(a\mathrm{d}_{P,2}\)（\(a:P_1\to N\)），其下降映射恰为 \(a|_{\Omega^2M}\)。任意 \(n\ge1\) 同样有 \(\ker\pi_n=\operatorname{im}\mathrm{d}_{P,n+1}\)，因此余圈唯一经过 \(\Omega^nM\)，余边则对应商公式分母中的限制映射。这里沿用\secref{b4:sec:0802}投射消解计算中的无符号微分 \(f\mapsto f\mathrm{d}_P\)。

给定 \(n\) 阶扩张，将 \(E_0\to M\) 记为 \(\mathrm{d}_{E,0}\)，并将 \(P_0\to M\) 记为 \(\mathrm{d}_{P,0}\)。依次利用 \(P_i\) 的投射性，构造与微分、增广交换的提升 \(t_i:P_i\to E_i\)，\(0\le i<n\)。最后 \(t_{n-1}|_{\Omega^nM}\) 被下一微分零化，故经最左端的核嵌入 \(i:N\to E_{n-1}\) 分解，得到 \(u:\Omega^nM\to N\)。

为证明其余类与提升无关，只需把两套提升在最左端的差写成一个 \(P_{n-1}\to N\) 映射的限制。不能直接断言这一差落入 \(N\)：较低次数的提升也可能不同，须先逐层消去由它们传来的误差。对两套提升 \(t,s\)，从最低次数起构造 \(h_i:P_i\to E_{i+1}\)，使
\[
 t_i-s_i=\mathrm{d}_Eh_i+h_{i-1}\mathrm{d}_P\qquad(0\le i<n-1),\qquad h_{-1}=0.
\]
零次的差被增广零化。若已构造前一步，令
\[
 e_i=t_i-s_i-h_{i-1}\mathrm{d}_{P,i}.
\]
由两套提升的交换等式及上一同伦等式，得
\[
\begin{aligned}
\mathrm{d}_{E,i}e_i
&=(t_{i-1}-s_{i-1})\mathrm{d}_{P,i}
  -\mathrm{d}_{E,i}h_{i-1}\mathrm{d}_{P,i}\\
&=h_{i-2}\mathrm{d}_{P,i-1}\mathrm{d}_{P,i}=0
\qquad(i\ge1).
\end{aligned}
\]
所以尚未消去的误差落在核中；中间正合性把这个核识别为 \(E_{i+1}\) 的像，\(P_i\) 投射才使它提升为 \(h_i\)。到 \(i=n-1\) 时，核已经是 \(i(N)\)，于是存在 \(a:P_{n-1}\to N\) 使
\[
 t_{n-1}-s_{n-1}-h_{n-2}\mathrm{d}_{P,n-1}=ia.
\]
限制到 \(\Omega^nM\) 后，微分项为零，再由 \(i\) 单射得到 \(u_t-u_s=a|_{\Omega^nM}\)，这恰在商公式的分母中。\(n=1\) 时直接用 \(t_0-s_0\) 即可。扩张态射把一套提升变成另一扩张的提升，因而整个箭头串都保持该余类。

对任意 \(u:\Omega^nM\to N\)，把截短消解
\[
 0\longrightarrow\Omega^nM\longrightarrow P_{n-1}\longrightarrow\cdots
 \longrightarrow P_0\longrightarrow M\longrightarrow0
\]
的最左端沿 \(u\) 推出。新项是
\(E_{n-1}(u)=(N\oplus P_{n-1})/\{(u(x),-x)\}\)，其余项保持 \(P_i\)。与一阶情形相同，核及像的计算证明新列正合。若 \(u-v=a|_{\Omega^nM}\)，在新项上用 \([(z,p)]\mapsto[(z+a(p),p)]\)，在其他项上用恒等映射，就得到两个新扩张的同构。

从任意原扩张取出的 \(u\) 又给出这个推出扩张到原扩张的态射：最左新项用 \([(z,p)]\mapsto i(z)+t_{n-1}(p)\)，其他项用 \(t_i\)。标准推出扩张的这些项取自 \(P_\bullet\)，原扩张的 \(E_i\) 却不必投射，所构造的中间箭头也不必可逆；它们之所以足以识别两扩张，正是因为高阶等价由扩张态射生成。因此两个方向互逆。特别地，每个扩张类都有一个由 \(u\in\operatorname{Hom}_R(\Omega^nM,N)\) 得到的代表；这个 Hom 集的上述商集便给出全部扩张类的集合参数化。

接着核对两变量的作用。若 \(v:N\to N'\)，沿 \(v\) 推出扩张，并把原来的提升复合到推出列中，最左端诱导的映射变成 \(vu:\Omega^nM\to N'\)。它对应消解余圈的后复合，故与 \(\operatorname{Ext}_R^n(M,v)\) 相容。

若 \(f:M'\to M\)，取投射消解 \(P'_\bullet\to M'\)，由\cref{b4:thm:resolution-comparison}选取提升 \(f\) 的比较映射 \(c:P'_\bullet\to P_\bullet\)。记
\[
 c_\Omega=c_{n-1}|_{\Omega^nM'}:\Omega^nM'\longrightarrow\Omega^nM.
\]
这里 \(f\) 决定要拉回的扩张，\(c\) 实现 Ext 计算中的预复合，\(c_\Omega\) 则记录它在关系模上的作用。拉回列的最低项为 \(E_0\times_M M'\)，其提升明确取为
\[
 P'_0\longrightarrow E_0\times_M M',\qquad
 x\longmapsto(t_0c_0(x),\varepsilon_{P'}(x));
\]
它落在纤维积中，因为 \(\mathrm{d}_{E,0}t_0c_0=f\varepsilon_{P'}\)。其余次数用原提升与 \(c_i\) 的复合，最左端映射因而成为 \(u c_\Omega\)。若余圈 \(b:P_n\to N\) 写为 \(b=u\pi_n\)，记相应满射为 \(\pi'_n:P'_n\to\Omega^nM'\)，则
\[
 b c_n=u\pi_nc_n=u c_\Omega\pi'_n,
\]
所以这个商类恰是 \(\operatorname{Ext}_R^n(f,N)\) 所给的类。若换另一比较提升 \(c'\)，由\cref{b4:prop:comparison-homotopy-unique}，存在同伦 \(h_i:P'_i\to P_{i+1}\)，使
\[
 c_i-c'_i=\mathrm{d}_{P,i+1}h_i+h_{i-1}\mathrm{d}_{P',i},\qquad h_{-1}=0.
\]
在 \(\Omega^nM'\) 上第二项为零，因此
\[
 u c_\Omega-u c'_\Omega
 =\left.(u\pi_nh_{n-1})\right|_{\Omega^nM'}.
\]
其中 \(h_{n-1}:P'_{n-1}\to P_n\)，故 \(u\pi_nh_{n-1}\) 定义在整个 \(P'_{n-1}\) 上，右侧恰属于商公式的分母。这既证明拉回与 Ext 的预复合相容，也证明更换比较提升不改变所得类。

最后把换消解与换扩张提升区分开来。若 \(P'_\bullet\) 也是 \(M\) 的投射消解，上面的构造取 \(f=1_M\)，就把两种商描述按 \([u]\mapsto[u c_\Omega]\) 比较。取反向比较后，两种复合分别与相应的恒等映射同伦，所以诱导互逆的同构；余圈等式也表明这正是\chapref{b4:ch:08}用于识别两种 Ext 计算的比较同构。因而不同消解给出的识别在该同构下相容。前面固定消解时已经证明换扩张中的提升、沿扩张态射改变代表均不改变商类，这里消去的则是投射消解本身的选择。

对扩张 \(\xi,\xi'\)，先逐项取直和 \(\xi\oplus\xi'\)，得到右端为 \(M\oplus M\)、左端为 \(N\oplus N\) 的列；再在右端沿 \(\Delta_M:M\to M\oplus M\) 拉回，最后在左端沿 \(+_N:N\oplus N\to N\) 推出。也就是
\[
 \xi+\xi'=(+_N)_*\Delta_M^*(\xi\oplus\xi').
\]
分别取 \(P_\bullet\) 到两列的提升，再用逐项对角映射 \(P_\bullet\to P_\bullet\oplus P_\bullet\) 提升到拉回列。在最左端，所得映射是 \((u,u'):\Omega^nM\to N\oplus N\)；推出后成为 \(+_N(u,u')=u+u'\)。所以 Baer 和对应商集中的加法。特别地，它不依赖代表扩张，并由此给出与 Ext 相同的交换群结构。
\end{proof}

对 \(n=0\)，约定 \(\operatorname{YExt}_R^0(M,N)=\operatorname{Hom}_R(M,N)\)。在没有足够投射或内射对象的一般交换范畴中，仍可讨论扩张及其箭头串；本节已证的与导出函子的比较则是在模范畴中进行的。

\subsection{Yoneda 乘积与消解乘积}
\label{b4:subsec:0901:03}

若一列以 \(N\) 为商，另一列以 \(N\) 为核，就能把它们首尾连接。例如两个一次扩张
\[
0\longrightarrow L\xrightarrow{i}E\xrightarrow{v}N\longrightarrow0,
\qquad
0\longrightarrow N\xrightarrow{j}F\xrightarrow{p}M\longrightarrow0
\]
拼接为
\[
0\longrightarrow L\xrightarrow{i}E\xrightarrow{jv}F
\xrightarrow{p}M\longrightarrow0.
\]
由于 \(j\) 单射，\(\ker(jv)=\ker v=\operatorname{im}i\)；由于 \(v\) 满射，\(\operatorname{im}(jv)=j(N)=\ker p\)，所以拼接列仍正合。中间箭头是明确的复合 \(E\to N\to F\)，而不是将两列中的 \(N\) 直接略去。这一操作给出 Ext 群之间并非来自普通数乘的乘法。

\begin{definition}\label{b4:def:yoneda-product}
设 \(R\) 为含幺环，\(L,N,M\) 为幺性左 \(R\) 模，\(m,n\ge1\) 为整数。对 \(\alpha\in\operatorname{YExt}_R^m(N,L)\)、\(\beta\in\operatorname{YExt}_R^n(M,N)\)，把代表扩张在共同的 \(N\) 处拼接，删除两处 \(N\)，以原来的满射到 \(N\) 与从 \(N\) 出发的单射之复合作为中间箭头。所得 \((m+n)\) 阶扩张类记为 \(\alpha\beta\)，称为\term[yonedachengji]{Yoneda 乘积}[Yoneda product]。约定 \(\operatorname{YExt}_R^0(X,Y)=\operatorname{Hom}_R(X,Y)\)：左因子为零次同态 \(N\to L\) 时沿该同态推出，右因子为零次同态 \(M\to N\) 时沿该同态拉回；两个因子均为零次时取同态复合。
\end{definition}
\symidx{\(\alpha\beta\)：Yoneda 扩张类的乘积}

\begin{proposition}\label{b4:prop:yoneda-resolution-product}
设 \(R\) 为含幺环，\(L,N,M\) 为幺性左 \(R\) 模。Yoneda 乘积良定义、双加性且结合，并在 Yoneda 扩张与 Ext 的自然对应下等于下述消解乘积。取投射消解 \(P_\bullet\to M\)、\(Q_\bullet\to N\)，以 \(\varepsilon_Q:Q_0\to N\) 表示增广，设 \(m,n\ge0\) 为整数。若 \(b:P_n\to N\)、\(a:Q_m\to L\) 分别是 \(\operatorname{Hom}_R(P_\bullet,N)\)、\(\operatorname{Hom}_R(Q_\bullet,L)\) 中的余圈，把 \(b\) 逐次提升为
\[
 t_i:P_{n+i}\longrightarrow Q_i,\qquad
 \varepsilon_Qt_0=b,\qquad \mathrm{d}_Qt_i=t_{i-1}\mathrm{d}_P\quad(i\ge1),
\]
则乘积由余圈 \(a t_m:P_{n+m}\to L\) 代表。
\end{proposition}
\begin{proof}
拼接在共同项处正合，因为像和核都由被删除的 \(N\) 描述。若一个因子沿某个扩张态射改变，把该态射与另一列的恒等映射拼接，仍得到扩张态射。因此拼接保持等价关系的每个生成关系，也就保持由正向或反向扩张态射组成的任意有限箭头串。零次因子的拉回或推出也把扩张态射送到扩张态射，故这些情形同样良定义。三个正次数因子依次拼接不改变项及其复合箭头。零次因子位于外端时，拉回或推出的泛性质使两种顺序相容；它位于中间时，拉回的投影与推出的结构映射给出两种拼接列之间的扩张态射。连续的零次因子则由同态复合的结合律处理。因而乘积结合。

对于消解公式，先沿满射 \(\varepsilon_Q:Q_0\to N\) 提升余圈 \(b:P_n\to N\)，得到 \(t_0:P_n\to Q_0\)。因此 \(Q\) 的零次已经对应 \(P\) 的第 \(n\) 次。由 \(b\mathrm{d}_{P,n+1}=0\)，\(t_0\mathrm{d}_{P,n+1}\) 落在 \(\ker\varepsilon_Q\)，可用 \(P_{n+1}\) 的投射性提升至 \(Q_1\)。以后每步同样先用 \(\mathrm{d}_P^2=0\) 核对落在核中，再提升；每向 \(Q\) 的高次推进一步，源也从 \(P_{n+i}\) 推进到 \(P_{n+i+1}\)。到第 \(m\) 步，恰可形成
\[
 P_{n+m}\xrightarrow{t_m}Q_m\xrightarrow{a}L,
 \qquad
 a t_m\mathrm{d}_{P,n+m+1}
 =a\mathrm{d}_{Q,m+1}t_{m+1}=0.
\]
这既说明所得是余圈，也说明乘积次数为何是 \(n+m\)。

对 \(m,n\ge1\)，要比较两种乘积，只须读取拼接扩张最左端的关系映射，证明其余圈为 \(at_m\)。用\cref{b4:thm:yoneda-ext-comparison}的推出构造代表 \(\alpha,\beta\)：把 \(\alpha\) 的最左新项记为 \(A\)，把 \(\beta\) 的最左新项记为 \(B\)，其余项来自两条消解。设 \(a,b\) 经合冲模的分解分别为
\[
 \bar a:\Omega^mN\to L,\qquad \bar b:\Omega^nM\to N,
\]
则这两个推出项为
\[
\begin{aligned}
 A&=(L\oplus Q_{m-1})/\{(\bar a(x),-x):x\in\Omega^mN\},\\
 B&=(N\oplus P_{n-1})/\{(\bar b(y),-y):y\in\Omega^nM\}.
\end{aligned}
\]
以 \(i_\alpha:L\to A\)、\(i_\beta:N\to B\) 表示单射，以 \(j_\alpha:Q_{m-1}\to A\)、\(j_\beta:P_{n-1}\to B\) 表示 \(q\mapsto[(0,q)]\) 型的映射。要将两段提升在共同的 \(N\) 处接起来，须使经 \(P_{n-1}\) 进入 \(B\) 与经 \(Q_0\to N\) 进入 \(B\) 的两条路径相同，所需等式正是下面的方块：
\[
\begin{tikzcd}[column sep=large,row sep=large]
 P_n \ar[r,"\mathrm{d}_{P,n}"] \ar[d,"t_0"']
   & P_{n-1} \ar[d,"j_\beta"] \\
 Q_0 \ar[r,"i_\beta\varepsilon_Q"'] & B.
\end{tikzcd}
\]
它交换，因为 \(j_\beta\mathrm{d}_{P,n}=i_\beta b=i_\beta\varepsilon_Qt_0\)。因此从 \(P_0\) 到 \(P_{n-2}\) 的恒等映射、\(P_{n-1}\to B\) 的 \(j_\beta\)，以及
\[
 P_{n+i}\xrightarrow{t_i}Q_i\quad(0\le i<m-1),\qquad
 P_{n+m-1}\xrightarrow{j_\alpha t_{m-1}}A
\]
组成到拼接列的逐项提升；当 \(m=1\) 或 \(n=1\) 时，相应的指标区间为空。若 \(m=1\)，拼接箭头 \(A\to B\) 由 \(i_\beta\varepsilon_Q:Q_0\to B\) 及 \(L\to B\) 的零映射在推出上诱导，所以上面的方块仍给出所需相容性。

这些映射给出到拼接列的提升以后，最左端的关系映射由 \(P_{n+m}\to P_{n+m-1}\to A\) 读出。下面的上方方块使用提升等式，下方方块使用 \(A\) 的推出关系，将这条路径与 \(P_{n+m}\to Q_m\to L\to A\) 比较：
\[
\begin{tikzcd}[column sep=large,row sep=large]
 P_{n+m} \ar[r,"\mathrm{d}_{P,n+m}"] \ar[d,"t_m"']
   & P_{n+m-1} \ar[d,"t_{m-1}"] \\
 Q_m \ar[r,"\mathrm{d}_{Q,m}"] \ar[d,"a"']
   & Q_{m-1} \ar[d,"j_\alpha"] \\
 L \ar[r,"i_\alpha"'] & A.
\end{tikzcd}
\]
下方方块交换，是因为推出关系给出 \(j_\alpha\mathrm{d}_{Q,m}=i_\alpha a\)。于是
\[
 j_\alpha t_{m-1}\mathrm{d}_{P,n+m}=i_\alpha a t_m.
\]
由于 \(\operatorname{im}\mathrm{d}_{P,n+m}=\Omega^{n+m}M\) 且 \(i_\alpha\) 单射，拼接列在最左端诱导的映射正是余圈 \(a t_m\) 经 \(\Omega^{n+m}M\) 的分解。因此拼接扩张对应所陈述的消解乘积。该对应也证明换 \(t_i\)、换余圈代表或换消解均不改变所得 Ext 类。

固定 \(b\) 的提升时，\((a+a')t_m=a t_m+a't_m\)；固定 \(a\) 时，两套 \(b,b'\) 的提升之和提升 \(b+b'\)。因此乘积双加性。若 \(m=0\)，写 \(a=h\varepsilon_Q\)，其中 \(h:N\to L\)，则 \(a t_0=h b\)，正是沿 \(h\) 推出的后复合。若 \(n=0\)，写 \(b=f\varepsilon_P\)，其中 \(f:M\to N\)，则 \(t_i\) 是提升 \(f\) 的比较映射，\(a t_m\) 正是自然性中沿 \(f\) 拉回的余圈。两个次数均为零时，这又化为同态复合。
\end{proof}

于是 \(\bigoplus_{n\ge0}\operatorname{Ext}_R^n(M,M)\) 成为带单位的分次环，单位是 \(1_M\)。这里没有一般的分次交换性断言：次数零部分已经可能是非交换环 \(\operatorname{End}_R(M)\)。例如在域 \(R=k\) 上取 \(M=k^2\)，这一部分就是矩阵环 \(M_2(k)\)。

\ProseParagraphBreak
即使次数零部分交换，正次数也未必满足分次交换律。取域 \(k\) 及交换环 \(R=k[\epsilon]/(\epsilon^2)\)，令 \(M=k=R/(\epsilon)\)。\secref{b4:sec:0701}已证明它有各项为 \(R\)、微分均为乘 \(\epsilon\) 的周期自由消解；\secref{b4:sec:0802}由施加 Hom 后微分为零得到每次 Ext 为 \(k\)。次数一类 \(t\) 由自然商映射 \(q:P_1=R\to k\) 表示，提升可取底层均为 \(1_R\) 的 \(t_i:P_{i+1}\to P_i\)。它们使次数降低一，而不是原消解中的零次恒等映射。上述乘积公式将 \(t^2\) 表示为 \(q t_1:P_2\to k\)，它仍是非零商映射；Hom 复形的微分为零，故它也不是余边。因此 \(t^2\ne0\)。若 \(\operatorname{char}k\ne2\)，分次交换律却会要求 \(t^2=-t^2\)，这说明正次数中的障碍也确实存在。\cref{b4:ex:09-dual-number-product}进一步计算这一例子的整个自扩张环。

\ProseParagraphBreak
消解乘积的计算采用 \(\delta f=f\mathrm{d}\) 的次数约定。与\cref{b4:def:hom-complex}的内部 Hom 约定比较时，把 \(P_i,Q_i\) 放在上链次数 \(-i\)，并记 \(c_r=(-1)^{r(r+1)/2}\)。\secref{b4:sec:0802}的重标度将内部 Hom 的余圈 \(c_nb\) 对应于这里的 \(b\)。相应的次数 \(n\) 提升在分量 \(P_{n+i}\to Q_i\) 上须取
\[
 T_i=c_n(-1)^{ni}t_i,\qquad
 \mathrm{d}_QT_i=(-1)^n\cdot T_{i-1}\mathrm{d}_P\quad(i\ge1),
\]
才能满足内部 Hom 的余圈条件。将另一个余圈写成 \(c_ma\)，其复合为 \(c_mc_n(-1)^{mn}at_m\)。由于
\[
 c_{m+n}=(-1)^{mn}c_mc_n,
\]
在次数 \(m+n\) 作回同一重标度后，仍得到 \(at_m\)。两种符号约定因此给出同一个乘积，但须同时调整余圈与全部提升分量。Ext 扩张的基本对应和 Baer 和可参阅 Weibel\cite[Theorem 3.4.3, Definition 3.4.4, Corollary 3.4.5]{Weibel1994HomologicalAlgebra}；高阶扩张的进一步讨论见该书\cite[Vista 3.4.6]{Weibel1994HomologicalAlgebra}。

% !TeX root = ../../Book4.tex
% 纲要标题：### 10.2 截短消解与同调维数估计
% 纲要位置：计划纲要.md，第 2719 行。

\section{截短消解与同调维数估计}
\label{b4:sec:1002}

高次导出群是否消失，可以转化为消解末端的合冲模是否投射。这个判据既判断消解能否终止，也说明局部化以后投射维数需要怎样核对。

% 纲要内容：
%
% * 截短消解及合冲模
% * 短正合列中的同调维数不等式
% * 第三卷局部投射维数所需的最低平移版本已在第8.5节证明；本节作一般化整理
% * 局部投射维数的讨论作为交换代数应用，可在学习第三卷后回读；不参与前两小节的证明
%

\subsection{截短消解及合冲模}
\label{b4:subsec:1002:01}

第八章已经证明一步维数平移：短正合列中的投射或内射对象使长正合列的相邻高次项同构。本节用它确定消解何时真正能够终止，而不再重复连接同态的构造。

\begin{proposition}\label{b4:prop:truncation-dimension}
设 \(R\) 为含幺环，\(M\) 为幺性左 \(R\) 模，\(d\ge1\) 为整数。给定左 \(R\) 模正合列
\[
 0\longrightarrow K_d\longrightarrow P_{d-1}\longrightarrow\cdots
 \longrightarrow P_0\longrightarrow M\longrightarrow0,
\]
其中各 \(P_i\) 投射。则 \(\operatorname{pd}_RM\le d\) 当且仅当 \(K_d\) 投射。若 \(0\to M\to I^0\to\cdots\to I^{d-1}\to L^d\to0\) 正合且各 \(I^i\) 内射，则 \(\operatorname{id}_RM\le d\) 当且仅当 \(L^d\) 内射。若第一列的 \(P_i\) 仅要求平坦，则 \(\operatorname{fd}_RM\le d\) 当且仅当 \(K_d\) 平坦。
\end{proposition}
\begin{proof}
若末端对象具有所需性质，给定列本身就是相应长度的消解。反向分别利用
\[
 \operatorname{Ext}_R^1(K_d,N)\simeq\operatorname{Ext}_R^{d+1}(M,N),\qquad
 \operatorname{Ext}_R^1(N,L^d)\simeq\operatorname{Ext}_R^{d+1}(N,M),
\]
以及 \(\operatorname{Tor}_1^R(T,K_d)\simeq\operatorname{Tor}_{d+1}^R(T,M)\)。投射和内射的式子是\cref{b4:prop:iterated-dimension-shift}；平坦版本对各短正合列用 Tor 长正合列，并用平坦项高次 Tor 为零。右侧由\cref{b4:thm:dimension-vanishing}消失，左侧的次数一判据给出所需性质。
\end{proof}

因此某个 \(d\) 步合冲模投射时，任何投射消解的 \(d\) 步合冲模都投射，但这些合冲模不必同构。给一个消解额外加入自由直和项，就会改变实际的核；不变的是上述消失条件及其表达的维数。

\subsection{短正合列中的同调维数不等式}
\label{b4:subsec:1002:02}

\begin{proposition}\label{b4:prop:dimension-short-exact}
设 \(R\) 为含幺环，\(0\to A\to B\to C\to0\) 为幺性左 \(R\) 模短正合列。以下所有维数均在 \(R\) 上计算。投射维数满足
\[
\begin{aligned}
 \operatorname{pd}B&\le\max\{\operatorname{pd}A,\operatorname{pd}C\},\\
 \operatorname{pd}A&\le\max\{\operatorname{pd}B,\operatorname{pd}C-1\},\\
 \operatorname{pd}C&\le\max\{\operatorname{pd}B,\operatorname{pd}A+1\}.
\end{aligned}
\]
平坦维数满足相同的三式。内射维数满足
\[
\begin{aligned}
 \operatorname{id}B&\le\max\{\operatorname{id}A,\operatorname{id}C\},\\
 \operatorname{id}A&\le\max\{\operatorname{id}B,\operatorname{id}C+1\},\\
 \operatorname{id}C&\le\max\{\operatorname{id}B,\operatorname{id}A-1\}.
\end{aligned}
\]
这里非零模的维数非负，零模的维数为 \(-\infty\)，并按通常扩展约定处理 \(\pm\infty\) 与有限整数的加减。
\end{proposition}
\begin{proof}
固定左模 \(N\)。Ext 的第一变量长正合列包含
\[
 \operatorname{Ext}^i(C,N)\to\operatorname{Ext}^i(B,N)\to\operatorname{Ext}^i(A,N)
 \to\operatorname{Ext}^{i+1}(C,N).
\]
第一式由两端的消失得到中间的消失；第二式用 \(\operatorname{Ext}^i(B,N)\) 及 \(\operatorname{Ext}^{i+1}(C,N)\)；第三式把次数改成 \(i-1,i\)，使用 \(\operatorname{Ext}^{i-1}(A,N)\) 及 \(\operatorname{Ext}^i(B,N)\)。对每个大于相应上界的正整数次数，这些群均为零，再用\cref{b4:thm:dimension-vanishing}。若上界为负，则只能来自零对象或 \(B=0\) 的平凡短列，结论直接成立。

平坦情形固定右模 \(T\)，用 \(\operatorname{Tor}_i(T,-)\) 的长正合列中三种相邻位置得到相同上界。内射情形改用第二变量长正合列
\(\operatorname{Ext}^i(N,A)\to\operatorname{Ext}^i(N,B)\to\operatorname{Ext}^i(N,C)\to\operatorname{Ext}^{i+1}(N,A)\)，逐一取相邻项，得到所写三式。
\end{proof}

这些不等式也给出精确的一步下降。例如 \(0\to K\to P\to M\to0\) 中 \(P\) 投射，若 \(\operatorname{pd}M=d\ge1\) 有限，则 \(\operatorname{pd}K=d-1\)：一式给出不超过 \(d-1\)，另一式排除更小值。若 \(d=0\)，短列分裂，\(K\) 投射或为零；不能强行套用数值 \(d-1=-1\)。

\subsection{交换代数应用：局部投射维数}
\label{b4:subsec:1002:03}

前两小节的截短判据与维数不等式已经由第八章的长正合列和维数平移、\secref{b4:sec:1001}的消失判据证明。本小节把这些一般结果用于交换 Noether 局部环；所调用的第三卷极小消解理论只参与这项应用，不构成前两小节的前提。

第三卷命题17.1.2已经在交换 Noether 局部环上构造有限生成模的极小自由消解，并证明其微分矩阵的元素都在极大理想中。定理17.2.3据此建立局部投射维数判据。这里用本节的一般理论解释其中“某一步消失就能截短”的依据。

\ProseParagraphBreak
具体地，设 \((R,\mathfrak m,k)\) 为交换 Noether 局部环，\(M\ne0\) 为有限生成的 \(R\)-模，\(F_\bullet\to M\) 为极小自由消解。第三卷命题17.1.3的极小消解与 Betti 数\footnote{贝蒂（Enrico Betti，1823-1892），意大利数学家，研究代数与拓扑；以其名字命名的数记录各维同调的秩。}结论给出
\[
 \operatorname{Tor}_i^R(k,M)\simeq k\otimes_R F_i,\qquad
 \operatorname{Ext}_R^i(M,k)\simeq\operatorname{Hom}_R(F_i,k).
\]
理由是两个复形的微分都变为零。有限性使 Nakayama 引理适用，所以 \(k\otimes F_i=0\) 就是 \(F_i=0\)，此后极小消解终止。结合本节的截短判据，投射维数等于这些 Tor 非零次数的上确界。这解释了为什么在这一特定环境中只需测试剩余域，而一般环上的判据必须量化所有模。

\ProseParagraphBreak
例如在 \(R=k[x,y]_{(x,y)}\) 上，剩余域的 Koszul 消解为
\[
 0\longrightarrow R\xrightarrow{\binom{-y}{x}}R^2
 \xrightarrow{(x\ \ y)}R\longrightarrow k\longrightarrow0.
\]
它极小，故 \(\operatorname{Tor}_2^R(k,k)=k\ne0\)，同时长度二给出高次消失，从而 \(\operatorname{pd}_Rk=2\)。相反，\cref{b4:exm:infinite-dimension-dual-numbers}的每一步都保留一个自由项，有限生成并不意味着有限投射维数。这里不重证第三卷的极小消解存在性或 Auslander--Buchsbaum 公式\footnote{奥斯兰德（Maurice Auslander，1926-1994），美国数学家，研究同调代数与表示论；布赫斯鲍姆（David Alvin Buchsbaum，1929-2021），美国数学家，研究同调代数及交换代数中的自由消解。}；它们分别依赖局部有限生成理论和深度理论。

% !TeX root = ../../Book4.tex
% 纲要标题：### 10.3 万有系数定理
% 纲要标签：b4:sec:0903
% 纲要位置：计划纲要.md，第 2697 行。

\section{万有系数定理}
\label{b4:sec:0903}

改变系数可能使原来的单射失效，因此同调与系数张量之间会多出 Tor 项。主理想整环上圈与边界自由，使这个差额形成自然的短正合列。

% 纲要内容：
%
% * 主理想整环上的万有系数短正合列
% * 自然短正合列与非典范分裂
% * 系数变换的计算
%

\subsection{主理想整环上的万有系数短正合列}
\label{b4:subsec:0903:01}

主理想整环上，自由模的子模仍自由，即使自由模的秩无限也成立，见\cref{b4:thm:pid-free-submodule}。因此自由链复形的圈和边界都是自由模。这个事实既使 Tor 项可由长度一的消解算出，也保证下面的短正合列能够分裂。

\begin{theorem}[同调的万有系数定理]\label{b4:thm:uct-homology}
设 \(R\) 为主理想整环，\(C\) 为自由 \(R\)-模链复形，\(A\) 为 \(R\)-模。对每个整数 \(n\)，有短正合列
\[
 0\longrightarrow H_n(C)\otimes_R A
 \xrightarrow{\kappa}H_n(C\otimes_R A)
 \xrightarrow{\rho}\operatorname{Tor}_1^R(H_{n-1}(C),A)
 \longrightarrow0.
\]
此列对自由链复形之间的链映射 \(C\to C'\) 和系数模同态 \(A\to A'\) 都协变自然。它分裂，但一般没有关于 \(C\) 自然的分裂。
\end{theorem}
\begin{proof}
令 \(Z_n=\ker \mathrm{d}_n\)、\(B_n=\operatorname{im}\mathrm{d}_{n+1}\)。由于 \(B_{n-1}\) 自由，序列 \(0\to Z_n\to C_n\to B_{n-1}\to0\) 分裂。张量后仍是短正合列，并拼成复形的短正合列，其中外侧复形的微分为零。其同调长正合列在所需位置给出
\[
 B_n\otimes A\longrightarrow Z_n\otimes A\longrightarrow H_n(C\otimes A)
 \longrightarrow B_{n-1}\otimes A\longrightarrow Z_{n-1}\otimes A.
\]
末箭头是包含 \(B_{n-1}\hookrightarrow Z_{n-1}\) 张量 \(A\)，这是由连接同态对提升的微分直接算出的。前一箭头的余核是 \(H_n(C)\otimes A\)；后一箭头的核由自由消解 \(0\to B_{n-1}\to Z_{n-1}\to H_{n-1}(C)\to0\) 识别为所写的 Tor 群。所有映射在选分裂之前已经定义，所以序列自然。

固定次数 \(n\)，选取投影 \(p_n:C_n\to Z_n\)，并记商映射为 \(q_n:Z_n\to H_n(C)\)。在 \(C\otimes A\) 的 \(n\) 次圈上定义
\[
 [x]\longmapsto(q_n\otimes1_A)(p_n\otimes1_A)(x).
\]
若 \(x\) 改变一个边界 \((\mathrm{d}_{n+1}\otimes1_A)(y)\)，由于 \(p_n\) 在 \(Z_n\) 上是恒等映射，而 \(\mathrm{d}_{n+1}\) 的像在 \(B_n\subseteq Z_n\) 中，这个差经 \(q_n\otimes1_A\) 映为零。因此公式在同调上良定义；它与 \(\kappa\) 的复合是恒等映射，故给出分裂。

这里的投影只在固定次数上使用。即使逐次选取 \(p_n\)，它们一般也不构成从 \(C\) 到零微分复形 \(Z\) 的链映射：\(p_{n-1}\mathrm{d}_n=\mathrm{d}_n\)，而链映射条件要求这个复合为零。上述分裂依赖补模的选择；下一小节给出无法随链复形自然选取分裂的例子。
\end{proof}

\begin{theorem}[上同调的万有系数定理]\label{b4:thm:uct-cohomology}
设 \(R\) 为主理想整环，\(C\) 为自由 \(R\)-模链复形，\(A\) 为 \(R\)-模。令上链复形 \(\operatorname{Hom}_R(C,A)^n=\operatorname{Hom}_R(C_n,A)\)，微分为 \(f\mapsto f\circ\mathrm{d}_{C,n+1}\)。则对每个整数 \(n\)，有短正合列
\[
 0\longrightarrow\operatorname{Ext}_R^1(H_{n-1}(C),A)
 \longrightarrow H^n(\operatorname{Hom}_R(C,A))
 \xrightarrow{r}\operatorname{Hom}_R(H_n(C),A)\longrightarrow0.
\]
此列对自由链复形之间的链映射 \(C\to C'\) 反变自然，对系数模同态 \(A\to A'\) 协变自然；前者在中间项诱导从 \(H^n(\operatorname{Hom}_R(C',A))\) 到 \(H^n(\operatorname{Hom}_R(C,A))\) 的映射。此列分裂，但一般没有关于 \(C\) 自然的分裂。
\end{theorem}
\begin{proof}
余圈 \(f:C_n\to A\) 零化 \(B_n\)，所以限制到 \(Z_n\) 后经 \(H_n(C)\) 分解；余边的限制为零，这定义了 \(r\)。由于 \(Z_n\) 在 \(C_n\) 中为直和因子，任意 \(H_n(C)\to A\) 都能先拉回 \(Z_n\) 再延拓至 \(C_n\)，故 \(r\) 满射。

\(r([f])=0\) 恰好表示 \(f|_{Z_n}=0\)，因而 \(f=t \mathrm{d}_n\)，其中 \(t:B_{n-1}\to A\)。这样的 \(f\) 是余边，当且仅当 \(t\) 能延拓至 \(C_{n-1}\)。因为 \(Z_{n-1}\) 也是直和因子，这又等价于 \(t\) 能延拓至 \(Z_{n-1}\)。于是
\[
 \ker r=\operatorname{Hom}_R(B_{n-1},A)/
 \operatorname{res}\operatorname{Hom}_R(Z_{n-1},A)
 =\operatorname{Ext}_R^1(H_{n-1}(C),A).
\]
链映射保持圈与边界，系数同态则与上述限制、分解及取商相容，故这个核的识别给出所述自然性。

选定投影 \(p_n:C_n\to Z_n\)，并记 \(q_n:Z_n\to H_n(C)\) 为商映射。对 \(\phi:H_n(C)\to A\)，取 \(n\) 次余链 \(\phi q_np_n\)。由于 \(p_n\) 在 \(Z_n\) 上是恒等映射，且 \(q_n\) 零化 \(B_n=\mathrm{d}_{n+1}(C_{n+1})\)，有
\[
 (\phi q_np_n)\mathrm{d}_{n+1}=0.
\]
所以它是余圈；限制到 \(Z_n\) 后为 \(\phi q_n\)，从而 \(\phi\mapsto[\phi q_np_n]\) 是 \(r\) 的右逆。这里构造的是固定次数上的余圈代表，无须把投影 \(p_n\) 本身视为链映射。
\end{proof}

上述 Hom 计算与\cref{b4:def:hom-complex}的内部 Hom 仅差每度的符号重标识，所得上同调相同。两个版本可与 Weibel\cite[Universal Coefficient Theorem for Homology 3.6.2, Universal Coefficient Theorem for Cohomology 3.6.5]{Weibel1994HomologicalAlgebra}对照；该书把同调版本先写在自由交换群上，本节使用同一证明得到主理想整环版本。

\subsection{自然短正合列与非典范分裂}
\label{b4:subsec:0903:02}

“自然短正合列”是指链映射 \(C\to C'\) 及系数同态 \(A\to A'\) 诱导整个正合列的交换图：同调版本对两变量都协变，上同调版本对 \(C\) 反变、对 \(A\) 协变。它不意味着左右两项在中间项中都已选定。自然的是第一项的像及第二项的商，而直接把中间项写成直和还要作额外选择。

\begin{example}\label{b4:exm:uct-no-natural-splitting}
下面固定系数 \(A=\mathbb F_2\)，说明分裂不能对链复形变量 \(C\) 自然。在 \(\mathbb Z\) 上取 \(C_1=\mathbb Za\oplus\mathbb Zb\)、\(C_0=\mathbb Zc\)，\(\mathrm{d}(a)=2c\)、\(\mathrm{d}(b)=0\)，其余项为零。链自同构
\[
 f_1(a)=a+b,\quad f_1(b)=b,\quad f_0(c)=c
\]
在 \(H_1(C)=\mathbb Zb\) 及 \(H_0(C)=\mathbb Z/2\) 上都诱导恒等映射，而在
\(H_1(C\otimes A)=A\bar a\oplus A\bar b\) 上是非平凡剪切。同调万有系数列左项是 \(A\bar b\)，商由 \(\bar a\) 表示。任何商的截面都把 \(1\) 送到 \(\bar a+\lambda\bar b\)，但 \(f\) 把它改成 \(\bar a+(\lambda+1)\bar b\)。所以没有与全部链映射相容的截面。

\ProseParagraphBreak
上同调版本也由同一例子说明。在 \(H^1(\operatorname{Hom}(C,A))\) 中，左端 Ext 项由 \(a^*\) 生成，右端 Hom 项由 \(b^*\) 表示。\(f^*(a^*)=a^*\)、\(f^*(b^*)=a^*+b^*\)，再一次排除了自然截面。
\end{example}

这个反例没有排除固定 \(C\) 后对系数变量的自然性。事实上，一旦固定各次数上的投影 \(p_n:C_n\to Z_n\)，前述两种分裂公式都与系数同态 \(A\to A'\) 相容。

\ProseParagraphBreak
如果 \(A\) 平坦，同调公式的 Tor 项消失，\(\kappa\) 就是自然同构；如果 \(H_{n-1}(C)\) 自由，两个版本的附加项都消失。这两种简化来自明确的消失条件，不需要选择分裂。

\subsection{系数变换的计算}
\label{b4:subsec:0903:03}

设 \(C\) 是有限秩自由整数链复形。若已由 Smith 标准形\footnote{史密斯（Henry John Stephen Smith，1826-1883），爱尔兰数学家，研究二次型、初等因子与数论。}算得
\[
 H_n(C)\simeq\mathbb Z^{r_n}\oplus\bigoplus_j\mathbb Z/d_{n,j},
 \qquad d_{n,j}>1,
\]
则改变系数不必重新化简整个张量复形。对 \(A=\mathbb Z/m\)，令 \(g_{n,j}=\gcd(d_{n,j},m)\)。长度一的循环模消解给出
\[
 (\mathbb Z/d)\otimes\mathbb Z/m\simeq\mathbb Z/\gcd(d,m),\qquad
 \operatorname{Tor}_1^{\mathbb Z}(\mathbb Z/d,\mathbb Z/m)
 \simeq\mathbb Z/\gcd(d,m).
\]
因此同调的同构类型是
\[
 H_n(C\otimes\mathbb Z/m)\simeq
 (\mathbb Z/m)^{r_n}\oplus\bigoplus_j\mathbb Z/g_{n,j}
 \oplus\bigoplus_j\mathbb Z/g_{n-1,j}.
\]
这个式子确定模的同构类型，并未指定内部直和因子。

\begin{example}\label{b4:exm:uct-integer-computation}
取 \(C_2=\mathbb Z\)、\(C_1=\mathbb Z^2\)、\(C_0=\mathbb Z\)，
\[
 \mathrm{d}_2(t)=(0,6t),\qquad \mathrm{d}_1(x,y)=4x.
\]
有 \(H_2(C)=0\)、\(H_1(C)=\mathbb Z/6\)、\(H_0(C)=\mathbb Z/4\)。取 \(A=\mathbb Z/2\)，两个微分都变成零，故新的同调依次为 \(A,A^2,A\)（次数 \(2,1,0\)）。在次数 \(1\)，其中一份 \(A\) 来自 \(H_1(C)\otimes A\)，另一份来自 \(\operatorname{Tor}_1(H_0(C),A)\)；次数 \(2\) 完全来自 \(\operatorname{Tor}_1(H_1(C),A)\)。

\ProseParagraphBreak
取上同调的整数系数，则 \(H^0(\operatorname{Hom}(C,\mathbb Z))=0\)、\(H^1=\mathbb Z/4\)、\(H^2=\mathbb Z/6\)。这些群都来自相邻次数的 Ext 项；直接对两个微分作转置也得到同样结果。因而“对偶以后只需对同调取 Hom”会漏去本例全部的非零上同调。
\end{example}

% !TeX root = ../../Book4.tex
% 纲要标题：### 10.4 Künneth 公式
% 纲要标签：b4:sec:0904
% 纲要位置：计划纲要.md，第 2703 行。

\section{Künneth 公式}
\label{b4:sec:0904}

张量两个复形的同调，由各自同调的张量积及一次 Tor 共同控制。域上 Tor 消失，而主理想整环上即使短正合列分裂，分裂也通常没有自然选择。

% 纲要内容：
%
% * 域和主理想整环的基本版本
% * Tor 项及成立假设
% * 与第三卷 Koszul 计算的联系
%
% ---
%

\subsection{域和主理想整环的基本版本}
\label{b4:subsec:0904:01}

\begin{theorem}[Künneth 公式]\label{b4:thm:kunneth-pid}
设 \(R\) 为主理想整环，\(C,D\) 为下有界的自由 \(R\)-模链复形。对每个整数 \(n\)，有对 \(C,D\) 的链映射自然的短正合列
\[
\begin{split}
0\longrightarrow&\ \bigoplus_{p+q=n}H_p(C)\otimes_R H_q(D)
 \xrightarrow{\kappa_n}H_n(C\otimes_R D)\\
 \longrightarrow&\ \bigoplus_{p+q=n-1}\operatorname{Tor}_1^R(H_p(C),H_q(D))
 \longrightarrow0.
\end{split}
\]
此列非典范地分裂。特别地，若 \(R=k\) 为域，则 \(\kappa_n\) 是自然同构。
\end{theorem}
\begin{proof}
把 \(Z_p=Z_p(C)\) 放在次数 \(p\)，组成零微分复形 \(Z\)；把 \(B_{p-1}=B_{p-1}(C)\) 放在次数 \(p\)，组成零微分复形 \(B'\)。短正合列 \(0\to Z\to C\to B'\to0\) 逐项分裂，因此与 \(D\) 张量后仍是短正合列。\(Z_p,B_p\) 都自由，由\cref{b4:prop:flat-coefficients-homology}，两侧复形的同调分别是
\[
 \bigoplus_{p+q=n}Z_p\otimes H_q(D),\qquad
 \bigoplus_{p+q=n}B_{p-1}\otimes H_q(D).
\]
在同调长正合列中，连接同态在每一项上由包含 \(B_{p-1}\hookrightarrow Z_{p-1}\) 张量 \(H_q(D)\) 给出。具体地，把边界元素提升到 \(C_p\)，再作总微分，来自 \(D\) 的一项对圈代表为零，剩下的就是该包含。这些包含张量映射的余核与核，由自由消解
\(0\to B_p\to Z_p\to H_p(C)\to0\) 分别识别为张量积及 \(\operatorname{Tor}_1\)，从而得到所写自然正合列。

还须证明分裂。记 \(i_p:B_p\hookrightarrow Z_p\) 为包含。因为 \(B_{p-1}\) 自由，可以为满射 \(\mathrm d_p:C_p\to B_{p-1}\) 选取截面 \(s_{p-1}:B_{p-1}\to C_p\)。于是
\[
 Z_p\oplus B_{p-1}\xrightarrow{\ \sim\ }C_p,
 \qquad (z,b)\longmapsto z+s_{p-1}(b).
\]
在这些坐标下，\(\mathrm d_p(z,b)=(i_{p-1}(b),0)\in Z_{p-1}\oplus B_{p-2}\)。因此 \(C\) 同构于两项复形的直和
\[
 C\simeq\bigoplus_p E(p),\qquad
 E(p):\quad B_p\xrightarrow{i_p}Z_p
 \quad(\text{次数 }p+1,p),
\]
其中 \(E(p)\) 的其余项为零。对 \(D\) 同样选取截面，得到
\[
 D\simeq\bigoplus_q F(q),\qquad
 F(q):\quad B_q(D)\hookrightarrow Z_q(D)
 \quad(\text{次数 }q+1,q),
\]
其中 \(F(q)\) 的其余项为零，且同调仅可能在次数 \(q\) 非零，其值为 \(H_q(D)\)。于是张量复形分解为 \(E(p)\otimes_R F(q)\) 的直和；下有界性保证每个总次数只涉及有限多个非零的 \((p,q)\)。

把 \(H_q(D)\) 看成只在次数 \(q\) 非零的复形。商映射 \(Z_q(D)\to H_q(D)\) 给出拟同构 \(F(q)\to H_q(D)\)。因为 \(E(p)\) 有界且自由，\cref{b4:prop:bounded-flat-tensor}说明
\[
 E(p)\otimes_R F(q)\longrightarrow E(p)\otimes_R H_q(D)
\]
也是拟同构。右侧是 \(B_p\otimes_R H_q(D)\to Z_p\otimes_R H_q(D)\)，放在次数 \(p+q+1,p+q\)。由 \(H_p(C)\) 的长度一自由消解，它在这两个次数的同调分别为 \(\operatorname{Tor}_1^R(H_p(C),H_q(D))\) 与 \(H_p(C)\otimes_R H_q(D)\)。因此在总次数 \(n\)，张量积部分来自 \(p+q=n\)，Tor 部分来自 \(p+q=n-1\)，得到直和分解。圈的张量代表元表明，前一部分恰是 \(\kappa_n\) 的像，后一部分通过自然正合列的商映射同构到其右端，因而给出一个分裂。

自然性来自前半段的圈模、边界模及同调长正合列：链映射保持这些结构，所以 \(\kappa_n\) 和右端商映射都是自然的。后半段所选截面却一般不与链映射相容，故所得直和分解没有自然选择。这与万有系数定理中自然短正合列和所选分裂的区别相同。

域上所有模平坦，Tor 项为零，故最后的结论成立。
\end{proof}

更一般的 Künneth 公式见 Weibel\cite[Theorem 3.6.3]{Weibel1994HomologicalAlgebra}：只要求一个复形的各项及边界模平坦，另一个复形可以任意，也不要求两者下有界。本节采用下有界自由复形这一较强的充分条件，使每条总次数对角线有限，并可直接沿用\cref{b4:prop:bounded-flat-tensor}的总化论证；下有界性并非 Künneth 公式本身的必要条件。

\subsection{Tor 项及成立假设}
\label{b4:subsec:0904:02}

\begin{example}\label{b4:exm:kunneth-two-cyclic}
设 \(m,n\) 为正整数，在 \(\mathbb Z\) 上分别取两项自由消解 \(C=(\mathbb Z\xrightarrow{m}\mathbb Z)\)、\(D=(\mathbb Z\xrightarrow{n}\mathbb Z)\)，次数均为 \(1,0\)。张量总复形是
\[
 0\longrightarrow\mathbb Z\xrightarrow{\binom{-n}{m}}
 \mathbb Z^2\xrightarrow{(m\ \ n)}\mathbb Z\longrightarrow0.
\]
令 \(g=\gcd(m,n)\)。第一个映射单射；第二个映射的像为 \(g\mathbb Z\)，核由 \((n/g,-m/g)\) 生成，而第一个映射的像是该核的 \(g\) 倍。因此
\[
 H_0(C\otimes D)=\mathbb Z/g,\qquad H_1(C\otimes D)=\mathbb Z/g,
 \qquad H_2(C\otimes D)=0.
\]
第二个非零群是 Tor 项；只取两个消解原有同调的张量积，会漏掉它。
\end{example}

定理中的自由性假设不能直接删去。仍取主理想整环 \(R=\mathbb Z\)，让 \(C,D\) 都是集中在次数零的 \(\mathbb Z/2\)。此时普通张量复形仅在次数零非零，但 \(\operatorname{Tor}_1^{\mathbb Z}(\mathbb Z/2,\mathbb Z/2)=\mathbb Z/2\)。若无条件套用上面的短正合列，在次数一会错误地要求零模满射到 \(\mathbb Z/2\)。这里应先取平坦消解再张量；普通张量复形本身没有包含该高阶信息。

\ProseParagraphBreak
主定理中的下有界性保证每条总次数对角线只有有限个非零项，因而上述有限对角总化工具适用。在一般环上，若要用无界复形替换张量积中的某个因子，还须检验它是否保持拟同构；逐项自由本身不能保证这一点，见\cref{b4:exm:unbounded-flat-tensor}。该反例的底环不是主理想整环，不能据此认定主定理中的下有界性是必要条件。

\subsection{与第三卷 Koszul 计算的联系}
\label{b4:subsec:0904:03}

第三卷第十五章的 Koszul 复形，可以看成两项复形的张量积。对交换环 \(R\) 中元素 \(x\)，写
\[
 K(x):\quad 0\longrightarrow R\xrightarrow{x}R\longrightarrow0
 \qquad(\text{次数 }1,0).
\]
多个元素的 Koszul 复形由 \(K(x_1)\otimes_R\cdots\otimes_R K(x_r)\) 给出。各因子次数一生成元的交换引入负号，这与外代数描述中的符号相同。第三卷命题15.2.2已经证明该张量识别，第15.3节证明正则序列时的正合性；这里说明换系数后如何读出新的同调。

\begin{proposition}\label{b4:prop:koszul-adjoin-element}
设 \(R\) 为交换含幺环，\(C\) 为 \(R\) 模链复形，\(x\in R\)。记 \(K(x)\) 为仅在次数 \(1,0\) 非零的链复形 \(R\xrightarrow{x}R\)，微分为乘 \(x\)。对每个整数 \(n\)，有自然短正合列
\[
 0\longrightarrow H_n(C)/xH_n(C)
 \longrightarrow H_n(C\otimes_R K(x))
 \longrightarrow\{\,z\in H_{n-1}(C):xz=0\,\}
 \longrightarrow0.
\]
这里不要求 \(R\) 为主理想整环，也不保证分裂。
\end{proposition}
\begin{proof}
张量复形在次数 \(n\) 是 \(C_n\oplus C_{n-1}\)，其微分为
\(\mathrm{d}(a,b)=(\mathrm{d}_Ca+(-1)^{n-1}xb,\mathrm{d}_Cb)\)。取包含 \(a\mapsto(a,0)\) 及投影 \(q_n(a,b)=(-1)^{n-1}b\)，便得到短正合列
\(0\to C\to C\otimes K(x)\xrightarrow{q}C[1]\to0\)，其中链记号下 \(C[1]_n=C_{n-1}\)、\(\mathrm{d}_{C[1]}=-\mathrm{d}_C\)。等式 \(q_{n-1}\mathrm{d}=-\mathrm{d}_Cq_n\) 直接保证投影是链映射。对于 \(C_{n-1}\) 中的圈 \(z\)，提升为 \((0,(-1)^{n-1}z)\)，微分是 \((xz,0)\)，故连接同态就是乘 \(x\)。因此同调长正合列的相关部分是
\[
 H_n(C)\xrightarrow{x}H_n(C)\longrightarrow H_n(C\otimes K(x))
 \longrightarrow H_{n-1}(C)\xrightarrow{x}H_{n-1}(C).
\]
取中间部分的核与余核即得结论。
\end{proof}

例如在 \(R=k[x]\) 上取两个相同元素 \(x,x\)。第一个复形 \(K(x)\) 只有 \(H_0=k\)，而 \(x\) 在该同调上作用为零，故 \(K(x,x)\) 的 \(H_0,H_1\) 都是 \(k\)。次数一的具体圈是 \(e_1-e_2\)，其倍数 \(x(e_1-e_2)\) 是次数二的边界（依生成元次序差一个负号）。这说明第二个 \(x\) 不是 \(R/(x)\) 上的非零因子，故重复元素不能形成正则序列。


\begin{chaptersummary}
扩张把 Ext 的余圈解释为正合列，投射消解的截短与推出使两种描述可以互相转换。Baer 和来自对角与加法，Yoneda 乘积来自拼接；它们分别对应余圈的相加和提升后的复合。

\ProseParagraphBreak
张量总复形中的符号保证微分平方为零，平坦性和有界性则承担保持同调的不同任务。在主理想整环上，自由子模的结构使万有系数与 Künneth 公式只需一次 Ext 或 Tor 修正。自然的是短正合列及其映射，分裂通常依赖选择。循环模的两项消解和 Koszul 复形都提供了能够直接核验这些差别的计算。
\end{chaptersummary}

\BookExercises

\begin{exercise}\label{b4:ex:09-extensions-parameters}
设 \(p\) 为素数，\(t\in\mathbb F_p\)。用生成元 \(x,y\) 和关系 \(px=0\)、\(py=tx\) 定义交换群 \(E_t\)。定义 \(i:\mathbb Z/p\mathbb Z\to E_t\) 为 \(i(\bar1)=x\)，定义 \(q:E_t\to\mathbb Z/p\mathbb Z\) 为 \(q(y)=\bar1\)、\(q(x)=\bar0\)。证明 \(0\to\mathbb Z/p\mathbb Z\xrightarrow{i}E_t\xrightarrow{q}\mathbb Z/p\mathbb Z\to0\) 是扩张，求各中间群的同构类型，并说明 Baer 和怎样作用于 \(t\)。
\end{exercise}
\begin{solution}
关系先把 \(b\) 降至 \(0\le b<p\)，再把 \(a\) 降至 \(0\le a<p\)，故每个元素都可写成 \(a x+b y\)。取 \(t\) 的整数代表 \(\widetilde t\)，关系子格由矩阵
\[
 M=\begin{pmatrix}p&-\widetilde t\\0&p\end{pmatrix}
\]
的两列生成。第一卷第 9.4 节的子格指数命题说明，非奇异整数方阵生成的子格在 \(\mathbb Z^2\) 中的指数等于行列式的绝对值；这里 \(|E_t|=|\det M|=p^2\)。上述 \(p^2\) 个代表因而互不相同，表示也唯一。因此 \(x\) 的阶为 \(p\)，核恰为它生成的群。当 \(t=0\) 时 \(E_t=(\mathbb Z/p\mathbb Z)^2\)；当 \(t\ne0\) 时 \(py=tx\ne0\)，所以 \(y\) 的阶为 \(p^2\)，且生成全群。取自由满射 \(\mathbb Z\to\mathbb Z/p\mathbb Z\)，将 \(1\) 提升到 \(y\)，在核 \(p\mathbb Z\) 上对应的映射为 \(p\mapsto t\)。故扩张类为 \(t\)，Baer 和把参数相加。中间群同构并不能消去两端指定映射的区别。
\end{solution}

\begin{exercise}\label{b4:ex:09-pushout-scalar}
设 \(R\) 为交换环，\(a\in R\) 为非零因子，\(N\) 为 \(R\)-模。证明在 \(\operatorname{Ext}_R^1(R/(a),N)\simeq N/aN\) 的识别下，沿乘 \(r:N\to N\) 推出扩张对应乘 \(r\)。扩张类 \(\bar n\) 在何时分裂？
\end{exercise}
\begin{solution}
自由消解为 \(0\to R\xrightarrow{a}R\to R/(a)\to0\)。类 \(\bar n\) 由把关系 \(a\) 送到 \(n\) 的映射表示；推出后这个映射变为把 \(a\) 送到 \(rn\)，故对应 \(\overline{rn}\)。由\cref{b4:thm:ext1-extensions}中分裂扩张对应零元的结论，原扩张分裂当且仅当 \(n\in aN\)。非零因子条件保证所写两项列确实是自由消解。
\end{solution}

\begin{exercise}\label{b4:ex:09-higher-not-isomorphism}
在域 \(k\) 上，比较二阶扩张
\(0\to k\to k^2\to k^2\to k\to0\)，其箭头依次为 \(a\mapsto(a,0)\)、\((x,y)\mapsto(y,0)\)、\((u,v)\mapsto v\)，与 \(0\to k\xrightarrow{1}k\xrightarrow{0}k\xrightarrow{1}k\to0\)。给出扩张态射，并说明为何不能要求中间项逐项同构。
\end{exercise}
\begin{solution}
前一列各位置的核与像都是所写坐标轴，所以正合。在两个中间项分别使用 \((x,y)\mapsto x\)、\((u,v)\mapsto v\)，两端使用恒等映射，得到前一列到后一列的扩张态射；中间方块的两条复合都是零。因此两扩张等价。但中间项的维数分别为二和一，不可能逐项同构。
\end{solution}

\begin{exercise}\label{b4:ex:09-dual-number-product}
设 \(k\) 为域，\(R=k[\varepsilon]/(\varepsilon^2)\)，把 \(k\) 识别为 \(R/(\varepsilon)\)。用周期自由消解计算 \(\operatorname{Ext}_R^*(k,k)\) 的 Yoneda 乘积，证明它是一个次数一元生成的多项式环。若 \(\operatorname{char}k\ne2\)，这说明了什么？
\end{exercise}
\begin{solution}
消解每项为 \(R\)，微分乘 \(\varepsilon\)。施加 \(\operatorname{Hom}_R(-,k)\) 后微分为零，每次 Ext 均为 \(k\)。次数一生成元 \(t\) 由商映射 \(P_1=R\to k\) 表示。为计算乘积，把它提升为
\[
 t_i:P_{i+1}\longrightarrow P_i,\qquad t_i=1_R\quad(i\ge0).
\]
这些映射把消解次数降低一；它们的底层模映射都是恒等，但不是消解的同次数恒等映射。因两侧微分都乘 \(\varepsilon\)，有 \(\mathrm{d}_i t_i=t_{i-1}\mathrm{d}_{i+1}\)（\(i\ge1\)）；底端与增广复合就是代表 \(t\) 的商映射。按\cref{b4:prop:yoneda-resolution-product}的提升复合法，\(t^n\) 由 \(P_n\to k\) 的商映射代表。它非零，且因 Hom 复形的微分为零而不可能是余边，故生成该次数的一维 Ext 群。因此自扩张环为 \(k[t]\)、\(\deg t=1\)。若特征不为二，则 \(t^2\ne0\) 而分次交换律会要求 \(t^2=-t^2\)。这里讨论的是自扩张环的 Yoneda 乘积；即使底环 \(R\) 交换，一般模的自扩张环也不必分次交换。群上同调等场合的分次交换性还使用了相应乘积的额外结构。
\end{solution}

\begin{exercise}\label{b4:ex:09-tensor-sign}
在特征不为二的域 \(k\) 上，令 \(C=D=(k\xrightarrow{1}k)\)，次数为 \(1,0\)。若在张量复形中漏去 \((-1)^p\)，证明“微分”的平方不为零；使用正确符号时证明总复形正合。
\end{exercise}
\begin{solution}
记次数一生成元为 \(u\in C_1\)、\(v\in D_1\)，其微分为次数零生成元 \(u_0,v_0\)。不带符号的第一次像是
\[
 u_0\otimes v+u\otimes v_0
 \in(C_0\otimes D_1)\oplus(C_1\otimes D_0).
\]
再作用一次，两项都映到 \(u_0\otimes v_0\in C_0\otimes D_0\)，故平方为 \(2u_0\otimes v_0\ne0\)。使用正确符号时，按次数一项 \((C_1\otimes D_0)\oplus(C_0\otimes D_1)\) 的坐标次序，微分为 \(k\xrightarrow{(-1,1)}k^2\xrightarrow{(1,1)}k\)，前箭头单射、后箭头满射，中间核由 \((-1,1)\) 生成，故正合。
\end{solution}

\begin{exercise}\label{b4:ex:09-rational-coefficients}
设整数链复形 \(C\) 满足每个 \(H_n(C)\) 都是挠群。证明 \(C\otimes\mathbb Q\) 正合。这里是否必须要求 \(C_n\) 自由？
\end{exercise}
\begin{solution}
\(\mathbb Q\) 作为整数模平坦，\cref{b4:prop:flat-coefficients-homology}给出 \(H_n(C\otimes\mathbb Q)=H_n(C)\otimes\mathbb Q\)。挠元 \(x\) 若被 \(m\ne0\) 零化，则 \(x\otimes q=mx\otimes(q/m)=0\)，故右侧为零。这个证明不要求原复形逐项自由。
\end{solution}

\begin{exercise}\label{b4:ex:09-uct-free-and-torsion}
设有限秩自由整数复形 \(C\) 在相邻次数满足 \(H_n(C)=\mathbb Z^2\oplus\mathbb Z/6\mathbb Z\)、\(H_{n-1}(C)=\mathbb Z\oplus\mathbb Z/10\mathbb Z\)。求 \(H_n(C\otimes\mathbb Z/4\mathbb Z)\) 的同构类型及 \(H^n(\operatorname{Hom}(C,\mathbb Z))\) 的同构类型。
\end{exercise}
\begin{solution}
同调公式的左项为 \((\mathbb Z/4\mathbb Z)^2\oplus\mathbb Z/2\mathbb Z\)，右项为 \(\mathbb Z/2\mathbb Z\)。单凭这两个端项还不能确定中间群的同构类型；这里可用\cref{b4:thm:uct-homology}，因为整数环是主理想整环且 \(C\) 逐项自由，其万有系数短正合列分裂。因此答案是 \((\mathbb Z/4\mathbb Z)^2\oplus(\mathbb Z/2\mathbb Z)^2\)。上同调公式的 Ext 项为 \(\mathbb Z/10\mathbb Z\)，Hom 项为 \(\mathbb Z^2\)；由\cref{b4:thm:uct-cohomology}的分裂，答案是 \(\mathbb Z^2\oplus\mathbb Z/10\mathbb Z\)。这些是同构类型，不指定自然的内部直和。
\end{solution}

\begin{exercise}\label{b4:ex:09-modp-betti}
设 \(C\) 为有限秩自由整数链复形，\(H_n(C)=\mathbb Z^{r_n}\oplus\bigoplus_j\mathbb Z/d_{n,j}\mathbb Z\)。固定素数 \(p\)，以 \(t_n\) 表示被 \(p\) 整除的 \(d_{n,j}\) 的个数。证明 \(\dim_{\mathbb F_p}H_n(C\otimes\mathbb F_p)=r_n+t_n+t_{n-1}\)。
\end{exercise}
\begin{solution}
对每个循环直和项，张量 \(\mathbb F_p\) 及一次 Tor 都是零或 \(\mathbb F_p\)，其非零条件恰为 \(p\mid d_{n,j}\)。自由部分张量贡献 \(r_n\)，而一次 Tor 为零。将这些维数代入万有系数短正合列，维数可加，得到公式；此步不需要选分裂。
\end{solution}

\begin{exercise}\label{b4:ex:09-kunneth-numbers}
取整数两项复形 \(C=(\mathbb Z\xrightarrow{6}\mathbb Z)\)、\(D=(\mathbb Z\xrightarrow{15}\mathbb Z)\)，次数均为 \(1,0\)。求总复形各次同调，并给出一次同调的圈生成元。
\end{exercise}
\begin{solution}
总微分为 \(\mathbb Z\xrightarrow{(-15,6)}\mathbb Z^2\xrightarrow{(6,15)}\mathbb Z\)。后者的像为 \(3\mathbb Z\)，核由 \((5,-2)\) 生成；前者的像由 \(-3(5,-2)\) 生成。因此 \(H_0=H_1=\mathbb Z/3\mathbb Z\)，\(H_2=0\)，一次类由 \([(5,-2)]\) 生成。这分别对应张量项及一次 Tor 项。
\end{solution}

\begin{exercise}\label{b4:ex:09-kunneth-flat-homology}
设 \(R\) 为主理想整环，\(C,D\) 为下有界自由复形，且每个 \(H_p(C)\) 平坦。证明 \(\kappa_n\) 为自然同构。若只知每个 \(C_p\) 平坦，为什么不能作此结论？
\end{exercise}
\begin{solution}
同调平坦使 Künneth 正合列的每个 \(\operatorname{Tor}_1(H_p(C),H_q(D))\) 为零，因此自然映射 \(\kappa_n\) 是同构。逐项平坦则只是复形各项的性质，商 \(Z_p/B_p\) 可能不平坦；上一题的 \(C_p,D_q\) 都自由，但一次同调全由 Tor 产生，\(\kappa_1\) 的定义域为零而值域非零。
\end{solution}

\begin{exercise}\label{b4:ex:09-euler-tensor}
设 \(C,D\) 为同一域上有限维向量空间的有界链复形。记
\[
\chi(C)=\sum_n(-1)^n\dim H_n(C).
\]
证明 \(\chi(C\otimes D)=\chi(C)\chi(D)\)。
\end{exercise}
\begin{solution}
域上的 Künneth 同构给出 \(\dim H_n(C\otimes D)=\sum_{p+q=n}\dim H_p(C)\dim H_q(D)\)。所有求和有限，乘 \((-1)^n\) 求和后以 \((-1)^{p+q}=(-1)^p(-1)^q\) 分离成两个有限和的乘积，得到等式。有界及有限维条件保证欧拉数作为普通整数有定义。
\end{solution}

\begin{exercise}\label{b4:ex:09-koszul-zero-factor}
设 \(R\) 为交换环，\(x\in R\) 为非零因子。求 \(K(x,0)=K(x)\otimes_R K(0)\) 的全部同调，并解释为什么不需要假设 \(R\) 为主理想整环。
\end{exercise}
\begin{solution}
令 \(C=K(x)\)。\(K(0)\) 只在次数 \(1,0\) 各有一个 \(R\)，微分为零，所以
\[
 (C\otimes_RK(0))_n=C_n\oplus C_{n-1},\qquad
 \mathrm{d}(a,b)=(\mathrm{d}_Ca,\mathrm{d}_Cb).
\]
本卷的链移位 \(C[1]_n=C_{n-1}\) 带微分 \(-\mathrm{d}_C\)。因此
\[
 (a,b)\longmapsto(a,(-1)^{n-1}b)
\]
给出 \(C\otimes_RK(0)\to C\oplus C[1]\) 的复形同构：第二分量在次数 \(n-1\) 的符号是 \((-1)^{n-2}\)，恰好等于对次数 \(n\) 的符号再乘负号。\(x\) 为非零因子使 \(C\) 在次数一的同调为零，在次数零的同调为 \(R/(x)\)，故 \(K(x,0)\) 的 \(H_0,H_1\) 都为 \(R/(x)\)，其他同调为零。这个直接复形分解在任意交换环上成立，不需要 Künneth 的主理想整环假设。
\end{solution}

\begin{exercise}\label{b4:ex:09-extension-zero-product}
设 \(R\) 为主理想整环，\(A,B,C\) 为 \(R\)-模。证明任意 \(\alpha\in\operatorname{Ext}_R^1(B,C)\)、\(\beta\in\operatorname{Ext}_R^1(A,B)\) 的 Yoneda 乘积为零，但两因子可以都非零。
\end{exercise}
\begin{solution}
任意 \(R\)-模有长度至多一的自由消解，所以 \(\operatorname{Ext}_R^2(A,C)=0\)，乘积必为零。取 \(R=\mathbb Z\)、\(A=B=C=\mathbb Z/2\mathbb Z\)，一次 Ext 群为 \(\mathbb Z/2\mathbb Z\)，两个因子都可取其非零元。这里的零指拼接所得二阶扩张在 Yoneda 等价下为零，并不要求其中间箭头为零，也不要求中间项与零扩张逐项同构；尤其不意味着两个短扩张至少有一个分裂。
\end{solution}

\begin{exercise}\label{b4:ex:09-uct-coefficient-map}
设 \(m,r\) 为正整数，\(C\) 为自由整数链复形，\(H_n(C)=0\)、\(H_{n-1}(C)=\mathbb Z/m\mathbb Z\)。用上同调万有系数定理求系数同态 \(\mathbb Z\to\mathbb Z/r\mathbb Z\) 在 \(H^n(\operatorname{Hom}(C,-))\) 上诱导的映射。
\end{exercise}
\begin{solution}
没有 Hom 项，故这些群自然等于 \(\operatorname{Ext}_{\mathbb Z}^1(\mathbb Z/m\mathbb Z,A)=A/mA\)。系数同态诱导的映射是 \(\mathbb Z/m\mathbb Z\to(\mathbb Z/r\mathbb Z)/m(\mathbb Z/r\mathbb Z)\)，把整数类送到同一整数的像。右侧同构于 \(\mathbb Z/\gcd(m,r)\mathbb Z\)，所以该映射是相应的自然取商。这里同构来自自然短正合列中另一端为零，不涉及非典范分裂。
\end{solution}

\begin{exercise}\label{b4:ex:09-uct-splitting-coefficients}
设 \(R\) 为主理想整环，\(C\) 为自由 \(R\)-模链复形，固定整数 \(n\) 及投影 \(r_n:C_n\to Z_n(C)\)。对每个 \(R\)-模 \(A\)，记同调万有系数列的映射为 \(\kappa_A,\rho_A\)，记由 \(r_n\) 构造的左逆为
\[
 \lambda_A:H_n(C\otimes_R A)\longrightarrow H_n(C)\otimes_R A.
\]
证明 \(\lambda_A\) 与所有系数同态 \(A\to A'\) 相容，从而固定 \(C\) 后可选取关于 \(A\) 自然的分裂。若换一个投影得到 \(\lambda'_A\)，证明存在唯一的关于 \(A\) 自然的同态
\[
 \theta_A:\operatorname{Tor}_1^R(H_{n-1}(C),A)
 \longrightarrow H_n(C)\otimes_R A,
 \qquad \lambda'_A-\lambda_A=\theta_A\rho_A.
\]
\end{exercise}
\begin{solution}
\(\lambda_A\) 由 \(r_n\otimes1_A\) 后接圈模到同调模的商映射诱导；这两个映射都与系数同态相容，所以所得左逆关于 \(A\) 自然。\(\rho_A\) 限制为 \(\ker\lambda_A\) 到右端 Tor 群的同构：它的核是 \(\ker\lambda_A\cap\operatorname{im}\kappa_A=0\)；若 \(\rho_A(x)=t\)，则 \(x-\kappa_A\lambda_A(x)\) 属于 \(\ker\lambda_A\) 且仍映到 \(t\)。这些限制同构与系数同态相容，其逆给出关于 \(A\) 自然的截面。

因为两个左逆在 \(\operatorname{im}\kappa_A=\ker\rho_A\) 上相同，\((\lambda'_A-\lambda_A)\kappa_A=0\)，差映射唯一经余核 \(\rho_A\) 分解，得到 \(\theta_A\)。对系数同态 \(g:A\to A'\)，记中间同调映射为 \(h_g=H_n(1_C\otimes g)\)，右端 Tor 映射为 \(T_g=\operatorname{Tor}_1(H_{n-1}(C),g)\)。由已有映射的自然性，
\[
 \begin{aligned}
 \theta_{A'}T_g\rho_A
 &=\theta_{A'}\rho_{A'}h_g
 =(\lambda'_{A'}-\lambda_{A'})h_g\\
 &=(1\otimes g)(\lambda'_A-\lambda_A)
 =(1\otimes g)\theta_A\rho_A.
 \end{aligned}
\]
再用 \(\rho_A\) 的满射性消去最右侧的 \(\rho_A\)，即得 \(\theta_{A'}T_g=(1\otimes g)\theta_A\)，这正是 \(\theta\) 的自然性。这里固定了 \(C\) 的投影，并未要求这些投影与不同复形之间的链映射相容，故不与正文的非自然分裂例子冲突。
\end{solution}

% !TeX root = ../../Book4.tex
% 纲要标题：## 第11章　同调维数与零调对象
% 纲要标签：b4:ch:10
% 纲要位置：计划纲要.md，第 3463 行。

\chapter{同调维数与零调对象}
\label{b4:ch:10}
\BookChapterNumber{b4:ch:10}

\begin{chapterabstract}
本章以消解长度和高次 Ext、Tor 的消失刻画同调维数，研究何时可截短消解，以及用零调对象计算导出函子的条件。最后研究复合函子的导出，区分正合性与内射像的零调条件在两种直接复合公式中的作用，并引出 Grothendieck 谱序列。
\end{chapterabstract}

\begin{learninggoals}
\begin{itemize}
\item 区分投射、内射、平坦维数及环的左、右整体维数。
\item 用 Ext、Tor 的消失判据证明消解可以截短，而不把某一消解的表面长度当成不变量。
\item 判断指定对象是否对指定函子零调，并用适当消解计算导出函子。
\item 分清复合中前函子正合、后函子正合及内射像零调三个条件的不同作用。
\end{itemize}
\end{learninggoals}

令 \(R=k[\varepsilon]/(\varepsilon^2)\)，其中 \(k\) 是域。商模 \(k=R/(\varepsilon)\) 由 \(R\twoheadrightarrow k\) 给出；核是 \(\varepsilon R\)，它又是乘 \(\varepsilon\) 的像，而乘 \(\varepsilon\) 的核仍是 \(\varepsilon R\)。因此沿着乘 \(\varepsilon\) 可以不断向左接出自由模。单凭这条无限消解，还不能排除另有一条有限投射消解；关键是对它取 \(\operatorname{Hom}_R(-,k)\) 后微分全为零，得到每个正次数的 \(\operatorname{Ext}_R^i(k,k)\) 都非零。有限投射消解必使足够高次的 Ext 消失，所以这里确实无法用有限投射消解替代。这与整数上的 \(\mathbb Z/n\mathbb Z\) 有长度一的自由消解形成对照。本章将说明消解长度怎样由 Ext、Tor 的高次消失检测，以及何时能换用对某个函子零调的对象。

\ProseParagraphBreak
本章直接使用\chapref{b4:ch:07}、\chapref{b4:ch:08}的消解、导出函子长正合列以及\secref{b4:sec:0805}建立的维数平移公式；\secref{b4:sec:1002}用这些公式判断何时能够截短消解，并估计同调维数。局部环上有限生成模的例子调用第三卷第17.1—17.2节已证明的极小自由消解和投射维数判据。\secref{b4:sec:1004}列出 Grothendieck 谱序列的适用条件，证明见\chapref{b4:ch:15}。

\ProseParagraphBreak
同调维数可对照 Weibel 的《An Introduction to Homological Algebra》\cite[Section 4.1]{Weibel1994HomologicalAlgebra}；零调对象与消解计算见\cite[F-Acyclic Objects 2.4.3, Exercise 2.4.3, Section 2.5]{Weibel1994HomologicalAlgebra}。这里按本书的左模约定书写 Ext，并在每个 Tor 判据中明确另一变量为右模。

% !TeX root = ../../Book4.tex
% 纲要标题：### 11.1 同调维数
% 纲要标签：b4:sec:1001
% 纲要位置：计划纲要.md，第 3465 行。

\section{同调维数}
\label{b4:sec:1001}

一个模的某条消解可能很长，但这未必来自模本身。我们要问的是它能否有更短的消解，以及导出群怎样检测消解能在何处结束。

% 纲要内容：
%
% * 投射、内射及平坦维数
% * 整体维数
% * 有限维数与消失判据
%

\subsection{投射、内射及平坦维数}
\label{b4:subsec:1001:01}

具体消解的长度按其最高非零次数计数，不把增广目标计作消解的一项。例如，非零投射模 \(P\) 有长度零的消解 \(0\to P\xrightarrow{1_P}P\to0\)。取非零投射模 \(Q\)，在任意两个相邻的正次数上附加可缩复形
\[
0\longrightarrow Q\xrightarrow{1_Q}Q\longrightarrow0,
\]
仍得到 \(P\) 的投射消解，却可使最高非零次数任意大。去掉这部分是对具体消解作同伦化简；要定义模本身的维数，则须在全部允许的消解中取能够达到的最短长度，而不能只看眼前一条消解有多少非零项。

\begin{definition}\label{b4:def:homological-dimensions}
设 \(R\) 为含幺环，\(M\ne0\) 为幺性左 \(R\) 模。对非负整数 \(n\)，考虑正合列
\[
 0\longrightarrow P_n\longrightarrow\cdots\longrightarrow P_0
 \longrightarrow M\longrightarrow0
\]
若每个 \(P_i\)（\(0\le i\le n\)）投射，则在这类列中取最小长度 \(n\)，称为 \(M\) 的\term[toushe-weishu]{投射维数}[projective dimension]，记为 \(\operatorname{pd}_RM\)。把投射模换为平坦左 \(R\) 模，得到\term[pingtan-weishu]{平坦维数}[flat dimension] \(\operatorname{fd}_RM\)。若改用正合列
\[
 0\longrightarrow M\longrightarrow I^0\longrightarrow\cdots\longrightarrow I^n\longrightarrow0
\]
且每个 \(I^j\)（\(0\le j\le n\)）内射，则最小长度称为\term[neishe-weishu]{内射维数}[injective dimension]，记为 \(\operatorname{id}_RM\)。没有有限长度消解时，对应维数为 \(\infty\)；对零模三者均约定为 \(-\infty\)。右模的维数定义在 \(R^{\mathrm{op}}\) 上。零模的约定不表示存在负无限长度的消解，而是使零项不影响维数比较中的最大值，例如 \(M\oplus0=M\)，相应的最大值也应保持不变。本章对有限整数 \(r\) 约定
\[
\begin{gathered}
 (-\infty)\pm r=-\infty,\qquad (+\infty)\pm r=+\infty,
 \\
 \max\{-\infty,d\}=d
 \quad\bigl(d\in\mathbb Z\cup\{\pm\infty\}\bigr),
\end{gathered}
\]
并约定空集的上确界为 \(\sup\varnothing=-\infty\)。
\end{definition}
\symidx{\(\operatorname{pd}_RM\)：模的投射维数}
\symidx{\(\operatorname{id}_RM\)：模的内射维数}
\symidx{\(\operatorname{fd}_RM\)：模的平坦维数}

这些方向来自\secref{b4:sec:0701}中的构造：投射消解从到 \(M\) 的满射开始，逐次对核取投射满射，向左延伸；内射消解从 \(M\) 到内射模的单射开始，逐次把余核嵌入内射模，向右延伸。平坦消解与投射消解同向，只减弱了各项的要求。长度零时，正合列中唯一的消解项与 \(M\) 同构，所以非零模的三种维数为零，分别等价于它本身投射、内射、平坦。每个投射模都平坦，故每条投射消解也是平坦消解，总有 \(\operatorname{fd}_RM\le\operatorname{pd}_RM\)。

\ProseParagraphBreak
这个不等式可能严格。令 \(S=\mathbb Z\setminus\{0\}\)，则 \(\mathbb Q=S^{-1}\mathbb Z\)。第二卷命题6.4.6、6.4.7给出模的局部化与张量积的自然同构，以及局部化的正合性，利用整数环的交换性得到
\[
A\otimes_{\mathbb Z}\mathbb Q\cong S^{-1}A
\]
对交换群 \(A\) 自然，\(-\otimes_{\mathbb Z}\mathbb Q\) 保持短正合列，故 \(\mathbb Q\) 平坦。它却不投射：若投射，第二卷定理7.1.2使它成为某个自由交换群的直和因子，而\cref{b4:thm:pid-free-submodule}使这个子群自由。非零自由交换群中任一基向量都不能被 \(2\) 整除，与 \(\mathbb Q\) 可除矛盾。因此 \(\operatorname{pd}_{\mathbb Z}\mathbb Q>0\)。另一方面，\(\mathbb Z\) 是主理想整环，同一定理给出 \(\mathbb Q\) 的长度至多一的自由消解，故其投射维数不超过一。两者合得
\[
\operatorname{fd}_{\mathbb Z}\mathbb Q=0,
\qquad \operatorname{pd}_{\mathbb Z}\mathbb Q=1.
\]

\ProseParagraphBreak
同一个交换群改变底环后也会改变维数。设 \(p\) 为素数，\(\mathbb Z/p\mathbb Z\) 作为 \(\mathbb F_p\)-模投射维数为零，作为整数模的投射维数为一，而作为 \(\mathbb Z/p^2\mathbb Z\)-模有无限投射维数。对最后一个底环，令 \(S=\mathbb Z/p^2\mathbb Z\)，乘 \(p\) 的核与像均为 \(pS\)，反复使用这个映射便得到自由消解 \(\cdots\xrightarrow{p}S\xrightarrow{p}S\to\mathbb Z/p\mathbb Z\to0\)。施加 \(\operatorname{Hom}_S(-,\mathbb Z/p\mathbb Z)\) 后微分为零，故任意正次数的 Ext 自群都是非零的 \(\mathbb Z/p\mathbb Z\)，与有限投射消解在高次数没有项矛盾。\cref{b4:ex:10-three-base-rings}进一步比较这三个底环下的完整计算；群的大小不能决定这里的差别。

\subsection{整体维数}
\label{b4:subsec:1001:02}

\begin{definition}\label{b4:def:global-dimension}
设 \(R\) 为非零含幺环。所有幺性左 \(R\) 模投射维数的上确界
\[
 \operatorname{l.gl.dim}R=\sup\{\,\operatorname{pd}_RM:M\;\text{为左}\;R\text{-模}\,\}
\]
称为 \(R\) 的\term[zuozhengtiweishu]{左整体维数}[left global dimension]。右整体维数定义为 \(\operatorname{r.gl.dim}R=\operatorname{l.gl.dim}R^{\mathrm{op}}\)。交换环时记为 \(\operatorname{gl.dim}R\)。
\end{definition}
\symidx{\(\operatorname{l.gl.dim}R,\ \operatorname{r.gl.dim}R\)：环的左整体维数与右整体维数}
\symidx{\(\operatorname{gl.dim}R\)：交换环的整体维数}

域及除环的所有模都有基，故相应整体维数为零。非交换环的左、右整体维数不能不加说明地混写。这里的上确界也包括无限生成模；只讨论有限生成模时，得到的是另一个受限制的问题。

\ProseParagraphBreak
单个左模 \(M\) 投射，只保证以它为目标的满射有截面；左整体维数为零则要求每个左模都投射，因而每条短正合列 \(0\to A\to B\to C\to0\) 的末项 \(C\) 都投射，使该列分裂。反过来，若整个左模范畴的短正合列都分裂，对任意左模 \(M\) 取自由满射 \(F\to M\)，所得短正合列分裂，使 \(M\) 成为自由模的直和因子，因而投射。这就证明左整体维数为零当且仅当每条短正合左模列分裂；它测量的是整个模范畴的性质。第二卷第十三章的半单模结构理论进一步把这种性质与半单环联系起来。

\subsection{有限维数与消失判据}
\label{b4:subsec:1001:03}

消解能在第 \(d\) 次结束，会使所有更高次导出群消失。反过来，并不需要逐个检验这些高次数：维数平移把第 \(d+1\) 次的消失化为第 \(d\) 个合冲模或余合冲模的一次消失，再用投射、内射、平坦的一次判据，就能使消解在这一位置结束。这里必须让测试模遍历相应的全部模。

\begin{theorem}\label{b4:thm:dimension-vanishing}
设 \(R\) 为含幺环，\(M\) 为幺性左 \(R\) 模，\(d\ge0\) 为整数。有以下等价条件：
\begin{enumerate}
\item \(\operatorname{pd}_RM\le d\)，当且仅当对每个幺性左 \(R\) 模 \(N\) 有 \(\operatorname{Ext}_R^{d+1}(M,N)=0\)；这又等价于所有 \(i>d\) 的这些 Ext 群为零。
\item \(\operatorname{id}_RM\le d\)，当且仅当对每个幺性左 \(R\) 模 \(N\) 有 \(\operatorname{Ext}_R^{d+1}(N,M)=0\)；这又等价于所有 \(i>d\) 的这些 Ext 群为零。
\item \(\operatorname{fd}_RM\le d\)，当且仅当对每个幺性右 \(R\) 模 \(T\) 有 \(\operatorname{Tor}_{d+1}^R(T,M)=0\)；这又等价于所有 \(i>d\) 的这些 Tor 群为零。
\end{enumerate}
\end{theorem}
\begin{proof}
投射情形，长度至多为 \(d\) 的投射消解在第 \(d+1\) 次及以上没有项，故施加 \(\operatorname{Hom}_R(-,N)\) 后，高于 \(d\) 次的 Ext 为零。反过来，取任意投射消解，令 \(K_d=\Omega^dM\)，其中 \(K_0=M\)。由\cref{b4:prop:iterated-dimension-shift}，对每个左模 \(N\) 都有
\[
\operatorname{Ext}_R^1(K_d,N)\cong\operatorname{Ext}_R^{d+1}(M,N).
\]
若右边对所有 \(N\) 消失，则\cref{b4:thm:ext-projective-injective}使 \(K_d\) 投射。\(d>0\) 时，截取的正合列为
\[
0\longrightarrow K_d\longrightarrow P_{d-1}\longrightarrow\cdots
\longrightarrow P_0\longrightarrow M\longrightarrow0.
\]
把 \(K_d\) 放在第 \(d\) 次，就得到长度至多为 \(d\) 的投射消解；\(d=0\) 时判据直接使 \(M\) 投射。这也使所有更高次 Ext 自动消失。

内射情形，有限内射消解直接给出全部更高次 Ext 的消失。反过来，固定内射消解 \(M\to I^\bullet\)，令 \(L^0=M\)，并逐次取
\[
L^{j+1}=\operatorname{coker}(L^j\hookrightarrow I^j).
\]
其中每个单射由消解的正合性给出，见\cref{b4:prop:iterated-dimension-shift}中的余合冲构造。同一命题给出
\[
\operatorname{Ext}_R^1(N,L^d)\cong\operatorname{Ext}_R^{d+1}(N,M).
\]
若右边对所有左模 \(N\) 消失，一次内射判据使 \(L^d\) 内射。\(d>0\) 时以它作为最后一项，得到
\[
0\longrightarrow M\longrightarrow I^0\longrightarrow\cdots
\longrightarrow I^{d-1}\longrightarrow L^d\longrightarrow0,
\]
这是长度至多为 \(d\) 的内射消解；\(d=0\) 时检测的对象就是 \(L^0=M\)。

平坦情形，\cref{b4:thm:tor-flat-criterion}给出平坦左模与任意右模的正次 Tor 全部为零。把长度至多为 \(d\) 的平坦消解拆成短正合列，由\cref{b4:thm:tor-long-exact}逐次平移：高于 \(d\) 次的 Tor 最终化为最左平坦项的正次 Tor，故为零。反过来，可先取一个自由消解，仍令 \(K_d=\Omega^dM\)、\(K_0=M\)。\cref{b4:prop:iterated-dimension-shift}给出
\[
\operatorname{Tor}_1^R(T,K_d)\cong\operatorname{Tor}_{d+1}^R(T,M)
\]
对每个右模 \(T\) 成立。一次 Tor 对所有右模消失，使 \(K_d\) 平坦；\(d>0\) 时用与投射情形相同的截短列，原自由项与最左的 \(K_d\) 均平坦，因而得到长度至多为 \(d\) 的平坦消解。\(d=0\) 时则直接检测 \(M\) 的平坦性。这里取的仍是向左的合冲模，与内射情形的余合冲模不同。
\end{proof}

\begin{corollary}\label{b4:cor:global-dimension-ext}
设 \(R\) 为非零含幺环。左整体维数也等于所有幺性左 \(R\) 模内射维数的上确界，并且等于
\[
 \sup\{\,n\ge0:\operatorname{Ext}_R^n(M,N)\ne0
       \;\text{对某些左}\;R\text{-模}\;M,N\,\}.
\]
\end{corollary}
\begin{proof}
对每个整数 \(d\ge0\)，由上一条定理得到
\[
\begin{aligned}
\bigl(\forall M,\ \operatorname{pd}_RM\le d\bigr)
&\Longleftrightarrow
\bigl(\forall M\,\forall N,\ \operatorname{Ext}_R^{d+1}(M,N)=0\bigr)\\
&\Longleftrightarrow
\bigl(\forall N,\ \operatorname{id}_RN\le d\bigr),
\end{aligned}
\]
其中 \(M,N\) 都遍历幺性左 \(R\) 模。因而两种维数的全体取值具有相同的有限上界；没有有限上界时，两边的上确界都为 \(\infty\)，且对每个 \(d\ge0\) 都有一对 \(M,N\) 使第 \(d+1\) 次 Ext 非零。若共同上确界是有限整数 \(g>0\)，在 \(d=g-1\) 时上述统一消失不能成立，故有第 \(g\) 次的非零 Ext，而高于 \(g\) 次的 Ext 全为零。\(R\ne0\) 保证 \(\operatorname{Hom}_R(R,R)\ne0\)，涵盖 \(g=0\) 的情形。这证明两个上确界及所列 Ext 次数的上确界相等，并不要求同一个模的投射维数与内射维数相等。
\end{proof}

因此非域主理想整环 \(R\) 的整体维数恰为一：\cref{b4:thm:pid-free-submodule}给出所有模长度至多一的自由消解；取非零非单位 \(a\)，由\cref{b4:exm:ext-cyclic}得 \(\operatorname{Ext}_R^1(R/(a),R)=R/(a)\ne0\)，给出下界。

\begin{example}\label{b4:exm:infinite-dimension-dual-numbers}
设 \(k\) 为域，\(R=k[\varepsilon]/(\varepsilon^2)\)。把 \(k\) 识别为商模 \(R/(\varepsilon)\)，并由 \(R\) 的交换性同时视为左模与右模。它的周期自由消解为
\[
 \cdots\xrightarrow{\varepsilon}R\xrightarrow{\varepsilon}R
 \xrightarrow{\varepsilon}R\longrightarrow k\longrightarrow0.
\]
乘 \(\varepsilon\) 的核与像均为 \(\varepsilon R\)，也等于商映射的核，所以该列正合。它无限长，还不能排除 \(k\) 另有有限长度的消解。对它施加 \(\operatorname{Hom}_R(-,k)\) 或 \(-\otimes_Rk\)，因为 \(\varepsilon\) 在 \(k\) 上作用为零，全部微分都为零，而每一项都同构于 \(k\)。因此
\[
\operatorname{Ext}_R^i(k,k)\cong k,
\qquad \operatorname{Tor}_i^R(k,k)\cong k
\quad(i\ge0).
\]
对任意有限 \(d\ge0\)，第 \(d+1\) 次仍非零，\cref{b4:thm:dimension-vanishing}的三个判据遂分别排除有限投射、内射、平坦消解。于是 \(\operatorname{pd}_Rk=\operatorname{id}_Rk=\operatorname{fd}_Rk=\infty\)，且 \(\operatorname{gl.dim}R=\infty\)。无限维数的依据是任意高次仍有非零导出群，与所选消解的表面长度不同。
\end{example}

维数与消失条件的系统对应可参阅 Weibel\cite[Section 4.1, pd Lemma 4.1.6, id Lemma 4.1.8, fd Lemma 4.1.10]{Weibel1994HomologicalAlgebra}。

% !TeX root = ../../Book4.tex
% 纲要标题：### 11.2 截短消解与同调维数估计
% 纲要标签：b4:sec:1002
% 纲要位置：计划纲要.md，第 3471 行。

\section{截短消解与同调维数估计}
\label{b4:sec:1002}

高次导出群是否消失，可以转化为消解末端的合冲模是否投射。这个判据既判断给定消解能否截短，也说明局部环上的投射维数可以怎样检测。

% 纲要内容：
%
% * 截短消解及合冲模
% * 短正合列中的同调维数不等式
% * 第三卷局部投射维数所需的最低平移版本已在第9.5节证明；本节作一般化整理
% * 局部投射维数的讨论作为交换代数应用，可在学习第三卷后回读；不参与前两小节的证明
%

\subsection{截短消解及合冲模}
\label{b4:subsec:1002:01}

维数的定义只要求存在一条足够短的消解；下面的判据则作用于任意一条已经给定的消解。若 \(\operatorname{pd}_RM\le d\)，无论前 \(d\) 步选用了哪些投射满射，得到的第 \(d\) 个合冲模都投射，因而都能在此截短。内射方向有同样的结论；对于平坦维数，当前给定的各项只需平坦，不必投射。

\begin{proposition}\label{b4:prop:truncation-dimension}
设 \(R\) 为含幺环，\(M\) 为幺性左 \(R\) 模，\(d\ge1\) 为整数。给定幺性左 \(R\) 模正合列
\[
 0\longrightarrow K_d\longrightarrow P_{d-1}\longrightarrow\cdots
 \longrightarrow P_0\longrightarrow M\longrightarrow0,
\]
其中各 \(P_i\) 投射。则 \(\operatorname{pd}_RM\le d\) 当且仅当 \(K_d\) 投射。若幺性左 \(R\) 模序列 \(0\to M\to I^0\to\cdots\to I^{d-1}\to L^d\to0\) 正合且各 \(I^i\) 内射，则 \(\operatorname{id}_RM\le d\) 当且仅当 \(L^d\) 内射。若第一列的 \(P_i\) 仅要求平坦，则 \(\operatorname{fd}_RM\le d\) 当且仅当 \(K_d\) 平坦。
\end{proposition}
\begin{proof}
若末端对象具有所需性质，将它作为第 \(d\) 个投射项、内射项或平坦项，给定列本身就是长度至多 \(d\) 的相应消解。

反向令 \(K_0=M\)，把第一列拆成 \(0\to K_{j+1}\to P_j\to K_j\to0\)（\(0\le j<d\)）；这里 \(P_j\to K_j\) 是原微分到其像的满射，最后的核就是给定的 \(K_d\)。若各 \(P_j\) 投射，对任意幺性左模 \(N\) 依次用\cref{b4:thm:dimension-shift-basic}，得到
\[
 \operatorname{Ext}_R^1(K_d,N)\simeq\operatorname{Ext}_R^{d+1}(M,N).
\]
\(\operatorname{pd}_RM\le d\) 使右侧对所有 \(N\) 消失；\cref{b4:thm:ext-projective-injective}的一次判据便使这个给定的 \(K_d\) 投射。

内射列则令 \(L^0=M\)，依次取余核，拆成 \(0\to L^j\to I^j\to L^{j+1}\to0\)（\(0\le j<d\)）。第二变量的平移\cref{b4:prop:injective-dimension-shift}给出
\[
 \operatorname{Ext}_R^1(N,L^d)\simeq\operatorname{Ext}_R^{d+1}(N,M).
\]
若 \(\operatorname{id}_RM\le d\)，右侧对所有幺性左模 \(N\) 消失，一次 Ext 判据使 \(L^d\) 内射。

平坦版本仍拆第一列，但不使用 \(P_j\) 的投射性。对任意幺性右模 \(T\)，\cref{b4:thm:tor-flat-criterion}使 \(\operatorname{Tor}_i^R(T,P_j)=0\) 对每个 \(i>0\) 成立；\cref{b4:thm:tor-long-exact}于是逐步给出
\[
 \operatorname{Tor}_1^R(T,K_d)\simeq\operatorname{Tor}_{d+1}^R(T,M).
\]
若 \(\operatorname{fd}_RM\le d\)，右侧由\cref{b4:thm:dimension-vanishing}对所有 \(T\) 消失，同一平坦判据使 \(K_d\) 平坦。
\end{proof}

因此某个 \(d\) 步合冲模投射时，任何投射消解的 \(d\) 步合冲模都投射。这些合冲模不必同构，但加入投射直和项后同构，即稳定同构。具体地，两个投射消解的第一步给出短正合列 \(0\to K_1\to P_0\to M\to0\) 与 \(0\to K'_1\to P'_0\to M\to0\)，由\cref{b4:lem:schanuel}有
\[
 K_1\oplus P'_0\cong K'_1\oplus P_0.
\]
对第二步，记原消解给出的满射为 \(\pi_1:P_1\to K_1\)、\(\pi'_1:P'_1\to K'_1\)。分别补入恒等映射，得到
\[
\begin{aligned}
0&\longrightarrow K_2\longrightarrow P_1\oplus P'_0
 \xrightarrow{\pi_1\oplus1}K_1\oplus P'_0\longrightarrow0,\\
0&\longrightarrow K'_2\longrightarrow P'_1\oplus P_0
 \xrightarrow{\pi'_1\oplus1}K'_1\oplus P_0\longrightarrow0.
\end{aligned}
\]
恒等分量的核为零，所以两条新列的核仍分别是 \(K_2,K'_2\)；中间项都是投射模。利用第一步的同构将两个右端识别，再用 Schanuel 引理，得到
\[
 K_2\oplus P'_1\oplus P_0\cong K'_2\oplus P_1\oplus P'_0.
\]
同样，若在第 \(j\) 步已有 \(K_j\oplus Q\cong K'_j\oplus Q'\)，其中 \(Q,Q'\) 投射，就比较 \(P_j\oplus Q\to K_j\oplus Q\) 与 \(P'_j\oplus Q'\to K'_j\oplus Q'\)，恒等分量仍不改变核。再次应用该引理，便得到第 \(j+1\) 步的稳定同构。归纳可得投射模 \(Q,Q'\)，使 \(K_d\oplus Q\cong K'_d\oplus Q'\)。同调维数检测的是这些合冲模何时投射，并不要求它们作为模相互同构。

\subsection{短正合列中的同调维数不等式}
\label{b4:subsec:1002:02}

\begin{proposition}\label{b4:prop:dimension-short-exact}
设 \(R\) 为含幺环，\(0\to A\to B\to C\to0\) 为幺性左 \(R\) 模短正合列。以下所有维数均在 \(R\) 上计算。投射维数满足
\[
\begin{aligned}
 \operatorname{pd}B&\le\max\{\operatorname{pd}A,\operatorname{pd}C\},\\
 \operatorname{pd}A&\le\max\{\operatorname{pd}B,\operatorname{pd}C-1\},\\
 \operatorname{pd}C&\le\max\{\operatorname{pd}B,\operatorname{pd}A+1\}.
\end{aligned}
\]
平坦维数满足相同的三式。内射维数满足
\[
\begin{aligned}
 \operatorname{id}B&\le\max\{\operatorname{id}A,\operatorname{id}C\},\\
 \operatorname{id}A&\le\max\{\operatorname{id}B,\operatorname{id}C+1\},\\
 \operatorname{id}C&\le\max\{\operatorname{id}B,\operatorname{id}A-1\}.
\end{aligned}
\]
这里维数取值及 \(\pm\infty\) 的运算遵循\cref{b4:def:homological-dimensions}中的约定。
\end{proposition}
\begin{proof}
若某一右端为 \(\infty\)，该不等式直接成立。右端为负的情形由短列的正合性迫使被估计的模为零，结论也直接成立。以下考虑非负有限上界；零次 Ext、Tor 分别按 Hom、张量积解释。

固定幺性左模 \(N\)，\cref{b4:thm:ext-long-exact}在第一变量给出五项片段
\[
\begin{aligned}
 \operatorname{Ext}^{i-1}(A,N)&\longrightarrow\operatorname{Ext}^i(C,N)
 \longrightarrow\operatorname{Ext}^i(B,N)\\
 &\longrightarrow\operatorname{Ext}^i(A,N)
 \longrightarrow\operatorname{Ext}^{i+1}(C,N)
 \qquad(i\ge1).
\end{aligned}
\]
估计 \(B\) 时，取 \(i>\max\{\operatorname{pd}A,\operatorname{pd}C\}\)，它两侧的同次数 Ext 均为零，故 \(\operatorname{Ext}^i(B,N)=0\)。估计 \(A\) 时，条件
\(i>\max\{\operatorname{pd}B,\operatorname{pd}C-1\}\)
同时给出 \(i>\operatorname{pd}B\) 与 \(i+1>\operatorname{pd}C\)，使夹住 \(\operatorname{Ext}^i(A,N)\) 的两项消失。估计 \(C\) 时，条件
\(i>\max\{\operatorname{pd}B,\operatorname{pd}A+1\}\)
则给出 \(i>\operatorname{pd}B\) 与 \(i-1>\operatorname{pd}A\)，使 \(\operatorname{Ext}^{i-1}(A,N)\) 与 \(\operatorname{Ext}^i(B,N)\) 消失。这说明三个最大值中的次数偏移恰好来自 \(i+1\) 和 \(i-1\)。以上对每个 \(N\) 成立，\cref{b4:thm:dimension-vanishing}便给出三条投射维数不等式。

平坦情形固定幺性右模 \(T\)，\cref{b4:thm:tor-long-exact}给出
\[
\begin{aligned}
 \operatorname{Tor}_{i+1}(T,C)&\longrightarrow\operatorname{Tor}_i(T,A)
 \longrightarrow\operatorname{Tor}_i(T,B)\\
 &\longrightarrow\operatorname{Tor}_i(T,C)
 \longrightarrow\operatorname{Tor}_{i-1}(T,A)
 \qquad(i\ge1).
\end{aligned}
\]
\(B\) 仍被 \(A,C\) 的同次数项夹住；\(A\) 被 \(C\) 的 \(i+1\) 次与 \(B\) 的 \(i\) 次项夹住；\(C\) 被 \(B\) 的 \(i\) 次与 \(A\) 的 \(i-1\) 次项夹住。消失所需的次数条件与投射情形相同，故得到同样的三式。

内射情形固定幺性左模 \(N\)，改用第二变量的片段
\[
\begin{aligned}
 \operatorname{Ext}^{i-1}(N,C)&\longrightarrow\operatorname{Ext}^i(N,A)
 \longrightarrow\operatorname{Ext}^i(N,B)\\
 &\longrightarrow\operatorname{Ext}^i(N,C)
 \longrightarrow\operatorname{Ext}^{i+1}(N,A)
 \qquad(i\ge1).
\end{aligned}
\]
估计 \(B\) 仍只用 \(A,C\) 的同次数项。估计 \(A\) 却要令 \(i>\operatorname{id}B\)、\(i-1>\operatorname{id}C\)，所以其上界是 \(\max\{\operatorname{id}B,\operatorname{id}C+1\}\)；这时左侧的 \(\operatorname{Ext}^{i-1}(N,C)\) 与右侧的 \(\operatorname{Ext}^i(N,B)\) 都为零，正合性使 \(\operatorname{Ext}^i(N,A)=0\)。估计 \(C\) 则要令 \(i>\operatorname{id}B\)、\(i+1>\operatorname{id}A\)，上界便是 \(\max\{\operatorname{id}B,\operatorname{id}A-1\}\)。因此两处偏移与投射情形不同；对所有 \(N\) 使用内射消失判据即得结论。
\end{proof}

这些不等式也给出精确的一步下降。在 \(0\to K\to P\to M\to0\) 中，若 \(P\) 投射且 \(\operatorname{pd}M=d\ge1\) 有限，则 \(P\ne0\)，故中间项的维数固定为 \(\operatorname{pd}P=0\)。两个不同的不等式于是能分别给出 \(K\) 的上界与下界。第二个投射维数不等式给出
\[
 \operatorname{pd}K\le\max\{\operatorname{pd}P,\operatorname{pd}M-1\}
 =\max\{0,d-1\}=d-1.
\]
第三个投射维数不等式则给出
\[
 d=\operatorname{pd}M\le\max\{\operatorname{pd}P,\operatorname{pd}K+1\}
 =\max\{0,\operatorname{pd}K+1\}.
\]
因 \(d\ge1\)，后一式迫使 \(\operatorname{pd}K\ge d-1\)，从而 \(\operatorname{pd}K=d-1\)。特别地，\(d=1\) 时两式分别给出 \(\operatorname{pd}K\le0\) 与 \(\operatorname{pd}K\ge0\)，所以 \(K\) 是非零投射模。若 \(d=0\)，短列分裂，\(K\) 投射或为零；不能强行套用数值 \(d-1=-1\)。

\subsection{交换代数应用：局部投射维数}
\label{b4:subsec:1002:03}

一般环上的消失判据需要对所有测试模作检验。能否只选一个剩余域作测试，取决于消解在这个域上能保留什么信息：若微分仍非零，单个同调群为零只表示该处的核等于像，并不能说明对应自由项为零。

\ProseParagraphBreak
% 跨卷引用目标：b3:prop:minimal-resolution-existence（17.1.4）。
% 跨卷引用目标：b3:prop:minimal-resolution-betti（17.1.5）。
设 \((R,\mathfrak m,k)\) 为交换 Noether 局部环，\(M\ne0\) 为有限生成的 \(R\) 模。取各项均为有限秩的极小自由消解 \(F_\bullet\to M\)，即每个正次微分满足 \(\mathrm{d}_i(F_i)\subseteq\mathfrak mF_{i-1}\)。因为 \(k=R/\mathfrak m\)，微分矩阵与 \(k\) 张量后全部变为零；向 \(k\) 作 Hom 时，\(\mathfrak m\) 的作用也为零。这样的消解由第三卷命题17.1.4构造，命题17.1.5据此给出
\[
 \operatorname{Tor}_i^R(k,M)\simeq k\otimes_R F_i,\qquad
 \operatorname{Ext}_R^i(M,k)\simeq\operatorname{Hom}_R(F_i,k).
\]
这里 \(R\) 交换，\(k\) 同时视为右 \(R\) 模，Tor 的平衡计算使我们可以用第二变量的自由消解 \(F_\bullet\)。两个复形已无非零微分，故每个整数 \(i\ge0\) 的同调就是该项本身，模 \(M\) 的第 \(i\) 个 Betti 数\footnote{贝蒂（Enrico Betti，1823-1892），意大利数学家，研究代数与拓扑。}为
\[
 \beta_i^R(M)=\operatorname{rank}_R F_i
 =\dim_k\operatorname{Tor}_i^R(k,M).
\]
% 跨卷引用目标：b3:thm:nakayama（3.2.3）。
每个 \(F_i\) 有限生成，第三卷定理3.2.3的 Nakayama 引理使 \(k\otimes_RF_i=F_i/\mathfrak mF_i=0\) 等价于 \(F_i=0\)。还须检查零项之后会不会重新出现非零项。若 \(F_i=0\)，则 \(\mathrm{d}_{i+1}:F_{i+1}\to F_i\) 的目标为零，因而整个 \(F_{i+1}\) 都是其核；消解在 \(F_{i+1}\) 处正合，把这个核识别为 \(\mathrm{d}_{i+2}\) 的像；极小性又使该像包含在 \(\mathfrak mF_{i+1}\) 中。因此
\[
F_{i+1}=\ker\mathrm{d}_{i+1}=\operatorname{im}\mathrm{d}_{i+2}
\subseteq\mathfrak mF_{i+1}.
\]
由 \(F_{i+1}\) 有限生成，再用 Nakayama 引理得到 \(F_{i+1}=0\)。递推后所有更高项都为零；一般正合复形在一个零项后仍可接上非零的两项恒等复形，极小性排除了这种情形。

\ProseParagraphBreak
% 跨卷引用目标：b3:thm:local-pd-vanishing（17.2.3）。
对整数 \(d\ge0\)，若 \(\operatorname{Tor}_{d+1}^R(k,M)=0\)，便有 \(F_{d+1}=0\)，后续各项也为零，因而这条极小消解的长度不超过 \(d\)。反过来，长度不超过 \(d\) 的投射消解也是平坦消解，所以 \(\operatorname{pd}_RM\le d\) 使所有高于 \(d\) 次的 Tor 消失，尤其使这个剩余域测试消失。Ext 的计算同样检测 \(F_i\) 是否为零。因此
\[
\begin{aligned}
\operatorname{pd}_RM
&=\sup\{\,i\ge0:\operatorname{Tor}_i^R(k,M)\ne0\,\}\\
&=\sup\{\,i\ge0:\operatorname{Ext}_R^i(M,k)\ne0\,\},
\end{aligned}
\]
这也解释了第三卷定理17.2.3中只检验剩余域的判据。这里的局部性、极小性及各项有限生成性共同保证检测有效。

\ProseParagraphBreak
例如，设 \(k\) 为域。在 \(R=k[x,y]_{(x,y)}\) 上，剩余域的 Koszul 消解为
\[
 0\longrightarrow R\xrightarrow{\binom{-y}{x}}R^2
 \xrightarrow{(x\ \ y)}R\longrightarrow k\longrightarrow0.
\]
它极小，故 \(\operatorname{Tor}_2^R(k,k)=k\ne0\)，同时长度二给出高次消失，从而 \(\operatorname{pd}_Rk=2\)。相反，\cref{b4:exm:infinite-dimension-dual-numbers}的每一步都保留一个自由项，有限生成并不意味着有限投射维数。在有限投射维数情形，第三卷第17.3节的 Auslander--Buchsbaum 公式\footnote{奥斯兰德（Maurice Auslander，1926-1994），美国数学家，研究同调代数与表示论；布赫斯鲍姆（David Alvin Buchsbaum，1929-2021），美国数学家，研究同调代数及交换代数中的自由消解。}进一步将投射维数与模和底环的深度联系起来。

% !TeX root = ../../Book4.tex
% 纲要标题：### 11.3 零调对象
% 纲要标签：b4:sec:1003
% 纲要位置：计划纲要.md，第 3478 行。

\section{零调对象}
\label{b4:sec:1003}

计算一个指定函子时，消解项不必全部内射，但必须对该函子没有正次导出贡献。仅有原列正合并不足够，逐项的零调条件还需证明。

% 纲要内容：
%
% * 相对于函子的零调性
% * 使用零调消解计算导出函子
% * 不能随意用任意正合复形替代消解
%

\subsection{相对于函子的零调性}
\label{b4:subsec:1003:01}

内射消解对每个左正合函子都可用，但有时它的各项过大，实际计算并不方便。若只计算一个指定函子 \(F\)，可以使用对这个函子没有高次导出贡献的对象。

\begin{definition}\label{b4:def:functor-acyclic}
设 \(\mathcal A,\mathcal B\) 为交换范畴，\(\mathcal A\) 有足够多内射对象，\(F:\mathcal A\to\mathcal B\) 为左正合加性函子。若对象 \(Q\in\mathcal A\) 满足 \(R^iF(Q)=0\) 对所有整数 \(i>0\) 成立，则称 \(Q\) 为 \(F\)-\term[lingtiao-duixiang]{零调对象}[acyclic object]，也称 \(F\)-无上同调对象\termidx[wushangtongdiaoduixiang]{无上同调对象}。给定 \(\mathcal A\) 中的对象 \(M\) 及非负增广上链复形
\[
 0\longrightarrow M\longrightarrow A^0\longrightarrow A^1\longrightarrow\cdots
\]
若此序列正合，且每个 \(A^j\)（\(j\ge0\)）都为 \(F\)-零调对象，则称为 \(M\) 的 \(F\)-\term[lingtiao-xiaojie]{零调消解}[acyclic resolution]。
\end{definition}

这里的零调性描述一个对象相对于指定函子 \(F\) 的高次导出行为，只要求正次数的导出函子消失，不要求 \(F(Q)=0\)。\cref{b4:def:acyclic-complex}中的零调复形则描述整条复形内部的同调，要求所有次数的同调对象为零。零调消解要求各项对 \(F\) 零调；去掉增广项后的复形有 \(H^0(A^\bullet)\cong M\)，而施加 \(F\) 后的正次数上同调也可能非零。

\ProseParagraphBreak
内射对象对每个左正合加性函子都是零调对象；对于固定的函子，零调对象却未必内射。例如取 \(F=\operatorname{Hom}_{\mathbb Z}(\mathbb Z/2\mathbb Z,-)\)。对任意交换群 \(Q\)，长度一自由消解 \(0\to\mathbb Z\xrightarrow{\times2}\mathbb Z\to\mathbb Z/2\mathbb Z\to0\) 给出
\[
 R^1F(Q)\cong Q/2Q,\qquad R^iF(Q)=0\quad(i>1).
\]
因此 \(\mathbb Z/3\mathbb Z\) 对 \(F\) 零调，却因不可除而不是内射交换群；不过此时 \(F(\mathbb Z/3\mathbb Z)=0\)。若改取
\[
 Q=(\mathbb Q/\mathbb Z)\oplus\mathbb Z/3\mathbb Z,
\]
则 \(2Q=Q\)，但 \(3Q\ne Q\)，所以它仍对 \(F\) 零调而非内射。同时 \(F(Q)\cong Q[2]\cong\mathbb Z/2\mathbb Z\ne0\)，明确说明原函子的取值不必消失。

\begin{proposition}\label{b4:prop:acyclic-closure}
设 \(\mathcal A,\mathcal B\) 为交换范畴，\(\mathcal A\) 有足够内射对象，\(F:\mathcal A\to\mathcal B\) 为左正合加性函子，且 \(0\to A\to B\to C\to0\) 为 \(\mathcal A\) 中的短正合列。若 \(A,C\) 都 \(F\)-零调，则 \(B\) 也如此；若 \(A,B\) 都 \(F\)-零调，则 \(C\) 也如此。若 \(B,C\) 都 \(F\)-零调，则 \(A\) 零调当且仅当 \(F(B)\to F(C)\) 为满态射。
\end{proposition}
\begin{proof}
前两项分别由长正合列中的
\(R^iF(A)\to R^iF(B)\to R^iF(C)\) 及
\(R^iF(B)\to R^iF(C)\to R^{i+1}F(A)\) 得到，其中 \(i\ge1\)。当 \(B,C\) 都零调时，片段 \(R^{i-1}F(C)\to R^iF(A)\to R^iF(B)\) 已使 \(R^iF(A)=0\) 对 \(i\ge2\) 成立，但低次部分仍为
\[
 0\longrightarrow F(A)\longrightarrow F(B)\longrightarrow F(C)
 \longrightarrow R^1F(A)\longrightarrow0.
\]
所以唯一可能剩下的正次导出是 \(R^1F(A)\cong\operatorname{coker}(F(B)\to F(C))\)。它消失恰好要求 \(F(B)\to F(C)\) 满，这就是第三种情形多出的条件。
\end{proof}

\subsection{零调消解与导出函子的计算}
\label{b4:subsec:1003:02}

\begin{theorem}\label{b4:thm:acyclic-resolution}
设 \(\mathcal A,\mathcal B\) 为交换范畴，\(\mathcal A\) 有足够内射对象，\(F:\mathcal A\to\mathcal B\) 为左正合加性函子。若 \(0\to M\to A^0\to A^1\to\cdots\) 是 \(F\)-零调消解，则对每个 \(n\ge0\)，有典范同构
\[
 R^nF(M)\simeq H^n(F(A^\bullet)).
\]
此同构对增广消解之间的态射自然。
\end{theorem}
\begin{proof}
令 \(K^0=M\)，从增广嵌入 \(\iota^0:M\hookrightarrow A^0\) 起，逐次取 \(K^{j+1}=\operatorname{coker}\iota^j\)。消解正合性把这个商对象识别为 \(\operatorname{im}(\mathrm{d}_A^j)\)，并给出下一项中的嵌入 \(\iota^{j+1}:K^{j+1}\hookrightarrow A^{j+1}\)。这样前两步是
\[
\begin{gathered}
 0\longrightarrow M\longrightarrow A^0\longrightarrow K^1\longrightarrow0,\\
 0\longrightarrow K^1\longrightarrow A^1\longrightarrow K^2\longrightarrow0,
\end{gathered}
\]
一般的一步写成
\[
 0\longrightarrow K^j\xrightarrow{\iota^j} A^j\xrightarrow{p^j} K^{j+1}\longrightarrow0
 \qquad(j\ge0).
\]
先看一次上同调。微分 \(\mathrm{d}_A^1=\iota^2p^1\)，左正合性使 \(F(\iota^2)\) 单，并把 \(\ker F(p^1)\) 识别为 \(F(K^1)\)。因而施加 \(F\) 后的一次余圈对象是 \(F(K^1)\)，一次余边来自 \(F(A^0)\to F(K^1)\)。第一条短正合列的长正合列又因 \(A^0\) 零调而给出 \(F(A^0)\to F(K^1)\to R^1F(M)\to0\)，所以
\[
 H^1(F(A^\bullet))
 \cong\operatorname{coker}\bigl(F(A^0)\to F(K^1)\bigr)
 \cong R^1F(M).
\]
一般次数的计算把同一机制用于后面的短正合列。长正合列和 \(A^j\) 的零调性给出
\[
\begin{gathered}
 0\to F(K^j)\to F(A^j)\to F(K^{j+1})\to R^1F(K^j)\to0,\\
 R^iF(K^{j+1})\simeq R^{i+1}F(K^j)\qquad(i\ge1).
\end{gathered}
\]
消解的微分分解为
\[
 \mathrm{d}_A^n:\quad A^n\xrightarrow{p^n}K^{n+1}
 \xrightarrow{\iota^{n+1}}A^{n+1}.
\]
对第 \(n+1\) 步的短正合列用左正合性，得到 \(F(\iota^{n+1})\) 仍为单态射，所以复合它不改变 \(F(p^n)\) 的核。再对第 \(n\) 步用左正合性，\(F(\iota^n)\) 把 \(F(K^n)\) 识别为这个核。因此
\[
\begin{aligned}
 \ker\bigl(F(\mathrm{d}_A^n)\bigr)
 &=\ker\bigl(F(p^n)\bigr)\\
 &=F(K^n),
\end{aligned}
\]
即 \(F(A^\bullet)\) 在次数 \(n\) 的余圈对象为 \(F(K^n)\)。当 \(n\ge1\) 时，同调就是
\[
 \operatorname{coker}\bigl(F(A^{n-1})\to F(K^n)\bigr)
 \simeq R^1F(K^{n-1})\simeq R^nF(M).
\]
次数零则由左正合性直接等于 \(F(M)\)。这些同构只使用给定的正合列和连接同态，未选分裂，因而对消解态射自然。

为说明这里的“典范”也与最初采用内射消解的定义一致，取 \(M\to I^\bullet\)。下面只使用目标 \(I^j\) 的内射性，源 \(A^j\) 不必内射：先把 \(M\to I^0\) 沿 \(M\hookrightarrow A^0\) 延拓为 \(f^0:A^0\to I^0\)。若已构造到 \(f^j\)，则 \(\mathrm{d}_I^jf^j\) 在 \(K^j\) 上为零；在零次这是与增广相容，在正次这是已有链映射等式及 \(\mathrm{d}_I^j\mathrm{d}_I^{j-1}=0\) 的结果。于是该复合经余核 \(p^j\) 分解为 \(K^{j+1}\to I^{j+1}\)，再由 \(I^{j+1}\) 内射沿 \(\iota^{j+1}\) 延拓为 \(f^{j+1}\)。逐次得到增广上链映射 \(f^\bullet:A^\bullet\to I^\bullet\)：
\[
\begin{tikzcd}[column sep=2.2em,row sep=large]
0\arrow[r]&M\arrow[r]\arrow[d,equal]&A^0\arrow[r]\arrow[d,"f^0"]&A^1\arrow[r]\arrow[d,"f^1"]&\cdots\\
0\arrow[r]&M\arrow[r]&I^0\arrow[r]&I^1\arrow[r]&\cdots
\end{tikzcd}
\]
这里的比较映射通过延拓选取，其同伦唯一性由\cref{b4:prop:comparison-homotopy-unique}的内射版本保证。具体地，设 \(f^\bullet,g^\bullet:A^\bullet\to I^\bullet\) 都延拓 \(1_M\)，约定 \(I^{-1}=0\)、\(h^0=0\)。由于 \(f^0-g^0\) 在 \(M\) 上为零，它经 \(p^0:A^0\twoheadrightarrow K^1\) 分解为 \(K^1\to I^0\)。利用 \(I^0\) 的内射性，把后者沿 \(\iota^1:K^1\hookrightarrow A^1\) 延拓为 \(h^1:A^1\to I^0\)，于是 \(f^0-g^0=h^1\mathrm{d}_A^0\)。

假定已构造到 \(h^j\)（\(j\ge1\)），且低于 \(j\) 次的同伦等式成立。已选的 \(h^j\) 解释了由前一次数传来的差别，在次数 \(j\) 尚未消去的误差是
\[
 e^j=f^j-g^j-\mathrm{d}_I^{j-1}h^j.
\]
链映射等式及前一次的同伦等式给出
\[
\begin{aligned}
 e^j\mathrm{d}_A^{j-1}
 &=\mathrm{d}_I^{j-1}
   \bigl(f^{j-1}-g^{j-1}-h^j\mathrm{d}_A^{j-1}\bigr)\\
 &=\mathrm{d}_I^{j-1}\mathrm{d}_I^{j-2}h^{j-1}=0.
\end{aligned}
\]
这个等式说明误差在前一微分的像上为零。由于 \(\mathrm{d}_A^{j-1}=\iota^jp^{j-1}\) 且 \(p^{j-1}\) 为满态射，\(e^j\iota^j=0\)，故 \(e^j\) 经余核 \(p^j\) 下降为 \(\bar e^j:K^{j+1}\to I^j\)。再由 \(I^j\) 内射，把 \(\bar e^j\) 沿 \(\iota^{j+1}\) 延拓为 \(h^{j+1}:A^{j+1}\to I^j\)。这样逐次得到
\[
 f^j-g^j=\mathrm{d}_I^{j-1}h^j+h^{j+1}\mathrm{d}_A^j
 \qquad(j\ge0),
\]
即 \(f-g=\mathrm{d}_Ih+h\mathrm{d}_A\)。施加加性函子后，这个等式仍成立，所以两种比较映射诱导同一个同调映射。上面的短正合列在比较映射下构成交换图，连接同态自然性说明其诱导映射正是已经构造的同构。因此不同零调消解都通过给定的 \(R^nF(M)\) 典范比较，而不要求它们之间任意方向都有链映射。
\end{proof}

逐项零调并不使 \(F(A^\bullet)\) 成为零调复形。它只使各 \(A^j\) 的正次导出消失，却未要求 \(F(A^{j-1})\to F(K^j)\) 满。上面计算的余核正是这份满射性失败；沿短正合列平移后，它记录的是 \(R^jF(M)\)。

\ProseParagraphBreak
对有足够投射对象的交换范畴及右正合加性函子，对定义域和值域同时取相反范畴，就得到左导出版本：若 \(Q_\bullet\to M\) 正合，且 \(L_iF(Q_j)=0\) 对所有 \(i>0,j\ge0\) 成立，则 \(L_nF(M)\simeq H_n(F(Q_\bullet))\)。在相反范畴中，核与余核互换，内射延拓成为投射提升；再把上链次数改记为链次数，便得到这些正合列与同构，而不需额外假设。具体地，设 \(R\) 为含幺环，\(N\) 为幺性左 \(R\) 模。在幺性右 \(R\) 模范畴中取 \(F=-\otimes_RN\)，平坦右模 \(Q\) 满足 \(L_iF(Q)=\operatorname{Tor}_i^R(Q,N)=0\)（\(i>0\)）。因此任一平坦消解 \(\cdots\to Q_1\to Q_0\to M\to0\) 都给出
\[
 \operatorname{Tor}_n^R(M,N)\cong H_n(Q_\bullet\otimes_RN)
 \qquad(n\ge0),
\]
而不必要求各 \(Q_j\) 投射。

\subsection{正合复形作为消解的条件}
\label{b4:subsec:1003:03}

应用定理时，要检查增广列正合、各项对指定函子零调，并且消解从次数零沿一个方向展开。前两项给出连接同态的平移，最后一项保证每个固定的 \(n\ge2\) 经过有限次就回到 \(M\)：
\[
 R^1F(K^{n-1})\cong R^2F(K^{n-2})\cong\cdots\cong R^nF(K^0)
 =R^nF(M).
\]
一条向两端无限延伸的正合复形若没有这样的增广起点，即使局部短正合列允许平移，也不能据此完成与指定对象导出的比较。

\begin{example}\label{b4:exm:arbitrary-resolution-fails}
取 \(F=\operatorname{Hom}_{\mathbb Z}(\mathbb Z/2\mathbb Z,-)\) 和 \(M=\mathbb Z/2\mathbb Z\)。正合列
\[
 0\longrightarrow M\longrightarrow M\oplus\mathbb Z
 \longrightarrow\mathbb Z\longrightarrow0
\]
用自然包含和投影，确实是一个长度一的右消解。但施加 \(F\) 后，非增广上链复形为 \(\mathbb Z/2\mathbb Z\to0\)，一次上同调为零，而 \(R^1F(M)=\operatorname{Ext}_{\mathbb Z}^1(\mathbb Z/2\mathbb Z,\mathbb Z/2\mathbb Z)=\mathbb Z/2\mathbb Z\)。第零项的零调条件已经失败，因为
\[
 R^1F(M\oplus\mathbb Z)
 \cong (M\oplus\mathbb Z)/2(M\oplus\mathbb Z)
 \cong (\mathbb Z/2\mathbb Z)^{\oplus2}\ne0.
\]
末项 \(\mathbb Z\) 也有 \(R^1F(\mathbb Z)\cong\mathbb Z/2\mathbb Z\ne0\)。原列虽正合，却不是 \(F\)-零调消解，不能用它计算导出函子。
\end{example}

对同一个 \(M=\mathbb Z/2\mathbb Z\)，取
\[
 0\longrightarrow\mathbb Z/2\mathbb Z\xrightarrow{\eta}\mathbb Q/\mathbb Z
 \xrightarrow{\times2}\mathbb Q/\mathbb Z\longrightarrow0,
 \qquad \eta(\bar1)=\dfrac12+\mathbb Z.
\]
乘二映射为满态射，其核恰为 \(\eta(\mathbb Z/2\mathbb Z)\)，而 \(\mathbb Q/\mathbb Z\) 是可除群，因而此列为内射消解。\(F(\mathbb Q/\mathbb Z)\) 由 \(\mathbb Q/\mathbb Z\) 中被 \(2\) 零化的元素确定，同构于 \(\mathbb Z/2\mathbb Z\)；乘二在这些元素上为零，所以施加同一 \(F\) 后，非增广上链复形为
\[
 \mathbb Z/2\mathbb Z\xrightarrow{0}\mathbb Z/2\mathbb Z
 \qquad\text{（次数 }0,1\text{）}.
\]
于是
\[
 R^0F(\mathbb Z/2\mathbb Z)\simeq\mathbb Z/2\mathbb Z,
 \qquad R^1F(\mathbb Z/2\mathbb Z)\simeq\mathbb Z/2\mathbb Z.
\]
这两个正合消解针对同一个对象；前一个含有非零调项，后一个逐项内射，因此后者能够计算该对象的导出函子。

\ProseParagraphBreak
零调对象及消解计算可参阅 Weibel\cite[F-Acyclic Objects 2.4.3, Exercise 2.4.3, Section 2.5]{Weibel1994HomologicalAlgebra}。该书在左导出部分提出此类对象，并以维数平移给出练习；本节把当前所需的右导出版本作为正文定理证明。

% !TeX root = ../../Book4.tex
% 纲要标题：### 11.4 复合函子的导出
% 纲要标签：b4:sec:1004
% 纲要位置：计划纲要.md，第 3484 行。

\section{复合函子的导出}
\label{b4:sec:1004}

把内射消解先送入一个函子以后，其像未必适合计算另一个函子。检查内射像的零调性，才能分清直接复合导出与需要谱序列的情形。

% 纲要内容：
%
% * 函子作用于内射对象后的条件
% * 复合导出的障碍
% * Grothendieck 谱序列的假设预告
%
% ---
%

\subsection{内射对象在函子下的像}
\label{b4:subsec:1004:01}

取一个内射消解后，先作用 \(F\) 再作用 \(G\)，看似能够计算复合函子。但 \(F(I)\) 不一定内射，也不一定对 \(G\) 零调。这一步是复合导出问题中必须单独检查的对象条件。

\begin{proposition}\label{b4:prop:right-adjoint-injectives}
设 \(\mathcal A,\mathcal B\) 为交换范畴，\(U:\mathcal A\to\mathcal B\) 是加性函子，具有正合左伴随 \(V:\mathcal B\to\mathcal A\)。则 \(U\) 把内射对象送到内射对象。
\end{proposition}
\begin{proof}
设 \(I\) 为 \(\mathcal A\) 的内射对象。给定 \(\mathcal B\) 中单态射 \(j:X\hookrightarrow Y\) 及 \(f:X\to U(I)\)，先用伴随双射
\[
 \alpha_X:\operatorname{Hom}_{\mathcal B}(X,U(I))
 \xrightarrow{\sim}\operatorname{Hom}_{\mathcal A}(V(X),I)
\]
把 \(f\) 转置为 \(\widetilde f:V(X)\to I\)。因为 \(V\) 正合，\(V(j):V(X)\hookrightarrow V(Y)\) 仍单；由 \(I\) 内射，可延拓为 \(\widetilde h:V(Y)\to I\)，使 \(\widetilde hV(j)=\widetilde f\)。再令 \(h=\alpha_Y^{-1}(\widetilde h):Y\to U(I)\)。伴随双射对源对象的自然性给出
\[
 \alpha_X(hj)=\alpha_Y(h)V(j)
 =\widetilde hV(j)=\widetilde f=\alpha_X(f).
\]
因 \(\alpha_X\) 为双射，\(hj=f\)，所以所求延拓确实存在，\(U(I)\) 内射。
\end{proof}

例如，设 \(R,S\) 为含幺环，\(\varphi:R\to S\) 为保单位元的环同态，以下模均为幺性左模。限制标量 \(\operatorname{Res}_{\varphi}:S\text{-}\mathbf{Mod}\to R\text{-}\mathbf{Mod}\) 同时出现在两种伴随中：
\[
 S\otimes_R-\;\dashv\;\operatorname{Res}_{\varphi}
 \;\dashv\;\operatorname{Hom}_R({}_RS,-).
\]
张量中的 \(S\) 以 \(t\cdot r=t\varphi(r)\) 接受右 \(R\) 作用；余诱导中的 \({}_RS\) 则以 \(r\cdot t=\varphi(r)t\) 接受左 \(R\) 作用，并在 \(\operatorname{Hom}_R({}_RS,J)\) 上规定左 \(S\) 作用 \((s\cdot h)(t)=h(ts)\)。这些伴随双射已在第二卷命题6.3.4中证明；当前两种保持内射性的判断使用不同的左伴随。

\ProseParagraphBreak
若 \(S\) 作为右 \(R\) 模平坦，则 \(S\otimes_R-\) 正合，所以作为它的右伴随，限制标量保持内射对象。没有平坦性时这可能失败：\(\mathbb Z/2\mathbb Z\) 作为 \(\mathbb F_2\) 模内射，作为整数模却不内射。另一方面，限制标量本身始终正合，因为它保留底交换群和映射；因此它的右伴随，即余诱导 \(\operatorname{Hom}_R({}_RS,-)\)，总把内射左 \(R\) 模送到内射左 \(S\) 模，无需平坦性。

\begin{theorem}\label{b4:thm:derived-composite-exact-first}
设 \(\mathcal A,\mathcal B,\mathcal C\) 为交换范畴，\(\mathcal A,\mathcal B\) 有足够多内射对象。设 \(F:\mathcal A\to\mathcal B\)、\(G:\mathcal B\to\mathcal C\) 为加性函子，\(F\) 正合，\(G\) 左正合，且每个内射对象 \(I\in\mathcal A\) 的像 \(F(I)\) 都对 \(G\) 零调。则对每个对象 \(M\in\mathcal A\) 及整数 \(n\ge0\)，有自然同构
\[
 R^n(GF)(M)\simeq R^nG(F(M)).
\]
\end{theorem}
\begin{proof}
取 \(M\to I^\bullet\) 的内射消解。同一个复形 \(GF(I^\bullet)\) 将用于两种计算。因为 \(I^\bullet\) 是 \(M\) 的内射消解，定义直接给出
\[
 R^n(GF)(M)=H^n(GF(I^\bullet)).
\]
另一方面，\(F\) 正合保证 \(F(M)\to F(I^\bullet)\) 仍为正合消解；内射像对 \(G\) 零调则保证这条消解能够计算 \(G\) 的导出。由\cref{b4:thm:acyclic-resolution}，
\[
 H^n(GF(I^\bullet))\cong R^nG(F(M)).
\]
两条识别复合就是所需同构。前一识别来自导出函子定义，后一识别来自零调消解定理；两者都与消解比较映射相容，所以对 \(M\) 自然。
\end{proof}

这里的内射像条件不能由 \(F\) 的正合性替代。取素数 \(p\)，令 \(F:\mathbb F_p\text{-}\mathbf{Mod}\to\mathbf{Ab}\) 为忘却函子，\(G=\operatorname{Hom}_{\mathbb Z}(\mathbb Z/p\mathbb Z,-)\)。\(F\) 正合，且 \(\mathbb F_p\) 在源向量空间范畴中内射，但
\[
 R^1G(F(\mathbb F_p))
 =\operatorname{Ext}_{\mathbb Z}^1(\mathbb Z/p\mathbb Z,\mathbb Z/p\mathbb Z)
 \cong\mathbb Z/p\mathbb Z\ne0,
\]
所以这个内射对象的像不是 \(G\)-零调对象。同时，在 \(\bar1\) 处取值把 \(GF(V)\) 自然识别为向量空间 \(V\) 的底交换群，故 \(GF\) 也正合，从而
\[
 R^1(GF)(\mathbb F_p)=0
 \quad\text{而}\quad R^1G(F(\mathbb F_p))\cong\mathbb Z/p\mathbb Z.
\]
复合公式在这里确实失败，不能只由前函子正合来保证。

\subsection{复合导出的障碍}
\label{b4:subsec:1004:02}

若 \(F\) 只有左正合性，\(F(I^\bullet)\) 的正次上同调正是 \(R^qF(M)\)，一般不是 \(F(M)\) 的消解。即使其各项都对 \(G\) 零调，也不能在上一个证明中删去 \(F\) 正合这一条件。

\begin{example}\label{b4:exm:composite-first-not-exact}
取素数 \(p\)，令
\[
 F:\mathbf{Ab}\longrightarrow\mathbb F_p\text{-}\mathbf{Mod},\quad F(A)=A[p],
 \qquad G:\mathbb F_p\text{-}\mathbf{Mod}\longrightarrow\mathbf{Ab}
\]
分别为 \(p\)-挠子群函子及忘却函子。\(F\) 左正合，\(G\) 正合，所以每个 \(F(I)\) 都对 \(G\) 零调。可是取 \(M=\mathbb Z\) 的内射消解 \(0\to\mathbb Z\to\mathbb Q\to\mathbb Q/\mathbb Z\to0\)，施加 \(F\) 后，非增广复形为
\[
 F(\mathbb Q)\longrightarrow F(\mathbb Q/\mathbb Z):
 \qquad 0\longrightarrow\mathbb F_p
 \qquad\text{（次数 }0,1\text{）}.
\]
它在一次有非零上同调，因此不是 \(F(\mathbb Z)=0\) 的正合消解。再作用 \(G\) 不会消去这份上同调，故
\[
 R^1(GF)(\mathbb Z)=\operatorname{Ext}_{\mathbb Z}^1(\mathbb Z/p\mathbb Z,\mathbb Z)
 =\mathbb Z/p\mathbb Z,\qquad R^1G(F(\mathbb Z))=0.
\]
非零项来自先作用 \(F\) 时产生的一次导出，而不是来自 \(G\) 的导出。
\end{example}

\begin{proposition}\label{b4:prop:derived-composite-exact-second}
设 \(\mathcal A,\mathcal B,\mathcal C\) 为交换范畴，\(\mathcal A\) 有足够多内射对象，\(F:\mathcal A\to\mathcal B\) 为左正合加性函子，\(G:\mathcal B\to\mathcal C\) 为正合加性函子。则对每个对象 \(M\in\mathcal A\)，有
\[
 R^n(GF)(M)\simeq G(R^nF(M))\qquad(n\ge0)
\]
自然成立。
\end{proposition}
\begin{proof}
这里使用的是 \(G\) 与取同调交换，不需要把 \(F(I^\bullet)\) 看成正合消解。对于任意上链复形 \(C^\bullet\)，正合性使
\[
 Z^n(G(C^\bullet))\cong G(Z^n(C^\bullet)),\qquad
 B^n(G(C^\bullet))\cong G(B^n(C^\bullet)),
\]
因为余圈是微分的核，余边是前一微分的像。再由 \(G\) 保持余核及
\(H^n(C^\bullet)=\operatorname{coker}(B^n(C^\bullet)\hookrightarrow Z^n(C^\bullet))\)，得到
\[
 H^n(G(C^\bullet))\cong G(H^n(C^\bullet)).
\]
取 \(M\to I^\bullet\) 的内射消解，令 \(C^\bullet=F(I^\bullet)\)，便有
\[
 R^n(GF)(M)=H^n(GF(I^\bullet))
 \cong G(H^n(F(I^\bullet)))=G(R^nF(M)).
\]
这些识别来自核、像及余核的自然态射，故与消解比较相容，给出对 \(M\) 的自然性。
\end{proof}

这两个结果处理的是不同的简化情形：前函子正合且把内射对象送到后函子的零调对象时，只剩后函子的导出；后函子正合时，它直接作用于前函子的导出，无需对内射像另加条件。两者都只有左正合性时，两类高次信息必须同时保留。

\subsection{Grothendieck 谱序列的假设}
\label{b4:subsec:1004:03}

设 \(\mathcal A,\mathcal B,\mathcal C\) 为交换范畴，前两者有足够内射对象，\(F:\mathcal A\to\mathcal B\)、\(G:\mathcal B\to\mathcal C\) 均为左正合加性函子。关键附加条件是
\[
 R^pG(F(I))=0\qquad(p>0,\ I\;\text{为}\;\mathcal A\;\text{的内射对象}).
\]
在这些条件下，\secref{b4:sec:1503}中的\cref{b4:thm:grothendieck-spectral-sequence}将用双复形建立 Grothendieck 谱序列，其第二页和目标分别为
\[
 E_2^{p,q}=R^pG(R^qF(M))
 \quad\Longrightarrow\quad R^{p+q}(GF)(M),\qquad p,q\ge0.
\]
指标 \(q\) 记录先作用 \(F\) 时产生的导出次数，\(p\) 记录再对 \(R^qF(M)\) 作用 \(G\) 时的导出次数，最终目标的次数是 \(p+q\)。箭头表示每个目标 \(R^n(GF)(M)\) 带有有限递减滤过
\[
\begin{aligned}
 0&=\mathcal F^{n+1}R^n(GF)(M)\subseteq\mathcal F^nR^n(GF)(M)\subseteq\cdots\\
 &\subseteq\mathcal F^0R^n(GF)(M)=R^n(GF)(M),
\end{aligned}
\]
而极限页描述相邻商：
\[
 E_\infty^{p,n-p}\cong
 \mathcal F^pR^n(GF)(M)/\mathcal F^{p+1}R^n(GF)(M)
 \qquad(0\le p\le n).
\]
因此第二页的各项不能直接读成目标的直和因子；还须经过后续各页，最终得到的也只是滤过的相邻商。谱序列的微分、收敛及滤过将在\chapref{b4:ch:14}建立。上述假设与公式可对照 Weibel\cite[Cohomological Setup 5.8.1, Grothendieck Spectral Sequence Theorem 5.8.3]{Weibel1994HomologicalAlgebra}；该书先作用 \(G\) 再作用 \(F\)，与这里的函子字母次序相反。

\ProseParagraphBreak
在内射像条件成立的前提下，若 \(F\) 正合，则 \(R^qF(M)=0\) 对 \(q>0\) 成立，只剩 \(q=0\) 行；总次数 \(n\) 只可能有 \(E_2^{n,0}=R^nG(F(M))\) 非零。若 \(G\) 正合，内射像条件自动成立，且 \(R^pG(-)=0\) 对 \(p>0\) 成立，只剩 \(p=0\) 列；总次数 \(n\) 只可能有 \(E_2^{0,n}=G(R^nF(M))\) 非零。这两种情形都没有可能的非零后续微分，每个目标滤过也只剩一个相邻商，所以分别恢复前两小节的自然同构。一般情形中，先对 \(F(I^\bullet)\) 作进一步消解，再比较两个取同调次序；内射像的零调条件保证其中一个次序计算 \(R^*(GF)\)。低阶正合列的证明见\cref{b4:prop:spectral-five-term}。

\ProseParagraphBreak
\Cref{b4:tab:derived-composite-conditions}比较这三种情形。表中设 \(\mathcal A,\mathcal B,\mathcal C\) 为交换范畴，\(\mathcal A,\mathcal B\) 有足够多内射对象，\(F:\mathcal A\to\mathcal B\)、\(G:\mathcal B\to\mathcal C\) 为加性左正合函子，\(M\in\mathcal A\)，且 \(n,p,q\ge0\)。

\begin{table}[htbp]
\centering
\caption{复合导出的条件与结论}
\label{b4:tab:derived-composite-conditions}
\begin{tabularx}{\textwidth}{@{}>{\raggedright\arraybackslash}X>{\raggedright\arraybackslash}X@{}}
\toprule
附加条件 & 结论 \\
\midrule
\(F\) 正合，且每个内射对象 \(I\in\mathcal A\) 的像 \(F(I)\) 都对 \(G\) 零调
& \(R^n(GF)(M)\simeq R^nG(F(M))\) \\
\(G\) 正合
& \(R^n(GF)(M)\simeq G(R^nF(M))\) \\
内射对象 \(I\in\mathcal A\) 的像 \(F(I)\) 均对 \(G\) 零调
& \(E_2^{p,q}=R^pG(R^qF(M))\)\newline
\(\Longrightarrow R^{p+q}(GF)(M)\) \\
\bottomrule
\end{tabularx}
\end{table}



\begin{chaptersummary}
有限同调维数可以从消解长度定义，也可以由统一的高次消失刻画。维数平移使一个次数的消失转化为合冲模的次数一判据，因而保证任意消解都能在相应位置截短。局部极小消解另有有限生成和 Nakayama 引理的帮助，所以在那里剩余域足以检测终止性。

\ProseParagraphBreak
零调消解依赖正在计算的函子；任意正合消解并不能替代内射消解。对复合函子 \(GF\)，若 \(F\) 正合且把内射对象送到 \(G\)-零调对象，则只需计算 \(G\) 在 \(F(M)\) 上的导出；若 \(G\) 正合，则可直接将 \(G\) 作用于 \(F\) 的导出。两者都只有左正合性时，在同一内射像零调条件下，Grothendieck 谱序列把前函子的高次导出与后函子的导出共同组织起来。
\end{chaptersummary}

\BookExercises

\begin{exercise}\label{b4:ex:10-dimension-summands}
设 \(R\) 为环，\(A,B\) 为左 \(R\)-模。证明 \(\operatorname{pd}(A\oplus B)=\max\{\operatorname{pd}A,\operatorname{pd}B\}\)，并证明内射维数和平坦维数的同类公式。
\end{exercise}
\begin{solution}
有限直和的消解可逐项直和，给出不超过右侧的上界。反向，\(A,B\) 都是 \(A\oplus B\) 的直和因子；对任意左 \(R\)-模 \(N\)，有
\[
\operatorname{Ext}_R^i(A\oplus B,N)
\simeq\operatorname{Ext}_R^i(A,N)\oplus\operatorname{Ext}_R^i(B,N).
\]
若总模的投射维数至多为 \(d\)，则左端对 \(i>d\) 消失，右侧两项也消失；投射维数判据给出 \(\operatorname{pd}A,\operatorname{pd}B\le d\)。内射维数则要检测第二变量，使用
\[
\operatorname{Ext}_R^i(N,A\oplus B)
\simeq\operatorname{Ext}_R^i(N,A)\oplus\operatorname{Ext}_R^i(N,B).
\]
平坦维数使用任意右 \(R\)-模 \(N\) 及 \(\operatorname{Tor}_i^R(N,A\oplus B)\) 的同类分解。若某个直和项维数无限，总模也不可能具有有限维数；零模约定则与最大值公式相容。
\end{solution}

\begin{exercise}\label{b4:ex:10-flat-not-projective}
固定素数 \(p\)。证明 \(\operatorname{fd}_{\mathbb Z}\mathbb Z[1/p]=0\)、\(\operatorname{pd}_{\mathbb Z}\mathbb Z[1/p]=1\)。
\end{exercise}
\begin{solution}
\(\mathbb Z[1/p]\) 是整数局部化，张量它等于将模中的 \(p\) 变为可逆，正合性可由清除分母直接验证，因此平坦。任意同态 \(f:\mathbb Z[1/p]\to\mathbb Z\) 满足 \(f(1)=p^n\cdot f(1/p^n)\) 对每个 \(n\) 成立，只能有 \(f(1)=0\)，再由整数群无挠得 \(f=0\)。若该模投射，则嵌入一个自由交换群；复合任意坐标投影都为零，嵌入也为零，矛盾。整数整体维数为一，于是其投射维数恰为一。
\end{solution}

\begin{exercise}\label{b4:ex:10-three-base-rings}
设 \(p\) 为素数。分别把 \(\mathbb Z/p\mathbb Z\) 看成 \(\mathbb F_p\)、\(\mathbb Z\)、\(\mathbb Z/p^2\mathbb Z\) 上的模，求其投射维数，并对最后一个底环写出全部正次数 Ext 自群。
\end{exercise}
\begin{solution}
在域上维数为零。整数自由消解 \(0\to\mathbb Z\xrightarrow{p}\mathbb Z\to\mathbb Z/p\mathbb Z\to0\) 给出上界一，\(\operatorname{Ext}_{\mathbb Z}^1(\mathbb Z/p\mathbb Z,\mathbb Z)=\mathbb Z/p\mathbb Z\ne0\) 给出下界一。令 \(S=\mathbb Z/p^2\mathbb Z\)，则
\(\cdots\xrightarrow{p}S\xrightarrow{p}S\to\mathbb Z/p\mathbb Z\to0\)
正合，因为乘 \(p\) 的核与像均为 \(pS\)。对它取 \(\operatorname{Hom}_S(-,\mathbb Z/p\mathbb Z)\) 后全部微分为零，故每个 \(\operatorname{Ext}_S^i(\mathbb Z/p\mathbb Z,\mathbb Z/p\mathbb Z)=\mathbb Z/p\mathbb Z\)，投射维数无限。
\end{solution}

\begin{exercise}\label{b4:ex:10-detect-last-degree}
设 \(R\) 为环，非零左模 \(M\) 的投射维数为有限整数 \(d\)。证明存在左模 \(N\) 使 \(\operatorname{Ext}_R^d(M,N)\ne0\)。内射维数有怎样的对应说法？
\end{exercise}
\begin{solution}
若 \(d=0\)，取 \(N=M\)，恒等映射非零。若 \(d>0\) 且对所有 \(N\) 都有 \(\operatorname{Ext}^d(M,N)=0\)，则消失判据给出 \(\operatorname{pd}M\le d-1\)，与最小性矛盾。内射版本把两变量交换：若 \(\operatorname{id}M=d\)，必有 \(N\) 使 \(\operatorname{Ext}^d(N,M)\ne0\)，证明完全由内射消失判据得到。
\end{solution}

\begin{exercise}\label{b4:ex:10-global-dimension-attained}
设 \(R\) 为非零环，\(J\) 为集合，\((M_j)_{j\in J}\) 为一族左 \(R\)-模。沿用 \(\sup\varnothing=-\infty\) 的约定，证明 \(\operatorname{pd}_R(\bigoplus_jM_j)=\sup_j\operatorname{pd}_RM_j\)。由此证明若左整体维数无限，就存在投射维数无限的单个左模。
\end{exercise}
\begin{solution}
若 \(J\) 为空或所有 \(M_j\) 为零，两边都为 \(-\infty\)。否则，若右侧为有限整数 \(d\ge0\)，为每个 \(M_j\) 选长度至多 \(d\) 的投射消解。逐项直和保持正合，投射模的直和仍投射，故得到长度至多 \(d\) 的消解。反向，每个 \(M_j\) 都是总模的直和因子，前一题的消失判据给出其维数不大于总模维数。因此上确界公式成立，右侧为 \(\infty\) 时也成立。若整体维数无限，或已经有无限维数模，或为每个 \(n\ge0\) 选择 \(\operatorname{pd}M_n>n\)。在后一情形令 \(M=\bigoplus_{n\ge0}M_n\)，则 \(\operatorname{pd}M\ge\operatorname{pd}M_n>n\) 对每个 \(n\) 成立。假如 \(\operatorname{pd}M=d<\infty\)，取 \(n=d\) 就得到 \(d>d\) 的矛盾。因此这些可能各自有限的维数，可以由一个直和合成无限投射维数。
\end{solution}

\begin{exercise}\label{b4:ex:10-syzygy-step}
设 \(R\) 为环，\(0\to K\to P\to M\to0\) 是左 \(R\)-模短正合列，\(P\) 投射。证明 \(\operatorname{pd}M=\infty\) 当且仅当 \(\operatorname{pd}K=\infty\)。当 \(M\) 投射时，\(K\) 是否必须为零？
\end{exercise}
\begin{solution}
如果 \(K\) 有有限长度的投射消解，把它与给定列拼接，就得到 \(M\) 的有限长度投射消解。如果 \(M\) 有有限长度的投射消解，则截短判据或短正合列的不等式使 \(K\) 也有这样的消解；\(M=0\) 时 \(K\cong P\) 投射，结论也直接成立。两者同时不存在有限长度的投射消解，便得到所述无限维数的等价。\(M\) 投射时仅能推出给定列分裂、\(K\) 投射；例如取 \(R=\mathbb Z\)，坐标投影 \(\mathbb Z\oplus\mathbb Z\to\mathbb Z\) 就具有非零核 \(\mathbb Z\)。
\end{solution}

\begin{exercise}\label{b4:ex:10-local-top-tor}
设 \(k\) 为域，\(R=k[x,y]_{(x,y)}\)。利用正文中的剩余域消解求 \(\operatorname{pd}_R(x,y)\)，并计算 \(\operatorname{Tor}_1^R((x,y),k)\)。
\end{exercise}
\begin{solution}
在 \(0\to(x,y)\to R\to k\to0\) 中，\(\operatorname{pd}_Rk=2\)，所以一次合冲模 \((x,y)\) 的投射维数为一。维数平移给出 \(\operatorname{Tor}_1^R((x,y),k)\simeq\operatorname{Tor}_2^R(k,k)=k\)。也可截取消解 \(0\to R\to R^2\to(x,y)\to0\)，张量 \(k\) 后首箭头为零，得到相同结果。
\end{solution}

\begin{exercise}\label{b4:ex:10-acyclic-n-divisible}
设 \(n>1\)，\(F=\operatorname{Hom}_{\mathbb Z}(\mathbb Z/n\mathbb Z,-)\)。证明交换群 \(A\) 对 \(F\) 零调当且仅当 \(nA=A\)。给出一个满足该条件而非内射的群。
\end{exercise}
\begin{solution}
由 \(\mathbb Z/n\mathbb Z\) 的长度一自由消解，\(R^1F(A)=A/nA\)，且 \(R^iF(A)=0\) 对 \(i>1\) 成立，因此零调条件只要求乘这个固定的 \(n\) 满射，即 \(nA=A\)。内射交换群则必须可除，即 \(mA=A\) 对每个非零整数 \(m\) 都成立。取 \(A=\mathbb Z[1/n]\)，它满足前一条件；选择不整除 \(n\) 的素数 \(\ell\)，方程 \(\ell x=1\) 在该群中无解，所以它不满足后一条件，因而不是内射整数模。
\end{solution}

\begin{exercise}\label{b4:ex:10-acyclic-kernel-warning}
对 \(F=\operatorname{Hom}_{\mathbb Z}(\mathbb Z/2\mathbb Z,-)\)，在 \(0\to\mathbb Z\to\mathbb Q\to\mathbb Q/\mathbb Z\to0\) 中找出两个零调项和一个非零调项，并用 \(F\) 后的映射说明区别。
\end{exercise}
\begin{solution}
\(\mathbb Q\)、\(\mathbb Q/\mathbb Z\) 都内射，因而对 \(F\) 零调。\(F(\mathbb Q)=0\)，而 \(F(\mathbb Q/\mathbb Z)\) 由 \(1/2+\mathbb Z\) 生成，是 \(\mathbb Z/2\mathbb Z\)。低次长正合列成为
\[
0\longrightarrow F(\mathbb Q/\mathbb Z)
\xrightarrow{\delta}R^1F(\mathbb Z)\longrightarrow0,
\]
所以连接映射 \(\delta\) 是非零同构 \(\mathbb Z/2\mathbb Z\xrightarrow{\sim}\mathbb Z/2\mathbb Z\)。这正是 \(F(\mathbb Q)\to F(\mathbb Q/\mathbb Z)\) 不满射留下的余核，也是零调对象类不无条件地对满射之核封闭的原因。
\end{solution}

\begin{exercise}\label{b4:ex:10-small-acyclic-resolution}
用 \(0\to\mathbb Z\to\mathbb Z[1/2]\to\mathbb Z[1/2]/\mathbb Z\to0\) 计算 \(F=\operatorname{Hom}_{\mathbb Z}(\mathbb Z/2\mathbb Z,-)\) 在 \(\mathbb Z\) 上的导出函子，并说明这不是内射消解。
\end{exercise}
\begin{solution}
两个非增广项都可被 \(2\) 整除，所以由\cref{b4:ex:10-acyclic-n-divisible}的计算对 \(F\) 零调。\(F(\mathbb Z[1/2])=0\)，而商群中被 \(2\) 零化的元素恰为 \(0\) 和 \(1/2+\mathbb Z\)。因此所得上链复形是 \(0\to\mathbb Z/2\mathbb Z\)，次数为 \(0,1\)，给出 \(R^0F(\mathbb Z)=0\)、\(R^1F(\mathbb Z)=\mathbb Z/2\mathbb Z\)，更高次为零。\(\mathbb Z[1/2]\) 不能被 \(3\) 整除，故不是内射群，却已足以参与计算。所得复形的一次上同调非零：零调消解保证其同调计算 \(R^nF\)，并不保证施加 \(F\) 后仍正合。
\end{solution}

\begin{exercise}\label{b4:ex:10-composite-image-condition}
设 \(p\) 为素数，考虑环同态 \(\mathbb Z\to\mathbb F_p\) 对应的限制标量函子 \(\operatorname{Res}\) 与余诱导函子 \(C=\operatorname{Hom}_{\mathbb Z}(\mathbb F_p,-)\)。取内射交换群 \(I=\mathbb Q/\mathbb Z\)，求 \(C(I)\) 的 \(\mathbb F_p\)-模结构，并计算 \(\operatorname{Ext}_{\mathbb Z}^1(\mathbb Z/p\mathbb Z,\operatorname{Res}C(I))\)。由此判断两个函子各是否保持内射对象，说明这与 \(\operatorname{Res}\) 正合并不矛盾。
\end{exercise}
\begin{solution}
在 \(1\) 处取值把 \(C(I)\) 识别为 \(I[p]\)，后者由 \(1/p+\mathbb Z\) 生成。因此 \(C(I)\simeq\mathbb F_p\) 是一维向量空间，也是内射 \(\mathbb F_p\)-模。限制标量后得到 \(\mathbb Z/p\mathbb Z\)，整数上的长度一自由消解给出
\[
\operatorname{Ext}_{\mathbb Z}^1(\mathbb Z/p\mathbb Z,\operatorname{Res}C(I))
\simeq\operatorname{Ext}_{\mathbb Z}^1(\mathbb Z/p\mathbb Z,\mathbb Z/p\mathbb Z)
\simeq\mathbb Z/p\mathbb Z\ne0.
\]
故 \(\operatorname{Res}\) 不保持内射对象。另一方面，\(\operatorname{Res}\dashv C\)，限制标量始终正合，所以由\cref{b4:prop:right-adjoint-injectives}，余诱导 \(C\) 保持内射对象。\(\operatorname{Res}\) 自身的左伴随却是 \(\mathbb F_p\otimes_{\mathbb Z}-\)，它不正合：张量 \(\mathbb Z\xrightarrow{p}\mathbb Z\) 得到非单射的零映射。这两种伴随对内射性的结论因而不同。
\end{solution}

\begin{exercise}\label{b4:ex:10-coinduction-injective}
设 \(R\to S\) 为环同态，\(I\) 为内射左 \(R\)-模。在 \(\operatorname{Hom}_R(S,I)\) 上规定 \((s\cdot f)(t)=f(ts)\)。证明它是内射左 \(S\)-模，并说明这里不要求 \(S\) 作为 \(R\)-模平坦。
\end{exercise}
\begin{solution}
该公式与左 \(R\)-线性相容，且 \(s_1(s_2f)(t)=f(ts_1s_2)\)，定义左 \(S\)-作用。对左 \(S\)-模 \(M\)，映射在 \(1\) 处取值给出
\(\operatorname{Hom}_S(M,\operatorname{Hom}_R(S,I))\simeq\operatorname{Hom}_R(M,I)\)，逆映射把 \(g\) 送到 \(m\mapsto(s\mapsto g(sm))\)。所以余诱导是限制标量的右伴随，而限制标量始终正合。由\cref{b4:prop:right-adjoint-injectives}，余诱导保持内射；此处左伴随是限制标量，不是张量诱导，因此没有平坦性要求。
\end{solution}

\begin{exercise}\label{b4:ex:10-finite-acyclic-length}
在\cref{b4:thm:acyclic-resolution}的环境中，若 \(M\) 有长度 \(d\) 的 \(F\)-零调消解，证明 \(R^nF(M)=0\) 对 \(n>d\) 成立。为什么这不能推出 \(\operatorname{id}M\le d\)？
\end{exercise}
\begin{solution}
该消解施加 \(F\) 后，在高于 \(d\) 的次数没有项，所以其同调为零，零调消解定理给出所需消失。这只针对一个 \(F\)，而内射维数要求全部 \(\operatorname{Ext}^{d+1}(N,M)\) 消失。例如取 \(F=\operatorname{Id}\)，它正合，故 \(R^iF(M)=0\) 对所有对象 \(M\) 及所有 \(i>0\) 成立。于是任意 \(M\) 都有长度零的零调消解 \(0\to M\xrightarrow{1_M}M\to0\)，对 \(M\) 是否内射没有任何限制；取 \(M=\mathbb Z\) 就不是内射群。因此这种消解的长度只针对所选函子，不能读作对象的内射维数。
\end{solution}

\begin{exercise}\label{b4:ex:10-injective-dimension-integers}
求 \(\mathbb Z\)、\(\mathbb Q\)、\(\mathbb Q/\mathbb Z\) 的整数模内射维数，并说明为什么 \(\mathbb Z\) 的投射维数与内射维数不同。
\end{exercise}
\begin{solution}
\(\mathbb Q\) 与 \(\mathbb Q/\mathbb Z\) 都非零且可除，所以内射维数为零。\(0\to\mathbb Z\to\mathbb Q\to\mathbb Q/\mathbb Z\to0\) 给出 \(\operatorname{id}_{\mathbb Z}\mathbb Z\le1\)，而 \(\operatorname{Ext}_{\mathbb Z}^1(\mathbb Z/2\mathbb Z,\mathbb Z)=\mathbb Z/2\mathbb Z\ne0\) 给出等号。\(\mathbb Z\) 自身自由，所以投射维数为零；两种维数分别限制 Ext 的第一、第二变量，没有一般的相等要求。
\end{solution}

\begin{exercise}\label{b4:ex:10-dimension-bounds-sharpness}
在整数模范畴中，对\cref{b4:prop:dimension-short-exact}的三条投射维数不等式及三条内射维数不等式，分别构造取等和严格不等的短正合列。可以让同一条短正合列同时检验多个不等式。
\end{exercise}
\begin{solution}
固定素数 \(p\)。先取三条短正合列
\[
\begin{aligned}
\xi_1:&\quad0\longrightarrow\mathbb Z\xrightarrow{\,p\,}\mathbb Z
\longrightarrow\mathbb Z/p\mathbb Z\longrightarrow0,\\
\xi_2:&\quad0\longrightarrow\mathbb Z\longrightarrow\mathbb Q
\longrightarrow\mathbb Q/\mathbb Z\longrightarrow0,\\
\xi_3:&\quad0\longrightarrow\mathbb Z/p\mathbb Z\longrightarrow\mathbb Z/p\mathbb Z\oplus\mathbb Z
\longrightarrow\mathbb Z\longrightarrow0.
\end{aligned}
\]
前两列使用通常的乘法、包含和商映射，第三列使用直和的包含与投影。三个非零项的投射维数依次为 \((0,0,1)\)、\((0,1,1)\)、\((1,1,0)\)。这里 \(\mathbb Q/\mathbb Z\) 不是投射整数模：含非零挠元的群不能嵌入自由交换群，而整数整体维数为一给出上界。在 \(\xi_1\) 中，估计 \(B\) 得到 \(0<\max\{0,1\}=1\)，而估计 \(A,C\) 分别得到 \(0=\max\{0,1-1\}\) 和 \(1=\max\{0,0+1\}\)。在 \(\xi_2\) 中，估计 \(A\) 得到 \(0<\max\{1,1-1\}=1\)，而估计 \(B,C\) 的两式都为 \(1=1\)。在 \(\xi_3\) 中，估计 \(C\) 得到 \(0<\max\{1,1+1\}=2\)，而估计 \(A,B\) 的两式都为 \(1=1\)。所以每一列恰有一条所考察的上界严格，另外两式取等。

内射维数方面，\(\xi_2\) 的三项维数为 \((1,0,0)\)。估计 \(B\) 得到 \(0<\max\{1,0\}=1\)，而估计 \(A,C\) 分别得到 \(1=\max\{0,0+1\}\) 和 \(0=\max\{0,1-1\}\)。\(\xi_1\) 的三项维数为 \((1,1,1)\)：\(\mathbb Z/p\mathbb Z\) 的内射维数为一，因为 \(\operatorname{Ext}_{\mathbb Z}^1(\mathbb Z/p\mathbb Z,\mathbb Z/p\mathbb Z)\ne0\)，而整数整体维数为一。此时估计 \(A\) 得到 \(1<\max\{1,1+1\}=2\)，估计 \(B,C\) 则都为 \(1=1\)。最后取分裂列
\[
0\longrightarrow\mathbb Z\longrightarrow\mathbb Z\oplus\mathbb Q
\longrightarrow\mathbb Q\longrightarrow0.
\]
其三项内射维数为 \((1,1,0)\)。估计 \(C\) 得到 \(0<\max\{1,1-1\}=1\)，估计 \(A,B\) 则都为 \(1=1\)。故六条不等式都可以取等，也都可能严格。
\end{solution}

